\documentclass[11pt]{article}
\usepackage[margin=1.1in]{geometry}
\usepackage{amsmath,amssymb,amsthm}
\usepackage{enumitem}
\usepackage{booktabs}
\usepackage{xcolor}
\definecolor{linkblue}{RGB}{25,60,120}
\usepackage[colorlinks=true,linkcolor=linkblue,citecolor=linkblue,urlcolor=linkblue]{hyperref}
\usepackage[capitalize]{cleveref}
\emergencystretch=2em

\newtheorem{theorem}{Theorem}[section]
\newtheorem{lemma}[theorem]{Lemma}
\newtheorem{proposition}[theorem]{Proposition}
\newtheorem{corollary}[theorem]{Corollary}
\newtheorem{fact}[theorem]{Fact}
\theoremstyle{definition}
\newtheorem{definition}[theorem]{Definition}
\newtheorem{problem}[theorem]{Open Problem}
\theoremstyle{remark}
\newtheorem{remark}[theorem]{Remark}
\crefname{problem}{Problem}{Problems}
\Crefname{problem}{Problem}{Problems}
\crefname{lemma}{Lemma}{Lemmas}
\Crefname{lemma}{Lemma}{Lemmas}
\crefname{proposition}{Proposition}{Propositions}
\Crefname{proposition}{Proposition}{Propositions}
\crefname{corollary}{Corollary}{Corollaries}
\Crefname{corollary}{Corollary}{Corollaries}
\crefname{definition}{Definition}{Definitions}
\Crefname{definition}{Definition}{Definitions}
\crefname{remark}{Remark}{Remarks}
\Crefname{remark}{Remark}{Remarks}

\newcommand{\F}{\mathbb{F}}
\newcommand{\B}{B}
\newcommand{\RS}{\mathrm{RS}}
\newcommand{\eca}{\varepsilon_{\mathrm{ca}}}
\newcommand{\emca}{\varepsilon_{\mathrm{mca}}}
\newcommand{\dmin}{\delta_{\min}}
\newcommand{\dist}{\Delta}
\newcommand{\distS}{\Delta_S}
\newcommand{\Lst}{\Lambda}
\newcommand{\Syn}{\operatorname{Syn}}
\newcommand{\Hb}{\mathrm{H}_2}
\newcommand{\eone}{e_1}
\newcommand{\UU}{\mathbb{U}}

\title{Universal Field-Size Caps and a Two-Sided Sandwich\\ for Mutual Correlated Agreement on Smooth Reed--Solomon Domains}
\author{Przemek Chojecki\\ \small ulam.ai}
\date{July 2, 2026}

\begin{document}
\maketitle

\begin{abstract}
We study the mutual correlated agreement (MCA), correlated agreement (CA), and list challenges of Arnon--Boneh--Fenzi \cite{ABF26} for Reed--Solomon codes $C=\RS[\F,D,k]$ on smooth evaluation domains.  The goal is to locate the threshold $\delta_C^*(\varepsilon)$ at cryptographic targets such as $2^{-128}$.

On the unsafe side we prove a field-size-universal cap.  For every finite field with $n<|\F|<2^{256}$ and every smooth multiplicative row of rate $\rho\in\{\frac12,\frac14,\frac18,\frac1{16}\}$ with $k\le2^{40}$,
\[
\emca(C,\delta)\;>\;\frac1{2k}\Bigl(1-\frac n{|\F|}\Bigr)\;\ge\;2^{-86}\;\gg\;2^{-128}
\]
throughout the printed near-capacity band, with the same bound for $\eca$ on the integer-admissible subinterval.  The proof is elementary: quotient-locator pigeonhole floors are composed with a self-contained deep-point conversion, with no value-set counting, exponential-sum estimate, or equidistribution input.  The same mechanism extends to map-smooth domains, Chebyshev line rounds, bivariate circle codes, and rational scale maps relevant to genus-one rows.

On the safe side we prove two self-contained reductions.  First, for every linear code, MCA has error at most $(\lfloor\delta n\rfloor+1)/q$ below one third of the minimum distance.  Second, below one half of the minimum distance, MCA is bounded by ordinary CA up to an explicit tangent term.  Combining this with the classical unique-decoding proximity gap of Ben-Sasson--Carmon--Ishai--Kopparty--Saraf \cite{BCIKS20} gives a conditional half-distance safe frontier; a self-contained half-Johnson CA bound gives an import-free alternative.

For the deployed KoalaBear-sextic and Mersenne-$31$ circle rows, these results give explicit two-sided threshold intervals at the natural target of each row ($2^{-128}$ and $2^{-100}$ respectively).  The remaining gap is an explicit aperiodic band problem, represented by finite rank-pivot and compressed-eliminant certificate language.  The paper also records explicit witness words, explicit certifying pairs in residue-line normal form, an optimized failure profile, and exact integer certificates for the printed deployed-row verdicts; the sole optional non-elementary input is the imported half-distance CA theorem, isolated explicitly.
\end{abstract}

\medskip
\noindent\textbf{Keywords:} Reed--Solomon codes, proximity gaps, correlated agreement, list decoding, FRI, STARKs.\\
\textbf{MSC 2020:} 94B35 (primary); 11T71, 68Q25 (secondary).

{\small\setcounter{tocdepth}{2}\tableofcontents}

\section{Introduction}\label{sec:intro}

Proximity testing for Reed--Solomon codes underlies the soundness of modern SNARKs, and the sharpest available soundness accounting is phrased through \emph{mutual correlated agreement} (MCA): a random linear combination $f_1+\gamma f_2$ being $\delta$-close to the code should force $f_1,f_2$ to be simultaneously explained by codewords on a \emph{common} agreement set, except with small probability $\emca(C,\delta)$ over $\gamma$. The survey of Arnon, Boneh and Fenzi \cite{ABF26} crystallizes the state of the art into a \emph{grand MCA challenge}: for $C=\RS[\F,L,k]$ over a smooth domain $L$, with rate $\rho\in\{1/2,1/4,1/8,1/16\}$, degree $k\le2^{40}$, and field size $|\F|<2^{256}$, determine the largest proximity parameter $\delta^*$ at which $\emca(C,\delta)\le2^{-128}$ can hold.

Positive results reach the Johnson bound $1-\sqrt\rho$. On the negative side, \cite{Cho26b} showed $\delta^*_C(2^{-128})<1-\rho-1/64$ for all prime fields of size up to roughly $2^{150}$ (per-rate reaches between $2^{150.6}$ and $2^{161.8}$), by exhibiting quotient-locator witnesses whose bad-slope sets are value sets of restricted symmetric sums over small subgroups; it further proved a \emph{conditional} universal cap over every field of size at most $2^{256}$ --- error $2^{-42}$ at gap $2^{-11}$, assuming $|\F|\ge2n$ --- by composing quotient-core lists with the same conversion isolated and proved directly here (its universal-cap theorem). Above the value-set reach, what remained open was the \emph{error-one} question, identified there with the \emph{per-fiber collision problem}: controlling, at a \emph{fixed} large prime, the collisions of characteristic-zero locator values under reduction. Extension fields, the setting actually deployed (\cite[\S6.3]{ABF26} works over a sextic extension of KoalaBear), posed an additional obstruction: by the subfield-confinement theorem of \cite{Cho26b}, refined in \cref{sec:subfield}, every explicit witness of \cite{Cho26a,Cho26b} is $\B$-rational and contributes only $p^{-5}<2^{-154}$ over the deployed extension.

This paper sharpens the universal cap --- gap $2^{-9}$ (resp.\ $2^{-10}$) in place of $2^{-11}$, error $2^{-86}$ with no condition relating $n$ and $|\F|$, and extension fields handled directly --- for every field below $2^{256}$ at once, by changing what is counted.


\begin{theorem}[universal unsafe cap, informal; see \cref{thm:main,cor:grand}]\label{thm:informal}
Let $\F$ be any finite field with $|\F|<2^{256}$, let $D\subseteq\F^\times$ be a smooth multiplicative evaluation domain of order $n$, and let $C=\RS[\F,D,k]$ have $k=\rho n\le2^{40}$ with $\rho\in\{1/2,1/4,1/8,1/16\}$. Assume $2^{10}\mid n$ for $\rho\in\{1/2,1/4,1/8\}$, and $2^{11}\mid n$ for $\rho=1/16$. Then
\[
\emca(C,\delta)>2^{-86}\gg2^{-128}
\]
for every $\delta\in[1-\rho-2^{-9},1-\rho)$ at rates $1/2,1/4,1/8$, and for every $\delta\in[1-\rho-2^{-10},1-\rho)$ at rate $1/16$. The same lower bound holds for $\eca(C,\delta)$ on the integer-admissible subinterval ending at $1-\rho-1/n$. If $|\F|\ge2n$, the lower bound is $\ge2^{-42}$. Hence $\delta^*_C(2^{-128})\le1-\rho-2^{-9}$ for $\rho\in\{1/2,1/4,1/8\}$ and $\delta^*_C(2^{-128})\le1-\rho-2^{-10}$ for $\rho=1/16$.
\end{theorem}

Combining the unsafe caps, the safe-side theorems, and the certificate refinements gives the following two-sided summary.

\begin{theorem}[two-sided deployed intervals, informal; see \cref{thm:sandwich2}, \cref{tab:scanner2}, \cref{cor:circle-anyfield}]\label{thm:informal-sandwich}
For every deployed row --- the KoalaBear-sextic multiplicative row and the Mersenne-$31$ circle rows, at their stated challenge targets --- the grand MCA threshold is pinned two-sidedly: at rate $\rho=\tfrac12$,
\begin{align*}
\delta^*_C\ &\in\ \bigl[\tfrac{1-\rho}3,\ \tfrac{15331}{32768}\bigr] &&\text{unconditionally, and}\\
\delta^*_C\ &\in\ \bigl[\tfrac{1-\rho}2,\ \tfrac{15331}{32768}\bigr] &&\text{modulo one classical import \cite{BCIKS20},}
\end{align*}
so the threshold is known within a factor smaller than two \textup(the upper edge is $\tfrac{15331}{32768}\approx0.46787$\textup). The same two-sided sandwich holds for genus-one \textup(ECFFT\textup) rows and for bivariate circle rows over every challenge field, whether or not it contains $i$.
\end{theorem}

The proof composes two ingredients. The first (\cref{lem:fiber}) is a pigeonhole lower bound on list sizes at a single received word: the quotient locators of \cite{Cho26a} attach to each $\ell$-subset $A$ of the order-$N$ quotient coset $Q=D^{a}$ a codeword $c_A$ at distance $\le1-\rho-t/N$ from the word $u_{z_A}=x^{k+ta}+z_Ax^{k+(t-1)a}$, where $z_A=-\eone(A)$; pigeonholing the $\binom N\ell$ subsets over their $\le|\B|$ slope values produces one word $u_z$ carrying a list of size $\binom N\ell/|\B|$. Two refinements of \cite[Main Thm.\,(d)]{Cho26a} matter here: we run the construction at \emph{slack two} ($t=2$), so that the listed codewords have degree $\le k$ and live in $\RS[\F,D,k+1]$ without any divisibility demand on $k+1$ (which is odd in all challenge instantiations, while $a$ is a $2$-power); and we pigeonhole over the \emph{subfield} $\B$ containing $D$ rather than over $\F$, which is what makes the extension-field case go through. The second ingredient is the direct deep-point conversion of \cref{thm:A}: in the integer-error range $f\le n-k-1$, if $\eca(C,\delta)$ is below $\eta(\frac1k-\frac n{k|\F|})$, then \emph{every} list of $\RS[\F,D,k+1]$ at radius $\delta$ has size at most $\lceil|\F|\,\eca/(1-\eta)\rceil$. Taking $\eta=\frac12$, the heavy fiber of \cref{lem:fiber} violates this whenever $\binom N\ell\ge|\B|(|\F|/k+1)$ --- an inequality with hundreds of bits of slack throughout the challenge envelope --- and the contrapositive yields $\eca(C,\delta)>\frac1{2k}(1-n/|\F|)$ throughout the integer-admissible range. The MCA cap then persists up to capacity by support-wise MCA monotonicity.

No value-set counting, exponential-sum estimate, or equidistribution hypothesis enters: pigeonhole fibers exist over every field unconditionally, and \cref{thm:A} converts list mass into correlated-agreement error directly. The cost of this route, relative to the error-one results of \cite{Cho26a,Cho26b} below $2^{150}$, is an error of order $1/k$ rather than $1-o(1)$, and a gap $2^{-9}$ at the first three prize rates (uniformly $2^{-10}$ across all four), rather than the fixed gap $1/64$; \cref{sec:discussion} quantifies the trade-off and what remains open.

Beyond the universal cap (\cref{cor:grand}, with a two-tier summary combining it with \cite{Cho26b}), we prove: a deployed-parameter corollary at the KoalaBear-sextic instantiation of \cite[\S6.3]{ABF26}, with $\emca>2^{-22}$ at gap $2^{-7}$ and the same $\eca$ lower bound on the integer-admissible subinterval (\cref{cor:deployed}), extended to every interleaved row of \cite[Tables~2--3]{ABF26} via an interleaving-transfer lemma (\cref{lem:inter,cor:rows}); an independent slacked variant through \cite[Thm.~1.9]{BCHKS25} for robustness (\cref{prop:slacked}); and a subfield-confinement lemma (\cref{lem:confine}) showing that for lines with $\B$-valued words over an extension $\F/\B$, every CA- or MCA-bad slope lies in $\B$. The latter has a structural consequence (\cref{cor:Fvalued}): at deployed parameters the lines certifying \cref{cor:deployed} are necessarily \emph{not} $\B$-rational, so the failure proved here is fiber-borne and currently non-constructive; exhibiting explicit certifying lines is posed as Open Problem~\ref{prob:explicit}.

A further axis is that the counting side of the composition is not tied to the multiplicative structure of $D$. \Cref{sec:isogeny} proves the fiber lemma for arbitrary polynomial folding maps with complete uniform fibers (\cref{def:map-smooth,lem:phi-fiber}), removes the divisibility hypothesis $a\mid k$, and instantiates the resulting universal cap (\cref{thm:phi-cap}) on the deployed isogeny-smooth family: twin-coset $x$-coordinate domains of the circle group over $p\equiv3\pmod4$ carry exact Chebyshev fibers at every dyadic scale, and over challenge fields containing $i$ the bivariate circle code is diagonally equivalent to a Reed--Solomon code on a torus twin coset. The Mersenne-$31$/QM31 circle family of \cite{HLP24} thereby acquires the same negative bands as the multiplicative rows --- error $>2^{-23}$ at gap $2^{-7}$ at the deployed sizes, gaps $2^{-9}$--$2^{-11}$ universally (\cref{cor:circle-grand,cor:circle-deployed}) --- answering the circle case of the isogeny-smooth problem with the elliptic (ECFFT) case reduced to a precise lacunarity obstruction (\cref{rem:ecfft-obstruction}, Open Problem~\ref{prob:isogeny}).

The paper has three main components.  First, it gives a field-size-universal unsafe construction: quotient-fiber and graded-prefix floors produce large lists at one received word, and a direct deep-point conversion turns that list mass into CA and MCA failure near capacity.  Second, it supplies the safe side needed to interpret these failures as a threshold statement: a self-contained deep-regime theorem proves MCA soundness below one third of the distance, while a pincer theorem reduces MCA to ordinary CA up to half the distance, with the half-distance endpoint obtained from a single imported proximity-gap theorem.  Third, it packages the deployed consequences as finite certificates: the scanner outputs the safe edge, unsafe edge, and remaining open band for each row, and every printed inequality is checkable by exact integer arithmetic.

After these components, the remaining mathematical question is no longer a list of bookkeeping tasks.  It is the aperiodic band problem: between the half-distance safe frontier and the entropy--subfield unsafe frontier, prove a correlated-agreement theorem crossing toward capacity, or find a new floor construction beating the graded-prefix envelope (Open Problem~\ref{prob:band}).

\paragraph{Roadmap.} \Cref{sec:prelim} fixes definitions, aligned with \cite[Defs.~4.1, 4.3]{ABF26}, proves the direct deep-point conversion used for the main cap, and records the separate BCHKS/ABF slacked fallback. \Cref{sec:fiber} proves the fiber and interleaving lemmas and the conservative quotient-support upper ledger, \cref{sec:main} the main theorem, \cref{sec:grand} the challenge-level and deployed corollaries, and \cref{sec:subfield} subfield confinement. \Cref{sec:isogeny} extends the cap beyond multiplicative smoothness: it proves the map-smooth, divisibility-free fiber lemma and universal cap, the exact Chebyshev-fiber and torus twin-coset geometry of circle domains, and the resulting circle-code caps at the deployed Mersenne-$31$/QM31 parameters. \Cref{sec:prize-program} formulates the integer staircase and the certificate interface used to compare safe and unsafe radii.  \Cref{sec:discharge} develops the structural ledgers, extension-line lift, residual-band language, scanner, and finite certificate core.  \Cref{sec:pincer} proves the safe-side pincer below half the distance and the scale-optimized unsafe floors, then assembles the narrowed two-sided sandwich.  \Cref{sec:answers} treats rational scales, circle-code transport over every challenge field, explicit witnesses, optimized failure profiles, and the deployed-row certificate table.  \Cref{sec:discussion} discusses scope, limitations, and the remaining open problems.  Throughout, companion results of \cite{Cho26a,Cho26b} are cited by their printed theorem names rather than numbers, since the numbering of those manuscripts is not yet frozen.

\section{Preliminaries}\label{sec:prelim}

\paragraph{Codes and distances.} Throughout, $\B\subseteq\F$ are finite fields, $q:=|\F|$, and $D=\alpha H\subseteq\B^\times$ is a multiplicative coset of a cyclic subgroup $H$ of order $n$; we write $D$ for the evaluation domain (denoted $L$ in \cite{ABF26}). The subgroup case is $\alpha=1$, and taking $\B=\F$ is always allowed. For $k\le n$,
\[
\RS[\F,D,k]\;:=\;\bigl\{(p(x))_{x\in D}\;:\;p\in\F[X],\ \deg p<k\bigr\}\subseteq\F^n,
\qquad \rho:=k/n .
\]
For $u,v\in\F^n$, $\dist(u,v)$ is the relative Hamming distance, $\dist(u,C):=\min_{c\in C}\dist(u,c)$, and the relative minimum distance of $C=\RS[\F,D,k]$ is $\dmin(C)=1-\rho+1/n$. For $S\subseteq D$, $\distS(u,v)$ and $\distS(u,C)$ denote the corresponding restricted relative distances on $S$. Lists are denoted
\[
\Lst(C,\delta,u):=\{c\in C:\dist(u,c)\le\delta\},\qquad \Lst(C,\delta):=\max_{u\in\F^n}|\Lst(C,\delta,u)| .
\]
For $s\ge1$ the $s$-interleaved code is $C^{\equiv s}:=\{(c^{(1)},\dots,c^{(s)}):c^{(i)}\in C\}\subseteq(\F^n)^s$ with the column distance $\dist_s(F,G):=\frac1n\,\bigl|\{j:\exists i,\ f^{(i)}_j\ne g^{(i)}_j\}\bigr|$, and $\dist_s(F,C^{\equiv s}):=\min_{\mathbf c\in C^{\equiv s}}\dist_s(F,\mathbf c)$.

\begin{definition}[correlated agreement error; {\cite[Def.~4.1]{ABF26}} up to notation]\label{def:ca}
For $0\le\delta_{\mathrm{fld}}\le\delta_{\mathrm{int}}<1$,
\[
\eca(C,\delta_{\mathrm{fld}},\delta_{\mathrm{int}})\;:=\;\max_{f_1,f_2\in\F^n}\ \Pr_{\gamma\leftarrow\F}\Bigl[\ \dist(f_1+\gamma f_2,\,C)\le\delta_{\mathrm{fld}}\ \wedge\ \dist_2\bigl((f_1,f_2),\,C^{\equiv2}\bigr)>\delta_{\mathrm{int}}\ \Bigr],
\]
and the no-proximity-loss variant is $\eca(C,\delta):=\eca(C,\delta,\delta)$. A slope $\gamma$ realizing the event for a fixed pair $(f_1,f_2)$ is called \emph{CA-bad} for that pair.
\end{definition}

\begin{definition}[mutual correlated agreement error, support-wise; {\cite[Def.~4.3]{ABF26}} up to notation]\label{def:mca}
For $\delta\in(0,1)$,
\[
\emca(C,\delta):=\max_{f_1,f_2\in\F^n}\Pr_{\gamma\leftarrow\F}\left[
\begin{array}{l}
\exists S\subseteq D,\ |S|\ge(1-\delta)n,\ \distS(f_1+\gamma f_2,C)=0,\\[2pt]
\text{and }\distS\bigl((f_1,f_2),C^{\equiv2}\bigr)>0
\end{array}
\right].
\]
Equivalently, $\gamma$ is MCA-bad for $(f_1,f_2)$ if some large agreement set $S$ explains the point $f_1+\gamma f_2$, but the same set $S$ does not simultaneously explain $f_1$ and $f_2$ by two codewords of $C$.
\end{definition}

\begin{definition}[challenge threshold]\label{def:dstar}
Following the grand MCA challenge of \cite[\S1]{ABF26}, for $\varepsilon^*\in(0,1)$ set
$\delta^*_C(\varepsilon^*):=\sup\{\delta\in(0,1-\rho):\emca(C,\delta)\le\varepsilon^*\}$, with $\sup\emptyset:=0$.
\end{definition}

All displayed threshold intervals use the integer-ball convention: distances are integral, so a radius $\delta$ corresponds to the agreement requirement $\lceil(1-\delta)n\rceil$ and to $\lfloor\delta n\rfloor$ permitted errors. When a theorem certifies that every $\delta\in[\delta_0,1-\rho)$ is unsafe at target $\varepsilon^*$, it follows that $\delta^*_C(\varepsilon^*)\le\delta_0$ under the supremum convention; when it certifies safety at a single radius $\delta_1<1-\rho$, it follows that $\delta^*_C(\varepsilon^*)\ge\delta_1$.

\begin{fact}[{cf.\ \cite[Fact~4.5]{ABF26}}]\label{fact:chain}
For every $\delta\in(0,1)$: $\eca(C,\delta)\le\emca(C,\delta)$.
\end{fact}

\begin{proof}
Fix $(f_1,f_2)$ and let $\gamma$ be CA-bad: $\dist(f_1+\gamma f_2,C)\le\delta$ and $\dist_2((f_1,f_2),C^{\equiv2})>\delta$. Choose $S\subseteq D$ of size at least $(1-\delta)n$ on which $f_1+\gamma f_2$ agrees with a codeword of $C$. Since $(f_1,f_2)$ is not $\delta$-close to $C^{\equiv2}$, there is no set of size at least $(1-\delta)n$ on which $f_1$ and $f_2$ are simultaneously explained by codewords of $C$; in particular this fails on the chosen $S$. Thus $\gamma$ is MCA-bad for the same pair. Taking maxima over pairs gives the claim.
\end{proof}

\begin{lemma}[support-wise MCA monotonicity]\label{lem:mca-monotone}
For every linear code $C\subseteq\F^D$, the function
\[
        \delta\longmapsto \emca(C,\delta)
\]
is nondecreasing.
\end{lemma}

\begin{proof}
Fix a pair $(f_1,f_2)$ and suppose that $\gamma$ is MCA-bad at radius $\delta$. Then there is a set $S\subseteq D$ with
\[
        |S|\ge (1-\delta)n
\]
such that $f_1+\gamma f_2$ is explained by $C$ on $S$, but $(f_1,f_2)$ is not simultaneously explained by two codewords of $C$ on the same $S$. If $\delta'\ge\delta$, then
\[
        |S|\ge (1-\delta)n\ge (1-\delta')n,
\]
so the same set $S$ witnesses that $\gamma$ is MCA-bad at radius $\delta'$. Taking the maximum over pairs gives the claim.
\end{proof}

\paragraph{Conversion theorem.} The main composition below uses the deep-point list-to-CA conversion proved here directly. This is the same consequence previously imported from Crites--Stewart, with the integer-radius admissibility condition made explicit; see \cref{rem:import}. We keep the slacked BCHKS/ABF fallback as a separate imported route.

\begin{theorem}[deep-point list-to-CA conversion]
\label{thm:A}
Let $C=\RS[\F,D,k]$, let $C^+:=\RS[\F,D,k+1]$, let
$q:=|\F|$ and $n:=|D|$, and assume $q>n$.  Let $\delta\in(0,1)$ and set
$f_\delta:=\lfloor \delta n\rfloor$.  Assume
\[
        f_\delta \le n-k-1 .
\]
For any $\eta\in[0,1)$, if
\[
        \eca(C,\delta)
        \le
        \eta\Bigl(\frac1k-\frac n{kq}\Bigr),
\]
then
\[
        \Lst(C^+,\delta)
        \le
        \Bigl\lceil
        \frac{q\,\eca(C,\delta)}{1-\eta}
        \Bigr\rceil .
\]
\end{theorem}

\begin{proof}
Put $\epsilon:=\eca(C,\delta)$ and $f=f_\delta$.
Let $U:D\to\F$ be any received word and suppose that
$P_1,\ldots,P_L\in\F[X]_{<k+1}$ are distinct polynomials whose
restrictions to $D$ all lie in the Hamming ball of relative radius
$\delta$ around $U$.  Since distances are integral, each $P_i$ agrees
with $U$ on at least
\[
        a:=n-f
\]
points of $D$.  The hypothesis $f\le n-k-1$ gives $a\ge k+1$, in particular
$a>k$.

Let $\Omega:=\F\setminus D$, so $|\Omega|=q-n$.  For each
$\alpha\in\Omega$ define the simple-pole line
\[
        f_\alpha(x)=\frac{U(x)}{x-\alpha},
        \qquad
        g_\alpha(x)=-\frac1{x-\alpha}
        \qquad (x\in D).
\]
If $P_i$ agrees with $U$ on a set $S_i$ of size at least $a$, then for
$z_i(\alpha):=P_i(\alpha)$ and every $x\in S_i$,
\[
        f_\alpha(x)+z_i(\alpha)g_\alpha(x)
        =\frac{P_i(x)-P_i(\alpha)}{x-\alpha}.
\]
The right-hand side is the restriction of a polynomial of degree $<k$.
Thus each distinct value among the $P_i(\alpha)$ gives a slope for
which the line point is $\delta$-close to $C$.

We next check the correlated-agreement far condition.  If a polynomial
$G\in\F[X]_{<k}$ agreed with $g_\alpha$ on more than $k$ points of $D$,
then
\[
        (X-\alpha)G(X)+1
\]
would be a polynomial of degree at most $k$ with more than $k$ roots,
but its value at $X=\alpha$ is $1$, a contradiction.  Hence
$(f_\alpha,g_\alpha)$ is not $\delta$-close to $C^{\equiv2}$, because
any common explanation on a support of size at least $a>k$ would in
particular explain $g_\alpha$ there.  Therefore the distinct values of
$P_i(\alpha)$ are CA-bad slopes.

It remains to lower-bound the number of distinct values for some
$\alpha$.  For $i\ne j$, the polynomial $P_i-P_j$ is nonzero of degree
at most $k$, so it has at most $k$ roots in $\Omega$.  Let $M(\alpha)$
be the number of distinct values among $P_1(\alpha),\ldots,P_L(\alpha)$.
Over all $\alpha\in\Omega$, the total number of colliding unordered pairs
$(i,j)$ with $P_i(\alpha)=P_j(\alpha)$ is at most $k\binom L2$.
Equivalently, if the fiber sizes of the map $i\mapsto P_i(\alpha)$ are
$m_v(\alpha)$, then
\[
        \sum_{\alpha\in\Omega}\sum_v \binom{m_v(\alpha)}2
        \le k\binom L2 .
\]
Choose $\alpha\in\Omega$ for which the inner collision count is at most
$k\binom L2/(q-n)$.  If the fiber sizes of the map $i\mapsto P_i(\alpha)$ are
$m_1,\ldots,m_{M(\alpha)}$, then
\[
        \sum_r m_r=L,
        \qquad
        \sum_r m_r^2=L+2\sum_r\binom{m_r}{2}
        \le L+\frac{kL(L-1)}{q-n}.
\]
Cauchy--Schwarz therefore gives
\[
        L^2\le M(\alpha)\sum_r m_r^2,
        \qquad
        M(\alpha)\ge
        \frac{L(q-n)}{q-n+k(L-1)}
        \ge
        \frac{L(q-n)}{q-n+kL} .
\]
Consequently
\[
        \epsilon
        \ge
        \frac1q\cdot \frac{L(q-n)}{q-n+kL}.
\]
Solving this inequality for $L$ gives
\[
        L
        \le
        \frac{\epsilon q(q-n)}{q-n-k\epsilon q},
\]
provided $\epsilon<(q-n)/(kq)$.  Under the stated hypothesis,
$\epsilon\le\eta(q-n)/(kq)$ with $\eta<1$, so the denominator is at
least $(1-\eta)(q-n)$ and hence
\[
        L\le \frac{q\epsilon}{1-\eta}.
\]
Since $U$ and the list were arbitrary, the displayed ceiling bounds the
maximum list size of $C^+$ at radius $\delta$.
\end{proof}

\begin{remark}[relation to Crites--Stewart]
The displayed conversion is the only consequence of Crites--Stewart
needed by the universal-cap argument.  The proof above is included to make the main cap
self-contained.  The separate slacked fallback through BCHKS/ABF remains
an independent imported route; it is not used in the final deployed certificate
grammar and is retained as a robustness check.
\end{remark}

\begin{proposition}[quantitative deep-list floor]
\label{prop:quantitative-deep-list-floor}
Let $C=\RS[\F,D,k]$, let $C^+:=\RS[\F,D,k+1]$, put
$q=|\F|$ and $n=|D|$, and assume $q>n$.  Let
$\delta_0\in[0,1-k/n)$, and suppose that for some received word
$U:D\to\F$ there are $L\ge1$ distinct codewords of $C^+$ in the
closed ball of radius $\delta_0$ around $U$.  Then, for every
$\delta\in[\delta_0,1-k/n)$,
\[
        \eca(C,\delta)
        \ge
        \frac{L(q-n)}{q(q-n+kL)} .
\]
For every such $\delta>0$,
\[
        \emca(C,\delta)
        \ge
        \frac{L(q-n)}{q(q-n+kL)}
\]
as well.  More precisely, some received line has at least
\[
        \left\lceil\frac{L(q-n)}{q-n+kL}\right\rceil
\]
CA-bad finite slopes at every such radius; for $\delta>0$, the same slopes are
MCA-bad.
\end{proposition}

\begin{proof}
Fix $\delta\in[\delta_0,1-k/n)$ and choose distinct
$P_1,\ldots,P_L\in\F[X]_{<k+1}$ in the $\delta_0$-ball around $U$.
They are also in the $\delta$-ball.  Since $\delta<1-k/n$, if
$f=\lfloor\delta n\rfloor$ then $f\le n-k-1$, so every codeword in this
$\delta$-ball agrees with $U$ on at least $a=n-f>k$ points.

Repeat the simple-pole construction in the proof of \cref{thm:A}.  For
$\alpha\in\Omega:=\F\setminus D$ set
\[
        f_\alpha(x)=\frac{U(x)}{x-\alpha},
        \qquad
        g_\alpha(x)=-\frac1{x-\alpha}
        \qquad (x\in D).
\]
For every distinct value of $P_i(\alpha)$, the word
$f_\alpha+P_i(\alpha)g_\alpha$ is $\delta$-close to $C$.  The pair
$(f_\alpha,g_\alpha)$ is not $\delta$-close to $C^{\equiv2}$, because a
common explanation on more than $k$ points would in particular explain
$g_\alpha$ there, and then $(X-\alpha)G(X)+1$ would be a polynomial of degree at most $k$ with more than $k$ roots and value $1$ at $\alpha$.
Thus the distinct values of the map $i\mapsto P_i(\alpha)$ are CA-bad
slopes.

For $i\ne j$, the polynomial $P_i-P_j$ has degree at most $k$ and is
nonzero, hence it has at most $k$ roots in $\Omega$.  Averaging over
$|\Omega|=q-n$ gives an $\alpha$ for which the number of colliding
unordered pairs among the $L$ values is at most $k\binom L2/(q-n)$.  If
$m_1,\ldots,m_M$ are the value multiplicities for this $\alpha$, then
\[
        \sum_r m_r=L,
        \qquad
        \sum_r m_r^2
        \le L+\frac{kL(L-1)}{q-n}.
\]
By Cauchy--Schwarz,
\[
        M\ge
        \frac{L^2}{L+kL(L-1)/(q-n)}
        =
        \frac{L(q-n)}{q-n+k(L-1)}
        \ge
        \frac{L(q-n)}{q-n+kL}.
\]
Since $M$ is an integer, $M\ge\lceil L(q-n)/(q-n+kL)\rceil$.  Dividing
by the $q$ possible finite slopes gives the CA lower bound.  For $\delta>0$, the MCA
bound follows from \cref{fact:chain}.
\end{proof}

\begin{corollary}[deep-list trigger and ceiling]
\label{cor:deep-list-trigger-ceiling}
In the setting of \cref{prop:quantitative-deep-list-floor}, define
\[
        \mathcal E_{q,k}(L)
        :=\frac{L(q-n)}{q(q-n+kL)},
        \qquad
        \mathcal N_{q,k}(L)
        :=\left\lceil\frac{L(q-n)}{q-n+kL}\right\rceil .
\]
Then $\mathcal E_{q,k}(L)$ is increasing in $L$, and
\[
        \mathcal E_{q,k}(L)>
        \frac1{2k}\left(1-\frac nq\right)
        \quad\Longleftrightarrow\quad
        L>\frac{q-n}{k}.
\]
Moreover
\[
        \mathcal E_{q,k}(L)<\frac1k\left(1-\frac nq\right)
\]
for every finite $L$, and the right-hand side is the limiting value as
$L\to\infty$.
\end{corollary}

\begin{proof}
The derivative of $L/(q-n+kL)$ is positive.  The trigger inequality is
\[
        \frac{L}{q-n+kL}>\frac1{2k},
\]
which is equivalent to $kL>q-n$.  The upper bound follows from
$q-n+kL>kL$.
\end{proof}

\begin{remark}[why this is stronger than the contrapositive]
\label{rem:quantitative-floor-vs-contrapositive}
\Cref{thm:A} is optimized for contradiction: once $L>(q-n)/k$, taking
$\eta=1/2$ already forces the printed $(2k)^{-1}(1-n/q)$ lower bound.
\Cref{prop:quantitative-deep-list-floor} keeps the whole interpolation count.
It is the form a staircase scanner should use when a quotient fiber is too small
to cross the $1/(2k)$ trigger but still contributes a nonzero explicit bad-slope
numerator, and also when a very large fiber improves the constant toward
$(1/k)(1-n/q)$.
\end{remark}

\begin{proposition}[punctured simple-pole deep-list floor]
\label{prop:punctured-simple-pole-floor}
Let $C=\RS[\F,D,k]$, let $C^+:=\RS[\F,D,k+1]$, put
$q=|\F|$ and $n=|D|$.  Let
\[
        \Omega_0\subseteq \F\setminus D,
        \qquad m:=|\Omega_0|>0 .
\]
Let $\delta_0\in[0,1-k/n)$, and suppose that for some received word
$U:D\to\F$ there are $L\ge1$ distinct codewords of $C^+$ in the closed ball of
radius $\delta_0$ around $U$.  Then, for every
$\delta\in[\delta_0,1-k/n)$,
\[
        \eca(C,\delta)
        \ge
        \frac{Lm}{q(m+kL)} .
\]
For every such $\delta>0$,
\[
        \emca(C,\delta)
        \ge
        \frac{Lm}{q(m+kL)} .
\]
More precisely, there is some $\alpha\in\Omega_0$ such that the simple-pole line
\[
        f_\alpha(x)=\frac{U(x)}{x-\alpha},
        \qquad
        g_\alpha(x)=-\frac1{x-\alpha}
        \qquad (x\in D)
\]
has at least
\[
        \left\lceil\frac{Lm}{m+kL}\right\rceil
\]
CA-bad finite slopes at every such radius; for $\delta>0$, the same slopes are
MCA-bad.
\end{proposition}

\begin{proof}
Repeat the proof of \cref{prop:quantitative-deep-list-floor}, but average only
over $\Omega_0$ instead of over all of $\F\setminus D$.  The simple-pole
argument works for every $\alpha\notin D$, so every distinct value among
$P_1(\alpha),\ldots,P_L(\alpha)$ gives a CA-bad slope.  For $i\ne j$, the
nonzero polynomial $P_i-P_j$ has degree at most $k$, hence has at most $k$ roots
inside the smaller set $\Omega_0$.  Therefore the total number of colliding
unordered pairs over all $\alpha\in\Omega_0$ is at most $k\binom L2$.  Choose
$\alpha\in\Omega_0$ for which the collision count is at most
$k\binom L2/m$.  If $M$ is the number of distinct values at this $\alpha$, the
same Cauchy--Schwarz calculation gives
\[
        M\ge
        \frac{Lm}{m+k(L-1)}
        \ge
        \frac{Lm}{m+kL} .
\]
Since $M$ is integral, this gives the displayed bad-slope numerator.  Dividing
by the $q$ finite slopes gives the CA lower bound, and the MCA statement for
$\delta>0$ follows from \cref{fact:chain}.
\end{proof}

\begin{corollary}[extension-pole deep-list floor]
\label{cor:extension-pole-deep-list-floor}
Assume in addition that $\B\subsetneq\F$, $D\subseteq\B$, and $|D|\ge2$.  Put
$b:=|\B|$.  In the setting of
\cref{prop:punctured-simple-pole-floor}, take
\[
        \Omega_0:=\F\setminus\B,
        \qquad m=q-b .
\]
Then for every $\delta\in[\delta_0,1-k/n)$,
\[
        \eca(C,\delta)
        \ge
        \mathcal E^{\rm ext}_{q,b,k}(L)
        :=
        \frac{L(q-b)}{q(q-b+kL)} .
\]
For $\delta>0$ the same bound holds for $\emca(C,\delta)$.  Moreover the
witnessing line may be chosen with pole $\alpha\in\F\setminus\B$, and its pole
vector $((x-\alpha)^{-1})_{x\in D}$ is not a scalar multiple over $\F$ of any
$\B$-valued vector.  In particular the witness line is genuinely
extension-valued, not merely a $\B$-rational line viewed over $\F$.
\end{corollary}

\begin{proof}
The lower bound is \cref{prop:punctured-simple-pole-floor} with
$\Omega_0=\F\setminus\B$.  It remains to prove genuine extension-valuedness.
Suppose that, for some $\lambda\in\F^\times$, all numbers
$\lambda/(x-\alpha)$ with $x\in D$ lay in $\B$.  Choose two distinct points
$x,y\in D$.  Then
\[
        \frac{y-\alpha}{x-\alpha}
        =
        \frac{\lambda/(x-\alpha)}{\lambda/(y-\alpha)}
        \in \B .
\]
Writing this ratio as $r\in\B$, we cannot have $r=1$ because $x\ne y$.  Hence
\[
        \alpha=\frac{rx-y}{r-1}\in\B,
\]
contradicting $\alpha\notin\B$.  Thus the pole vector is not projectively
$\B$-valued.  In particular the displayed affine parametrization has no
$\B$-valued direction vector after rescaling.
\end{proof}

\begin{corollary}[extension-pole quotient-remainder floor]
\label{cor:extension-pole-quotient-remainder-floor}
Assume $\B\subsetneq\F$, $D\subseteq\B^\times$ is a multiplicative coset of order
$n\ge2$, put $b:=|\B|$, and let $C=\RS[\F,D,k]$.
Fix a quotient-remainder profile as in
\cref{lem:heaviest-prefix-locator-floor} with $K=k+1$ and agreement $A_0>k$.
Let
\[
        L:=H_{c,m,s}^{k+1}
\]
be the certified heaviest prefix-fiber size, or replace $L$ by either fallback
\[
        \left\lceil \frac{M_{c,m,s}}{I_{c,m,s}^{k+1}}\right\rceil,
        \qquad
        \left\lceil\frac{M_{c,m,s}}{|\B|^{w_c(s,A_0-k-1)}}\right\rceil .
\]
Then, for every
\[
        \delta\in[1-A_0/n,1-k/n),
\]
there is a genuinely extension-valued simple-pole line with at least
\[
        \mathcal N^{\rm ext}_{q,b,k}(L)
        :=
        \left\lceil\frac{L(q-b)}{q-b+kL}\right\rceil
\]
CA-bad finite slopes.  For $\delta>0$, these slopes are MCA-bad.  Equivalently,
\[
        \eca(C,\delta)
        \ge
        \frac{L(q-b)}{q(q-b+kL)},
        \qquad
        \emca(C,\delta)
        \ge
        \frac{L(q-b)}{q(q-b+kL)}
        \quad(\delta>0).
\]
If $L>(q-b)/k$, this gives the clean trigger
\[
        \eca(C,\delta)>
        \frac1{2k}\left(1-\frac bq\right),
        \qquad
        \emca(C,\delta)>
        \frac1{2k}\left(1-\frac bq\right)
        \quad(\delta>0).
\]
\end{corollary}

\begin{proof}
\Cref{lem:heaviest-prefix-locator-floor} gives a $\B$-valued received word with
at least $L$ codewords of $C^+$ in the closed ball of radius $1-A_0/n$.  Apply
\cref{cor:extension-pole-deep-list-floor}.  The trigger inequality is
\[
        \frac{L(q-b)}{q(q-b+kL)}>
        \frac{q-b}{2kq},
\]
which is equivalent to $kL>q-b$.
\end{proof}

\begin{definition}[extension-pole certificate]
\label{def:extension-pole-certificate}
An extension-pole certificate for a list $P_1,\ldots,P_L\in\F[X]_{<k+1}$ and a
received word $U$ consists of:
\begin{enumerate}[label=\textup{(\roman*)},leftmargin=2.2em]
\item a declared proper subfield $\B\subsetneq\F$ containing $D$ and a pole
$\alpha\in\F\setminus\B$;
\item the agreement supports showing that each $P_i$ lies in the declared ball
around $U$;
\item the bucket table for the values $P_i(\alpha)$, with duplicate values
coalesced;
\item the simple-pole line $(f_\alpha,g_\alpha)$ and the proof that
$g_\alpha$ is not projectively $\B$-valued;
\item the final bad-slope numerator, namely the number of distinct buckets in
\textup{(iii)}.
\end{enumerate}
The verifier checks $\alpha\notin\B$, the support agreements, the value buckets,
and the simple-pole far condition.  For quotient-remainder floors, the list in
\textup{(ii)} may itself be supplied by a prefix-fiber certificate from
\cref{def:prefix-fiber-certificate}.
\end{definition}

\begin{remark}[what the extension-pole construction changes]
\label{rem:extension-pole-status}
The subfield-confinement theorem says that $\B$-rational lines are inert over a
proper extension.  \Cref{cor:extension-pole-deep-list-floor} gives the matching
existence theorem on the failure side: any large list can be converted using
poles restricted to $\F\setminus\B$, so the resulting bad line is genuinely
extension-valued.  Thus the extension-line bucket in a prize scanner is no
longer only an abstract warning; it has an explicit simple-pole source and a
stable numerator with $q-b$ replacing $q-n$.  What remains open is not existence
of extension-valued witnesses, but their sharp classification inside the
aperiodic/residue-line ledgers and, for a fully explicit counterexample, printing
a concrete pole $\alpha$ with its value-bucket table.
\end{remark}

\begin{theorem}[{\cite[Thm.~1.9]{BCHKS25}}, as stated in {\cite[Thm.~5.2]{ABF26}}]\label{thm:B}
Let $C=\RS[\F,D,k]$, let $\rho=k/n$, and let $\delta\in(0,1-\rho)$ satisfy $\delta+2/n\le 1-\rho-1/n$. If
\[
\eca\Bigl(C,\ \delta+\tfrac2n,\ 1-\rho-\tfrac1n\Bigr)\;<\;\frac1{2n},
\qquad\text{then}\qquad
\Lst(C,\delta)\;<\;|\F| .
\]
\end{theorem}

\begin{remark}[on imports and radius conventions]\label{rem:import}
\Cref{thm:A} is the only Crites--Stewart-type consequence used by the main cap, and it is now proved above by the simple-pole deep-point argument.  The important admissibility condition is the integer-radius condition
\[
        \lfloor \delta n\rfloor \le n-k-1 .
\]
Consequently, the no-loss CA lower bound proved below should be read only in this admissible range. The Proximity-Prize MCA cap is unaffected: we prove the lower bound at the endpoint
\[
        \delta_N:=1-\rho-\frac2N
\]
and then use support-wise MCA monotonicity to extend the MCA failure to all larger sub-capacity radii.

\Cref{thm:B} remains a separate slacked fallback imported through the BCHKS/ABF route. It is not needed for the main field-size-universal MCA cap.
\end{remark}

\section{Locator fibers and interleaving transfer}\label{sec:fiber}

Fix $N\mid n$, set $a:=n/N$, and let $Q:=D^{a}:=\{x^a:x\in D\}\subseteq\B^\times$, a multiplicative coset of order $N$. The $a$-th power map $D\to Q$ is surjective and its fibers $S_b:=\{x\in D:x^a=b\}$ have exactly $a$ elements each. Since $a\mid n\mid|\B|-1$, the characteristic of $\F$ does not divide $a$, so $X^a-b$ is separable for $b\ne0$ and its full root set inside $D$ is $S_b$ whenever $b\in Q$. For a set $A$, $e_j(A)$ denotes the $j$-th elementary symmetric function of its elements.

\begin{lemma}[locator fibers are lists]\label{lem:fiber}
Let $\B\subseteq\F$, let $D\subseteq\B^\times$ be a multiplicative coset of order $n$, and let $k$ satisfy $a\mid k$ and $\rho=k/n$. Write $\ell_1:=\rho N+1$ and $\ell_2:=\rho N+2$.
\begin{enumerate}[label=\textup{(\roman*)}]
\item If $\ell_1\le N$, there exists $z\in\B$ such that, with $u_z:=\bigl(x^{k+a}+z\,x^{k}\bigr)_{x\in D}\in\F^n$,
\[
\bigl|\Lst\bigl(\RS[\F,D,k],\ 1-\rho-\tfrac1N,\ u_z\bigr)\bigr|\;\ge\;\binom{N}{\ell_1}\Big/|\B| .
\]
\item If $\ell_2\le N$, there exists $z\in\B$ such that, with $u_z:=\bigl(x^{k+2a}+z\,x^{k+a}\bigr)_{x\in D}\in\F^n$,
\[
\bigl|\Lst\bigl(\RS[\F,D,k+1],\ 1-\rho-\tfrac2N,\ u_z\bigr)\bigr|\;\ge\;\binom{N}{\ell_2}\Big/|\B| .
\]
\end{enumerate}
In both cases the words $u_z$ are $\B$-valued and all listed codewords lie in $\RS[\B,D,k]$ \textup(resp.\ $\RS[\B,D,k+1]$\textup).
\end{lemma}

\begin{proof}
We prove (ii); part (i) is the identical computation one degree rung lower and is \cite[Main Thm.\,(d)]{Cho26a} sharpened by pigeonholing over $\B$ in place of $\F$.

Let $A\subseteq Q$ with $|A|=\ell_2$, and put
\[
L_A(X)\;:=\;\prod_{b\in A}\bigl(X^a-b\bigr)\;=\;\sum_{j=0}^{\ell_2}(-1)^j e_j(A)\,X^{a(\ell_2-j)}\;=\;X^{k+2a}-e_1(A)\,X^{k+a}+R_A(X),
\]
where $a\ell_2=k+2a$, $a(\ell_2-1)=k+a$, and $R_A$ collects the terms with $j\ge2$, so $\deg R_A\le a(\ell_2-2)=k$. The polynomial $L_A$ is squarefree with root set $S_A:=\bigcup_{b\in A}S_b\subseteq D$ of size $a\ell_2=k+2a$.

Set $z_A:=-e_1(A)\in\B$ and $c_A:=\bigl(-R_A(x)\bigr)_{x\in D}$. Since $\deg R_A\le k$, we have $c_A\in\RS[\B,D,k+1]\subseteq\RS[\F,D,k+1]$. For every $x\in S_A$,
\[
0=L_A(x)=x^{k+2a}+z_A\,x^{k+a}+R_A(x),\qquad\text{i.e.}\qquad u_{z_A}(x)=c_A(x),
\]
so $u_{z_A}$ and $c_A$ agree on at least $k+2a$ points of $D$ and $\dist(u_{z_A},c_A)\le1-\rho-2/N$.

The map $A\mapsto c_A$ is injective on the $\ell_2$-subsets of $Q$ sharing a common slope value: if $z_A=z_{A'}$ and $c_A=c_{A'}$, then $R_A$ and $R_{A'}$ are polynomials of degree $\le k<n$ agreeing on all $n$ points of $D$, hence equal as polynomials; then $L_A=L_{A'}$, so $S_A=S_{A'}$, and $A=\{x^a:x\in S_A\}=A'$.

Finally, $z_A=-e_1(A)$ takes at most $|\B|$ values as $A$ ranges over the $\binom{N}{\ell_2}$ subsets; some $z\in\B$ is attained by at least $\binom{N}{\ell_2}/|\B|$ of them, and the corresponding codewords $c_A$ are pairwise distinct elements of $\Lst(\RS[\F,D,k+1],\,1-\rho-2/N,\,u_z)$.
\end{proof}

\begin{remark}[why slack two]\label{rem:slacktwo}
Part (i) applied to $C^+=\RS[\F,D,k+1]$ directly would require $a\mid k+1$; in every challenge instantiation $a$ is a power of two and $k+1$ is odd, so this fails. At slack two the listed codewords have degree $\le k$, i.e.\ lie in $C^+$, while the divisibility demand stays on $k$. This is exactly the interface needed by \cref{thm:A}, whose list conclusion concerns $C^+$ while its hypothesis concerns $C$.
\end{remark}

\begin{remark}[the subfield pigeonhole elsewhere]\label{rem:subpigeon}
The same one-line strengthening --- pigeonholing locator data over the field of definition $\B$ rather than over $\F$ --- also upgrades the coefficient-pigeonhole bound of \cite{Cho26b} to the generated field: the prefix map $\Phi_\sigma$ there has image in $\B^\sigma$ whenever $D\subseteq\B$, since the elementary symmetric prefixes of subsets of $D$ are $\B$-valued. This is the form consumed by the ``necessary'' half of the list-profile conjecture of \cite{Cho26b}, which states the entropy floor at the generated field $\tau^{*}(\rho,q_D)$; \cref{lem:fiber} supplies the missing subfield step.
\end{remark}

\begin{remark}[lists, not lines]\label{rem:lines}
The words $u_z$ of \cref{lem:fiber} lie on the $\B$-rational line $f+zg$ with $f=(x^{k+ta})_{x\in D}$, $g=(x^{k+(t-1)a})_{x\in D}$, but the lemma is purely a statement about lists at a received word and no line structure is used in \cref{sec:main}. This matters: by \cref{cor:Fvalued} below, over a proper extension $\F\supsetneq\B$ the lines that end up certifying the CA failure cannot be these (or any) $\B$-rational lines.
\end{remark}

\subsection{Quotient-profile floors with remainder supports}
\label{sec:quotient-remainder}

The fiber lemma above uses supports that are unions of complete quotient fibers.
For threshold work one also needs agreements that are not multiples of the fiber
size.  The following locator-prefix pigeonhole gives a finite lower ledger for
such remainder supports.  It is a lower-floor theorem only; the matching upper
ledger over all divisor scales remains open; it is the upper-ledger side of the residual band problem.

Let $c\mid n$ and put $Q_c:=D^c$, so the map
\[
        \pi_c:D\to Q_c,\qquad x\mapsto x^c
\]
has fibers of size $c$ and $|Q_c|=N=n/c$.  For $0\le s<c$ and an integer
$\sigma\ge0$, define the prefix weight
\[
        w_c(s,\sigma)
        :=\#\{h:1\le h\le\sigma,
              \ h\bmod c\in\{0,1,\ldots,s\}\},
\]
where residues are taken in $\{0,1,\ldots,c-1\}$.  Thus
$w_c(0,\sigma)=\lfloor \sigma/c\rfloor$ for full-fiber supports.

\begin{lemma}[quotient-remainder locator-prefix lower floor]
\label{lem:quotient-remainder-prefix}
Let $\B\subseteq\F$, let $D\subseteq\B^\times$ be a multiplicative coset of
order $n$, and let $K<n$.  Fix $c\mid n$, write $N=n/c$, and let
\[
        A_0=mc+s,
        \qquad 0\le s<c,
        \qquad 0\le m\le N,
        \qquad A_0\ge K .
\]
If $s>0$, assume $mc+s\le n$.  Put $\sigma=A_0-K$ and
\[
        M_{c,m,s}:=\binom Nm\binom{n-mc}{s}.
\]
Then there is a $\B$-valued received word $U:D\to\B$ such that
\[
        \bigl|\Lst(\RS[\F,D,K],\,1-A_0/n,\,U)\bigr|
        \ge
        \left\lceil \frac{M_{c,m,s}}{|\B|^{w_c(s,\sigma)}}\right\rceil .
\]
The listed codewords may be chosen in $\RS[\B,D,K]$.
\end{lemma}

\begin{proof}
For $M\subseteq Q_c$ with $|M|=m$, write
\[
        P_M(X)=\prod_{b\in M}(X^c-b).
\]
For $R\subseteq D\setminus \pi_c^{-1}(M)$ with $|R|=s$, put
\[
        E_R(X)=\prod_{r\in R}(X-r),
        \qquad
        L_{M,R}(X)=P_M(X)E_R(X).
\]
Then $L_{M,R}$ is a squarefree locator of degree $A_0$ whose roots are exactly
\(\pi_c^{-1}(M)\cup R\).  Write
\[
        L_{M,R}(X)=X^{A_0}+\sum_{h=1}^{A_0}\lambda_h(M,R)X^{A_0-h}.
\]
The coefficients lie in $\B$.  Moreover, because $P_M$ has monomials only in
degrees congruent to $0$ modulo $c$ and $E_R$ has degree $s$, the coefficient
$\lambda_h(M,R)$ is zero unless $h\bmod c\in\{0,1,\ldots,s\}$.
Consequently the high-degree prefix
\[
        \bigl(\lambda_h(M,R):1\le h\le\sigma,
              \ h\bmod c\in\{0,1,\ldots,s\}\bigr)
\]
takes at most $|\B|^{w_c(s,\sigma)}$ values.

Pigeonhole over these prefix values.  There is a fixed prefix vector $\xi$ and a
subfamily $\mathcal T_\xi$ of at least
$M_{c,m,s}/|\B|^{w_c(s,\sigma)}$ pairs $(M,R)$ with this prefix.  Let
\[
        U_\xi(X)=X^{A_0}+\sum_{1\le h\le\sigma}\lambda_h(\xi)X^{A_0-h}
\]
be the corresponding high-degree part, restricted to $D$.  For every
$(M,R)\in\mathcal T_\xi$, the difference
\[
        R_{M,R}(X):=L_{M,R}(X)-U_\xi(X)
\]
has degree $<K$.  Hence the codeword $-R_{M,R}$ agrees with $U_\xi$ on all
roots of $L_{M,R}$, a set of size $A_0$.

It remains only to check that the listed codewords are distinct.  If two pairs
$(M,R)$ and $(M',R')$ in the same prefix class gave the same codeword on $D$,
then the two remainders of degree $<K<n$ would be equal as polynomials.  Since
the high prefixes are also equal, $L_{M,R}=L_{M',R'}$.  Their root sets in $D$
are therefore equal.  A full fiber has size $c$ while $|R|=s<c$, so the full
fibers contained in this root set are exactly $M$, and then the residual set is
forced as well.  Thus $(M,R)=(M',R')$.
\end{proof}

\begin{definition}[prefix fibers and prefix-image certificates]
\label{def:prefix-fiber-certificate}
In the notation of \cref{lem:quotient-remainder-prefix}, let
$\Phi_{c,m,s}^{K}$ be the map from pairs $(M,R)$ to the high-degree prefix
\[
        \bigl(\lambda_h(M,R):1\le h\le A_0-K,
              \ h\bmod c\in\{0,1,\ldots,s\}\bigr)
        \in \B^{w_c(s,A_0-K)} .
\]
If $A_0=K$ this is the unique empty prefix.  For a prefix value $\xi$, put
\[
        \mu_{c,m,s}^{K}(\xi)
        :=\bigl| (\Phi_{c,m,s}^{K})^{-1}(\xi)\bigr|.
\]
Define the image size and the heaviest prefix fiber by
\[
        I_{c,m,s}^{K}:=\bigl|\operatorname{im}\Phi_{c,m,s}^{K}\bigr|,
        \qquad
        H_{c,m,s}^{K}:=\max_{\xi}\mu_{c,m,s}^{K}(\xi).
\]
Thus
\[
        H_{c,m,s}^{K}
        \ge
        \left\lceil\frac{M_{c,m,s}}{I_{c,m,s}^{K}}\right\rceil
        \ge
        \left\lceil\frac{M_{c,m,s}}{|\B|^{w_c(s,A_0-K)}}\right\rceil .
\]
A prefix-fiber certificate is an orbit decomposition, hash table, or other
verifiable bucket table for the finite map $\Phi_{c,m,s}^{K}$, including the
image size and the maximum fiber size.  If only $I_{c,m,s}^{K}$ is certified,
the middle lower bound is still valid; if no image certificate is supplied, the
ambient-box lower bound is valid.
\end{definition}

\begin{lemma}[heaviest-prefix locator floor]
\label{lem:heaviest-prefix-locator-floor}
Under the hypotheses of \cref{lem:quotient-remainder-prefix}, there is a
$\B$-valued received word $U:D\to\B$ such that
\[
        \bigl|\Lst(\RS[\F,D,K],\,1-A_0/n,\,U)\bigr|
        \ge H_{c,m,s}^{K}.
\]
Consequently
\[
        \bigl|\Lst(\RS[\F,D,K],\,1-A_0/n,\,U)\bigr|
        \ge
        \left\lceil \frac{M_{c,m,s}}{I_{c,m,s}^{K}}\right\rceil
        \ge
        \left\lceil \frac{M_{c,m,s}}{|\B|^{w_c(s,A_0-K)}}\right\rceil
\]
when only the coarser certificates are available.  The listed codewords may be
chosen in $\RS[\B,D,K]$.
\end{lemma}

\begin{proof}
Choose a prefix value $\xi$ with
$\mu_{c,m,s}^{K}(\xi)=H_{c,m,s}^{K}$ and form the high-degree word
\[
        U_\xi(X)=X^{A_0}+\sum_{1\le h\le A_0-K}
        \lambda_h(\xi)X^{A_0-h},
\]
where $\lambda_h(\xi)=0$ for the high coefficients outside the allowed
residue classes.  For every pair $(M,R)$ in this heaviest prefix fiber, the difference
$L_{M,R}(X)-U_\xi(X)$ has degree $<K$, so the negative of this remainder is a
codeword of $\RS[\B,D,K]$ agreeing with $U_\xi$ on the $A_0$ roots of
$L_{M,R}$.  If two pairs in the same prefix fiber gave the same codeword on
$D$, then their degree-$<K<n$ remainders would be equal as polynomials; together
with the common prefix this gives $L_{M,R}=L_{M',R'}$.  Their root sets in $D$
are equal.  A full quotient fiber has size $c$ while the residual set has size
$s<c$, so the full fibers contained in the root set determine $M$, and then the
residual set determines $R$.  Thus the codewords are pairwise distinct.  The two
coarser displayed bounds follow from \cref{def:prefix-fiber-certificate}.
\end{proof}

For $A\in\{k+1,
\ldots,n\}$ and $c\mid n$, write
\[
        A=m_c(A)c+s_c(A),
        \qquad 0\le s_c(A)<c,
\]
and set
\[
\begin{aligned}
        M_c(A)&:=
        \binom{n/c}{m_c(A)}
        \binom{n-cm_c(A)}{s_c(A)},\\
        L_{c,{\rm heavy}}^+(A)&:=
        H_{c,m_c(A),s_c(A)}^{k+1},\\
        L_{c,{\rm image}}^+(A)&:=
        \left\lceil\frac{M_c(A)}{I_{c,m_c(A),s_c(A)}^{k+1}}\right\rceil,\\
        L_{c,{\rm box}}^+(A)&:=
        \left\lceil
        \frac{M_c(A)}{|\B|^{w_c(s_c(A),A-k-1)}}
        \right\rceil .
\end{aligned}
\]
The three levels satisfy
\[
        L_{c,{\rm heavy}}^+(A)
        \ge L_{c,{\rm image}}^+(A)
        \ge L_{c,{\rm box}}^+(A).
\]
Here ``heavy'' uses a full prefix-fiber certificate, ``image'' uses only the
prefix-image size, and ``box'' is the certificate-free bound already implicit in
\cref{lem:quotient-remainder-prefix}.

\begin{theorem}[quotient-remainder deep-point lower ledger]
\label{thm:quotient-remainder-deep-floor}
Let $\B\subseteq\F$, let $q=|\F|>n$, let
$D\subseteq\B^\times$ be a multiplicative coset of order $n$, and let
$C=\RS[\F,D,k]$.  Fix an agreement value $A\in\{k+1,\ldots,n\}$ and a
divisor $c\mid n$.  For each
\[
        \star\in\{{\rm heavy},{\rm image},{\rm box}\}
\]
and every
\[
        \delta\in[1-A/n,\,1-k/n)
\]
one has
\[
        \eca(C,\delta)
        \ge \mathcal E_{q,k}\bigl(L_{c,\star}^+(A)\bigr).
\]
If $\delta>0$, then also
\[
        \emca(C,\delta)
        \ge \mathcal E_{q,k}\bigl(L_{c,\star}^+(A)\bigr).
\]
Here $\mathcal E_{q,k}$ is as in \cref{cor:deep-list-trigger-ceiling}.  More
precisely, some received line has at least
\[
        \mathcal N_{q,k}\bigl(L_{c,\star}^+(A)\bigr)
\]
CA-bad finite slopes at that radius, and for $\delta>0$ these slopes are
MCA-bad.

Consequently, for a closed agreement threshold $a\ge k+1$, define
\[
        B_{Q,\star}^-(a)
        :=\max_{\substack{c\mid n\\ A\in\{a,a+1,\ldots,n\}}}
          \mathcal N_{q,k}\bigl(L_{c,\star}^+(A)\bigr).
\]
Then $B_{Q,\star}^-(a)$ is a lower numerator for $q\,\eca(C,\delta)$ for
every $\delta\in[1-a/n,1-k/n)$, and for $q\,\emca(C,\delta)$ on the
same interval with $\delta>0$.  The strongest
certified ledger uses $\star={\rm heavy}$; the certificate-free ledger uses
$\star={\rm box}$.
\end{theorem}

\begin{proof}
Apply \cref{lem:heaviest-prefix-locator-floor} with $K=k+1$ and $A_0=A$.  This
produces a $C^+=\RS[\F,D,k+1]$ list of size at least
$L_{c,{\rm heavy}}^+(A)$ at radius $1-A/n$, with the image and ambient-box
variants giving the corresponding smaller certified list sizes.  Then apply
\cref{prop:quantitative-deep-list-floor}.  If $A\ge a$ and
$\delta\ge1-a/n$, then $\delta\ge1-A/n$, so the same converted floor applies.
The maximum over $(c,A)$ is the correct lower-ledger operation: the definitions
of $\eca$ and $\emca$ maximize over received lines, so lower witnesses from
different divisor scales or exact agreement values need not share one line and
must not be summed.
\end{proof}

\begin{corollary}[quotient-remainder trigger]
\label{cor:quotient-remainder-trigger}
In the setting of \cref{thm:quotient-remainder-deep-floor}, if
$L_{c,\star}^+(A)>(q-n)/k$ for any
$\star\in\{{\rm heavy},{\rm image},{\rm box}\}$, then, for every
$\delta\in[1-A/n,1-k/n)$,
\[
        \eca(C,\delta)>
        \frac1{2k}\left(1-\frac nq\right),
\]
and for every such $\delta>0$,
\[
        \emca(C,\delta)>
        \frac1{2k}\left(1-\frac nq\right).
\]
\end{corollary}

\begin{proof}
This is \cref{cor:deep-list-trigger-ceiling} applied inside
\cref{thm:quotient-remainder-deep-floor}.
\end{proof}

\begin{corollary}[quantitative first-grid floor]
\label{cor:quantitative-first-grid-floor}
Let $C=\RS[\F,D,k]$, let $q=|\F|>n$, and assume only that $D\subseteq\F$ has
$n$ distinct points.  Then for every
$\delta\in[1-(k+1)/n,1-k/n)$,
\[
        \eca(C,\delta)
        \ge
        \mathcal E_{q,k}\left(\binom n{k+1}\right),
\]
and for every such $\delta>0$,
\[
        \emca(C,\delta)
        \ge
        \mathcal E_{q,k}\left(\binom n{k+1}\right).
\]
In particular, if $\binom n{k+1}>(q-n)/k$, then the printed
$(2k)^{-1}(1-n/q)$ lower bound holds on this entire first grid interval without
any multiplicative smoothness assumption.
\end{corollary}

\begin{proof}
Use the ordinary locator identity for all $(k+1)$-subsets of $D$; this is the
$c=1$ argument behind \cref{lem:heaviest-prefix-locator-floor}, but it does not
use multiplicative smoothness.  The prefix is empty, so the list size is at
least
$\binom n{k+1}$.  Apply \cref{prop:quantitative-deep-list-floor} and then
\cref{cor:deep-list-trigger-ceiling}.
\end{proof}

\begin{definition}[quotient lower-ledger certificate]
\label{def:quotient-lower-ledger-certificate}
A quotient-remainder lower-ledger certificate for a row and a threshold $a$
contains, for each divisor scale $c$ and each exact agreement $A\ge a$ that is
used in the maximum:
\begin{enumerate}[label=\textup{(\roman*)},leftmargin=2.2em]
\item the decomposition $A=m_c(A)c+s_c(A)$;
\item the support count $M_c(A)$ and prefix weight $w_c(s_c(A),A-k-1)$;
\item either a full prefix-fiber bucket table certifying
$H_{c,m_c(A),s_c(A)}^{k+1}$, or an image-size certificate for
$I_{c,m_c(A),s_c(A)}^{k+1}$, or the declaration that the ambient-box denominator
is being used;
\item the resulting list floor $L_{c,\star}^+(A)$;
\item the converted slope numerator
$\mathcal N_{q,k}(L_{c,\star}^+(A))$ and error floor
$\mathcal E_{q,k}(L_{c,\star}^+(A))$.
\end{enumerate}
The verifier checks the locator-prefix formula, distinctness of codewords inside
a prefix fiber, the integer arithmetic, and the final maximum over divisor
scales and exact agreement values.
\end{definition}

\begin{remark}[quotient lower ledger versus quotient upper ledger]
\label{rem:T1-lower-ledger-status}
The quotient-profile lower side is now scanner-ready at every agreement value
$A=mc+s$: the generated-field contribution is the actual heaviest prefix fiber
$H_{c,m,s}^{k+1}$ when certified, the exact prefix-image denominator
$I_{c,m,s}^{k+1}$ when only image data are certified, and the ambient box
$|\B|^{w_c(s,A-k-1)}$ when no certificate is supplied.  The challenge-field size
$q$ enters only through the independent conversion functions
$\mathcal E_{q,k}$ and $\mathcal N_{q,k}$.  Across divisor scales the lower-side
operation is a maximum, not a union count, because the witness line may change.
The remaining upper-ledger problem is sharper and unchanged in logical status:
replace the safe support ledger and the exact declared-branch image ledger by an
exhaustive quotient-periodic classification, with duplicate slopes coalesced
across all divisor scales.
\end{remark}

\begin{corollary}[augmented slack-one quotient floor]
\label{cor:augmented-slack-one}
Let $\B\subseteq\F$ and let $D\subseteq\B^\times$ be a multiplicative coset
of order $n$.  Put $C=\RS[\F,D,k]$ and $C^+=\RS[\F,D,k+1]$.  If
$c\mid\gcd(n,k)$, then there is a $\B$-valued received word $U_c$ such that
\[
        \bigl|\Lst(C^+,\,1-(k+c)/n,\,U_c)\bigr|
        \ge
        \binom{n/c}{k/c+1} .
\]
One may take $U_c(x)=x^{k+c}$.  In particular, for $c=1$ this gives the
universal first-grid floor
\[
        \bigl|\Lst(C^+,\,1-(k+1)/n,\,x^{k+1})\bigr|
        \ge
        \binom n{k+1},
\]
with no quotient denominator.  The $c=1$ statement in fact holds for any set
$D\subseteq\B$ of $n$ distinct points, not only for multiplicative cosets.
\end{corollary}

\begin{proof}
For $c>1$, apply Lemma~\ref{lem:quotient-remainder-prefix} with $K=k+1$,
$s=0$, and $m=k/c+1$.  Then $A_0=k+c$ and $\sigma=A_0-K=c-1$, so
$w_c(0,c-1)=0$.  The high-degree prefix is only the monomial $X^{k+c}$, giving
the displayed word.  For $c=1$, the same argument is just the ordinary locator
identity for all $(k+1)$-subsets of $D$: the high-degree prefix of a monic
locator of degree $k+1$ is $X^{k+1}$, and its remainder has degree at most $k$.
\end{proof}

\begin{corollary}[first-grid deep-point cap]
\label{cor:first-grid-cap}
Let $C=\RS[\F,D,k]$, $q=|\F|$, and suppose $q>n$.  Let
$c\mid\gcd(n,k)$.  If $c>1$, assume $D\subseteq\B^\times$ is a
multiplicative coset as above; for $c=1$, no smoothness is needed.  If
\[
        \binom{n/c}{k/c+1}\ge q/k+1,
\]
then
\[
        \eca\bigl(C,1-(k+c)/n\bigr)
        >\frac1{2k}\left(1-\frac nq\right),
        \qquad
        \emca\bigl(C,1-(k+c)/n\bigr)
        >\frac1{2k}\left(1-\frac nq\right).
\]
For $c=1$ this is the first closed grid point below capacity.
\end{corollary}

\begin{proof}
Corollary~\ref{cor:augmented-slack-one} gives a $C^+$-list of size
$L\ge q/k+1$ at agreement $k+c$.  The proof of the deep-point conversion
Theorem~\ref{thm:A} applies whenever the agreement is $>k$, so it applies here
for every $c\ge1$.  If the displayed CA lower bound failed, then
\[
        \eca(C,1-(k+c)/n)\le \frac{q-n}{2kq}.
\]
Theorem~\ref{thm:A} with $\eta=1/2$ would give
\[
        \Lst(C^+,1-(k+c)/n)
        \le
        \left\lceil \frac{q-n}{k}\right\rceil
        < q/k+1,
\]
contradicting the list lower bound.  Every no-loss CA-bad slope is also
support-wise MCA-bad on the same explaining support, giving the MCA inequality.
\end{proof}

\begin{corollary}[first-grid cap for large challenge-envelope rows]
\label{cor:first-grid-grand}
Let $\rho\in\{\frac12,\frac14,\frac18,\frac1{16}\}$, let
$C=\RS[\F,D,k]$ with $|D|=n$, $k=\rho n$, $q=|\F|<2^{256}$, and
$q>n$.  Assume
\[
\begin{array}{c|cccc}
\rho & \frac12 & \frac14 & \frac18 & \frac1{16}\\ \hline
k\ge & 127 & 78 & 58 & 47
\end{array}
\]
for the corresponding rate.  Then the first closed grid point below capacity is
already CA- and MCA-unsafe:
\[
        \eca\bigl(C,1-\rho-1/n\bigr)
        >\frac1{2k}\left(1-\frac nq\right),
        \qquad
        \emca\bigl(C,1-\rho-1/n\bigr)
        >\frac1{2k}\left(1-\frac nq\right).
\]
If also $k\le2^{40}$, then this lower bound is $>2^{-86}$ and hence is far
above the prize target $2^{-128}$.
\end{corollary}

\begin{proof}
By Corollary~\ref{cor:first-grid-cap} with $c=1$, it is enough to check
\[
        \binom n{k+1}\ge q/k+1.
\]
Since $q<2^{256}$, the stronger sufficient condition is
$k(\binom n{k+1}-1)\ge2^{256}$.  The four printed lower bounds on $k$ are exactly
small finite checks of this inequality at the corresponding rate; the left side
is increasing thereafter.  For the error size, $q\ge n+1$, $n=k/\rho\le16k$,
and $k\le2^{40}$ give
\[
        \frac1{2k}\left(1-\frac nq\right)
        \ge \frac1{2k(n+1)}
        >2^{-86}.
\]
\end{proof}


\subsection{A conservative quotient-support upper ledger}
\label{sec:quotient-support-upper}

The remainder lower floor above supplies the unsafe side of the quotient-profile
ledger.  For a prize threshold one also needs a safe-side accounting rule: if a
later argument declares a bad parameter to be quotient-periodic, the declaration
must come with a divisor-lattice numerator that is counted with overlaps under
control.  The next statements give such a conservative upper ledger.  They do
not prove that every bad line parameter is quotient-periodic; they only say that
once a permitted support family has been specified, every ledger term witnessed
inside that family is bounded by the size of the family, with a factor $d$ for
finite-parameter degree-$d$ power curves.

Let $\mathcal S$ be a family of subsets of $D$, all of size at least $a$.  For a
received line $f+zg$, say that a finite slope is
\emph{$\mathcal S$-witnessed support-wise bad} if it has a support-wise
noncontained explanation on some $S\in\mathcal S$.  Define the analogous
$\mathcal S$-witnessed no-loss CA slopes by requiring that $f+zg$ be explained
on some $S\in\mathcal S$ while $(f,g)$ is not globally explained on any support
of size at least $a$.

\begin{lemma}[one support pays for at most one line parameter]
\label{lem:one-support-one-line}
Let $C=\RS[\F,D,k]$ and let $a\ge k$.  Fix a support
$S\subseteq D$ with $|S|\ge a$.  For any received line $f+zg$, at most one finite
slope $z\in\F$ can be support-wise noncontained on the fixed support $S$.
Likewise, at most one finite slope can be no-loss CA-bad while being explained on
this fixed support $S$.  Consequently, for every support family $\mathcal S$,
\[
        \#\{\mathcal S\text{-witnessed support-wise MCA-bad finite slopes}\}
        \le |\mathcal S|,
\]
and the same bound holds for no-loss CA slopes witnessed by $\mathcal S$.
\end{lemma}

\begin{proof}
Suppose two distinct finite slopes $z_1\ne z_2$ are both explained on the same
support $S$, by codewords $p_1,p_2\in C$:
\[
        f+z_i g=p_i\qquad\text{on }S,\quad i=1,2.
\]
Then on $S$,
\[
        g=\frac{p_1-p_2}{z_1-z_2},\qquad
        f=p_1-z_1g.
\]
Both right-hand sides are codewords of $C$.  Thus $(f,g)$ is simultaneously
explained on $S$.  This contradicts support-wise noncontainment on $S$.  It also
contradicts no-loss CA badness at agreement $a$, because $|S|\ge a$ gives a
global simultaneous explanation support of the forbidden size.  Hence each fixed
support contributes at most one bad finite line parameter, and summing over
supports gives the displayed bound.
\end{proof}

For a finite-parameter power curve
\[
        W_\gamma=f_0+\gamma f_1+\cdots+\gamma^d f_d,
        \qquad \gamma\in\F,
\]
say that a parameter is support-wise curve-bad on $S$ if $W_\gamma$ is explained
by a codeword on $S$, but the coefficient tuple $(f_0,\ldots,f_d)$ is not
simultaneously explained by $C^{\equiv(d+1)}$ on $S$.  The no-loss curve-CA
variant replaces this support-wise noncontainment by the corresponding global
far condition at agreement $a$.

\begin{lemma}[one support pays for at most $d$ curve parameters]
\label{lem:one-support-d-curve}
Let $C=\RS[\F,D,k]$, let $a\ge k$, and fix a support $S\subseteq D$ with
$|S|\ge a$.  For a degree-$d$ power curve $W_\gamma$, at most $d$ finite
parameters $\gamma\in\F$ can be support-wise curve-bad on $S$.  The same bound
holds for no-loss curve-CA parameters explained on $S$.  Consequently any
support family $\mathcal S$ witnesses at most $d|\mathcal S|$ bad finite curve
parameters.
\end{lemma}

\begin{proof}
If $d+1$ distinct parameters $\gamma_0,\ldots,\gamma_d$ were explained on the
same support $S$, with codewords $P_j\in C$, then Lagrange interpolation in the
parameter variable gives codewords $Q_0,\ldots,Q_d\in C$ such that
\[
        \sum_{i=0}^d \gamma_j^i Q_i=P_j
        \qquad (0\le j\le d).
\]
On every point of $S$, the same interpolation applied to the values of
$W_{\gamma_j}$ gives $f_i=Q_i$ for every coefficient $i$.  Hence
$(f_0,\ldots,f_d)$ is simultaneously explained by $C^{\equiv(d+1)}$ on $S$.
This contradicts support-wise curve-noncontainment on $S$, and also contradicts
no-loss curve-CA badness at agreement $a$.  Thus every support contributes at
most $d$ parameters.
\end{proof}

\begin{lemma}[one support pays for at most one interleaved tuple]
\label{lem:one-support-one-interleaved}
Let $C=\RS[\F,D,k]$, let $a\ge k$, and let $\mu\ge1$.  Fix a $\mu$-row received
word $U=(U_1,\ldots,U_\mu)$.  For a fixed support $S\subseteq D$ with
$|S|\ge a$, there is at most one tuple $(P_1,\ldots,P_\mu)\in C^\mu$ agreeing
with $U$ on every row over $S$.  Hence the number of interleaved list tuples whose
common agreement support belongs to a support family $\mathcal S$ is at most
$|\mathcal S|$.
\end{lemma}

\begin{proof}
For each row $i$, two degree-$<k$ polynomials agreeing with $U_i$ on the fixed
support $S$ agree with each other on at least $k$ points, so they are equal.  Thus
each row has at most one codeword compatible with $S$, and therefore the whole
interleaved tuple is unique if it exists.
\end{proof}

We now instantiate the support-family ledger on quotient-remainder supports.  For
$c\mid n$ and an integer $A$ with $0\le A\le n$, write
\[
        A=m_c(A)c+s_c(A),\qquad 0\le s_c(A)<c,
\]
and define
\[
\begin{aligned}
        \mathcal S_c(A)
        :=\{&\pi_c^{-1}(M)\cup R:
          M\subseteq Q_c,
          |M|=m_c(A),\\
        & R\subseteq D\setminus\pi_c^{-1}(M),
          |R|=s_c(A)\} .
\end{aligned}
\]
For a divisor set $\mathcal C\subseteq\{c:c\mid n\}$ put
\[
        \mathcal S_{\mathcal C}^{\ge a}
        :=\bigcup_{c\in\mathcal C}\bigcup_{A=a}^n \mathcal S_c(A)
\]
and
\[
        U_{\mathcal C}^{\rm supp}(a)
        :=\left|\mathcal S_{\mathcal C}^{\ge a}\right|.
\]
When the exact union is inconvenient, use the safe sum
\[
\label{eq:quotient-support-safe-sum}
        U_{\mathcal C}^{\rm sum}(a)
        :=\sum_{c\in\mathcal C}\sum_{A=a}^n
          \binom{n/c}{\lfloor A/c\rfloor}
          \binom{n-c\lfloor A/c\rfloor}{A-c\lfloor A/c\rfloor},
\]
so that $U_{\mathcal C}^{\rm supp}(a)\le U_{\mathcal C}^{\rm sum}(a)$.

\begin{proposition}[divisor-union quotient support ledger]
\label{prop:divisor-union-support-ledger}
Let $C=\RS[\F,D,k]$ and let $a\ge k$.  Fix a divisor set $\mathcal C$ and let
$\mathcal S_{\mathcal C}^{\ge a}$ be the quotient-remainder support family above.
Then, for every received line, the number of finite support-wise MCA-bad slopes
whose witness support lies in $\mathcal S_{\mathcal C}^{\ge a}$ is at most
\[
        U_{\mathcal C}^{\rm supp}(a)
        \le U_{\mathcal C}^{\rm sum}(a).
\]
The same bound holds for no-loss CA slopes witnessed by this family.  For a
finite-parameter degree-$d$ power curve, the corresponding support-wise curve-MCA
and no-loss curve-CA numerator is at most
\[
        d\,U_{\mathcal C}^{\rm supp}(a)
        \le d\,U_{\mathcal C}^{\rm sum}(a).
\]
For every interleaving arity $\mu\ge1$, the number of interleaved list tuples whose
common agreement support belongs to $\mathcal S_{\mathcal C}^{\ge a}$ is at most
$U_{\mathcal C}^{\rm supp}(a)$.
\end{proposition}

\begin{proof}
Apply Lemmas~\ref{lem:one-support-one-line}, \ref{lem:one-support-d-curve}, and
\ref{lem:one-support-one-interleaved} to the support family
$\mathcal S_{\mathcal C}^{\ge a}$.  The inequality
$U_{\mathcal C}^{\rm supp}(a)\le U_{\mathcal C}^{\rm sum}(a)$ is only the union
bound over the divisor and agreement-size indices.
\end{proof}

\begin{remark}[what this does and does not finish]
\label{rem:quotient-upper-status}
Proposition~\ref{prop:divisor-union-support-ledger} supplies a certified
safe-side quotient ledger at the support-family level, and it is already enough
for a staircase scanner to avoid double-counting divisor scales: one may compute
the exact union $U_{\mathcal C}^{\rm supp}$ or use the printed safe sum.  This is
not yet the sharp quotient upper theorem needed for a full prize solution.  Near
capacity the support count can be exponentially larger than the number of
distinct bad slopes, so the next refinement is a slope-collision quotient ledger:
classify when many quotient-remainder supports map to the same line or curve
parameter, and replace the support count by the corresponding distinct-parameter
union count.
\end{remark}


\begin{definition}[block support-profile polynomial]
\label{def:block-support-profile-polynomial}
Let $D$ be a multiplicative coset of order $n$, let
$\mathcal C$ be a nonempty finite set of divisors of $n$, and put
\[
        h:=\operatorname{lcm}\mathcal C,
        \qquad N_h:=n/h .
\]
For each $c\in\mathcal C$ let $H_c\le H_h$ denote the subgroup of the cyclic
kernel of order $h$ with $|H_c|=c$.  If $T\subseteq H_h$, define
\[
        \phi_c(T):=\#\{\text{cosets }uH_c\subseteq H_h:\ uH_c\subseteq T\},
\]
the number of complete $c$-subfibers contained in the block subset $T$.
The block support-profile polynomial of $\mathcal C$ is
\[
        G_{\mathcal C,h}(Y,(Z_c)_{c\in\mathcal C})
        :=\sum_{T\subseteq H_h}
          Y^{|T|}\prod_{c\in\mathcal C}Z_c^{\phi_c(T)} .
\]
For a vector $\nu=(\nu_c)_{c\in\mathcal C}$ write
$Z^\nu:=\prod_{c\in\mathcal C}Z_c^{\nu_c}$, and define
$N_{\mathcal C,h}(A;\nu)$ by
\[
        G_{\mathcal C,h}(Y,Z)^{N_h}
        =\sum_{A=0}^{n}\sum_{\nu}
          N_{\mathcal C,h}(A;\nu)Y^A Z^\nu .
\]
\end{definition}

\begin{lemma}[block polynomial counts exact quotient profiles]
\label{lem:block-polynomial-counts-profiles}
For a support $S\subseteq D$ and $c\mid n$, let
\[
        F_c(S):=\#\{\text{$c$-fibers of }D\to D^c\text{ contained in }S\}.
\]
Then $N_{\mathcal C,h}(A;\nu)$ is exactly the number of supports
$S\subseteq D$ such that
\[
        |S|=A,
        \qquad F_c(S)=\nu_c\quad(c\in\mathcal C).
\]
\end{lemma}

\begin{proof}
The $h$-fibers partition $D$ into $N_h$ blocks, and every $c$-fiber with
$c\in\mathcal C$ is contained in a unique $h$-fiber because $c\mid h$.  After
choosing an identification of each $h$-fiber with the cyclic kernel $H_h$, the
number of selected points and the number of complete internal $c$-subfibers are
recorded by the monomial
\[
        Y^{|T|}\prod_{c\in\mathcal C}Z_c^{\phi_c(T)}
\]
for the block subset $T$.  The choices in distinct $h$-fibers are independent,
and the global statistics $|S|$ and $F_c(S)$ are the sums of the block statistics.
Multiplying the block polynomial over the $N_h$ blocks and extracting the
coefficient gives the asserted count.
\end{proof}

\begin{lemma}[membership by full-fiber profile]
\label{lem:quotient-support-membership-profile}
Fix $A$ and $c\mid n$, and write
\[
        A=m_c(A)c+s_c(A),
        \qquad 0\le s_c(A)<c .
\]
For a support $S\subseteq D$ of size $A$, one has
\[
        S\in\mathcal S_c(A)
        \qquad\Longleftrightarrow\qquad
        F_c(S)=m_c(A)=\lfloor A/c\rfloor .
\]
\end{lemma}

\begin{proof}
If $S\in\mathcal S_c(A)$, then by definition it is the union of
$m_c(A)$ complete $c$-fibers and $s_c(A)<c$ residual points outside those fibers;
there is no room for any further complete $c$-fiber.  Hence
$F_c(S)=m_c(A)$.  Conversely, if $F_c(S)=m_c(A)$, take the complete $c$-fibers
contained in $S$ as the set $M$ in the definition of $\mathcal S_c(A)$; the
remaining $A-m_c(A)c=s_c(A)$ selected points form the residual set.
\end{proof}

\begin{theorem}[exact divisor-lattice support-union formula]
\label{thm:exact-divisor-block-support-ledger}
Let $\mathcal C$ be a nonempty finite set of divisors of $n$, and define
$G_{\mathcal C,h}$ and $N_{\mathcal C,h}(A;\nu)$ as in
\cref{def:block-support-profile-polynomial}.  Then the exact support-family
union from \cref{prop:divisor-union-support-ledger} is
\[
\boxed{
        U_{\mathcal C}^{\rm supp}(a)
        =\sum_{A=a}^{n}\sum_{\nu}
          N_{\mathcal C,h}(A;\nu)\,
          \mathbf 1\!
          \left[\exists c\in\mathcal C:
                 \nu_c=\left\lfloor A/c\right\rfloor\right].
      }
\]
Consequently the quotient-support branch numerators in
\cref{prop:divisor-union-support-ledger} may be computed by this coefficient
formula instead of by the safe sum
$U_{\mathcal C}^{\rm sum}(a)$.
\end{theorem}

\begin{proof}
For a fixed exact size $A$, the union
$\bigcup_{c\in\mathcal C}\mathcal S_c(A)$ consists precisely of those supports
$S$ of size $A$ for which $S\in\mathcal S_c(A)$ for at least one divisor
$c\in\mathcal C$.  By \cref{lem:quotient-support-membership-profile}, this is
exactly the condition
\[
        F_c(S)=\lfloor A/c\rfloor
\]
for at least one $c\in\mathcal C$.  By
\cref{lem:block-polynomial-counts-profiles}, the number of supports with a fixed
profile vector $\nu=(F_c(S))_{c\in\mathcal C}$ is
$N_{\mathcal C,h}(A;\nu)$.  Summing over all exact sizes $A\ge a$ and over the
profile vectors satisfying the displayed predicate gives the formula.
\end{proof}

\begin{corollary}[single-divisor sanity check]
\label{cor:single-divisor-support-ledger}
For $\mathcal C=\{c\}$ the block polynomial is
\[
        G_{\{c\},c}(Y,Z)=(1+Y)^c+(Z-1)Y^c,
\]
and \cref{thm:exact-divisor-block-support-ledger} gives
\[
        \left|\mathcal S_c(A)\right|
        =\binom{n/c}{\lfloor A/c\rfloor}
          \binom{n-c\lfloor A/c\rfloor}{A-c\lfloor A/c\rfloor},
\]
the summand used in the safe ledger.
\end{corollary}

\begin{proof}
Inside one $c$-fiber, every proper subset contributes only $Y^{|T|}$, while the
full subset contributes $Y^cZ$.  Thus the block polynomial is
$(1+Y)^c-Y^c+Y^cZ$.  Raising to $n/c$ and selecting exactly
$m=\lfloor A/c\rfloor$ full blocks and $s=A-mc$ residual points in the remaining
blocks gives the displayed count.
\end{proof}

\begin{definition}[divisor-block support-ledger certificate]
\label{def:divisor-block-support-ledger-certificate}
A divisor-block support-ledger certificate for a divisor set $\mathcal C$ and a
threshold $a$ consists of:
\begin{enumerate}[label=\textup{(\roman*)},leftmargin=2.2em]
\item the lcm block size $h=\operatorname{lcm}\mathcal C$ and the block count
$N_h=n/h$;
\item the internal coset table of every subgroup $H_c\le H_h$ for
$c\in\mathcal C$;
\item the block polynomial $G_{\mathcal C,h}$, either as a full coefficient table
or as an orbit-compressed table under translations of $H_h$;
\item the coefficient extraction from $G_{\mathcal C,h}^{N_h}$ for every exact
size $A\ge a$;
\item the predicate table
$\mathbf 1[\exists c\in\mathcal C:\nu_c=\lfloor A/c\rfloor]$ and the final
integer $U_{\mathcal C}^{\rm supp}(a)$.
\end{enumerate}
The verifier checks the block enumeration, the subgroup containment relation
$c\mid h$, coefficient arithmetic, and the final predicate sum.  This is an exact
support-union certificate; it does not by itself coalesce distinct line or curve
parameters that arise from different supports.
\end{definition}

\begin{remark}[support ledger versus slope ledger]
\label{rem:support-vs-slope-ledger}
\Cref{thm:exact-divisor-block-support-ledger} closes the divisor-lattice overlap
problem at the support-family level: intersections between quotient descriptions
at incomparable divisor scales are now counted by one coefficient extraction in
the common $h$-block.  This is still weaker than the desired distinct-slope
quotient ledger.  To finish that stronger ledger one must apply the exact
support-image maps of
\cref{prop:distinct-parameter-line-ledger,prop:distinct-parameter-curve-ledger}
to the support set counted here and then prove or certify the resulting parameter
image size after duplicate coalescing.
\end{remark}

\subsection{Distinct-parameter quotient image ledger}
\label{sec:quotient-distinct-parameter}

The support-union ledger of \cref{sec:quotient-support-upper} is deliberately
safe but often very loose: many quotient-remainder supports can certify the same
line slope or the same curve parameter.  The next refinement is to count the
image of the actual parameter map.  This subsection gives an exact image ledger
for the quotient branch.  It is still a branch ledger, not an aperiodic upper
bound: it counts only witnesses whose agreement support lies in a declared
quotient-remainder support family.

We first record the support-to-parameter map in syndrome form.  Let
$C=\RS[\F,D,k]$, $|D|=n$, and put $r=n-k$.  For each $x\in D$ set
\[
        \lambda_x:=\left(\prod_{y\in D,\ y\ne x}(x-y)\right)^{-1}.
\]
For a word $Y:D\to\F$, define its syndrome vector
\[
        \Syn(Y)_m:=\sum_{x\in D}\lambda_x x^mY(x),
        \qquad 0\le m<r .
\]
Thus $\Syn(Y)=0$ if and only if $Y\in C$.  If
$T\subseteq D$ has size $j\le r$, write
\[
        L_T(X)=\prod_{x\in T}(X-x)=\ell_0+\ell_1X+\cdots+\ell_jX^j,
        \qquad
        \boldsymbol\ell_T=(\ell_0,\ldots,\ell_j)^t,
\]
and put $t=r-j$.  For $w\in\F^r$ let $H_{t,j}(w)$ be the $t\times(j+1)$ Hankel
window
\[
        H_{t,j}(w)_{a,b}=w_{a+b},
        \qquad 0\le a<t,\quad 0\le b\le j .
\]
When $S=D\setminus T$, the integer $t$ is exactly $|S|-k$.


\begin{lemma}[support-locator syndrome recurrence]
\label{lem:support-locator-syndrome-recurrence}
Let $Y:D\to\F$ and let $T\subseteq D$ with $|T|=j\le r$.  Put
$S=D\setminus T$ and $t=r-j$.  Then $Y$ is explained by $C$ on $S$ if and only if
\[
        H_{t,j}(\Syn(Y))\boldsymbol\ell_T=0 .
\]
\end{lemma}

\begin{proof}
Let $W_T\subseteq\F^r$ be the span of the parity-check columns indexed by $T$,
\[
        W_T=\operatorname{span}_\F\{\lambda_x(1,x,\ldots,x^{r-1}):x\in T\}.
\]
The word $Y$ is explained by $C$ on $S$ if and only if $Y-c$ is supported on $T$
for some $c\in C$, equivalently $\Syn(Y)\in W_T$.  If
$w\in W_T$, write $w_m=\sum_{x\in T}a_xx^m$.  Then, for $0\le h<t$,
\[
        \bigl(H_{t,j}(w)\boldsymbol\ell_T\bigr)_h
        =\sum_{b=0}^j\ell_b w_{h+b}
        =\sum_{x\in T}a_xx^hL_T(x)=0.
\]
Conversely, because $L_T$ is monic, the displayed recurrence determines all
coordinates $w_j,\ldots,w_{r-1}$ from the first $j$ coordinates, so its solution
space has dimension at most $j$.  The $j$ parity-check columns indexed by the
distinct points of $T$ have Vandermonde rank $j$, hence $W_T$ has dimension $j$
and is exactly the recurrence space.
\end{proof}

\begin{lemma}[exact affine and projective support-image map]
\label{lem:exact-line-image-map}
Let $f,g:D\to\F$, put $u=\Syn(f)$ and $v=\Syn(g)$, and fix
$T\subseteq D$ with $|T|=j\le r$ and $S=D\setminus T$.  Define
\[
        A_T:=H_{t,j}(u)\boldsymbol\ell_T,
        \qquad
        B_T:=H_{t,j}(v)\boldsymbol\ell_T .
\]
A finite slope $z\in\F$ is explained on the fixed support $S$ and is
support-wise noncontained there if and only if
\[
        A_T+zB_T=0,
        \qquad B_T\ne0 .
\]
Consequently the fixed support contributes either no finite bad slope or the
unique value
\[
        \theta_T=-A_{T,m}/B_{T,m}
\]
for any coordinate $m$ with $B_{T,m}\ne0$, provided all nonzero coordinates give
the same value.

For the projective line $\alpha f+\beta g$, the fixed support contributes a
support-wise noncontained point $[\alpha:\beta]\in\mathbb P^1(\F)$ if and only if
$A_T$ and $B_T$ are linearly dependent and not both zero.  In that case the
contributed projective parameter is the unique point satisfying
\[
        \alpha A_T+\beta B_T=0 .
\]
\end{lemma}

\begin{proof}
A word $Y$ is explained by $C$ on $S=D\setminus T$ if and only if
$Y-c$ is supported on $T$ for some $c\in C$, which is equivalent to
$\Syn(Y)$ lying in the span of the parity-check columns indexed by $T$.  The
standard locator recurrence says that this is equivalent to
\[
        H_{t,j}(\Syn(Y))\boldsymbol\ell_T=0 .
\]
Applying this to $Y=f+zg$ gives $A_T+zB_T=0$.  Under this incidence, simultaneous
explanation of $f$ and $g$ on $S$ is equivalent to $B_T=0$, because
$A_T=(A_T+zB_T)-zB_T$.  This proves the finite-slope statement.  The projective
statement is the same calculation with $Y=\alpha f+\beta g$; if
$A_T=B_T=0$ then both coefficient words are already explained on $S$, so the
point is contained rather than bad.
\end{proof}

Let $\mathcal S_{\mathcal C}^{\ge a}$ be the quotient-remainder support family
from \cref{sec:quotient-support-upper}, and let
\[
        \mathcal T_{\mathcal C}^{\ge a}
        :=\{D\setminus S:S\in\mathcal S_{\mathcal C}^{\ge a}\}
\]
be the corresponding co-support family.  For a received line define
\[
        \Theta_{\mathcal C}^{\rm aff}(f,g;a)
        :=\{\theta_T:T\in\mathcal T_{\mathcal C}^{\ge a},
              \theta_T\text{ is defined by }\cref{lem:exact-line-image-map}\}
\]
and define $\Theta_{\mathcal C}^{\rm proj}(f,g;a)$ analogously using the
projective point of \cref{lem:exact-line-image-map}.

\begin{proposition}[distinct-parameter quotient ledger for lines]
\label{prop:distinct-parameter-line-ledger}
For every received line $f+zg$, the finite affine slopes that are support-wise
MCA-bad and have a witness support in
$\mathcal S_{\mathcal C}^{\ge a}$ are exactly the elements of
$\Theta_{\mathcal C}^{\rm aff}(f,g;a)$.  Hence the exact quotient-branch line
numerator is
\[
        |\Theta_{\mathcal C}^{\rm aff}(f,g;a)|
        \le U_{\mathcal C}^{\rm supp}(a).
\]
The corresponding projective-slope numerator is
\[
        |\Theta_{\mathcal C}^{\rm proj}(f,g;a)|
        \le U_{\mathcal C}^{\rm supp}(a).
\]
The same image ledgers upper-bound no-loss CA parameters whose explaining
support lies in the same quotient family.
\end{proposition}

\begin{proof}
This is only \cref{lem:exact-line-image-map} applied to every co-support in the
family, followed by taking a union of the resulting parameter values.  The CA
claim follows because no-loss CA-badness is stronger than having an explained
slope on the support and forbids simultaneous explanation on any support of size
at least $a$.
\end{proof}

The curve analogue is just as explicit.  For a degree-$d$ power curve
\[
        W_\gamma=f_0+\gamma f_1+\cdots+\gamma^d f_d,
        \qquad \gamma\in\F,
\]
put $u_i=\Syn(f_i)$ and, for $T\subseteq D$, define
\[
        V_{i,T}:=H_{t,j}(u_i)\boldsymbol\ell_T\in\F^t .
\]
Let
\[
        \Gamma_T^{(d)}
        :=\{\gamma\in\F:
             V_{0,T}+\gamma V_{1,T}+\cdots+\gamma^dV_{d,T}=0,
             \ (V_{0,T},\ldots,V_{d,T})\ne0\} .
\]
The nonzero tuple condition is exactly support-wise noncontainment of the
coefficient tuple on the fixed support.

\begin{lemma}[exact support-image map for finite-parameter power curves]
\label{lem:exact-curve-image-map}
For a fixed support $S=D\setminus T$, the set of support-wise noncontained
curve-bad finite parameters is exactly $\Gamma_T^{(d)}$.  In particular
$|\Gamma_T^{(d)}|\le d$.
\end{lemma}

\begin{proof}
The same syndrome-locator recurrence applied to $W_\gamma$ gives the displayed
vector equation.  If every $V_{i,T}$ is zero, then every coefficient word
$f_i$ is already explained on $S$, so the parameter is contained.  Otherwise at
least one scalar coordinate of
$V_{0,T}+\gamma V_{1,T}+\cdots+\gamma^dV_{d,T}$ is a nonzero polynomial in
$\gamma$ of degree at most $d$, so it has at most $d$ roots.
\end{proof}

Define the exact quotient curve image
\[
        \Gamma_{\mathcal C}^{(d)}(f_0,\ldots,f_d;a)
        :=\bigcup_{T\in\mathcal T_{\mathcal C}^{\ge a}}\Gamma_T^{(d)} .
\]

\begin{proposition}[distinct-parameter quotient ledger for curves]
\label{prop:distinct-parameter-curve-ledger}
For a finite-parameter degree-$d$ power curve, the quotient-branch support-wise
curve-MCA numerator is exactly
\[
        |\Gamma_{\mathcal C}^{(d)}(f_0,
        \ldots,f_d;a)| .
\]
It satisfies the safe bounds
\[
        |\Gamma_{\mathcal C}^{(d)}(f_0,\ldots,f_d;a)|
        \le d\,U_{\mathcal C}^{\rm supp}(a)
        \le d\,U_{\mathcal C}^{\rm sum}(a).
\]
The same image ledger upper-bounds no-loss curve-CA parameters whose explaining
support lies in the quotient family.
\end{proposition}

\begin{proof}
Apply \cref{lem:exact-curve-image-map} to each support and take the union of the
resulting parameter sets.  The displayed inequalities follow from
$|\Gamma_T^{(d)}|\le d$ and from the support-union bound of
\cref{prop:divisor-union-support-ledger}.
\end{proof}

\begin{definition}[fixed-support parameter gcd]
\label{def:fixed-support-parameter-gcd}
Let $C=\RS[\F,D,k]$, let $q=|\F|$, and let
\[
        W_\Gamma=f_0+\Gamma f_1+\cdots+\Gamma^d f_d,
        \qquad \Gamma\in\F,
\]
be a finite-parameter degree-$d$ power curve of received words.  Fix a
co-support $T\subseteq D$, put $S=D\setminus T$, and assume
$|S|>k$.  With the notation of \cref{lem:exact-curve-image-map}, write
\[
        P_{T,m}(\Gamma)
        :=\sum_{i=0}^d \Gamma^i (V_{i,T})_m,
        \qquad 0\le m<t,
        \qquad t=|S|-k .
\]
Define the fixed-support contribution polynomial
$g_T^{(d)}(\Gamma)\in\F[\Gamma]$ as follows.  If all coordinate polynomials
$P_{T,m}$ are identically zero, set $g_T^{(d)}=1$.  Otherwise let
$g_T^{(d)}$ be the monic greatest common divisor of the nonzero polynomials
among $P_{T,0},\ldots,P_{T,t-1}$.
\end{definition}

\begin{lemma}[fixed-support gcd equals the noncontained parameter set]
\label{lem:fixed-support-gcd-exact}
In the setting of \cref{def:fixed-support-parameter-gcd}, the finite parameters
$\gamma\in\F$ for which $W_\gamma$ is explained on $S=D\setminus T$ and the
coefficient tuple $(f_0,\ldots,f_d)$ is not simultaneously explained on $S$ are
exactly the roots in $\F$ of $g_T^{(d)}$.  In particular
\[
        \deg g_T^{(d)}\le d,
\]
with equality possible only before taking common factors across different
supports.
\end{lemma}

\begin{proof}
By \cref{lem:exact-curve-image-map}, the fixed support explains $W_\gamma$ exactly
when
\[
        P_{T,m}(\gamma)=0\qquad(0\le m<t).
\]
If every coordinate polynomial is identically zero, then every coefficient
vector $V_{i,T}$ is zero, so every coefficient word $f_i$ is already explained on
$S$; this is the contained branch and contributes no noncontained parameter.
The convention $g_T^{(d)}=1$ gives the empty root set, as required.

Otherwise the common zeros of the coordinate polynomials are exactly the roots
of their monic gcd after zero coordinates are discarded.  Since at least one
coordinate polynomial is nonzero of degree at most $d$, the gcd also has degree
at most $d$.  The nonzero tuple condition in \cref{lem:exact-curve-image-map} is
precisely the assertion that the coordinate polynomials are not all zero.
\end{proof}

\begin{definition}[quotient image lcm polynomial]
\label{def:quotient-image-lcm-polynomial}
Let $\mathcal T$ be any finite family of co-supports $T\subseteq D$ with
$|D\setminus T|>k$.  For a degree-$d$ finite-parameter power curve define
\[
        R_{\mathcal T}^{(d)}(\Gamma)
        :=\operatorname{rad}\operatorname{lcm}_{T\in\mathcal T} g_T^{(d)}(\Gamma),
\]
where factors equal to $1$ are ignored and the empty lcm is $1$.  For the quotient branch at agreement
threshold $a>k$ and divisor set $\mathcal C$, we use
\[
        \mathcal T=\mathcal T_{\mathcal C}^{\ge a}
        :=\{D\setminus S:S\in\mathcal S_{\mathcal C}^{\ge a}\}.
\]
The finite-field root-count polynomial is
\[
        R_{\mathcal T,F}^{(d)}(\Gamma)
        :=\gcd\bigl(R_{\mathcal T}^{(d)}(\Gamma),\Gamma^q-\Gamma\bigr).
\]
Because $R_{\mathcal T}^{(d)}$ is squarefree, the degree of
$R_{\mathcal T,F}^{(d)}$ is the number of $\F$-rational roots of
$R_{\mathcal T}^{(d)}$.
\end{definition}

\begin{theorem}[exact finite-parameter quotient image ledger]
\label{thm:exact-quotient-image-lcm-ledger}
Let $a>k$, let $\mathcal C$ be a finite divisor set, and let
$\mathcal T_{\mathcal C}^{\ge a}$ be the quotient co-support family.  For every
finite-parameter degree-$d$ power curve, the exact quotient-branch support-wise
curve-MCA numerator is
\[
        \boxed{
        \left|\Gamma_{\mathcal C}^{(d)}(f_0,\ldots,f_d;a)\right|
        =\deg R_{\mathcal T_{\mathcal C}^{\ge a},F}^{(d)} .
        }
\]
Consequently
\[
        \left|\Gamma_{\mathcal C}^{(d)}(f_0,\ldots,f_d;a)\right|
        \le \deg R_{\mathcal T_{\mathcal C}^{\ge a}}^{(d)}
        \le d\,U_{\mathcal C}^{\rm supp}(a),
\]
and the final support count may be evaluated by the exact divisor-block formula
of \cref{thm:exact-divisor-block-support-ledger}.  The same lcm ledger
upper-bounds no-loss curve-CA parameters whose explaining support lies in the
same quotient family.
\end{theorem}

\begin{proof}
By \cref{lem:fixed-support-gcd-exact}, for each fixed co-support $T$ the
noncontained parameters witnessed on $D\setminus T$ are exactly the $\F$-roots of
$g_T^{(d)}$.  Taking the union over
$T\in\mathcal T_{\mathcal C}^{\ge a}$ is therefore the same as taking the
$\F$-rational root set of the squarefree lcm
$R_{\mathcal T_{\mathcal C}^{\ge a}}^{(d)}$.  Over a finite field, the
$\F$-rational roots of a squarefree polynomial are counted by the degree of its
gcd with $\Gamma^q-\Gamma$, proving the displayed equality.

The degree bound follows from $\deg g_T^{(d)}\le d$ and from the fact that a
squarefree lcm has degree at most the sum of the degrees of the factors before
coalescing.  The no-loss curve-CA claim follows because a no-loss CA parameter
with an explaining support in the quotient family is, on that support,
noncontained; otherwise the coefficient tuple would have a simultaneous
explanation on a support of size at least $a$.
\end{proof}

\begin{corollary}[affine lines as the degree-one lcm ledger]
\label{cor:affine-line-lcm-ledger}
For an affine received line $f+Zg$, take $d=1$, $f_0=f$, and $f_1=g$.  Then
$R_{\mathcal T_{\mathcal C}^{\ge a}}^{(1)}(Z)$ is the squarefree product of the
linear factors $Z-z$ for the distinct finite quotient-witnessed support-wise
MCA-bad slopes.  Hence
\[
        \#\{\text{finite quotient-witnessed MCA-bad slopes}\}
        =\deg R_{\mathcal T_{\mathcal C}^{\ge a}}^{(1)} .
\]
The same polynomial upper-bounds no-loss CA finite slopes whose explaining
support lies in the declared quotient family.
\end{corollary}

\begin{proof}
For $d=1$, each nonconstant fixed-support gcd has degree at most one.  Its root,
if present, is exactly the unique finite slope supplied by
\cref{lem:exact-line-image-map}.  Taking the squarefree lcm over the declared
co-support family therefore multiplies one copy of $Z-z$ for each distinct
finite slope and no copies for supports that contribute no finite noncontained
slope.
\end{proof}

\begin{definition}[projective line lcm ledger]
\label{def:projective-line-lcm-ledger}
For a projective received line $Z_0f+Z_1g$ and a co-support $T$ with
$|D\setminus T|>k$, form the homogeneous linear coordinate forms
\[
        L_{T,m}(Z_0,Z_1):=Z_0(A_T)_m+Z_1(B_T)_m,
        \qquad 0\le m<t .
\]
If all $L_{T,m}$ are zero, set $h_T^{\rm proj}=1$.  Otherwise let
$h_T^{\rm proj}$ be the gcd of the nonzero linear forms, normalized by a fixed
monomial order.  Thus $h_T^{\rm proj}$ is either $1$ or a homogeneous linear
form.  For a co-support family $\mathcal T$ define
\[
        R_{\mathcal T}^{\rm proj}(Z_0,Z_1)
        :=\operatorname{rad}\operatorname{lcm}_{T\in\mathcal T}h_T^{\rm proj}.
\]
\end{definition}

\begin{corollary}[exact projective quotient image ledger]
\label{cor:projective-line-lcm-ledger}
For projective slopes, the exact quotient-branch support-wise MCA numerator is
\[
        \deg R_{\mathcal T_{\mathcal C}^{\ge a}}^{\rm proj}.
\]
The same degree upper-bounds no-loss projective-CA parameters whose explaining
support lies in the quotient family.
\end{corollary}

\begin{proof}
A nonzero family of homogeneous linear forms on $\mathbb P^1$ has a common zero
if and only if the forms have a common homogeneous linear factor; this factor
cuts out exactly one $\F$-rational projective point.  The all-zero case is the
contained branch and contributes no bad point, by convention.  The squarefree
lcm therefore contains one linear factor for each distinct projective parameter
and no duplicates.
\end{proof}

\begin{definition}[lcm quotient-image certificate]
\label{def:lcm-quotient-image-certificate}
An lcm quotient-image certificate for a line or finite-parameter curve consists
of:
\begin{enumerate}[label=\textup{(\roman*)},leftmargin=2.2em]
\item the divisor set $\mathcal C$, the agreement threshold $a$, and either the
exact support-union certificate of \cref{def:divisor-block-support-ledger-certificate}
or an explicit hash/orbit enumeration of the co-support family
$\mathcal T_{\mathcal C}^{\ge a}$;
\item for every represented co-support $T$, the locator vector
$\boldsymbol\ell_T$ and the syndrome-window coefficient vectors $V_{i,T}$;
\item the coordinate polynomials $P_{T,m}(\Gamma)$, the fixed-support gcd
$g_T^{(d)}$, and a proof that the all-zero case has been assigned to the
contained branch;
\item the squarefree lcm $R_{\mathcal T}^{(d)}$ and, over the finite field, the
root-count factor $R_{\mathcal T,F}^{(d)}=\gcd(R_{\mathcal T}^{(d)},\Gamma^q-\Gamma)$;
\item the declared numerator $\deg R_{\mathcal T,F}^{(d)}$ for finite curves or
$\deg R_{\mathcal T}^{\rm proj}$ for projective lines.
\end{enumerate}
The verifier checks the locator split condition, the syndrome recurrence, the
gcd and lcm arithmetic, squarefreeness, and the final finite-field root count.
This certificate replaces a support count by an exact distinct-parameter count
inside the declared quotient branch.
\end{definition}

\begin{remark}[status after the lcm image ledger]
\label{rem:lcm-image-ledger-status}
The quotient ledgers now separate three tasks that were previously entangled.
First, \cref{thm:quotient-remainder-deep-floor} supplies lower floors by
heaviest generated-field prefix fibers.  Second,
\cref{thm:exact-divisor-block-support-ledger} counts the declared quotient
support union with divisor overlaps handled exactly.  Third,
\cref{thm:exact-quotient-image-lcm-ledger,cor:affine-line-lcm-ledger,cor:projective-line-lcm-ledger}
push that support family through the fixed-locator equations and coalesce
duplicate finite, projective, and curve parameters by gcd/lcm arithmetic.  The
remaining quotient-side obstruction is not duplicate coalescing inside a declared
quotient family; it is the structural safe-side theorem that every bad parameter
in the intended regime either belongs to this quotient image, to the tangent
branch, or to the aperiodic/extension buckets bounded elsewhere.
\end{remark}

\begin{definition}[quotient image certificate]
\label{def:quotient-image-certificate}
A quotient image certificate for a line or curve ledger consists of:
\begin{enumerate}[label=\textup{(\roman*)}]
\item the divisor set $\mathcal C$ and the agreement threshold $a$;
\item a deterministic description, hash, or orbit decomposition of
$\mathcal T_{\mathcal C}^{\ge a}$;
\item for every represented co-support $T$, the locator coefficients
$\boldsymbol\ell_T$ and the syndrome-window vectors $A_T,B_T$ for a line, or
$V_{0,T},\ldots,V_{d,T}$ for a curve;
\item the declared parameter value, parameter bucket, or the declaration that no
parameter is contributed;
\item a bucket table proving the number of distinct parameters after duplicate
coalescing.
\end{enumerate}
The verifier checks the locator split condition, the Hankel equations, the
noncontainment vector, and the bucket consistency.  This certificate is the
preferred replacement for a raw quotient-support count whenever the support count
is much larger than the actual parameter image.
\end{definition}

\begin{remark}[status of the refinement]
\label{rem:distinct-ledger-status}
The propositions above are exact for the declared quotient branch and are a real
refinement over the support ledger: the numerator is an image size, not a support
count.  They still do not prove that all bad parameters are quotient-periodic.
The remaining safe-side theorem for the full prize is the aperiodic Hankel-packing
bound: after tangent and quotient image branches are removed, one must bound the
parameters coming from co-supports outside the declared quotient families.
\end{remark}


\subsection{Aperiodic Hankel chart atlas}
\label{sec:aperiodic-hankel-certificates}

The quotient image ledger of \cref{sec:quotient-distinct-parameter} counts bad
parameters whose witness support lies in a declared quotient-remainder family.
The complementary safe-side problem is to control all other supports.  This
subsection is the contributor-facing formulation of that problem.  It does not claim
that the needed aperiodic eliminants always exist.  It gives a finite chart
atlas, with precise ideals and pivots, so that a verifier can check a supplied
eliminant and convert it into a bad-parameter upper bound.

The main point is to avoid one huge saturated incidence ideal.  We split the
aperiodic branch into:
\begin{enumerate}[label=\textup{(\roman*)},leftmargin=2.2em]
\item a regular overdetermined Hankel bucket, where a single maximal minor in the
line parameter already gives an eliminant;
\item finite affine pivot charts, one for each noncontainment coordinate;
\item a separate projective-infinity chart, if projective slopes are being
counted;
\item finite-parameter curve pivot charts, one for a nonzero coefficient of one
coordinate polynomial;
\item a singular bucket, where all regular minors vanish or the projection to the
parameter line remains positive-dimensional after the previous splits.
\end{enumerate}
Only the last bucket is genuinely unresolved.  A staircase scanner should keep it
as a named obstruction rather than hiding it inside the word ``aperiodic.''

Fix an exact agreement value $A\ge k$ and put
\[
        j=j_A:=n-A,
        \qquad
        t=t_A:=A-k,
        \qquad
        R=n-k .
\]
Let
\[
        P_D(X):=\prod_{x\in D}(X-x).
\]
A monic degree-$j$ co-support locator is written
\[
        L_{\ell}(X)=X^j+\ell_{j-1}X^{j-1}+\cdots+\ell_0,
        \qquad \ell=(\ell_0,\ldots,\ell_{j-1})\in\mathbb A^j .
\]
Let $R_{D,j}(\ell;X)$ be the remainder of $P_D(X)$ modulo $L_{\ell}(X)$, and let
$I_{\rm split,j}$ be the ideal generated by the $j$ coefficients of
$R_{D,j}$.  Since $P_D$ is squarefree, the $F$-points of
$V(I_{\rm split,j})$ are exactly the squarefree $D$-split monic degree-$j$
locators.

For a word $y:D\to F$ write $\Syn(y)\in F^R$ for the parity-check syndrome used
above.  Let
\[
        \widetilde\ell=(\ell_0,\ldots,\ell_{j-1},1)^t .
\]
For an affine line $f+zg$, put $u=\Syn(f)$ and $v=\Syn(g)$ and define
\[
        A(\ell):=H_{t,j}(u)\widetilde\ell,
        \qquad
        B(\ell):=H_{t,j}(v)\widetilde\ell
        \qquad\in F^t .
\]
For a fixed split locator, the support-wise finite-slope condition is exactly
\[
        A(\ell)+ZB(\ell)=0,
        \qquad
        B(\ell)\ne0 .
\]
The inequality is the noncontainment condition.

\subsubsection{Regular overdetermined bucket}

Assume first that
\[
        t_A\ge j_A+1,
        \qquad\text{equivalently}\qquad
        2A\ge n+k+1 .
\]
Set
\[
        M_A(Z):=H_{t_A,j_A}(u)+Z H_{t_A,j_A}(v).
\]
This is a $t_A\times(j_A+1)$ matrix with entries affine-linear in $Z$.

\begin{definition}[regular Hankel minor certificate]
\label{def:hankel-regularity-certificate}
A regular Hankel minor certificate at exact agreement $A$ is a row set
$I_A\subseteq\{0,\ldots,t_A-1\}$ of size $j_A+1$ such that
\[
        \Delta_A(Z)
        :=\det M_A(Z)_{I_A,\{0,\ldots,j_A\}}
\]
is not the zero polynomial in $F[Z]$.  If no such row set exists, the line is
called \emph{Hankel-singular at $A$}.
\end{definition}

\begin{lemma}[regular exact-agreement eliminant]
\label{lem:regular-exact-agreement-eliminant}
Assume $t_A\ge j_A+1$, and suppose $\Delta_A(Z)$ is a regular Hankel minor
certificate.  Every finite support-wise MCA-bad slope with a witness support of
exact size $A$ is a root of $\Delta_A$.  Hence the number of such slopes is at
most
\[
        \deg\Delta_A\le j_A+1=n-A+1 .
\]
The same root containment holds for no-loss CA-bad slopes with exact agreement
$A$.
\end{lemma}

\begin{proof}
Let $z$ be a bad finite slope with witness support $S$ of size $A$, and put
$T=D\setminus S$.  Then $|T|=j_A$, and the locator vector
$\widetilde\ell_T$ is a nonzero kernel vector of $M_A(z)$.  Therefore the column
rank of $M_A(z)$ is less than $j_A+1$, so every $(j_A+1)\times(j_A+1)$ maximal
minor vanishes at $z$, including $\Delta_A(z)$.  The degree bound follows because
$\Delta_A$ is the determinant of a $(j_A+1)\times(j_A+1)$ matrix whose entries
are affine-linear in $Z$.
\end{proof}

\begin{theorem}[regular closed-range Hankel packing]
\label{thm:regular-closed-ball-hankel-packing}
Let $a\ge k$ and suppose that for each exact agreement
$A\in\{a,a+1,\ldots,n\}$ with $t_A\ge j_A+1$ a regular Hankel minor certificate
$\Delta_A$ is supplied.  Then all bad finite slopes whose exact witness size is
in this regular bucket lie in the root set of
\[
        \Delta_{\ge a}(Z):=\prod_{A=a}^{n}\Delta_A(Z),
\]
with factors omitted for exact agreements outside the overdetermined range.
Consequently their number is at most
\[
        \sum_A \deg\Delta_A .
\]
If exact root sets are supplied, roots already charged to the tangent or
quotient-image ledgers may be removed from this count by duplicate coalescing.
\end{theorem}

\begin{proof}
Apply \cref{lem:regular-exact-agreement-eliminant} at the exact size of the
witness support and take the union of the resulting root sets.  The degree sum is
a safe union bound.
\end{proof}

\paragraph{Canonical regular rank-drop gcds.}
The regular-minor certificate below can be made canonical.  Instead of choosing
one nonzero maximal minor, take the greatest common divisor of all nonzero
maximal minors.  This gives exactly the finite-parameter values at which the
Hankel pencil can drop rank, before the split-locator and noncontainment tests
are imposed.  Thus it is an upper ledger for the regular bucket, and often a much
smaller one than the degree of a single selected minor.

Fix an exact agreement value $A\ge k$, and keep the notation
\[
        j=j_A:=n-A,
        \qquad
        t=t_A:=A-k,
        \qquad
        r=n-k .
\]
The regular overdetermined range is
\[
        t_A>j_A,
        \qquad\text{i.e.}\qquad
        t_A\ge j_A+1,
        \qquad\text{equivalently}\qquad
        2A\ge n+k+1 .
\]

\begin{definition}[canonical affine rank-drop gcd]
\label{def:canonical-affine-rankdrop-gcd}
Fix an affine received line $f+Zg$, put
$u=\Syn(f)$ and $v=\Syn(g)$, and assume $t_A\ge j_A+1$.  Set
\[
        M_A(Z):=H_{t_A,j_A}(u)+Z H_{t_A,j_A}(v).
\]
Let $\mathfrak M_A(Z)$ be the finite set of all $(j_A+1)\times(j_A+1)$ maximal
minors of $M_A(Z)$.  If every element of $\mathfrak M_A(Z)$ is the zero
polynomial, the exact-agreement bucket is called \emph{rank-drop singular}.
Otherwise define
\[
        G_A^{\rm aff}(Z)
        :=\operatorname{rad}\gcd\{\Delta(Z):
              \Delta\in\mathfrak M_A(Z),\ \Delta\ne0\},
\]
with the gcd normalized to be monic.  Over the finite field $\F$ of size $q$ set
\[
        G_{A,F}^{\rm aff}(Z):=
        \gcd\bigl(G_A^{\rm aff}(Z),Z^q-Z\bigr).
\]
Thus $G_{A,F}^{\rm aff}$ is the squarefree polynomial whose roots are precisely
those $F$-rational parameters at which all nonzero maximal minors vanish.
\end{definition}

\begin{theorem}[canonical exact-agreement rank-drop ledger]
\label{thm:canonical-affine-rankdrop-ledger}
Assume $t_A\ge j_A+1$ and the exact-agreement bucket is not rank-drop singular.
Every finite support-wise MCA-bad slope with an exact witness support of size
$A$ is a root of $G_A^{\rm aff}$.  Consequently the number of such slopes is at
most
\[
        \deg G_{A,F}^{\rm aff}
        \le \deg G_A^{\rm aff}
        \le j_A+1=n-A+1 .
\]
The same bound holds for no-loss CA-bad slopes with exact agreement $A$.
\end{theorem}

\begin{proof}
Let $z\in F$ be support-wise bad with exact witness support $S$ of size $A$, and
put $T=D\setminus S$.  By \cref{lem:support-locator-syndrome-recurrence}, the
locator vector $\boldsymbol\ell_T$ is a nonzero kernel vector of $M_A(z)$.
Therefore $\operatorname{rank}M_A(z)<j_A+1$, so every maximal minor of $M_A(Z)$
vanishes at $Z=z$.  In particular every nonzero maximal minor vanishes at $z$,
and hence their gcd $G_A^{\rm aff}$ vanishes at $z$.

The finite-field root count is exactly $\deg G_{A,F}^{\rm aff}$ because
$G_{A,F}^{\rm aff}$ is squarefree and divides $Z^q-Z$.  Each maximal minor has
degree at most $j_A+1$, since it is the determinant of a matrix whose entries are
affine-linear in $Z$; the gcd cannot have larger degree.  The no-loss CA claim is
identical, because no-loss CA-badness with exact agreement $A$ also supplies an
explaining support of size $A$ and hence the same locator-kernel condition.
\end{proof}

\begin{remark}[why this improves the one-minor certificate]
\label{rem:gcd-rankdrop-improves-one-minor}
\Cref{lem:regular-exact-agreement-eliminant} used one certified nonzero maximal
minor and therefore counted the roots of a single determinant.  The polynomial
$G_A^{\rm aff}$ uses all maximal minors and counts their common finite-field root
set.  It is never worse than any one-minor certificate and can be much smaller;
it is also canonical, so two scanners that use the same syndrome convention print
the same regular-bucket polynomial.
\end{remark}

\begin{definition}[closed-ball rank-drop lcm]
\label{def:closed-ball-rankdrop-lcm}
For an agreement threshold $a>k$, let
\[
        \mathcal A_{\rm reg}(a):=
        \{A\in\{a,a+1,\ldots,n\}: t_A\ge j_A+1
          \text{ and the bucket }A\text{ is not rank-drop singular}\}.
\]
Define
\[
        G_{\ge a}^{\rm aff}(Z)
        :=\operatorname{rad}\operatorname{lcm}_{A\in\mathcal A_{\rm reg}(a)}
          G_A^{\rm aff}(Z),
        \qquad
        G_{\ge a,F}^{\rm aff}(Z):=
        \gcd\bigl(G_{\ge a}^{\rm aff}(Z),Z^q-Z\bigr),
\]
with the empty lcm equal to $1$.  Exact agreement values
$A\in\{a,\ldots,n\}\setminus\mathcal A_{\rm reg}(a)$ are recorded as residual
buckets: they are either underdetermined ($t_A<j_A+1$) or rank-drop singular.
\end{definition}

\begin{corollary}[canonical regular closed-ball Hankel packing]
\label{cor:canonical-regular-closed-ball-hankel-packing}
All finite support-wise MCA-bad slopes whose exact witness size lies in
$\mathcal A_{\rm reg}(a)$ are roots of $G_{\ge a}^{\rm aff}$.  Hence their number
is at most
\[
        \deg G_{\ge a,F}^{\rm aff}
        \le \sum_{A\in\mathcal A_{\rm reg}(a)}(n-A+1).
\]
The same bound holds for no-loss CA-bad slopes in the same regular buckets.
\end{corollary}

\begin{proof}
Apply \cref{thm:canonical-affine-rankdrop-ledger} at the exact agreement size of
the witness and take the union over $A$.  The squarefree lcm coalesces duplicate
roots across exact-agreement buckets; the final gcd with $Z^q-Z$ counts only
finite-field roots.  The displayed degree sum is the safe bound obtained before
coalescing.
\end{proof}

\begin{definition}[paid-root removal for the regular branch]
\label{def:paid-root-removal-regular-branch}
Suppose a squarefree polynomial $P_{\rm paid,F}(Z)$ is supplied whose roots are
exactly the finite-field parameters already charged to tangent/common-line and
quotient-image ledgers.  Define the regular residual polynomial
\[
        G_{{\rm reg}\setminus{\rm paid},F}^{\rm aff}(Z)
        :=\frac{G_{\ge a,F}^{\rm aff}(Z)}
        {\gcd(G_{\ge a,F}^{\rm aff}(Z),P_{\rm paid,F}(Z))} .
\]
Its degree is the number of regular rank-drop parameters not already paid by the
previous ledgers.
\end{definition}

\begin{corollary}[quotient/tangent-removed regular aperiodic numerator]
\label{cor:quotient-tangent-removed-regular-numerator}
If $P_{\rm paid,F}$ is exact in the sense of
\cref{def:paid-root-removal-regular-branch}, then the number of finite slopes in
the regular overdetermined branch that remain after removing the paid
quotient/tangent parameters is at most
\[
        \deg G_{{\rm reg}\setminus{\rm paid},F}^{\rm aff} .
\]
This is a safe aperiodic numerator for the regular buckets.  The underdetermined
and rank-drop singular exact-agreement buckets must still be handled by pivot
charts or by a separately named residual obstruction.
\end{corollary}

\begin{proof}
By \cref{cor:canonical-regular-closed-ball-hankel-packing}, every regular-bucket
bad slope is a root of $G_{\ge a,F}^{\rm aff}$.  Removing the roots of the exact
paid polynomial deletes precisely those already charged parameters from this
finite root set.  The remaining root set is counted by the displayed quotient
because all polynomials are squarefree divisors of $Z^q-Z$.
\end{proof}

\begin{definition}[finite-parameter curve rank-drop gcd]
\label{def:curve-rankdrop-gcd}
For a finite-parameter degree-$d$ curve
\[
        W_\Gamma=f_0+\Gamma f_1+\cdots+\Gamma^d f_d,
\]
put $u_i=\Syn(f_i)$ and, at exact agreement $A$, define
\[
        M_A^{(d)}(\Gamma):=
        \sum_{i=0}^{d}\Gamma^i H_{t_A,j_A}(u_i).
\]
When $t_A\ge j_A+1$ and not all maximal minors of $M_A^{(d)}$ vanish
identically, set
\[
        G_A^{(d)}(\Gamma)
        :=\operatorname{rad}\gcd\{\Delta(\Gamma):
              \Delta\text{ is a nonzero maximal minor of }M_A^{(d)}\},
\]
and
\[
        G_{A,F}^{(d)}(\Gamma):=
        \gcd\bigl(G_A^{(d)}(\Gamma),\Gamma^q-\Gamma\bigr).
\]
\end{definition}

\begin{theorem}[canonical curve rank-drop ledger]
\label{thm:canonical-curve-rankdrop-ledger}
In a non-singular overdetermined exact-agreement bucket, every finite
curve-MCA-bad parameter is a root of $G_A^{(d)}$.  Consequently the number of
such parameters is at most
\[
        \deg G_{A,F}^{(d)}
        \le \deg G_A^{(d)}
        \le d(j_A+1).
\]
The same bound holds for no-loss curve-CA parameters in the same bucket.  The
closed-ball form is obtained by taking the squarefree lcm of the polynomials
$G_A^{(d)}$ over the regular exact-agreement buckets and then taking the gcd with
$\Gamma^q-\Gamma$.
\end{theorem}

\begin{proof}
A curve-bad parameter with exact witness support $S$ gives, for
$T=D\setminus S$, a nonzero locator vector in the kernel of
$M_A^{(d)}(\gamma)$.  Hence every maximal minor vanishes at $\gamma$, and the
same gcd argument as in \cref{thm:canonical-affine-rankdrop-ledger} applies.
Each matrix entry has degree at most $d$ in $\Gamma$, so each maximal minor has
degree at most $d(j_A+1)$.
\end{proof}

\begin{definition}[regular rank-drop certificate]
\label{def:regular-rankdrop-certificate}
A regular rank-drop certificate for a row, a received affine line or finite
curve, and an agreement threshold $a$ consists of:
\begin{enumerate}[label=\textup{(\roman*)},leftmargin=2.2em]
\item the syndrome vectors of the coefficient words and the exact-agreement range
$A\in\{a,\ldots,n\}$;
\item for every overdetermined exact agreement $A$, either a declaration that all
maximal minors vanish identically, or the canonical gcd polynomial
$G_A^{\rm aff}$, respectively $G_A^{(d)}$, together with ideal-membership or raw
minor arithmetic proving that it is the gcd of all nonzero maximal minors;
\item the squarefree lcm over all non-singular regular buckets and the finite-root
factor obtained by gcd with $Z^q-Z$ or $\Gamma^q-\Gamma$;
\item if paid branches are removed, the exact paid-root polynomial and the finite
root-set quotient of \cref{def:paid-root-removal-regular-branch};
\item a residual-bucket table listing all underdetermined and rank-drop singular
exact agreements that are not covered by the regular gcd ledger.
\end{enumerate}
The verifier checks the syndrome convention, the Hankel matrices, the maximal
minor gcds, squarefree lcms, finite-field root factors, paid-root removal, and the
residual-bucket table.
\end{definition}

\begin{remark}[status after the rank-drop ledger]
\label{rem:rankdrop-ledger-status}
The regular overdetermined part of the aperiodic Hankel problem is now canonical:
one no longer has to choose a minor or accept a degree sum when the common root
set of all maximal minors is smaller.  This still does not solve the near-capacity
underdetermined buckets, which form the hard aperiodic residual class.
Those buckets remain assigned to the affine pivot charts, projective-infinity
charts, curve coefficient pivots, or named singular residual obstructions.
\end{remark}

\subsubsection{Finite affine pivot charts}

The regular minor test is intentionally coarse and only works directly in the
overdetermined range.  The next atlas is the one contributors should use to
construct actual aperiodic eliminants.

\begin{definition}[aperiodic locator chart]
\label{def:aperiodic-locator-chart}
An aperiodic locator chart at exact agreement $A$ is a pair
\[
        X=(E_X,\Delta_X),
\]
where $E_X$ is a finite list of polynomial equations in the locator coefficients
$\ell$, and $\Delta_X(\ell)$ is a product of declared nonzero tests.  The chart
represents the constructible set
\[
        V(I_{\rm split,j}+\langle E_X\rangle)\setminus V(\Delta_X).
\]
The equations and inequations must include, or reference, the removal of already
paid tangent/common-code-line and quotient-image branches.  A family of charts
$\mathfrak X_A$ is a valid aperiodic chart cover if the union of their
constructible sets contains every remaining split locator of exact co-support
size $j_A$ not assigned to those earlier ledgers.
\end{definition}

For a coordinate $h\in\{0,\ldots,t-1\}$ define the affine pivot consistency ideal
\[
\begin{aligned}
        J^{\rm aff}_{X,h}:={}&
        I_{\rm split,j}+\langle E_X\rangle
        +\langle A_m(\ell)B_h(\ell)-A_h(\ell)B_m(\ell):0\le m<t\rangle \\
        &+\langle ZB_h(\ell)+A_h(\ell)\rangle
        \subseteq F[\ell_0,\ldots,\ell_{j-1},Z],
\end{aligned}
\]
and saturate it by the pivot denominator:
\[
        \widehat J^{\rm aff}_{X,h}:=
        J^{\rm aff}_{X,h}:(\Delta_X B_h)^\infty .
\]
The equations $A_mB_h-A_hB_m=0$ force $A$ and $B$ to be collinear on the pivot
chart, and $ZB_h+A_h=0$ records the unique finite slope.

\begin{theorem}[pivoted affine-line eliminant certificate]
\label{thm:aperiodic-affine-eliminant}
Let $X$ be an aperiodic locator chart and let $h$ be a noncontainment pivot.  If
there is a nonzero polynomial
\[
        Q_{X,h}(Z)\in \widehat J^{\rm aff}_{X,h}\cap F[Z]
\]
of degree $D_{X,h}$, then the number of finite support-wise MCA-bad slopes in
that chart with $B_h\ne0$ is at most $\min\{|F|,D_{X,h}\}$.  The same bound holds
for no-loss CA-bad slopes in that pivot chart.
\end{theorem}

\begin{proof}
A bad slope in the chart has a split locator with $\Delta_X\ne0$ and, on this
pivot chart, $B_h\ne0$.  The Hankel incidence equation gives
$A+ZB=0$, hence the point $(\ell,Z)$ satisfies all generators of
$J^{\rm aff}_{X,h}$ and survives the saturation.  Therefore $Q_{X,h}(Z)=0$.
A nonzero degree-$D_{X,h}$ polynomial over $F$ has at most $D_{X,h}$ roots.
\end{proof}

\begin{remark}[why the pivot split is useful]
\label{rem:pivot-split-useful}
The saturation by the whole vector $B$ in the unsplit incidence ideal is replaced
by the explicit open cover $B_h\ne0$.  A scanner can try all $h$, produce small
univariate eliminants on successful pivots, and report the residual locus
$B_0=\cdots=B_{t-1}=0$ as contained for support-wise MCA.  For no-loss CA, that
residual locus also contributes no bad slope at the same support, because if
$A=B=0$ then both $f$ and $g$ are explained on the witness support.
\end{remark}

\subsubsection{Projective-slope split}

For projective parameters $[Z_0:Z_1]$ the incidence equation is
\[
        Z_0A(\ell)+Z_1B(\ell)=0 .
\]
The finite patch $Z_0\ne0$ is exactly the affine pivot atlas above with
$Z=Z_1/Z_0$.  The extra point at infinity is $[0:1]$.

\begin{definition}[projective infinity chart]
\label{def:projective-infinity-chart}
For a chart $X$, the projective-infinity chart is the constructible locus
\[
        I_{\rm split,j}+\langle E_X\rangle+
        \langle B_0,\ldots,B_{t-1}\rangle,
        \qquad
        \Delta_X\ne0,
        \quad A\ne0 .
\]
If it is nonempty, it contributes the single projective parameter $[0:1]$.
If it is empty, the infinity point contributes nothing on $X$.
\end{definition}

\begin{corollary}[projective eliminant certificate]
\label{cor:projective-eliminant-certificate}
For a chart cover $\mathfrak X_A$, the projective-slope numerator is bounded by
the union of the finite-patch root sets from
\cref{thm:aperiodic-affine-eliminant} together with at most one extra point
$[0:1]$, provided the infinity charts are certified either empty or nonempty.
Equivalently, a homogeneous nonzero eliminant
\[
        Q_X(Z_0,Z_1)
\]
of degree $D_X$ bounds the number of projective bad parameters in that chart by
$D_X$.
\end{corollary}

\begin{proof}
The finite patch is the affine theorem.  The complement of that patch is the
single point $[0:1]$, which is governed by \cref{def:projective-infinity-chart}.
For the homogeneous formulation, every bad projective parameter is a zero of the
homogeneous eliminant, and a nonzero homogeneous polynomial of degree $D_X$ on
$\mathbb P^1(F)$ has at most $D_X$ zeros.
\end{proof}

\subsubsection{Finite-parameter curve charts}

For a degree-$d$ finite power curve
\[
        W_\Gamma=f_0+\Gamma f_1+\cdots+\Gamma^d f_d,
\]
put
\[
        V_i(\ell):=H_{t,j}(\Syn(f_i))\widetilde\ell\in F^t,
        \qquad 0\le i\le d .
\]
The curve incidence equation is
\[
        V_0(\ell)+\Gamma V_1(\ell)+\cdots+\Gamma^d V_d(\ell)=0 .
\]
Curve-MCA noncontainment means that not all coefficient vectors $V_i(\ell)$ are
zero.  Thus the noncontained branch is covered by coefficient pivots
$(i,h)$ with $(V_i)_h\ne0$.

For a chart $X$ and a coefficient pivot $(i,h)$, define
\[
\begin{aligned}
        J^{\rm curve,d}_{X,i,h}:={}&
        I_{\rm split,j}+\langle E_X\rangle \\
        &+\left\langle
          \sum_{r=0}^d \Gamma^r (V_r(\ell))_m:
          0\le m<t
          \right\rangle
        \subseteq F[\ell_0,\ldots,\ell_{j-1},\Gamma],
\end{aligned}
\]
and set
\[
        \widehat J^{\rm curve,d}_{X,i,h}
        :=J^{\rm curve,d}_{X,i,h}:(\Delta_X (V_i)_h)^\infty .
\]

\begin{theorem}[pivoted curve eliminant certificate]
\label{thm:aperiodic-curve-eliminant}
If a chart $X$ and coefficient pivot $(i,h)$ have a nonzero eliminant
\[
        Q_{X,i,h}(\Gamma)\in
        \widehat J^{\rm curve,d}_{X,i,h}\cap F[\Gamma]
\]
of degree $D_{X,i,h}$, then the number of finite degree-$d$ curve-MCA bad
parameters in that chart and pivot is at most $\min\{|F|,D_{X,i,h}\}$.  The same
bound applies to no-loss curve-CA parameters in the same pivot chart.
\end{theorem}

\begin{proof}
Every bad curve parameter on the pivot chart gives a split locator satisfying all
curve-incidence equations and the nonzero pivot condition $(V_i)_h\ne0$.  It
therefore survives the saturation and is a root of the supplied eliminant.
\end{proof}

% -----------------------------------------------------------------------------
\paragraph{Residual rank-pivot charts.}
The gcd ledger above removes the regular overdetermined branch.  What remains is
not a formless exceptional set: every remaining bad parameter is still a point of
a finite collection of determinantal rank charts.  The point of the next
statements is to make this residual bucket verifier-facing.  A certificate may
now say exactly which rank chart it uses, which noncontainment coordinate is
open, and which univariate or homogeneous eliminant cuts out the parameter
projection.  If no such eliminant is supplied, the chart is a named residual
obstruction rather than an unaccounted term.

Fix an exact agreement value $A\ge k$ and keep
\[
        j=j_A:=n-A,
        \qquad
        t=t_A:=A-k,
        \qquad
        r=n-k .
\]
For a projective received line $Z_0 f+Z_1 g$, put
\[
        U:=H_{t,j}(\Syn(f)),
        \qquad
        V:=H_{t,j}(\Syn(g)),
        \qquad
        M^{\rm proj}_A(Z_0,Z_1):=Z_0U+Z_1V .
\]
For a locator vector $\widetilde\ell$ write
\[
        A(\ell):=U\widetilde\ell,
        \qquad
        B(\ell):=V\widetilde\ell .
\]
Thus the projective support equation is
\[
        Z_0A(\ell)+Z_1B(\ell)=0,
\]
and projective noncontainment is $(A(\ell),B(\ell))\ne(0,0)$.

Let $X=(E_X,\Delta_X)$ be a locator chart in the sense of
\cref{def:aperiodic-locator-chart}.  For row and column sets
$I\subseteq\{0,\ldots,t-1\}$ and
$J\subseteq\{0,\ldots,j\}$ with $|I|=|J|=s$, write
\[
        \Delta_{I,J}(Z_0,Z_1):=
        \det M^{\rm proj}_A(Z_0,Z_1)_{I,J},
\]
with the convention that the unique $0\times0$ minor is $1$.  Let
$\mathcal M_{s+1}(M^{\rm proj}_A)$ denote the set of all $(s+1)\times(s+1)$
minors, empty when $s=\min\{t,j+1\}$.

For a noncontainment pivot
\[
        \pi\in\{A_h:0\le h<t\}\cup\{B_h:0\le h<t\},
\]
let $\Pi_\pi(\ell)$ denote the corresponding scalar coordinate.  Define the
saturated projective rank-pivot ideal
\[
\begin{aligned}
        \widehat J^{\rm proj}_{X,I,J,\pi}:=
        \bigl(&I_{{\rm split},j}+\langle E_X\rangle
        +\langle M^{\rm proj}_A(Z_0,Z_1)\widetilde\ell\rangle\\
        &+\langle \Delta: \Delta\in\mathcal M_{s+1}(M^{\rm proj}_A)\rangle
        \bigr):
        \bigl(\Delta_X\Delta_{I,J}\Pi_\pi\bigr)^\infty
        \subseteq F[\ell_0,\ldots,\ell_{j-1},Z_0,Z_1].
\end{aligned}
\]
The open condition $\Delta_{I,J}\ne0$ and the vanishing of all larger minors put
$M^{\rm proj}_A$ on the exact-rank-$s$ stratum on this chart.  The saturation by
$\Pi_\pi$ chooses one of the finitely many open noncontainment coordinates.

\begin{theorem}[residual rank-pivot cover]
\label{thm:residual-rank-pivot-cover}
Fix an exact agreement value $A$ and a chart cover $\mathfrak X_A$ of the
remaining locator branch after tangent/common-line and quotient-image removals.
Every projective support-wise MCA-bad parameter with exact witness size $A$ in
that remaining branch lies on one of the rank-pivot charts
\[
        \widehat J^{\rm proj}_{X,I,J,\pi}
\]
for some $X\in\mathfrak X_A$, some rank $s$, some nonzero $s\times s$ minor
$\Delta_{I,J}$ at the parameter, and some nonzero noncontainment coordinate
$\Pi_\pi$.  The same cover contains every no-loss projective CA-bad parameter in
that branch.
\end{theorem}

\begin{proof}
Let $[z_0:z_1]$ be bad with witness support $S$, co-support $T$, and locator
vector $\widetilde\ell_T$.  Since the chart family covers the remaining branch,
$\widetilde\ell_T$ lies in some $X\in\mathfrak X_A$ with $\Delta_X\ne0$.  The
support equation gives
$M^{\rm proj}_A(z_0,z_1)\widetilde\ell_T=0$.  Let
$s=\operatorname{rank}M^{\rm proj}_A(z_0,z_1)$.  Choose a nonzero
$s\times s$ minor at this parameter; then every $(s+1)\times(s+1)$ minor
vanishes.  Finally, support-wise noncontainment says that
$(A(\ell_T),B(\ell_T))\ne(0,0)$, so some coordinate pivot $\Pi_\pi$ is nonzero.
Thus the point survives the displayed saturation and belongs to the indicated
rank-pivot chart.  The no-loss CA statement uses the same explaining support and
same noncontainment conclusion.
\end{proof}

\begin{theorem}[projective residual eliminant certificate]
\label{thm:projective-residual-eliminant-certificate}
If a rank-pivot chart has a nonzero homogeneous eliminant
\[
        Q_{X,I,J,\pi}(Z_0,Z_1)
        \in
        \widehat J^{\rm proj}_{X,I,J,\pi}\cap F[Z_0,Z_1]
\]
of degree $D_{X,I,J,\pi}$, then the number of projective line parameters in that
chart is at most $D_{X,I,J,\pi}$.  More exactly, over a finite field $F$ of size
$q$, the number of finite-patch roots is
\[
        \deg\gcd\bigl(Q_{X,I,J,\pi}(1,Z),Z^q-Z\bigr),
\]
and the point at infinity contributes precisely when
$Q_{X,I,J,\pi}(0,1)=0$.  Thus the exact finite-field count is
\[
        \deg\gcd\bigl(Q_{X,I,J,\pi}(1,Z),Z^q-Z\bigr)
        +\mathbf 1_{Q_{X,I,J,\pi}(0,1)=0}.
\]
\end{theorem}

\begin{proof}
Every parameter represented by the chart gives a point of the saturated
incidence variety, hence is a zero of every element of the elimination ideal,
including $Q_{X,I,J,\pi}$.  A nonzero homogeneous polynomial of degree $D$ on
$\mathbb P^1(F)$ has at most $D$ projective zeros.  On the affine patch
$Z_0=1$ the finite parameters are exactly the roots of $Q(1,Z)$ in $F$, which
are counted by the gcd with $Z^q-Z$ after squarefree reduction.  The complement
of the finite patch is the single point $[0:1]$, giving the displayed indicator.
\end{proof}

\begin{corollary}[finite and infinity patches]
\label{cor:finite-infinity-rank-pivot-patches}
The finite affine rank-pivot certificate is obtained from
\cref{thm:projective-residual-eliminant-certificate} by setting $Z_0=1$ and
using only noncontainment pivots $B_h\ne0$.  The infinity chart is obtained by
setting $Z_0=0$, $Z_1=1$, imposing
\[
        B(\ell)=0,
        \qquad
        A_h(\ell)\ne0
\]
for some $h$, and applying the same saturated rank-pivot test to the constant
matrix $V$.  Consequently the projective sampler has no hidden extra case: every
finite or infinite bad point is represented by one of these two patch types.
\end{corollary}

\begin{proof}
On the finite patch $[1:Z]$, the support equation is
$A(\ell)+ZB(\ell)=0$.  If $B(\ell)=0$, then also $A(\ell)=0$, so the fixed
support is contained; hence every finite noncontained point has some
$B_h\ne0$.  At infinity $[0:1]$, the support equation is $B(\ell)=0$ and
noncontainment is exactly $A(\ell)\ne0$.  These alternatives exhaust
$\mathbb P^1(F)$.
\end{proof}

\paragraph{Curve coefficient-pivot rank charts.}
For a finite-parameter degree-$d$ power curve
\[
        W_\Gamma=f_0+\Gamma f_1+\cdots+\Gamma^d f_d,
\]
put $U_i:=H_{t,j}(\Syn(f_i))$ and
\[
        M_A^{(d)}(\Gamma):=
        \sum_{i=0}^{d}\Gamma^i U_i,
        \qquad
        V_i(\ell):=U_i\widetilde\ell .
\]
For row and column sets $I,J$ of size $s$, write
$\Delta_{I,J}^{(d)}(\Gamma)$ for the corresponding $s\times s$ minor of
$M_A^{(d)}(\Gamma)$, and let $\mathcal M_{s+1}(M_A^{(d)})$ be the set of all
larger minors.  For a coefficient noncontainment pivot $(i,h)$ define
\[
\begin{aligned}
        \widehat J^{\rm curve,d}_{X,I,J,i,h}:=
        \bigl(&I_{{\rm split},j}+\langle E_X\rangle
        +\langle M_A^{(d)}(\Gamma)\widetilde\ell\rangle\\
        &+\langle \Delta: \Delta\in\mathcal M_{s+1}(M_A^{(d)})\rangle
        \bigr):
        \bigl(\Delta_X\Delta_{I,J}^{(d)}(V_i)_h\bigr)^\infty .
\end{aligned}
\]

\begin{theorem}[curve residual rank-pivot eliminant]
\label{thm:curve-residual-rank-pivot-eliminant}
The curve-bad parameters in a residual exact-agreement branch are covered by the
charts $\widehat J^{\rm curve,d}_{X,I,J,i,h}$.  If such a chart has a nonzero
eliminant
\[
        Q_{X,I,J,i,h}(\Gamma)
        \in \widehat J^{\rm curve,d}_{X,I,J,i,h}\cap F[\Gamma]
\]
of degree $D_{X,I,J,i,h}$, then that chart contributes at most
\[
        \deg\gcd\bigl(Q_{X,I,J,i,h}(\Gamma),\Gamma^q-\Gamma\bigr)
        \le D_{X,I,J,i,h}
\]
finite parameters over $F$.  The same bound applies to no-loss curve-CA
parameters in the same chart.
\end{theorem}

\begin{proof}
A curve-bad parameter with locator $\widetilde\ell_T$ satisfies
$M_A^{(d)}(\gamma)\widetilde\ell_T=0$.  Taking $s$ to be the rank of
$M_A^{(d)}(\gamma)$ gives a nonzero $s\times s$ minor and vanishing of all larger
minors.  Noncontainment of the coefficient tuple means that some coordinate
$(V_i(\ell_T))_h$ is nonzero.  Thus the point lies on one of the saturated
rank-pivot charts and any univariate eliminant for that chart must vanish at
$\gamma$.  The finite-field root count is the gcd with $\Gamma^q-\Gamma$.
\end{proof}

\paragraph{Extension-line Weil restriction.}
The previous charts are intrinsic over the line field.  For extension-field
samplers one also needs a generated-field ledger: a low-dimensional set over the
base field $\B$ can be negligible after division by $|F|$, even when it is large
as a subset of $\B$.  The following is the certificate form of that accounting.

Let $F/\B$ have degree $m$, choose a $\B$-basis
$\omega_0,\ldots,\omega_{m-1}$ of $F$, and write a finite affine line parameter
as
\[
        Z=Z_0\omega_0+\cdots+Z_{m-1}\omega_{m-1},
        \qquad Z_i\in\B .
\]
Assume $D\subseteq\B$ and write every locator coefficient in the same base field
coordinates when necessary.  Replacing each $F$-valued equation in a rank-pivot
chart by its $m$ coordinate equations over $\B$ gives a Weil-restricted saturated
ideal
\[
        \widehat J^{\rm WR}_{Y}
        \subseteq
        \B[\ell\text{-coordinates},Z_0,\ldots,Z_{m-1}]
\]
for a chart $Y$; noncontainment is split by one nonzero $\B$-coordinate of the
chosen $F$-valued pivot and included in the saturation.
Let $\pi_Z$ denote projection to the parameter-coordinate space
$\mathbb A^m_{\B}$.

\begin{theorem}[extension-line dimension--degree ledger]
\label{thm:extension-line-dimension-degree-ledger}
Suppose that for a Weil-restricted residual chart $Y$ a certificate gives a
closed set
\[
        Y_Z\supseteq \pi_Z\bigl(V(\widehat J^{\rm WR}_{Y})\bigr)
        \subseteq \mathbb A^m_{\B}
\]
of degree at most $\Delta_Y$ and dimension at most $e_Y$.  Then the number of
extension-field affine parameters contributed by this chart is at most
\[
        \Delta_Y |\B|^{e_Y}.
\]
If $Y_Z$ is zero-dimensional and radical, this bound may be replaced by the
exact number of its $\B$-rational points.  After normalization by the line field,
this chart contributes at most
\[
        \frac{\Delta_Y |\B|^{e_Y}}{|F|}
        =\Delta_Y |\B|^{e_Y-m}
\]
to the MCA or CA error numerator.
\end{theorem}

\begin{proof}
The chosen basis identifies $F$ with $\B^m$, so affine line parameters over $F$
are exactly the $\B$-rational points of $\mathbb A^m_{\B}$.  Every bad parameter
represented by the chart lies in the projection of the Weil-restricted incidence
variety and hence in $Y_Z(\B)$.  The standard affine degree point-count bound gives
$|Y_Z(\B)|\le\Delta_Y|\B|^{e_Y}$.  The normalized statement divides by
$|F|=|\B|^m$.
\end{proof}

\begin{corollary}[base-rational line inertness as a chart certificate]
\label{cor:base-rational-line-inertness-chart}
If $f$ and $g$ are $\B$-valued, then every finite affine bad slope on every
rank-pivot chart lies in $\B$.  Equivalently, in the above coordinates all
extension-valued charts with $(Z_1,\ldots,Z_{m-1})\ne0$ are empty.
\end{corollary}

\begin{proof}
For a fixed support and locator over $\B$, the vectors $A(\ell)$ and $B(\ell)$
are $\B$-valued.  On a finite noncontained chart some coordinate $B_h$ is
nonzero, and the slope is forced by
$Z=-A_h/B_h\in\B$.  Hence no genuinely extension-valued finite parameter can
survive the chart equations.
\end{proof}

\paragraph{Residual aperiodic certificate theorem.}
We can now state the safe-side theorem in the form a scanner should print.  Let
$P_{\rm paid,F}$ be the exact squarefree finite-root polynomial already charged
to tangent/common-line and quotient-image branches, and let
$G_{{\rm reg}\setminus{\rm paid},F}^{\rm aff}$ be the regular residual polynomial
of \cref{def:paid-root-removal-regular-branch}.  For the remaining exact
agreement levels, let $\mathfrak R_A$ be a finite list of rank-pivot residual
charts, each equipped with either a finite/projective/curve eliminant or a
Weil-restricted extension projection certificate.

\begin{theorem}[scanner-checkable residual aperiodic ledger]
\label{thm:scanner-checkable-residual-aperiodic-ledger}
Assume that the chart families $\mathfrak R_A$, for $A\ge a$, cover every
remaining locator after tangent/common-line and quotient-image removals.  Suppose
each residual chart $Y$ is certified by one of the following data:
\begin{enumerate}[label=\textup{(\roman*)},leftmargin=2.2em]
\item an affine eliminant $Q_Y(Z)$, contributing
$\deg\gcd(Q_Y,Z^q-Z)$;
\item a homogeneous projective eliminant $Q_Y(Z_0,Z_1)$, contributing
$\deg\gcd(Q_Y(1,Z),Z^q-Z)+\mathbf 1_{Q_Y(0,1)=0}$;
\item a finite curve eliminant $Q_Y(\Gamma)$, contributing
$\deg\gcd(Q_Y,\Gamma^q-\Gamma)$;
\item an extension Weil-restricted projection certificate $(\Delta_Y,e_Y)$,
contributing $\Delta_Y|\B|^{e_Y}$ parameters.
\end{enumerate}
Then the finite affine aperiodic numerator after quotient/tangent removals is
bounded by
\[
        \deg G_{{\rm reg}\setminus{\rm paid},F}^{\rm aff}
        +\sum_{A\ge a}\sum_{Y\in\mathfrak R_A} B_Y,
\]
where $B_Y$ is the corresponding contribution above.  The projective and curve
variants are obtained by using the projective or curve contributions for all
charts.  If exact root sets or exact zero-dimensional projections are printed,
the displayed safe sum may be replaced by their union after duplicate
coalescing.
\end{theorem}

\begin{proof}
The regular nonsingular overdetermined branch is covered by
\cref{cor:quotient-tangent-removed-regular-numerator}.  Every remaining bad
parameter is covered by the rank-pivot cover theorem, or by its curve analogue,
after applying the declared chart cover.  Each chart contribution is an upper
bound by \cref{thm:projective-residual-eliminant-certificate,thm:curve-residual-rank-pivot-eliminant,thm:extension-line-dimension-degree-ledger}.  Summing the
chart bounds is a safe union bound; replacing the sum by an exact union is valid
when the certificate supplies the corresponding finite root or projection sets.
\end{proof}

\begin{definition}[singular residual certificate]
\label{def:singular-residual-certificate}
A singular or underdetermined residual aperiodic certificate consists of:
\begin{enumerate}[label=\textup{(\roman*)},leftmargin=2.2em]
\item the paid tangent/common-line and quotient-image root ledgers already
removed;
\item the regular rank-drop gcd/lcm ledger and its paid-root removal;
\item for each remaining exact agreement $A$, a chart cover of the residual
locator set after the paid branches are removed;
\item for every chart, its rank $s$, pivot minor $\Delta_{I,J}$, noncontainment
pivot, and either an affine, projective, or curve eliminant, or a
Weil-restricted projection degree/dimension certificate;
\item a residual table for every chart on which no eliminant or projection bound
is supplied, labelled as quotient, tangent/common-line, base-field exceptional,
genuinely extension-valued, or new singular obstruction.
\end{enumerate}
The verifier checks chart coverage, rank-stratum equations, saturations by the
pivot and noncontainment coordinates, eliminant membership, finite-field root
counts, extension projection bounds, and duplicate coalescing when claimed.
\end{definition}

\begin{remark}[status after the singular-bucket charts]
\label{rem:singular-bucket-chart-status}
This does not prove that the hard aperiodic residual buckets are small.  It does remove
one ambiguity: underdetermined and rank-drop singular buckets now have a precise
finite certificate language, including finite slopes, projective infinity,
finite-parameter curves, and genuinely extension-valued parameters.  The next
mathematical step is to construct low-degree eliminants, or prove structural
classification, for these rank-pivot charts in the near-capacity regimes where
$t_A=A-k$ is small.
\end{remark}

% === Combined Insert A.2: kernel-compressed residual eliminants and exact residual images ===

\paragraph{Kernel-compressed rank-pivot eliminants.}
The rank-pivot charts above are already finite certificate objects, but the
variables are still the full locator coordinates.  The next reduction removes the
linear kernel equations on a pivot chart by Cramer's rule.  It is the algebraic
form a scanner should use before attempting Gr\"obner elimination or resultants:
all remaining equations are split-locator equations, chart-removal equations,
rank-stratum equations, and noncontainment pivots in the free kernel coordinates
and the line parameter.

We write the statement first on the finite affine patch.  Fix an exact agreement
value $A$, put $j=n-A$ and $t=A-k$, and set
\[
        M_A(Z):=U+ZV,
        \qquad
        U:=H_{t,j}(\Syn(f)),\quad V:=H_{t,j}(\Syn(g)).
\]
Let $I\subseteq\{0,\ldots,t-1\}$ and $J\subseteq\{0,\ldots,j\}$ have the same
cardinality $s$, and put
\[
        \Delta_{I,J}(Z):=\det M_A(Z)_{I,J}.
\]
Let $C_J:=\{0,\ldots,j\}\setminus J$ and introduce free kernel coordinates
$Y=(Y_c)_{c\in C_J}$.  On the open set $\Delta_{I,J}\ne0$, define the polynomial
vector
\[
        \Lambda^{\#}_{I,J}(Z,Y)\in F[Z,Y]^{j+1}
\]
by
\[
        \bigl(\Lambda^{\#}_{I,J}\bigr)_{C_J}
        :=\Delta_{I,J}Y,
        \qquad
        \bigl(\Lambda^{\#}_{I,J}\bigr)_{J}
        :=-\operatorname{adj}\bigl(M_A(Z)_{I,J}\bigr)
             M_A(Z)_{I,C_J}Y .
\]
Thus
\[
        M_A(Z)_{I,*}\Lambda^{\#}_{I,J}(Z,Y)=0.
\]
If, in addition, all $(s+1)\times(s+1)$ minors of $M_A(Z)$ vanish, then every row
of $M_A(Z)$ lies in the span of the pivot rows, and hence
\[
        M_A(Z)\Lambda^{\#}_{I,J}(Z,Y)=0 .
\]

To encode splitting without dividing by the leading coefficient, use homogeneous
locator coordinates.  For
$\lambda=(\lambda_0,\ldots,\lambda_j)$ set
\[
        L_\lambda(X):=\sum_{b=0}^{j}\lambda_bX^b .
\]
Let $Q(X)=\sum_{r=0}^{n-j}q_rX^r$ be an auxiliary quotient polynomial and let
$I_{{\rm hsplit},j}$ be the homogeneous split incidence
\[
        P_D(X)=L_\lambda(X)Q(X),
        \qquad \lambda_j\ne0,
\]
viewed as the coefficient equations saturated by $\lambda_j$.  Equivalently,
its points with $\lambda_j\ne0$ are precisely the nonzero scalar multiples of the
monic locators $\prod_{x\in T}(X-x)$ with $T\subseteq D$ and $|T|=j$.
For a monic locator chart $X=(E_X,\Delta_X)$, write $E_X^{\rm hom}$ and
$\Delta_X^{\rm hom}$ for the usual homogenizations in the coordinates
$\lambda_0,\ldots,\lambda_j$.

For a finite noncontainment pivot $B_h\ne0$, set
\[
        \Pi_h^{\#}(Z,Y):=\bigl(V\Lambda^{\#}_{I,J}(Z,Y)\bigr)_h .
\]
The \emph{kernel-compressed affine rank-pivot ideal} is
\[
\begin{aligned}
        \mathcal K^{\rm aff}_{X,I,J,h}:=
        \bigl(&\langle \operatorname{coeff}_X(P_D-L_{\Lambda^{\#}_{I,J}}Q)\rangle
        +\langle E_X^{\rm hom}(\Lambda^{\#}_{I,J})\rangle \\
        &+\langle \Delta: \Delta\in\mathcal M_{s+1}(M_A)\rangle\bigr):
        \bigl(\Delta_{I,J}\,\Delta_X^{\rm hom}(\Lambda^{\#}_{I,J})\,
              (\Lambda^{\#}_{I,J})_j\,\Pi_h^{\#}\bigr)^\infty
\end{aligned}
\]
in the ring $F[Z,Y,(q_r)]$.  Here
$L_{\Lambda^{\#}_{I,J}}$ means $L_\lambda$ after substituting
$\lambda=\Lambda^{\#}_{I,J}(Z,Y)$.

\begin{theorem}[kernel-compressed affine eliminant certificate]
\label{thm:kernel-compressed-affine-eliminant}
Every finite affine bad slope represented by the rank-pivot chart
$(X,I,J,h)$ is in the projection of
$V(\mathcal K^{\rm aff}_{X,I,J,h})$ to the $Z$-line.  Consequently, if
\[
        Q_{X,I,J,h}(Z)\in
        \mathcal K^{\rm aff}_{X,I,J,h}\cap F[Z]
\]
is nonzero, then the chart contributes at most
\[
        \deg\gcd(Q_{X,I,J,h}(Z),Z^q-Z)
\]
finite slopes over $F_q$.
\end{theorem}

\begin{proof}
Let $z$ be a bad slope in the chart, and let $\lambda$ be a nonzero scalar
multiple of the corresponding locator vector.  Since the pivot minor is nonzero,
put $Y_c=\lambda_c/\Delta_{I,J}(z)$ for $c\in C_J$.  Cramer's rule and the pivot
row equations give
\[
        \Lambda^{\#}_{I,J}(z,Y)=\lambda .
\]
The locator is split, lies in the chart $X$, has nonzero leading coefficient,
and has the chosen nonzero noncontainment coordinate.  The larger minors vanish
on the exact-rank stratum.  Thus the point survives the displayed saturation and
lies in the compressed incidence.  Every element of the elimination ideal must
vanish on its $Z$-coordinate, and the finite-field count is obtained by
intersecting with $Z^q-Z$.
\end{proof}

\begin{remark}[why the compression is useful]
The original chart has $j$ locator variables, one parameter variable, and the
split equations.  The compressed chart has only $j+1-s$ free kernel coordinates
before imposing the leading-coordinate and chart opens.  In the overdetermined
rank-drop case this is often a large saving.  In the genuinely underdetermined
near-capacity case the saving may be small, but the construction still separates
three causes of difficulty: many split locators, a dominant parameter projection,
or a high-degree finite eliminant.
\end{remark}

The projective and curve variants are identical after replacing the parameter
ring.  For projective lines, replace $M_A(Z)$ by
$M_A^{\rm proj}(Z_0,Z_1)$, use homogeneous pivot minors, and take
homogeneous eliminants in $F[Z_0,Z_1]$.  On the infinity patch $[0:1]$, this is
just the constant matrix $V$ with equations $B(\ell)=0$ and an open pivot
$A_h\ne0$.  For a degree-$d$ finite power curve, replace $M_A(Z)$ by
\[
        M_A^{(d)}(\Gamma)=\sum_{i=0}^{d}\Gamma^i
        H_{t,j}(\Syn(f_i))
\]
and use coefficient pivots $(V_i)_h\ne0$.

\begin{corollary}[kernel-compressed projective and curve eliminants]
\label{cor:kernel-compressed-projective-curve-eliminants}
The projective rank-pivot chart obtained from the compressed construction has the
same parameter projection as the uncompressed rank-pivot chart on the declared
opens.  A nonzero homogeneous eliminant $Q(Z_0,Z_1)$ contributes at most
\[
        \deg\gcd(Q(1,Z),Z^q-Z)+\mathbf 1_{Q(0,1)=0}
\]
projective parameters over $F_q$.  For a finite degree-$d$ curve chart, a nonzero
eliminant $Q(\Gamma)$ contributes at most
\[
        \deg\gcd(Q(\Gamma),\Gamma^q-\Gamma)
\]
parameters.
\end{corollary}

\begin{proof}
The same Cramer parametrization applies to the projective pencil and to the
curve pencil.  The root-count statements are the finite-field root counts for a
homogeneous polynomial on $\mathbb P^1(F_q)$ and a univariate polynomial on
$\mathbb A^1(F_q)$.
\end{proof}

\paragraph{Exact residual support-image lcm ledger.}
The compressed eliminant is a constructive way to find a small parameter
polynomial.  There is also an exact, definition-level residual image ledger that
works for any finite co-support family, quotient or not.  This is the clean
object against which scanner-produced eliminants should be audited.

Let $\mathcal T$ be any family of co-supports $T\subseteq D$ of size $j=n-A$.
For a line $f+Zg$, put
\[
        A_T:=H_{t,j}(\Syn(f))\boldsymbol\ell_T,
        \qquad
        B_T:=H_{t,j}(\Syn(g))\boldsymbol\ell_T .
\]
Define the finite affine fixed-support polynomial
\[
        g_T^{\rm aff}(Z):=
        \begin{cases}
        \operatorname{monic}\gcd\{(A_T)_h+Z(B_T)_h:0\le h<t\},
             & (A_T,B_T)\ne(0,0),\\
        1,   & (A_T,B_T)=(0,0),
        \end{cases}
\]
with the convention that the gcd of constants with no common root is $1$.  Thus
$g_T^{\rm aff}$ has degree either $0$ or $1$.  The contained case
$(A_T,B_T)=(0,0)$ is assigned the polynomial $1$, because it contributes no
support-wise bad slope.
Set
\[
        R_{\mathcal T}^{\rm aff}(Z):=
        \operatorname{rad}\operatorname{lcm}_{T\in\mathcal T}g_T^{\rm aff}(Z).
\]
For projective lines, define $g_T^{\rm proj}(Z_0,Z_1)$ as the monic homogeneous
gcd of the linear forms
\[
        Z_0(A_T)_h+Z_1(B_T)_h,
        \qquad 0\le h<t,
\]
again setting it equal to $1$ in the contained case, and put
\[
        R_{\mathcal T}^{\rm proj}:=
        \operatorname{rad}\operatorname{lcm}_{T\in\mathcal T}g_T^{\rm proj}.
\]
For a finite degree-$d$ curve, use the vector polynomials
\[
        \sum_{i=0}^{d}\Gamma^i(V_{i,T})_h
\]
and define $R_{\mathcal T}^{(d)}(\Gamma)$ by the same radical lcm rule, with the
all-zero coefficient tuple treated as contained.

\begin{theorem}[exact arbitrary-residual image ledger]
\label{thm:exact-arbitrary-residual-image-ledger}
Let $\mathcal T$ be any finite co-support family.  The exact finite affine line
parameters witnessed by $\mathcal T$ are the roots in $F$ of
$R_{\mathcal T}^{\rm aff}$.  Hence their number is
\[
        \deg\gcd(R_{\mathcal T}^{\rm aff}(Z),Z^q-Z).
\]
The exact projective parameters witnessed by $\mathcal T$ are the projective
zeros of $R_{\mathcal T}^{\rm proj}$, counted over $F_q$ by
\[
        \deg\gcd(R_{\mathcal T}^{\rm proj}(1,Z),Z^q-Z)
        +\mathbf 1_{R_{\mathcal T}^{\rm proj}(0,1)=0}.
\]
For finite degree-$d$ curves, the exact finite parameter count is
\[
        \deg\gcd(R_{\mathcal T}^{(d)}(\Gamma),\Gamma^q-\Gamma).
\]
The same ledgers upper-bound no-loss CA parameters witnessed by the same
co-support family.
\end{theorem}

\begin{proof}
For a fixed co-support $T$, the syndrome-locator recurrence says that the line
point is explained on $D\setminus T$ precisely when
$A_T+ZB_T=0$.  If $A_T=B_T=0$, then both coefficient words are explained on the
same support and the fixed support is contained, so it is correctly charged as
$1$.  Otherwise the common roots of the scalar coordinates are exactly the
finite bad slopes contributed by $T$, and this common-root set is represented by
$g_T^{\rm aff}$.  Taking the radical lcm over $T\in\mathcal T$ coalesces duplicate
slopes.  The projective and curve statements are identical, using the homogeneous
or degree-$d$ coordinate polynomials.  No-loss CA badness is stronger than
existence of an explaining support with noncontainment, so the same image set is
an upper ledger.
\end{proof}


\begin{corollary}[closed-ball residual exact image after paid-branch removal]
\label{cor:residual-exact-image-after-paid-removal}
Fix a closed-ball agreement threshold $a$.  For each exact agreement
$A\in\{a,a+1,\ldots,n\}$, let $\mathcal T_{{\rm all},A}$ be the family of
co-supports of size $n-A$, let $\mathcal T_{{\rm paid},A}\subseteq
\mathcal T_{{\rm all},A}$ be a certified paid support family, and put
\[
        \mathcal T_{{\rm res},A}
        :=\mathcal T_{{\rm all},A}\setminus\mathcal T_{{\rm paid},A}.
\]
Let $R_{{\rm res},A}^{\rm aff}$ be the polynomial of
\cref{thm:exact-arbitrary-residual-image-ledger} for the fixed-size family
$\mathcal T_{{\rm res},A}$.  Define the closed-ball residual polynomial
\[
        R_{{\rm res},\ge a}^{\rm aff}(Z)
        :=\operatorname{rad}\operatorname{lcm}_{A=a}^{n}
          R_{{\rm res},A}^{\rm aff}(Z),
        \qquad
        R_{{\rm res},\ge a,F}^{\rm aff}(Z)
        :=\gcd\bigl(R_{{\rm res},\ge a}^{\rm aff}(Z),Z^q-Z\bigr).
\]
Then the exact finite affine residual numerator after paid-support removal is
\[
        \deg R_{{\rm res},\ge a,F}^{\rm aff}.
\]
The projective and finite-curve residual numerators are obtained by the same
exact-agreement radical-lcm construction, using
$R_{{\rm res},A}^{\rm proj}$ and $R_{{\rm res},A}^{(d)}$ respectively, followed by
the finite-field root count of \cref{thm:exact-arbitrary-residual-image-ledger}.
If a scanner also prints an exact paid root polynomial $P_{\rm paid,F}$ for the
paid branch, the total distinct finite affine numerator is
\[
        \deg\operatorname{rad}\operatorname{lcm}
        \bigl(P_{\rm paid,F},R_{{\rm res},\ge a,F}^{\rm aff}\bigr),
\]
not the sum of the paid and residual degrees.
\end{corollary}

\begin{proof}
Apply \cref{thm:exact-arbitrary-residual-image-ledger} separately at each exact
agreement size $A$, where the co-support size $j=n-A$ and the Hankel window
height $t=A-k$ are fixed.  The radical lcm over $A$ is exactly union and duplicate
coalescing across exact agreement levels.  Intersecting with $Z^q-Z$ counts only
$F$-rational finite slopes.  The final displayed lcm performs the same duplicate
coalescing between paid roots and residual roots.
\end{proof}

\paragraph{Elimination trichotomy for residual charts.}
The exact image lcm exists because split locators are a finite set.  A compressed
rank-pivot certificate is useful when it finds a much smaller polynomial without
enumerating all residual supports.  The following trichotomy is the scanner
contract.

\begin{proposition}[compressed-chart elimination trichotomy]
\label{prop:compressed-chart-elimination-trichotomy}
For a compressed finite affine residual chart over the algebraic closure, exactly
one of the following occurs:
\begin{enumerate}[label=\textup{(\roman*)},leftmargin=2.2em]
\item the saturated compressed ideal is the unit ideal, and the chart is empty;
\item the saturated compressed ideal is proper and the elimination ideal in $F[Z]$ contains a nonzero polynomial, and its
squarefree generator gives a finite parameter envelope containing the exact
support-image lcm of that chart;
\item the saturated compressed ideal is proper and the elimination ideal in $F[Z]$ is zero, equivalently the projection of the
relaxed chart to the $Z$-line is Zariski dense.
\end{enumerate}
If the full homogeneous split equations are present, case \textup{(iii)} can occur
only through a positive-dimensional relaxation or through a contained component
not removed by the noncontainment saturation; an exact split, noncontained chart
has finite parameter image.
\end{proposition}

\begin{proof}
The first three alternatives are the standard elimination theorem for the
projection of an affine variety to a line.  If the full split equations are
present and the leading coordinate is nonzero, the locator coordinate belongs to
the finite set of scalar multiples of the finitely many degree-$j$ locators
supported on $D$.  After projectivizing the locator scale, each fixed support
contributes at most one finite noncontained slope by
\cref{lem:one-support-one-line}.  Hence the exact noncontained split image is a
finite subset of the parameter line.  A dense projection can therefore only come
from a relaxation of the split condition or from an unremoved contained component
on which every parameter is explained but no noncontainment pivot should survive.
\end{proof}

\begin{remark}[compressed-residual status]
The singular bucket has now been reduced to two concrete certificate tasks.  The
first is computational-algebraic: use the kernel-compressed ideals to find small
univariate or homogeneous eliminants.  The second is enumerative but exact: when
compression is not yet sharp enough, print the residual support-image lcm for a
finite residual family and coalesce it with the paid quotient/tangent roots.
Neither task proves the final aperiodic packing theorem by itself, but both
replace the phrase ``singular residual bucket'' by checkable objects with exact
root-count semantics.
\end{remark}

\subsubsection{Fallback and final certificate}

The eliminant certificates above are strongest when the projection to the
parameter line is finite.  A failed eliminant search should not be discarded; it
should be split further or recorded as a singular bucket.  The following fallback
is safe but usually too weak for a final $2^{-128}$ threshold.

\begin{proposition}[dimension-degree fallback]
\label{prop:dimension-degree-fallback}
Let $Y\subseteq\mathbb A^N_F$ be an affine variety of degree at most $\Delta$ and
Krull dimension at most $e$.  Then
\[
        |Y(F)|\le \Delta |F|^e .
\]
Consequently, if a saturated Hankel incidence chart for line or curve parameters
has degree $\Delta_X$ and dimension $e_X$, then the number of bad parameters in
that chart is at most $\Delta_X|F|^{e_X}$.
\end{proposition}

\begin{proof}
This is the standard affine degree point-count bound, proved by induction on
dimension using generic affine hyperplane sections.  The parameter image is no
larger than the incidence set.
\end{proof}

For a threshold statement, use finitely many charts and agreement levels.  Let
$\mathfrak X_A$ be a chart cover of the declared aperiodic branch at exact
agreement $A$.  If each finite affine pivot chart has a univariate eliminant, set
\[
        B_{\rm ap}^{\rm aff,elim}(a)
        :=\sum_{A=a}^n\sum_{X\in\mathfrak X_A}
          \sum_{h=0}^{t_A-1}\deg Q_{X,h} .
\]
For projective slopes add the certified infinity contribution, and for degree-$d$
curves replace the inner sum by the coefficient-pivot sum
$\sum_{i=0}^d\sum_h\deg Q_{X,i,h}$.  Exact root-union tables may replace these
safe sums whenever duplicate parameters across charts are certified.

\begin{definition}[aperiodic Hankel-packing certificate]
\label{def:aperiodic-hankel-certificate}
An aperiodic Hankel-packing certificate for a row, a line or curve family, and an
agreement threshold $a$ consists of:
\begin{enumerate}[label=\textup{(\roman*)},leftmargin=2.2em]
\item the tangent/common-code-line and quotient-image ledgers that have already
been removed;
\item for each exact agreement $A\in\{a,\ldots,n\}$, either a regular minor
certificate $\Delta_A$ or a declaration that the exact-agreement bucket is
Hankel-singular;
\item for every singular or near-capacity bucket, a finite constructible locator
chart cover $\mathfrak X_A$ as in \cref{def:aperiodic-locator-chart};
\item for affine lines, a list of pivot eliminants $Q_{X,h}(Z)$, or a failed-pivot
record explaining which residual locus remains unsplit;
\item for projective slopes, the finite-patch eliminants plus the infinity-chart
empty/nonempty certificate;
\item for curves, coefficient-pivot eliminants $Q_{X,i,h}(\Gamma)$;
\item a root-union table or a safe degree-sum table producing the declared
aperiodic numerator $B_{\rm ap}(a)$;
\item if any chart uses the dimension-degree fallback, its degree, dimension, and
an explicit statement that the resulting bound is only a fallback, not a sharp
threshold certificate.
\end{enumerate}
The verifier checks locator divisibility by $P_D$, chart coverage of the declared
aperiodic complement, nonzero pivot tests, ideal membership of eliminants after
saturation, and final root-count arithmetic.
\end{definition}

\begin{remark}[contributor workflow]
\label{rem:aperiodic-contributor-workflow}
For each row and exact agreement $A$, contributors should try the following in
order.  First test the overdetermined regular minors.  If a nonzero minor is
found, record its root set and remove roots already paid by tangent or quotient
ledgers.  If all maximal minors vanish identically, move the instance to the
singular bucket.  On the singular or near-capacity bucket, build locator charts,
then split by affine pivots $B_h\ne0$; for projective samplers also test the
infinity chart $B=0$, $A\ne0$; for curves split by coefficient pivots
$(V_i)_h\ne0$.  A chart is complete only when it has either a nonzero eliminant,
a certified empty locus, or a named residual obstruction.  Residual obstructions
should be labelled as quotient-periodic, tangent/common-code-line,
subfield/extension exceptional, or candidate new obstruction.
\end{remark}

\begin{remark}[status]
\label{rem:aperiodic-certificate-status}
\Cref{thm:regular-closed-ball-hankel-packing,thm:aperiodic-affine-eliminant,cor:projective-eliminant-certificate,thm:aperiodic-curve-eliminant}
are proved certificate theorems.  They do not prove that all rows admit small
aperiodic eliminants.  They make the remaining MCA safe-side task concrete: after
tangent and quotient branches are removed, construct small eliminants for the
pivot charts, or identify the singular chart that prevents such an eliminant.
\end{remark}


\begin{lemma}[interleaving transfer]\label{lem:inter}
For every linear code $C\subseteq\F^n$, every $s\ge1$ and every $\delta\in(0,1)$:
\[
\eca\bigl(C^{\equiv s},\delta\bigr)\ \ge\ \eca(C,\delta),
\qquad
\emca\bigl(C^{\equiv s},\delta\bigr)\ \ge\ \emca(C,\delta).
\]
\end{lemma}

\begin{proof}
For $h\in\F^n$ write $h^{\Delta}:=(h,\dots,h)\in(\F^n)^s$. For any codeword tuple $\mathbf c=(c^{(1)},\dots,c^{(s)})\in C^{\equiv s}$, the columns where $h^\Delta$ and $\mathbf c$ agree are $\bigcap_i\{j:c^{(i)}_j=h_j\}\subseteq\{j:c^{(1)}_j=h_j\}$, so $\dist_s(h^\Delta,C^{\equiv s})\ge\dist(h,C)$; taking $c^{(1)}=\dots=c^{(s)}$ gives equality:
$\dist_s(h^\Delta,C^{\equiv s})=\dist(h,C)$.

Fix $(f_1,f_2)$ attaining $\eca(C,\delta)$ (resp.\ $\emca$) and consider the pair $(f_1^\Delta,f_2^\Delta)$ for $C^{\equiv s}$, with the slope $\gamma$ acting row-wise, so $f_1^\Delta+\gamma f_2^\Delta=(f_1+\gamma f_2)^\Delta$. By the displayed equality, $\dist_s\bigl((f_1+\gamma f_2)^\Delta,C^{\equiv s}\bigr)\le\delta$ iff $\dist(f_1+\gamma f_2,C)\le\delta$. For the correlated condition, the pair $(f_1^\Delta,f_2^\Delta)$ viewed as a $2s$-row interleaved word has rows $f_1$ ($s$ times) and $f_2$ ($s$ times); by the same column argument its distance to $(C^{\equiv s})^{\equiv2}=C^{\equiv2s}$ equals $\dist_2((f_1,f_2),C^{\equiv2})$. For the MCA condition, a common set $S$ explains $f_1^\Delta$ and $f_2^\Delta$ iff it explains $f_1$ and $f_2$. Hence $\gamma$ is CA-bad (resp.\ MCA-bad) for $(f_1^\Delta,f_2^\Delta)$ iff it is so for $(f_1,f_2)$, and the maxima over pairs can only grow.
\end{proof}

\section{The main theorem}\label{sec:main}

\begin{theorem}[universal cap]\label{thm:main}
Let $\B\subseteq\F$ be finite fields, let $q=|\F|$, let $D\subseteq\B^\times$ be a multiplicative coset of order $n$, and let $k$ have $\rho=k/n\in(0,1)$; set $C:=\RS[\F,D,k]$. Let $N\mid n$ satisfy $a:=n/N\mid k$ and $(1-\rho)N\ge3$. If
\begin{equation}\label{eq:hyp}
        \binom{N}{\rho N+2}
        \ge
        |\B|\Bigl(\frac qk+1\Bigr),
\end{equation}
then, writing
\[
        \delta_N:=1-\rho-\frac2N,
\]
we have
\[
\eca(C,\delta)
>
\frac1{2k}\Bigl(1-\frac nq\Bigr)
\]
for every integer-admissible radius
\[
        \delta_N\le \delta \le 1-\rho-\frac1n .
\]
In particular,
\[
        \emca(C,\delta)
        >
        \frac1{2k}\Bigl(1-\frac nq\Bigr)
\]
for every
\[
        \delta_N\le \delta < 1-\rho .
\]
Consequently,
\[
        \delta^*_C(\varepsilon^*)
        \le
        1-\rho-\frac2N
\]
for every
\[
        \varepsilon^*
        \le
        \frac1{2k}\Bigl(1-\frac nq\Bigr).
\]
\end{theorem}

\begin{proof}
Note first that $\rho N=k/a\in\mathbb Z$, $\ell_2=\rho N+2\le N-1$ by $(1-\rho)N\ge3$, and $\delta_N=1-\rho-2/N>0$ for the same reason. Also $k+1\le n$.

Fix
\[
        \delta\in[\delta_N,\,1-\rho-\tfrac1n].
\]
Then
\[
        \lfloor \delta n\rfloor \le n-k-1,
\]
so \cref{thm:A} is applicable.

Suppose toward a contradiction that
\[
        \eca(C,\delta)
        \le
        \frac1{2k}\Bigl(1-\frac nq\Bigr).
\]
Apply \cref{thm:A} with $\eta=\frac12$. Its hypothesis holds, so
\[
        \Lst(C^+,\delta)
        \le
        \Bigl\lceil 2q\cdot \eca(C,\delta)\Bigr\rceil
        \le
        \Bigl\lceil\frac{q-n}{k}\Bigr\rceil
        <
        \frac{q-n}{k}+1
        \le
        \frac qk+1 .
\]
On the other hand, by \cref{lem:fiber}(ii) and monotonicity of lists in the radius, since $\delta\ge\delta_N$,
\[
        \Lst(C^+,\delta)
        \ge
        \bigl|\Lst(C^+,\delta_N,u_z)\bigr|
        \ge
        \binom{N}{\ell_2}\Big/|\B|
        \ge
        \frac qk+1
\]
by \eqref{eq:hyp}. This is a contradiction. Hence
\[
        \eca(C,\delta)
        >
        \frac1{2k}\Bigl(1-\frac nq\Bigr)
\]
for every $\delta\in[\delta_N,1-\rho-1/n]$.

In particular, the above holds at $\delta=\delta_N$. By \cref{fact:chain},
\[
        \emca(C,\delta_N)
        \ge
        \eca(C,\delta_N)
        >
        \frac1{2k}\Bigl(1-\frac nq\Bigr).
\]
By support-wise MCA monotonicity, \cref{lem:mca-monotone}, the same lower bound for $\emca(C,\delta)$ holds for every
\[
        \delta\in[\delta_N,1-\rho).
\]
Thus no sub-capacity radius
\[
        \delta\ge 1-\rho-\frac2N
\]
is $\varepsilon^*$-admissible whenever
\[
        \varepsilon^*
        \le
        \frac1{2k}\Bigl(1-\frac nq\Bigr),
\]
which proves the claimed bound on $\delta^*_C(\varepsilon^*)$.
\end{proof}

\begin{remark}[constants]\label{rem:eta}
Choosing $\eta=1-\theta$ in \cref{thm:A} improves the conclusion to $\eca>\frac{1-\theta}k(1-\frac nq)$ at the price of strengthening \eqref{eq:hyp} to $\binom N{\rho N+2}\ge|\B|\bigl(\frac q{\theta k}+1\bigr)$; given the hundreds of bits of slack in \eqref{eq:hyp} throughout the challenge envelope (\cref{tab:numerics}), the error constant is effectively $(1-o(1))/k$. We fix $\eta=\frac12$ for cleanliness.
\end{remark}

\begin{remark}[scope]\label{rem:scope}
No smoothness is assumed beyond the existence of a single divisor $N\mid n$ with $a=n/N\mid k$ and \eqref{eq:hyp}. The theorem applies verbatim to $2$-adic, mixed-radix, or any other FFT-friendly multiplicative domain, and --- via the subfield pigeonhole --- to any extension field $\F\supseteq\B$ with $D\subseteq\B$.

The deep-point conversion is an integer-radius theorem with the condition $\lfloor\delta n\rfloor\le n-k-1$. Therefore the no-loss CA lower bound is stated for
\[
        \delta \le 1-\rho-\frac1n .
\]
The MCA threshold cap is stronger in the range needed for the Proximity Prize: once the support-wise MCA lower bound is proved at
\[
        \delta_N=1-\rho-\frac2N,
\]
monotonicity of support-wise MCA extends it to every larger $\delta<1-\rho$.
\end{remark}

\section{Consequences for the grand MCA challenge}\label{sec:grand}

\begin{corollary}[universal field-size cap for the challenge envelope]\label{cor:grand}
Let $\rho\in\{\frac12,\frac14,\frac18,\frac1{16}\}$ and set
\[
N_\rho:=1024\quad(\rho\in\{\tfrac12,\tfrac14,\tfrac18\}),\qquad N_{1/16}:=2048 .
\]
Let $\F$ be any finite field with $q<2^{256}$, let $\B\subseteq\F$ be any subfield, let $D\subseteq\B^\times$ be a multiplicative coset of order $n$ with $N_\rho\mid n$, and let $k=\rho n\le2^{40}$. Then $C=\RS[\F,D,k]$ satisfies
\[
\emca(C,\delta)
>
\frac1{2k}\Bigl(1-\frac nq\Bigr)
\ge
2^{-86}
\gg
2^{-128}
\]
for every
\[
        \delta\in
        \Bigl[1-\rho-\frac2{N_\rho},\,1-\rho\Bigr).
\]
Moreover, the same lower bound holds for $\eca(C,\delta)$ throughout the integer-admissible subinterval
\[
        \delta\in
        \Bigl[1-\rho-\frac2{N_\rho},\,1-\rho-\frac1n\Bigr].
\]
If $q\ge2n$, the lower bound improves to $\ge2^{-42}$. In particular,
\[
\delta^*_C(2^{-128})\;\le\;1-\rho-2^{-9}\quad(\rho\in\{\tfrac12,\tfrac14,\tfrac18\}),
\qquad
\delta^*_C(2^{-128})\;\le\;1-\rho-2^{-10}\quad(\rho=\tfrac1{16}).
\]
\end{corollary}

\begin{proof}
We verify the hypotheses of \cref{thm:main} with $N=N_\rho$. Divisibility: $a=n/N_\rho$ divides $k$ because $k/a=\rho N_\rho\in\mathbb Z$ (equal to $512,256,128,128$ respectively); $(1-\rho)N_\rho\ge128\ge3$. For \eqref{eq:hyp}, the entropy bound $\binom N\ell\ge2^{N\Hb(\ell/N)}/(N+1)$ together with the conservative evaluations of $\Hb$ recorded in \cref{tab:numerics} gives $\log_2\binom{N_\rho}{\rho N_\rho+2}\ge552$ in all four cases, while
\[
|\B|\Bigl(\frac qk+1\Bigr)\;\le\;q\Bigl(\frac qk+1\Bigr)\;\le\;q^2+q\;<\;2^{512}+2^{256}\;<\;2^{513}\;\le\;2^{552},
\]
so \eqref{eq:hyp} holds with at least $39$ bits to spare (and over $500$ bits at rate $\frac12$). Since $D\subseteq\F^\times$, we have $n\le q-1$; for fixed $n$, the function $1-n/q$ is minimized at $q=n+1$, hence $1-n/q\ge1/(n+1)$. Also $n=k/\rho\le16k\le2^{44}$. Therefore
\[
\frac1{2k}\Bigl(1-\frac nq\Bigr)\ge\frac1{2k(n+1)}>2^{-86}.
\]
If $q\ge2n$ then $1-n/q\ge1/2$, giving the sharper $1/(4k)\ge2^{-42}$. The gap is $2/N_\rho=2^{-9}$ for $\rho\in\{1/2,1/4,1/8\}$ and $2/N_\rho=2^{-10}$ for $\rho=1/16$. \Cref{thm:main} gives the MCA lower bound throughout the stated sub-capacity interval, and gives the CA lower bound on the integer-admissible subinterval ending at $1-\rho-1/n$.
\end{proof}

\begin{table}[ht]
\centering
\begin{tabular}{ccccccc}
\toprule
$\rho$ & $N_\rho$ & $\ell_2=\rho N_\rho+2$ & $\Hb(\ell_2/N_\rho)\ge$ & $\log_2\binom{N_\rho}{\ell_2}\ge$ & needed & gap $2/N_\rho$\\
\midrule
$1/2$  & $1024$ & $514$ & $0.99998$ & $1013$ & $513$ & $2^{-9}$\\
$1/4$  & $1024$ & $258$ & $0.8143$  & $823$  & $513$ & $2^{-9}$\\
$1/8$  & $1024$ & $130$ & $0.5490$  & $552$  & $513$ & $2^{-9}$\\
$1/16$ & $2048$ & $130$ & $0.34105$ & $687$  & $513$ & $2^{-10}$\\
\bottomrule
\end{tabular}
\caption{Numerics for \cref{cor:grand}: $\log_2\binom{N}{\ell}\ge N\Hb(\ell/N)-\log_2(N+1)$, rounded down. The ``needed'' column is the universal upper bound $\log_2\bigl(|\B|(q/k+1)\bigr)<513$ valid throughout the envelope $q<2^{256}$, $k\ge1$.}
\label{tab:numerics}
\end{table}

\begin{remark}[hypothesis $N_\rho\mid n$]
For the $2$-adic domains of the challenge this holds whenever $n\ge2^{10}$ (resp.\ $2^{11}$), i.e.\ for all but very small instances. For mixed-radix smooth $n$, replace $N_\rho$ by any divisor $N\mid n$ of comparable size with $\rho N\in\mathbb Z$; the entropy count in \cref{tab:numerics} is insensitive to the exact value of $N$ at fixed $\log_2 N$. The divisibility-free fiber lemma removes the integrality requirement $\rho N\in\mathbb Z$ altogether: \cref{lem:phi-fiber}, \cref{cor:mixed-radix} caps every mixed-radix row possessing \emph{any} divisor in the printed dyadic windows.
\end{remark}

\begin{remark}[two-tier answer to the grand challenge]\label{rem:twotier}
Combining \cref{cor:grand} with the grand-challenge cap of \cite[Prop.\ ``Grand-challenge cap'']{Cho26b} (gap $\frac1{64}$, prime fields up to $\approx2^{150}$, via value-set coverage) yields the current state of the negative side for smooth multiplicative domains:
\[
\delta^*_C(2^{-128})\;\le\;
\begin{cases}
1-\rho-\frac1{64}, & \F\ \text{prime},\ q\lesssim2^{150},\\[2pt]
1-\rho-2^{-9}, & \rho\in\{\frac12,\frac14,\frac18\},\ \text{every}\ \F,\ q<2^{256},\\[2pt]
1-\rho-2^{-10}, & \rho=\frac1{16},\ \text{every}\ \F,\ q<2^{256}.
\end{cases}
\]
Together with the positive results collected in \cite[Table~1]{ABF26}, this pins the grand MCA threshold between the Johnson region and the explicit capacity-edge caps above, throughout the stated challenge envelope. The second and third tiers supersede the conditional cap of \cite[Thm.\ ``Universal cap'']{Cho26b} ($\delta^*<1-\rho-2^{-11}$ under $|\F|\ge2n$), which was the previous state in this band. Circle rows --- line rounds and bivariate circle codes --- now carry an analogous negative band; see \cref{cor:circle-grand} and \cref{rem:circle-transfer}.
\end{remark}

\begin{corollary}[deployed parameters: KoalaBear sextic]\label{cor:deployed}
Let $\B=\F_p$ with $p=2^{31}-2^{24}+1$, let $\F=\F_{p^6}$ (so $q=p^6\approx2^{185.93}$), let $D\subseteq\B^\times$ be the subgroup of order $n=2^{21}$, and let $k=2^{20}$, $\rho=\frac12$, as in \cite[\S6.3]{ABF26}. Then
\[
\emca\bigl(\RS[\F_{p^6},D,2^{20}],\,\delta\bigr)
>
\bigl(1-2^{-164}\bigr)\,2^{-21}
>
2^{-22}
\]
for every
\[
        \delta\in\Bigl[\frac12-2^{-7},\,\frac12\Bigr).
\]
Moreover, the same lower bound holds for $\eca$ throughout the integer-admissible subinterval
\[
        \delta\in
        \Bigl[\frac12-2^{-7},\,\frac12-2^{-21}\Bigr).
\]
\end{corollary}

\begin{proof}
Apply \cref{thm:main} with $N=256$, $a=2^{13}\mid k$, $(1-\rho)N=128\ge3$, gap $2/N=2^{-7}$. Hypothesis \eqref{eq:hyp}: $\log_2\binom{256}{130}\ge256\cdot\Hb(130/256)-\log_2257\ge256\cdot0.99982-8.01>247$, while $|\B|(q/k+1)<2^{31}\,(2^{166}+1)<2^{198}$. The error bound is $\frac1{2k}(1-n/q)$ with $n/q=2^{21}/p^6<2^{-164}$. \Cref{thm:main} gives the stated MCA bound on $[\frac12-2^{-7},\frac12)$ and the stated CA bound on $[\frac12-2^{-7},\frac12-2^{-21})$.
\end{proof}

\begin{corollary}[all interleaved rows of {\cite[Tables 2--3]{ABF26}}]\label{cor:rows}
In the setting of \cref{cor:deployed}, for each $s=2^j$ with $0\le j\le12$ let $C_s:=\RS[\F_{p^6},D_s,k_s]$ with $|D_s|=n_s=2^{21}/s$ and $k_s=2^{20}/s$ (rate $\frac12$, $s\,n_s=2^{21}$ as deployed). Then for every $\delta\in\bigl[\frac12-2^{-7},\,\frac12\bigr)$,
\[
\emca\bigl(C_s^{\equiv s},\delta\bigr)
>
\frac s{2^{21}}\bigl(1-2^{-164}\bigr)
\ge
2^{-21}\bigl(1-2^{-164}\bigr).
\]
Moreover, the same lower bound holds for $\eca(C_s^{\equiv s},\delta)$ on the integer-admissible subinterval
\[
        \delta\in
        \Bigl[\frac12-2^{-7},\,\frac12-\frac1{n_s}\Bigr).
\]
Every row of the deployed parameter family therefore fails the $2^{-128}$ target at gap $2^{-7}$ below capacity.
\end{corollary}

\begin{proof}
For each $s\le2^{12}$: $256\mid n_s$ (as $n_s\ge2^9$), $a_s=n_s/256=2^{13}/s$ divides $k_s=2^{20}/s$, and $(1-\rho)\cdot256=128\ge3$. Hypothesis \eqref{eq:hyp} for $C_s$: $|\B|(q/k_s+1)<2^{31}(2^{166}\,s+1)<2^{198}\,s\le2^{210}\le2^{247}\le\binom{256}{130}$. \Cref{thm:main} gives the MCA lower bound for $C_s$ throughout $[\frac12-2^{-7},\frac12)$, and the CA lower bound throughout $[\frac12-2^{-7},\frac12-1/n_s]$. \Cref{lem:inter} transfers the corresponding lower bounds to the $s$-interleaved code.
\end{proof}

\begin{proposition}[independent slacked variant via \cref{thm:B}]\label{prop:slacked}
In the setting of \cref{thm:main}, assume instead $(1-\rho)N\ge2$, $n\ge3N$, and
\[
\binom{N}{\rho N+1}\;\ge\;|\B|\cdot q .
\]
Then
\[
\eca\Bigl(C,\ 1-\rho-\frac1N+\frac2n,\ 1-\rho-\frac1n\Bigr)\;\ge\;\frac1{2n}.
\]
In particular, with $N=N_\rho$ as in \cref{cor:grand} and $n\ge3N_\rho$, every challenge-envelope instantiation satisfies this slacked CA lower bound, and $\frac1{2n}\ge2^{-45}\gg2^{-128}$.
\end{proposition}

\begin{proof}
Set $\delta_0:=1-\rho-\frac1N$. By $(1-\rho)N\ge2$, $\delta_0\in(0,1-\rho)$. By $n\ge3N$, we also have $\delta_0+2/n\le1-\rho-1/n$, so the parameters in \cref{thm:B} are admissible. By \cref{lem:fiber}(i), some $u_z$ carries at least $\binom N{\rho N+1}/|\B|\ge q$ codewords of $C$ within radius $\delta_0$, so $\Lst(C,\delta_0)\ge q$. The contrapositive of \cref{thm:B} gives
\[
\eca\Bigl(C,\delta_0+\frac2n,1-\rho-\frac1n\Bigr)\ge\frac1{2n}.
\]
For the numerical claim, the worst case of $\log_2\binom{N_\rho}{\rho N_\rho+1}$ over the four rates is at $\rho=1/8$, where the entropy bound used in \cref{tab:numerics} certifies $\log_2\binom{1024}{129}\ge1024\,\Hb(129/1024)-\log_2 1025\ge549$ (the exact value is $\approx554.7$); hence $\binom{N_\rho}{\rho N_\rho+1}\ge2^{549}>2^{512}>|\B|q$ throughout $q<2^{256}$. Finally $n=k/\rho\le16k\le2^{44}$, so $1/(2n)\ge2^{-45}$.
\end{proof}

\begin{remark}
\Cref{prop:slacked} is logically independent of \cref{thm:A} and yields a weaker error ($\frac1{2n}$ versus $\frac1{2k}$) in the \emph{maximally slacked} form of correlated agreement --- the internal radius sits at capacity minus $\frac1n$ --- which is the weakest failure mode one can certify. That even this fails by $83$ orders of magnitude (in bits) above the $2^{-128}$ target gives a separate fallback route through the imported BCHKS/ABF theorem.
\end{remark}

\section{Subfield confinement and the shape of certifying lines}\label{sec:subfield}

The deployed instantiations run the protocol over an extension $\F/\B$ while keeping the domain in the base: $D\subseteq\B$. The subfield-confinement theorem of \cite[Thm.\ ``Subfield confinement'']{Cho26b} pins down where bad slopes can live for lines whose words are base-rational: all support-wise MCA-bad slopes lie in $\B$. The following lemma refines this by tracking two radii, so that it also covers the proximity-loss form of correlated agreement (\cref{def:ca}); part (b) is the statement of \cite{Cho26b}, reproved here for completeness in the present notation.

\begin{lemma}[subfield confinement; refines {\cite{Cho26b}}]\label{lem:confine}
Let $\B\subseteq\F$ be finite fields, $D\subseteq\B^\times$, $C=\RS[\F,D,k]$, and let $f,g\in\B^n\subseteq\F^n$ be $\B$-valued words. Then:
\begin{enumerate}[label=\textup{(\alph*)}]
\item for any $0\le\delta_{\mathrm{fld}}\le\delta_{\mathrm{int}}<1$, every CA-bad slope for $(f,g)$ at radii $(\delta_{\mathrm{fld}},\delta_{\mathrm{int}})$ lies in $\B$;
\item for any $\delta\in(0,1)$, every MCA-bad slope for $(f,g)$ at radius $\delta$ lies in $\B$.
\end{enumerate}
Consequently the bad-slope set of any $\B$-valued pair has density at most $|\B|/q$ in $\F$. The same holds for pairs $(\lambda f,\lambda g)$ with $\lambda\in\F^\times$ and $f,g$ $\B$-valued, by $\F$-linearity of $C$.
\end{lemma}

\begin{proof}
Fix a $\B$-basis $1=\beta_0,\beta_1,\dots,\beta_{m-1}$ of $\F$. Every word $w\in\F^n$ decomposes uniquely as $w=\sum_i\beta_i w_i$ with $w_i\in\B^n$, and every codeword $c\in\RS[\F,D,k]$ decomposes as $c=\sum_i\beta_i c_i$ with $c_i\in\RS[\B,D,k]$: write the coefficient vector of the underlying polynomial in the basis and evaluate at $D\subseteq\B$.

Let $z\in\F\setminus\B$, say $z=\sum_iz_i\beta_i$ with $z_{i^*}\ne0$ for some $i^*\ge1$. We first record the support-wise consequence. Suppose that $c\in C$ agrees with $f+zg$ on some set $S\subseteq D$. Since $f_j,g_j\in\B$, the basis components of $(f+zg)_j$ are $f_j+z_0g_j$ (component $0$) and $z_ig_j$ (components $i\ge1$). Agreement on $S$ holds componentwise, so on $S$:
\[
f+z_0g=c_0,\qquad z_{i^*}\,g=c_{i^*}.
\]
Hence $g$ agrees on $S$ with $z_{i^*}^{-1}c_{i^*}\in C$, and then $f=(f+z_0g)-z_0g$ agrees on $S$ with $c_0-z_0z_{i^*}^{-1}c_{i^*}\in C$. Thus every support on which the line point $f+zg$ is explained by $C$ also explains the pair $(f,g)$ on that same support.

For CA, if $\dist(f+zg,C)\le\delta_{\mathrm{fld}}$, choose such an $S$ with $|S|\ge(1-\delta_{\mathrm{fld}})n$; the preceding paragraph gives $\dist_2((f,g),C^{\equiv2})\le\delta_{\mathrm{fld}}\le\delta_{\mathrm{int}}$, so $z$ is not CA-bad. For MCA, the preceding paragraph rules out every possible bad support $S$ in the support-wise definition. Hence $z$ is not MCA-bad either.
\end{proof}

\begin{corollary}[certifying lines over extensions are genuinely $\F$-valued]\label{cor:Fvalued}
At the parameters of \cref{cor:deployed}, $|\B|/q=p^{-5}<2^{-154}$. Hence any pair $(f,g)$ whose CA- or MCA-bad-slope density at some radius
\[
        \delta\in
        \Bigl[\frac12-2^{-7},\,\frac12\Bigr)
\]
exceeds $2^{-154}$ contains a word that is not $\B$-valued, even after removing a common scalar factor. In particular, any pair attaining the MCA lower bound $>2^{-22}$ of \cref{cor:deployed} is genuinely $\F$-valued. On the integer-admissible subinterval
\[
        \delta\in
        \Bigl[\frac12-2^{-7},\,\frac12-2^{-21}\Bigr),
\]
the same statement applies to pairs attaining the CA lower bound. All explicit failure witnesses currently known on these domains (the locator lines of \cite{Cho26a,Cho26b} and the fiber words of \cref{lem:fiber}) are $\B$-rational and are therefore inert here: the failure certified by \cref{cor:deployed} is fiber-borne and, at present, non-constructive.
\end{corollary}

\begin{proof}
Immediate from \cref{lem:confine} and \cref{cor:deployed}: a $\B$-valued pair, up to common scalar, has CA- and MCA-bad-slope density at most $|\B|/q=p^{-5}<2^{-154}<2^{-22}$. Thus any pair witnessing the deployed lower bound cannot be of this form.
\end{proof}

\begin{remark}
\Cref{lem:confine,cor:Fvalued} sharpen the division of labor between the two grand challenges of \cite{ABF26} over deployed extensions: the \emph{list} side does not deflate ($\RS[\B,D,k]\subseteq\RS[\F,D,k]$, so all base-field list lower bounds, including \cref{lem:fiber}, persist verbatim), while every sub-$2^{-128}$ \emph{MCA} attack by explicit base-rational witnesses is killed by confinement --- and yet MCA still fails at $2^{-22}$.  The failure can also be made constructive once a large list is supplied: \cref{cor:extension-pole-deep-list-floor,cor:extension-pole-quotient-remainder-floor} allow one to choose a simple pole in $\F\setminus\B$ and print its value-bucket table. The deflation and decoupling corollaries of \cite{Cho26b} record the same division at the no-loss radius, using the same corrected rounding $p^{-5}=2^{-154.94\ldots}<2^{-154}$ as in \cref{cor:Fvalued}. Exhibiting deployed-size bucket tables for such lines is Open Problem~\ref{prob:explicit}.
\end{remark}


\section{Beyond multiplicative smoothness: map-smooth domains, Chebyshev fibers, and the circle-code cap}\label{sec:isogeny}

The fiber lemma \cref{lem:fiber} uses exactly two features of the power map $x\mapsto x^a$ on a multiplicative coset: its fibers inside the domain are complete and of uniform size $a$, and the locators $\prod_{b\in A}(X^a-b)$ are \emph{lacunary} --- expanded in powers of $X^a$, their two top terms occupy degrees $a\ell$ and $a(\ell-1)$, leaving everything from degree $a(\ell-2)$ downward to the listed codeword. This section isolates exactly these two features. \Cref{lem:phi-fiber} proves the fiber lemma for an arbitrary polynomial folding map $\varphi$ whose fibers inside the domain are complete and uniform, and simultaneously removes the divisibility hypothesis $a\mid k$ by rounding the locator length; \cref{thm:phi-cap} is the corresponding universal cap, containing \cref{thm:main} as the special case $\varphi=X^a$ with $a\mid k$. We then verify both features for the deployed isogeny-smooth family, the circle-group domains of Circle STARKs over $p=2^{31}-1$ \cite{HLP24}: twin-coset $x$-coordinate domains carry exact Chebyshev fibers at every dyadic scale (\cref{lem:cheb-fibers}), and over any challenge field containing $i$ the bivariate circle code is diagonally equivalent to a Reed--Solomon code on a torus twin coset, an equivalence preserving lists and both correlated-agreement errors exactly (\cref{lem:diag-invariance,lem:circle-rs}). The resulting universal caps (\cref{cor:circle-grand}) and deployed instantiations (\cref{cor:circle-deployed}) settle the circle-group case of the isogeny-smooth problem posed below; what remains of that problem is isolated, with a precise obstruction, in \cref{rem:ecfft-obstruction} and restated as Open Problem~\ref{prob:isogeny}.

\subsection{Map-smooth domains and a divisibility-free fiber lemma}
\label{sec:map-smooth}

\begin{definition}[map-smooth domain]\label{def:map-smooth}
Let $\B\subseteq\F$ be finite fields, let $\varphi\in\B[X]$ have degree $a\ge1$, and let $D\subseteq\B$ be a finite set of size $n$. The set $D$ is \emph{$(\varphi,a)$-smooth} if every nonempty fiber $\varphi^{-1}(w)\cap D$, $w\in\varphi(D)$, has exactly $a$ elements. In that case $Q:=\varphi(D)\subseteq\B$ has size $N:=n/a$, and $D$ is the disjoint union of the $N$ fibers $S_b:=\{x\in D:\varphi(x)=b\}$, $b\in Q$.
\end{definition}

Two remarks on scope. First, a multiplicative coset $D\subseteq\B^\times$ of order $n$ is $(X^a,a)$-smooth for every $a\mid n$: the fibers are the cosets of the order-$a$ subgroup contained in $D$'s underlying group. Second, nothing is assumed about the roots of $\varphi-w$ outside $D$: the definition constrains fibers inside $D$ only, so separability or squarefreeness of $\varphi-w$ plays no role below.

\begin{lemma}[map-smooth fiber lemma, divisibility-free]\label{lem:phi-fiber}
Let $\B\subseteq\F$ and let $D\subseteq\B$ be a $(\varphi,a)$-smooth domain of size $n$, with $Q=\varphi(D)$ and $N=n/a$ as above. Let $1\le k<n$ and $\rho:=k/n$. Put
\[
\ell_1:=\Bigl\lfloor\frac{k-1}a\Bigr\rfloor+2,\qquad
\ell_2:=\Bigl\lfloor\frac{k}a\Bigr\rfloor+2,\qquad
A_1:=a\ell_1,\qquad A_2:=a\ell_2,
\]
so that $k+a\le A_1\le k+2a-1$ and $k+a+1\le A_2\le k+2a$, with $A_1=k+a$ and $A_2=k+2a$ when $a\mid k$.
\begin{enumerate}[label=\textup{(\roman*)}]
\item If $\ell_1\le N-1$, there exists $z\in\B$ such that, with $u_z:=\bigl(\varphi(x)^{\ell_1}+z\,\varphi(x)^{\ell_1-1}\bigr)_{x\in D}\in\B^n$,
\[
\bigl|\Lst\bigl(\RS[\F,D,k],\,1-\tfrac{A_1}n,\,u_z\bigr)\bigr|\ \ge\ \binom N{\ell_1}\Big/|\B| .
\]
\item If $\ell_2\le N-1$, there exists $z\in\B$ such that, with $u_z:=\bigl(\varphi(x)^{\ell_2}+z\,\varphi(x)^{\ell_2-1}\bigr)_{x\in D}\in\B^n$,
\[
\bigl|\Lst\bigl(\RS[\F,D,k+1],\,1-\tfrac{A_2}n,\,u_z\bigr)\bigr|\ \ge\ \binom N{\ell_2}\Big/|\B| .
\]
\end{enumerate}
In both cases $u_z$ is $\B$-valued, and the listed codewords lie in $\RS[\B,D,a(\ell_1-2)+1]\subseteq\RS[\B,D,k]$ \textup(resp.\ $\RS[\B,D,a(\ell_2-2)+1]\subseteq\RS[\B,D,k+1]$\textup).
\end{lemma}

\begin{proof}
First the bracketing of $A_1,A_2$: writing $k=a\lfloor k/a\rfloor+s$ with $0\le s\le a-1$ gives $A_2=k-s+2a\in[k+a+1,k+2a]$; writing $k-1=a\lfloor(k-1)/a\rfloor+s_1$ with $0\le s_1\le a-1$ gives $A_1=k-1-s_1+2a\in[k+a,k+2a-1]$; and $a\mid k$ forces $s=0$ and $s_1=a-1$.

We prove (ii); part (i) is the identical computation one degree rung lower, with $k+1$ replaced by $k$ throughout. Let $A\subseteq Q$ with $|A|=\ell_2$, and put
\[
P_A(Y):=\prod_{b\in A}(Y-b)=Y^{\ell_2}-\eone(A)\,Y^{\ell_2-1}+R_A(Y),
\qquad
R_A(Y):=\sum_{j=2}^{\ell_2}(-1)^je_j(A)\,Y^{\ell_2-j},
\]
so $\deg_YR_A\le\ell_2-2$, and set $L_A(X):=P_A(\varphi(X))\in\B[X]$. The polynomial $L_A$ vanishes on
\[
S_A:=\varphi^{-1}(A)\cap D=\bigsqcup_{b\in A}S_b,
\qquad
|S_A|=a\ell_2=A_2,
\]
since the fibers $S_b$ are complete, pairwise disjoint, and of size exactly $a$.

Set $z_A:=-\eone(A)\in\B$ and $c_A:=\bigl(-R_A(\varphi(x))\bigr)_{x\in D}$. Since $\deg_XR_A(\varphi(X))\le a(\ell_2-2)=A_2-2a\le k$, we have $c_A\in\RS[\B,D,a(\ell_2-2)+1]\subseteq\RS[\B,D,k+1]\subseteq\RS[\F,D,k+1]$. For every $x\in S_A$,
\[
0=L_A(x)=\varphi(x)^{\ell_2}+z_A\,\varphi(x)^{\ell_2-1}+R_A(\varphi(x)),
\qquad\text{i.e.}\qquad
u_{z_A}(x)=c_A(x),
\]
so $u_{z_A}$ and $c_A$ agree on at least $A_2$ points of $D$ and $\dist(u_{z_A},c_A)\le1-A_2/n$.

The assignment $A\mapsto c_A$ is injective among subsets sharing a common value of $z_A$. Indeed, suppose $z_A=z_{A'}$ and $c_A=c_{A'}$. Then $R_A(\varphi(X))$ and $R_{A'}(\varphi(X))$ are polynomials of degree $\le k<n$ agreeing at all $n$ points of $D$, hence equal as polynomials. Composition with the nonconstant $\varphi$ is injective on $\F[Y]$: if $R:=R_A-R_{A'}\ne0$ has degree $d\ge1$, then $R(\varphi(X))$ has degree $ad\ge1$ with leading coefficient $\mathrm{lc}(R)\,\mathrm{lc}(\varphi)^d\ne0$, and if $d=0$ then $R(\varphi(X))=R$ is a nonzero constant; either way $R(\varphi(X))\ne0$. Hence $R_A=R_{A'}$, so $P_A=Y^{\ell_2}+z_AY^{\ell_2-1}+R_A(Y)=P_{A'}$, and the root multisets give $A=A'$.

Finally, $z_A=-\eone(A)$ lies in $\B$ because $A\subseteq Q\subseteq\B$, so it takes at most $|\B|$ values as $A$ ranges over the $\binom N{\ell_2}$ subsets of $Q$ of size $\ell_2$ (here $\ell_2\le N-1\le N$ is used to have this many subsets). Some value $z\in\B$ is attained by at least $\binom N{\ell_2}/|\B|$ subsets, and the corresponding codewords $c_A$ are pairwise distinct elements of $\Lst(\RS[\F,D,k+1],1-A_2/n,u_z)$.
\end{proof}

\begin{remark}[the divisibility hypothesis is removed]\label{rem:no-divisibility}
With $\varphi=X^a$, $D$ a multiplicative coset, and $a\mid k$, \cref{lem:phi-fiber} is \cref{lem:fiber} verbatim: $A_1=k+a$, $A_2=k+2a$, and $u_z$ has the printed monomial shape. The rounding $\ell_2=\lfloor k/a\rfloor+2$ removes the hypothesis $a\mid k$ at the cost of a certified agreement $A_2\in[k+a+1,k+2a]$ rather than exactly $k+2a$; the certified radius therefore satisfies
\[
1-\rho-\frac2N\ \le\ 1-\frac{A_2}n\ \le\ 1-\rho-\frac1N-\frac1n,
\]
so the listed ball always sits at least one quotient step $\approx1/N$ below capacity. This matters precisely where $a\mid k$ genuinely fails: odd-dimension circle codes (\cref{cor:circle-grand}) and mixed-radix multiplicative rows without an integral $\rho N$ (\cref{cor:mixed-radix}). \Cref{rem:slacktwo} already used slack two to move the divisibility demand from $k+1$ onto $k$; the present rounding removes it altogether.
\end{remark}

\begin{theorem}[universal cap for map-smooth domains]\label{thm:phi-cap}
Let $\B\subseteq\F$ be finite fields, $q:=|\F|$, and let $D\subseteq\B$ be a $(\varphi,a)$-smooth domain of size $n<q$. Let $1\le k<n$, $\rho:=k/n$, $C:=\RS[\F,D,k]$, and let $N=n/a$, $\ell_2=\lfloor k/a\rfloor+2$, $A_2=a\ell_2$ be as in \cref{lem:phi-fiber}. Assume $\ell_2\le N-1$ and
\begin{equation}\label{eq:hyp-phi}
\binom N{\ell_2}\ \ge\ |\B|\Bigl(\frac qk+1\Bigr).
\end{equation}
Then
\[
\eca(C,\delta)\;>\;\frac1{2k}\Bigl(1-\frac nq\Bigr)
\qquad\text{for every }\ \delta\in\Bigl[1-\frac{A_2}n,\ 1-\rho-\frac1n\Bigr],
\]
and
\[
\emca(C,\delta)\;>\;\frac1{2k}\Bigl(1-\frac nq\Bigr)
\qquad\text{for every }\ \delta\in\Bigl[1-\frac{A_2}n,\ 1-\rho\Bigr).
\]
Consequently
\[
\delta^*_C(\varepsilon^*)\ \le\ 1-\frac{A_2}n\ \le\ 1-\rho-\frac{a+1}n
\qquad\text{for every }\ \varepsilon^*\le\frac1{2k}\Bigl(1-\frac nq\Bigr).
\]
\end{theorem}

\begin{proof}
Since $\ell_2\le N-1$ we have $A_2\le n-a<n$, and since $A_2\ge k+a+1\ge k+2$ both stated intervals are nonempty, with $1-A_2/n\le1-\rho-(a+1)/n\le1-\rho-2/n$.

Fix $\delta\in[1-A_2/n,\,1-\rho-1/n]$; then $\lfloor\delta n\rfloor\le n-k-1$, so \cref{thm:A} is applicable to $C$ at radius $\delta$. Suppose toward a contradiction that $\eca(C,\delta)\le\frac1{2k}(1-\frac nq)$. Applying \cref{thm:A} with $\eta=\frac12$ gives, for $C^+:=\RS[\F,D,k+1]$,
\[
\Lst(C^+,\delta)\ \le\ \bigl\lceil2q\,\eca(C,\delta)\bigr\rceil
\ \le\ \Bigl\lceil\frac{q-n}k\Bigr\rceil
\ <\ \frac{q-n}k+1\ \le\ \frac qk+1 .
\]
On the other hand, by \cref{lem:phi-fiber}(ii) and monotonicity of lists in the radius (using $\delta\ge1-A_2/n$),
\[
\Lst(C^+,\delta)\ \ge\ \bigl|\Lst\bigl(C^+,1-\tfrac{A_2}n,u_z\bigr)\bigr|\ \ge\ \binom N{\ell_2}\Big/|\B|\ \ge\ \frac qk+1
\]
by \eqref{eq:hyp-phi}, a contradiction. Hence the $\eca$ lower bound holds throughout the stated integer-admissible interval. In particular it holds at $\delta_0:=1-A_2/n$, where \cref{fact:chain} gives the same lower bound for $\emca(C,\delta_0)$, and support-wise MCA monotonicity (\cref{lem:mca-monotone}) extends it to every $\delta\in[\delta_0,1-\rho)$. The bound on $\delta^*_C$ follows exactly as in \cref{thm:main}, using $A_2\ge k+a+1$.
\end{proof}

\begin{remark}[relation to \cref{thm:main}]\label{rem:phi-cap-scope}
With $\varphi=X^a$, $D$ a multiplicative coset, and $a\mid k$, \cref{thm:phi-cap} is \cref{thm:main} verbatim: $A_2=k+2a$, the left endpoint is $\delta_N=1-\rho-2/N$, hypothesis \eqref{eq:hyp-phi} is \eqref{eq:hyp}, and $\ell_2\le N-1$ is the printed condition $(1-\rho)N\ge3$. The generalization is strict in two independent directions --- the folding map and the removal of $a\mid k$ --- at no cost in constants. As in \cref{rem:eta}, the constant $\frac1{2k}$ can be pushed to $\frac{1-o(1)}k$ given slack in \eqref{eq:hyp-phi}.
\end{remark}

\begin{corollary}[mixed-radix multiplicative rows, divisibility-free]\label{cor:mixed-radix}
Let $\rho\in\{\frac12,\frac14,\frac18,\frac1{16}\}$, let $\F$ be any finite field with $q:=|\F|<2^{256}$, let $\B\subseteq\F$ be any subfield, let $D\subseteq\B^\times$ be a multiplicative coset of order $n$, and let $k=\rho n\le2^{40}$, $C=\RS[\F,D,k]$. Suppose $n$ has \emph{some} divisor $N$ with
\[
2^{10}\le N\le2^{11}\ \ \bigl(\rho\in\{\tfrac12,\tfrac14,\tfrac18\}\bigr),
\qquad\text{resp.}\qquad
2^{11}\le N\le2^{12}\ \ \bigl(\rho=\tfrac1{16}\bigr),
\]
with no divisibility condition relating $N$ and $k$ \textup(in particular $\rho N$ need not be an integer\textup). Then
\[
\emca(C,\delta)\;>\;\frac1{2k}\Bigl(1-\frac nq\Bigr)\;\ge\;2^{-86}\;\gg\;2^{-128}
\qquad\text{for every }\ \delta\in\Bigl[1-\rho-\frac1N,\ 1-\rho\Bigr),
\]
and the same lower bound holds for $\eca(C,\delta)$ on the integer-admissible subinterval $[1-\frac{A_2}n,\,1-\rho-\frac1n]$, with $A_2$ as in \cref{lem:phi-fiber} for $a=n/N$. In particular
\[
\delta^*_C(2^{-128})\ \le\ 1-\rho-2^{-11}\ \ \bigl(\rho\in\{\tfrac12,\tfrac14,\tfrac18\}\bigr),
\qquad
\delta^*_C(2^{-128})\ \le\ 1-\rho-2^{-12}\ \ \bigl(\rho=\tfrac1{16}\bigr).
\]
\end{corollary}

\begin{proof}
Apply \cref{thm:phi-cap} with $\varphi=X^{n/N}$ and $a=n/N$; the coset $D$ is $(X^a,a)$-smooth, and $n<q$ as in \cref{cor:grand}. Since $\lfloor k/a\rfloor\in(k/a-1,k/a]$, we have $\ell_2/N\in(\rho+\frac1N,\rho+\frac2N]$, so $\ell_2\le\rho N+2\le N/2+2\le N-1$. For \eqref{eq:hyp-phi}: the interval $(\rho+\frac1N,\rho+\frac2N]$ is contained in $[\frac18,\,0.502]$ for the first three rates and in $[\frac1{16},\,0.0645]$ for $\rho=\frac1{16}$; on each interval the concave function $\Hb$ attains its minimum at an endpoint, and $\Hb(0.502)\ge0.999\ge\Hb(\frac18)$ while $\Hb(0.0645)\ge\Hb(\frac1{16})$, so $\Hb(\ell_2/N)\ge\Hb(\frac18)\ge0.5435$, resp.\ $\Hb(\ell_2/N)\ge\Hb(\frac1{16})\ge0.3372$. Hence, using the worst $N$ in each window,
\[
\log_2\binom N{\ell_2}\ \ge\ N\,\Hb(\ell_2/N)-\log_2(N+1)\ \ge\
\begin{cases}
2^{10}\cdot0.5435-\log_2(2^{11}+1)\ \ge\ 545,\\[2pt]
2^{11}\cdot0.3372-\log_2(2^{12}+1)\ \ge\ 678,
\end{cases}
\]
while $|\B|(q/k+1)\le q(q/k+1)\le q^2+q<2^{513}$, exactly as in \cref{cor:grand}. \Cref{thm:phi-cap} gives the lower bounds on $[1-\frac{A_2}n,1-\rho-\frac1n]$ and $[1-\frac{A_2}n,1-\rho)$; since $A_2\ge k+a+1$, the MCA interval contains $[1-\rho-\frac1N,1-\rho)$, and $\frac1N\ge2^{-11}$, resp.\ $2^{-12}$. The error numerics $\frac1{2k}(1-\frac nq)\ge\frac1{2k(n+1)}>2^{-86}$ (using $n\le16k\le2^{44}$) are those of \cref{cor:grand}.
\end{proof}

For FFT-friendly mixed-radix orders $n=2^{\alpha}3^{\beta}5^{\gamma}\cdots$ the divisor set is dense on the logarithmic scale, so the window hypothesis is mild; and taking the smallest available $N$ in the window recovers the smallest gap.

\subsection{Twin cosets on the circle, Chebyshev fibers, and torus uniformization}
\label{sec:circle-geometry}

Throughout this subsection $p$ is a prime with $p\equiv3\pmod4$, so that $X^2+1$ is irreducible over $\F_p$ and $\F_{p^2}=\F_p(i)$ with $i^2=-1$. Let
\[
\UU\ :=\ \{u\in\F_{p^2}^{\times}:\ u^{p+1}=1\}
\]
be the norm-one torus, a cyclic group of order $p+1$; for the deployed Mersenne prime $p=2^{31}-1$ one has $p+1=2^{31}$. The circle $\mathcal C(\F_p):=\{(x,y)\in\F_p^2:x^2+y^2=1\}$ is in bijection with $\UU$ via
\[
\iota(x,y)\ :=\ x+iy,
\qquad
\iota^{-1}(u)\ =\ \Bigl(\frac{u+u^{-1}}2,\ \frac{u-u^{-1}}{2i}\Bigr),
\]
using $u^p=u^{-1}$ for $u\in\UU$ and $i^p=-i$. Write $\chi(u):=\frac{u+u^{-1}}2\in\F_p$ for the $x$-coordinate map; $\chi(u)=\chi(v)$ holds if and only if $v\in\{u,u^{-1}\}$. The Chebyshev polynomials of the first kind, normalized by $T_0=1$, $T_1=X$, $T_{a+1}=2XT_a-T_{a-1}$, have integer coefficients and degree $a$, reduce modulo $p$ to polynomials in $\F_p[X]$, and satisfy the semiconjugacy
\[
T_a(\chi(u))=\chi(u^a)\quad(u\in\UU,\ a\ge0),
\qquad\qquad
T_{ab}=T_a\circ T_b .
\]
Following the twin-coset geometry of circle FFTs \cite{HLP24}, for a subgroup $H\le\UU$ of order $M$ and $g\in\UU$ with $g^2\notin H$ we call
\[
\mathcal D\ :=\ gH\ \cup\ g^{-1}H\ \subseteq\ \UU
\]
a \emph{twin coset}. It has size $2M$, is closed under inversion, and contains no self-inverse element: if $u\in gH$ satisfied $u=u^{-1}$, then $u\in gH\cap(gH)^{-1}=gH\cap g^{-1}H$, forcing $g^2\in H$; likewise for $u\in g^{-1}H$. Consequently $\chi$ is exactly two-to-one on $\mathcal D$, and the $x$-coordinate domain $D:=\chi(\mathcal D)\subseteq\F_p$ has size exactly $M$. For $a\mid M$ write $H^{(a)}:=\{h^a:h\in H\}$ for the image subgroup, of order $M/a$, and $\mathcal D^{(a)}:=\{u^a:u\in\mathcal D\}$.

\begin{lemma}[power fibers on torus twin cosets]\label{lem:torus-fibers}
Let $\mathcal D=gH\cup g^{-1}H$ be a twin coset of size $2M$ and let $a\mid M$.
\begin{enumerate}[label=\textup{(\alph*)}]
\item If $g^{2a}\notin H^{(a)}$, then $\mathcal D$ is $(X^a,a)$-smooth \textup(over any field containing $\F_{p^2}$\textup), and $\mathcal D^{(a)}=g^aH^{(a)}\sqcup g^{-a}H^{(a)}$ is again a twin coset, of size $2M/a$.
\item If $g^{2a}\in H^{(a)}$, then $\mathcal D$ is $(X^a,2a)$-smooth, with $\mathcal D^{(a)}=g^aH^{(a)}$ a single coset of size $M/a$.
\end{enumerate}
Moreover, $g^{2a}\notin H^{(a)}$ holds for every $a\mid M$ simultaneously if and only if $g^{2M}\ne1$, i.e.\ $\operatorname{ord}(g)\nmid2M$.
\end{lemma}

\begin{proof}
Since $a\mid M\mid p+1$ and $\UU$ is cyclic of order $p+1$, the kernel $K_a$ of $u\mapsto u^a$ on $\UU$ has order $a$ and, being the unique subgroup of $\UU$ of that order, is contained in $H$. Hence the $a$-th power map sends the coset $gH$ onto $g^aH^{(a)}$ with every fiber a coset of $K_a$, i.e.\ exactly $a$-to-one, and likewise sends $g^{-1}H$ onto $g^{-a}H^{(a)}$. The two image cosets of the subgroup $H^{(a)}$ are disjoint or equal according to whether $g^{2a}\notin H^{(a)}$ or $g^{2a}\in H^{(a)}$. In the disjoint case, every $v\in\mathcal D^{(a)}$ lies in exactly one of the two image cosets and has exactly $a$ preimages in $\mathcal D$; the twin-coset condition for $\mathcal D^{(a)}$, namely $(g^a)^2\notin H^{(a)}$, is the hypothesis itself. In the coincident case, every $v\in g^aH^{(a)}$ has exactly $a$ preimages in $gH$ and exactly $a$ in $g^{-1}H$, hence exactly $2a$ in $\mathcal D$. For the last claim: at $a=M$ we have $H^{(M)}=\{1\}$, so the scale-$M$ condition reads $g^{2M}\ne1$; conversely, if $g^{2a}=h^a$ for some $a\mid M$ and $h\in H$, then $g^{2M}=(g^{2a})^{M/a}=h^{M}=1$.
\end{proof}

\begin{lemma}[exact Chebyshev fibers on $x$-coordinate twin-coset domains]\label{lem:cheb-fibers}
Let $\mathcal D=gH\cup g^{-1}H$ be a twin coset of size $2M$, let $D:=\chi(\mathcal D)\subseteq\F_p$, $|D|=M$, and let $a\mid M$ satisfy $g^{2a}\notin H^{(a)}$. Then $D$ is $(T_a,a)$-smooth over $\F_p$, and
\[
T_a(D)\ =\ \chi\bigl(\mathcal D^{(a)}\bigr)
\]
is the $x$-coordinate domain of the twin coset $\mathcal D^{(a)}$, of size $M/a$. In particular, if $\operatorname{ord}(g)\nmid2M$, this holds at every scale $a\mid M$ simultaneously, and for $M=2^m$ the tower
\[
D\ \xrightarrow{\ T_2\ }\ T_2(D)\ \xrightarrow{\ T_2\ }\ T_4(D)\ \xrightarrow{\ T_2\ }\ \cdots
\]
consists of $x$-coordinate twin-coset domains of sizes $M,M/2,M/4,\dots$ --- the line-round tower of circle FRI \cite{HLP24}.
\end{lemma}

\begin{proof}
The semiconjugacy gives $T_a(D)=T_a(\chi(\mathcal D))=\chi(\mathcal D^{(a)})$; by \cref{lem:torus-fibers}(a), $\mathcal D^{(a)}$ is a twin coset, so $\chi$ is exactly two-to-one on it and $|T_a(D)|=|\mathcal D^{(a)}|/2=M/a$.

Fix $w\in T_a(D)$ and write $w=\chi(v)$ with $v\in\mathcal D^{(a)}=g^aH^{(a)}\sqcup g^{-a}H^{(a)}$. Since $\bigl(g^aH^{(a)}\bigr)^{-1}=g^{-a}H^{(a)}$, the set $\mathcal D^{(a)}$ is closed under inversion, and after replacing $v$ by $v^{-1}$ if necessary (which does not change $w$) we may assume $v\in g^aH^{(a)}$. Then $v\ne v^{-1}$: otherwise $v$ would lie in both of the disjoint cosets $g^aH^{(a)}$ and $g^{-a}H^{(a)}$. For $x=\chi(u)\in D$ with $u\in\mathcal D$,
\[
T_a(x)=w
\iff
\chi(u^a)=\chi(v)
\iff
u^a\in\{v,v^{-1}\}.
\]
By the fiber structure in the proof of \cref{lem:torus-fibers}, the equation $u^a=v$ has exactly $a$ solutions $u\in gH$ and none in $g^{-1}H$ (whose $a$-th powers fill the disjoint coset $g^{-a}H^{(a)}$), and the equation $u^a=v^{-1}$ has exactly $a$ solutions in $g^{-1}H$ --- the inverses of the former --- and none in $gH$. Hence
\[
E_w\ :=\ \{u\in\mathcal D:\ u^a\in\{v,v^{-1}\}\}
\]
has exactly $2a$ elements, is closed under inversion, and, being a subset of $\mathcal D$, contains no self-inverse element. Since $\chi(u)=\chi(u')$ forces $u'\in\{u,u^{-1}\}$, the image $\chi(E_w)$ --- which is precisely the fiber $T_a^{-1}(w)\cap D$ --- has exactly $|E_w|/2=a$ elements. As $w\in T_a(D)$ was arbitrary, $D$ is $(T_a,a)$-smooth with $|T_a(D)|=M/a=|D|/a$, as required. The tower statement follows from the all-scales equivalence in \cref{lem:torus-fibers} together with $T_{2^{j+1}}=T_2\circ T_{2^j}$.
\end{proof}

\begin{remark}[standard-position cosets are twin cosets with $\operatorname{ord}(g)\nmid2M$]\label{rem:standard-position}
A standard-position coset of circle FFT \cite{HLP24} is a coset $gH_0$, $H_0:=\langle g^2\rangle$, of size $2M$ with $\operatorname{ord}(g)=4M$ (so $|H_0|=2M$). Let $K\le H_0$ be the index-two subgroup, of order $M$. Then $g^{-1}K=g\,g^{-2}K$, and $g^{-2}$ generates $H_0$, hence $g^{-2}\notin K$, so
\[
gH_0\ =\ gK\ \sqcup\ g^{-1}K
\]
is exactly the twin coset attached to $(K,g)$: the condition $g^2\notin K$ holds because $g^2$ generates $H_0\supsetneq K$, and $\operatorname{ord}(g)=4M\nmid2M$. Thus every standard-position domain of \cite{HLP24}, at every folding level, satisfies the hypotheses of \cref{lem:torus-fibers,lem:cheb-fibers} at every dyadic scale simultaneously.
\end{remark}

\begin{lemma}[diagonal invariance of lists and correlated agreement]\label{lem:diag-invariance}
Let $C\subseteq\F^n$ be a linear code and $t\in(\F^{\times})^n$; write $t\circ v$ for the coordinatewise product and $t\circ C:=\{t\circ c:c\in C\}$. Then for every received word $U$, every radius $\delta$, and every pair of radii,
\[
\Lst(t\circ C,\delta,t\circ U)=t\circ\Lst(C,\delta,U),
\qquad
\eca(t\circ C,\cdot,\cdot)=\eca(C,\cdot,\cdot),
\qquad
\emca(t\circ C,\cdot)=\emca(C,\cdot).
\]
\end{lemma}

\begin{proof}
The map $v\mapsto t\circ v$ is an $\F$-linear bijection of $\F^n$ that preserves coordinatewise agreement: $(t\circ v)_j=(t\circ w)_j\iff v_j=w_j$. Hence it preserves $\dist$ and every restricted distance $\distS$, maps $C$ onto $t\circ C$, acts row-wise on pairs, and commutes with the line structure: $t\circ(f_1+\gamma f_2)=(t\circ f_1)+\gamma\,(t\circ f_2)$. It therefore carries every witness of the defining events of $\Lst$, $\eca$, and $\emca$ for $C$ bijectively onto a witness for $t\circ C$ with the same slope $\gamma$ and the same agreement supports, and conversely via $t^{-1}$.
\end{proof}

\begin{lemma}[torus uniformization of circle codes]\label{lem:circle-rs}
Let $p\equiv3\pmod4$, let $\F\supseteq\F_{p^2}$ be a finite field \textup(so $i\in\F$\textup), let $E\subseteq\mathcal C(\F_p)$ be a set of circle points, and let $E':=\iota(E)\subseteq\UU\subseteq\F_{p^2}^{\times}$. For $w\ge0$ define the $\F$-coefficient circle code
\[
\mathcal C_w(\F,E)\ :=\ \bigl\{(f(P))_{P\in E}\ :\ f\in\F[x,y]/(x^2+y^2-1),\ \deg f\le w\bigr\}.
\]
Then, under the coordinate bijection $\iota:E\to E'$ and with the diagonal vector $t:=(u^{-w})_{u\in E'}$,
\[
\mathcal C_w(\F,E)\ =\ t\circ\RS[\F,E',2w+1].
\]
Consequently, by \cref{lem:diag-invariance}, the circle code and the Reed--Solomon code $\RS[\F,E',2w+1]$ of odd dimension $2w+1$ on the torus domain $E'$ have identical list sizes at every radius and identical $\eca$ and $\emca$ at every (pair of) radii.
\end{lemma}

\begin{proof}
Since $i\in\F$ and $x+iy$ is invertible in the quotient ring (with inverse $x-iy$, as $(x+iy)(x-iy)=x^2+y^2=1$), the substitution $u:=x+iy$ defines a ring isomorphism
\[
\F[x,y]/(x^2+y^2-1)\ \xrightarrow{\ \sim\ }\ \F[u,u^{-1}],
\qquad
x\mapsto\frac{u+u^{-1}}2,\quad y\mapsto\frac{u-u^{-1}}{2i},
\]
with inverse $u\mapsto x+iy$; it matches evaluation at $P\in E$ on the left with evaluation at $\iota(P)\in E'$ on the right. Every class of degree $\le w$ has the canonical form $f_0(x)+y\,f_1(x)$ with $\deg f_0\le w$ and $\deg f_1\le w-1$ (reduce $y^2=1-x^2$), and $\{h_0(x)+y\,h_1(x)\}$ is a free-module canonical form of the quotient ring over $\F[x]$, so this space has dimension exactly $2w+1$. Under the isomorphism,
\[
x^j=2^{-j}u^{-j}(u^2+1)^j\in u^{-j}\,\F[u]_{\le2j},
\qquad
y\,x^j\in u^{-(j+1)}\,\F[u]_{\le2j+2},
\]
so the degree-$\le w$ subspace maps into $u^{-w}\F[u]_{\le2w}$; both spaces have dimension $2w+1$ and the map is injective, hence the image is exactly $u^{-w}\F[u]_{\le2w}$. Evaluating on $E'$ gives
\[
\mathcal C_w(\F,E)
=\bigl\{(u^{-w}h(u))_{u\in E'}\ :\ h\in\F[u],\ \deg h\le2w\bigr\}
=t\circ\RS[\F,E',2w+1].
\qedhere
\]
\end{proof}

Note the parity: the uniformized dimension $k_c=2w+1$ is always odd, while the useful folding scales $a$ on a dyadic torus twin coset are even, so $a\nmid k_c$ \emph{always}. The divisibility-free form of \cref{lem:phi-fiber} is therefore not a luxury here: the original \cref{lem:fiber} can never be applied to a uniformized circle code directly.

\subsection{Universal caps for circle codes and circle-FRI line rounds}
\label{sec:circle-caps}

\begin{corollary}[universal circle-row cap]\label{cor:circle-grand}
Let $p\equiv3\pmod4$ be prime, let $\mathcal D=gH\cup g^{-1}H\subseteq\UU$ be a twin coset of size $2M$ with $\operatorname{ord}(g)\nmid2M$ \textup(e.g.\ any standard-position coset of \cite{HLP24}, by \cref{rem:standard-position}\textup), and let $\F$ be any finite field with $q:=|\F|<2^{256}$.
\begin{enumerate}[label=\textup{(\alph*)}]
\item \textup{(Line rounds.)} Assume $\F\supseteq\F_p$. Let $D:=\chi(\mathcal D)$, $n:=M$, and $C:=\RS[\F,D,k]$ with $k\le2^{40}$ and $\rho:=k/n\in[\tfrac1{16},\tfrac12]$. If $\rho\ge\frac18$, assume $2^{10}\mid M$ and set $N:=2^{10}$; if $\rho<\frac18$, assume $2^{11}\mid M$ and set $N:=2^{11}$. Then
\[
\emca(C,\delta)\;>\;\frac1{2k}\Bigl(1-\frac nq\Bigr)\;\ge\;2^{-86}\;\gg\;2^{-128}
\qquad\text{for every }\ \delta\in\Bigl[1-\rho-\frac1N,\ 1-\rho\Bigr),
\]
and the same lower bound holds for $\eca(C,\delta)$ on the integer-admissible subinterval $[1-\frac{A_2}n,\,1-\rho-\frac1n]$, with $A_2$ as in \cref{lem:phi-fiber} for $a=n/N$. If moreover $\frac nN\mid k$ --- as in every $2$-power instantiation at the four challenge rates --- the MCA band extends to $[1-\rho-\frac2N,\,1-\rho)$. In particular
\[
\delta^*_C(2^{-128})\le1-\rho-2^{-10}\ \ (\rho\ge\tfrac18),
\qquad
\delta^*_C(2^{-128})\le1-\rho-2^{-11}\ \ (\rho<\tfrac18),
\]
improving to $1-\rho-2^{-9}$ \textup(resp.\ $1-\rho-2^{-10}$\textup) under $\frac nN\mid k$.
\item \textup{(Circle codes.)} Assume $\F\supseteq\F_{p^2}$. Let $C:=\mathcal C_w(\F,\iota^{-1}(\mathcal D))$ with $k_c:=2w+1\le2^{41}$, $n_c:=2M$, and $\rho_c:=k_c/n_c\in[\tfrac1{16},\tfrac12]$. If $\rho_c\ge\frac18$, assume $2^{9}\mid M$ and set $N:=2^{10}$; if $\rho_c<\frac18$, assume $2^{10}\mid M$ and set $N:=2^{11}$. Then
\[
\emca(C,\delta)\;>\;\frac1{2k_c}\Bigl(1-\frac{n_c}q\Bigr)\;\ge\;2^{-88}\;\gg\;2^{-128}
\qquad\text{for every }\ \delta\in\Bigl[1-\rho_c-\frac1N,\ 1-\rho_c\Bigr),
\]
with the corresponding $\eca$ bound on $[1-\frac{A_2}{n_c},\,1-\rho_c-\frac1{n_c}]$ \textup(here $a:=n_c/N\nmid k_c$ always, since $k_c$ is odd\textup). In particular $\delta^*_C(2^{-128})\le1-\rho_c-2^{-10}$ for $\rho_c\ge\frac18$, and $\delta^*_C(2^{-128})\le1-\rho_c-2^{-11}$ for $\rho_c<\frac18$.
\end{enumerate}
\end{corollary}

\begin{proof}
(a) Put $a:=n/N=M/N$, so $a\mid M$, and $g^{2a}\notin H^{(a)}$ by the all-scales equivalence of \cref{lem:torus-fibers}; \cref{lem:cheb-fibers} then makes $D$ a $(T_a,a)$-smooth domain over $\B:=\F_p\subseteq\F$. We check the hypotheses of \cref{thm:phi-cap}. First, $n<q$: since $M\mid p+1$ and $\operatorname{ord}(g)\mid p+1$, the case $M=p+1$ would give $\operatorname{ord}(g)\mid p+1\mid2M$, contradicting $\operatorname{ord}(g)\nmid2M$; hence $M\le\frac{p+1}2$ and $n\le\frac{p+1}2<p\le q$. Second, $\ell_2\le\rho N+2\le N/2+2\le N-1$. Third, for \eqref{eq:hyp-phi}: $\ell_2/N\in(\rho+\frac1N,\rho+\frac2N]$, which is contained in $[\frac18,\,0.502]$ when $\rho\ge\frac18$ and in $[\frac1{16},\,0.126]$ when $\rho<\frac18$; on each interval the concave $\Hb$ is minimized at an endpoint, and both right endpoints have larger $\Hb$-value than the left, so $\Hb(\ell_2/N)\ge\Hb(\frac18)\ge0.5435$, resp.\ $\Hb(\ell_2/N)\ge\Hb(\frac1{16})\ge0.3372$. Hence
\[
\log_2\binom N{\ell_2}\ \ge\ N\,\Hb(\ell_2/N)-\log_2(N+1)\ \ge\ 546,\ \ \text{resp.}\ \ 679,
\]
while $|\B|(q/k+1)\le q(q/k+1)\le q^2+q<2^{513}$. The error numerics are those of \cref{cor:grand}: $\rho\ge\frac1{16}$ gives $n\le16k\le2^{44}$, so $\frac1{2k}(1-\frac nq)\ge\frac1{2k(n+1)}>2^{-86}$. \Cref{thm:phi-cap} now yields the $\eca$ and $\emca$ lower bounds on $[1-\frac{A_2}n,1-\rho-\frac1n]$ and $[1-\frac{A_2}n,1-\rho)$ respectively; since $A_2\ge k+a+1$, the MCA interval contains $[1-\rho-\frac1N,1-\rho)$, and when $\frac nN\mid k$ one has $A_2=k+2a$, so it equals $[1-\rho-\frac2N,1-\rho)$.

(b) By \cref{lem:circle-rs}, all list sizes, $\eca$, and $\emca$ of $C$ coincide with those of $C':=\RS[\F,\mathcal D,k_c]$ on the torus twin coset $\mathcal D\subseteq\F_{p^2}$. Put $a:=n_c/N=2M/N$; then $a\mid M$ (indeed $M/a=N/2\in\mathbb Z$), and $g^{2a}\notin H^{(a)}$ as before, so $\mathcal D$ is $(X^a,a)$-smooth over $\B:=\F_{p^2}\subseteq\F$ by \cref{lem:torus-fibers}(a). Also $n_c=2M<p+1\le p^2\le q$: the case $2M=p+1$ would again give $\operatorname{ord}(g)\mid p+1=2M$. The verification of \eqref{eq:hyp-phi} is identical to (a), with the same two entropy cases and $|\B|=p^2\le q$; the error numerics use $n_c\le16k_c\le2^{45}$, so $\frac1{2k_c}(1-\frac{n_c}q)\ge\frac1{2k_c(n_c+1)}>2^{-88}$. \Cref{thm:phi-cap} applied to $C'$, transported back through \cref{lem:circle-rs,lem:diag-invariance}, gives the stated bands; $A_2\ge k_c+a+1$ again yields the containment of $[1-\rho_c-\frac1N,1-\rho_c)$.
\end{proof}

\begin{corollary}[deployed circle parameters: Mersenne-31 with QM31]\label{cor:circle-deployed}
Let $p=2^{31}-1$, let $\F=\F_{p^4}$ be the deployed QM31 extension \textup(so $q=p^4>2^{123.9}$ and $\F\supseteq\F_{p^2}$\textup), and let $\mathcal D$ be a standard-position twin coset with $M=2^{21}$ \textup(so $\operatorname{ord}(g)=2^{23}\nmid2^{22}=2M$\textup).
\begin{enumerate}[label=\textup{(\alph*)}]
\item \textup{(Line round, rate $\frac12$.)} With $D=\chi(\mathcal D)$, $|D|=2^{21}$, and $C=\RS[\F_{p^4},D,2^{20}]$,
\[
\emca(C,\delta)\;>\;\bigl(1-2^{-102}\bigr)\,2^{-21}\;>\;2^{-22}
\qquad\text{for every }\ \delta\in\Bigl[\frac12-2^{-7},\ \frac12\Bigr),
\]
and the same lower bound holds for $\eca(C,\delta)$ for every $\delta\in[\frac12-2^{-7},\,\frac12-2^{-21}]$.
\item \textup{(Circle code.)} With $C=\mathcal C_w(\F_{p^4},\iota^{-1}(\mathcal D))$, $w=2^{20}$, $k_c=2^{21}+1$, $n_c=2^{22}$, $\rho_c=\frac12+2^{-22}$,
\[
\emca(C,\delta)\;>\;2^{-23}
\qquad\text{for every }\ \delta\in\Bigl[\frac12-2^{-7},\ 1-\rho_c\Bigr),
\]
and the same lower bound holds for $\eca(C,\delta)$ for every $\delta\in[\frac12-2^{-7},\,\frac12-2^{-21}]$.
\end{enumerate}
\end{corollary}

\begin{proof}
Both parts use $N=256$ and, as in \cref{cor:deployed}, $\log_2\binom{256}{130}\ge247$.

(a) Here $a=n/N=2^{13}$ divides $k=2^{20}$, so $\ell_2=\rho N+2=130$, $A_2=k+2a=2^{20}+2^{14}$, and $1-A_2/n=\frac12-2^{-7}$. The standard-position hypothesis gives exact Chebyshev fibers at scale $a$ (\cref{rem:standard-position,lem:cheb-fibers}), with $\B=\F_p$. Hypothesis \eqref{eq:hyp-phi}: $|\B|(q/k+1)=p\,(p^4/2^{20}+1)<2^{31}\cdot2^{105}=2^{136}\le2^{247}\le\binom{256}{130}$. The error bound is $\frac1{2k}(1-n/q)$ with $n/q=2^{21}/p^4<2^{-102}$. \Cref{thm:phi-cap} gives the MCA band $[\frac12-2^{-7},\frac12)$ and the $\eca$ band $[\frac12-2^{-7},\frac12-2^{-21}]$, since $1-\rho-\frac1n=\frac12-2^{-21}$.

(b) Here $a=n_c/N=2^{14}\mid M=2^{21}$, $\ell_2=\lfloor(2^{21}+1)/2^{14}\rfloor+2=2^7+2=130$, $A_2=2^{14}\cdot130=2^{21}+2^{15}$, and $1-A_2/n_c=\frac12-2^{-7}$. \Cref{lem:circle-rs} transports everything to $C'=\RS[\F_{p^4},\mathcal D,2^{21}+1]$ on the torus twin coset, which is $(X^a,a)$-smooth over $\B=\F_{p^2}$ by \cref{lem:torus-fibers,rem:standard-position}. Hypothesis \eqref{eq:hyp-phi}: $|\B|(q/k_c+1)=p^2\,\bigl(p^4/(2^{21}+1)+1\bigr)<2^{62}\cdot2^{104}=2^{166}\le2^{247}$. The error bound: $n_c/q=2^{22}/p^4<2^{-101}$, so
\[
\frac1{2k_c}\Bigl(1-\frac{n_c}q\Bigr)\;>\;\frac{1-2^{-101}}{2^{22}+2}\;>\;2^{-22}\bigl(1-2^{-21}\bigr)\bigl(1-2^{-101}\bigr)\;>\;2^{-23}.
\]
\Cref{thm:phi-cap} gives the MCA band $[\frac12-2^{-7},\,1-\rho_c)$ and the $\eca$ band $[\frac12-2^{-7},\,1-\rho_c-\frac1{n_c}]=[\frac12-2^{-7},\,\frac12-2^{-21}]$.
\end{proof}

\begin{corollary}[circle certifying lines are genuinely extension-valued]\label{cor:circle-Fvalued}
The subfield-confinement lemma \cref{lem:confine} holds verbatim for an arbitrary evaluation set $D\subseteq\B$: its proof uses only that the evaluation points lie in $\B$, not that they are invertible. Consequently, at the parameters of \cref{cor:circle-deployed}:
\begin{enumerate}[label=\textup{(\alph*)}]
\item any pair of words on $D=\chi(\mathcal D)$ that is $\F_p$-valued up to a common scalar factor has MCA-bad-slope density at most $p/q=p^{-3}<2^{-92}$ at every radius; in particular, every pair attaining the bound $>2^{-22}$ of \cref{cor:circle-deployed}\textup{(a)} is genuinely $\F_{p^4}$-valued;
\item the diagonal vector $t=(u^{-w})_{u}$ of \cref{lem:circle-rs} is $\F_{p^2}$-valued, so any pair of circle words whose $\iota$-transports are $\F_{p^2}$-valued up to a common scalar factor has MCA-bad-slope density at most $p^2/q=p^{-2}<2^{-61}$ at every radius; in particular, every circle pair attaining the bound $>2^{-23}$ of \cref{cor:circle-deployed}\textup{(b)} contains a word whose $\iota$-transport is not $\F_{p^2}$-valued, even after removing a common scalar.
\end{enumerate}
As in \cref{rem:extension-pole-status}, the extension-pole floors \cref{cor:extension-pole-deep-list-floor,cor:extension-pole-quotient-remainder-floor} supply the matching constructive mechanism: the fiber received words of \cref{lem:phi-fiber} on circle domains are $\B$-valued with $\B$-valued lists, and any such list converts through a simple pole $\alpha\in\F\setminus\B$ into genuinely extension-valued bad lines, with denominator $q-b$ ($b=|\B|\in\{p,p^2\}$) in place of $q-n$.
\end{corollary}

\begin{proof}
Inspecting the proof of \cref{lem:confine}: it decomposes words and codewords over a $\B$-basis of $\F$ and evaluates polynomials at the points of $D\subseteq\B$; no step uses $D\subseteq\B^{\times}$, so the two-radius statement holds for arbitrary $D\subseteq\B$, in particular for the Chebyshev domain $\chi(\mathcal D)\subseteq\F_p$ and the torus domain $\mathcal D\subseteq\F_{p^2}$. (a) is then \cref{lem:confine} with $\B=\F_p$ inside $q=p^4$: the bad-slope density of a $\B$-confined pair is at most $|\B|/q=p^{-3}<2^{-92}<2^{-22}$. (b) A circle pair whose $\iota$-transports are $\F_{p^2}$-valued up to a common scalar corresponds, after multiplying by the $\F_{p^2}$-valued diagonal $t$, to an $\F_{p^2}$-confined pair for $\RS[\F_{p^4},\mathcal D,k_c]$; \cref{lem:diag-invariance} preserves the bad-slope set, and \cref{lem:confine} with $\B=\F_{p^2}$ bounds its density by $p^2/q=p^{-2}<2^{-61}<2^{-23}$. The final claim is \cref{cor:extension-pole-deep-list-floor} applied with $\B\in\{\F_p,\F_{p^2}\}$ to the $\B$-valued received words and lists of \cref{lem:phi-fiber}.
\end{proof}

\begin{remark}[interleaving, slacked fallback, and the two-tier picture for circle rows]\label{rem:circle-transfer}
The interleaving-transfer lemma \cref{lem:inter} is stated for arbitrary linear codes, so every bound of \cref{cor:circle-grand,cor:circle-deployed} transfers verbatim to $s$-interleaved circle rows, exactly as in \cref{cor:rows}. The slacked BCHKS/ABF fallback also carries over: \cref{lem:phi-fiber}(i) supplies a list of size $\ge\binom N{\ell_1}/|\B|$ in $\RS[\F,D,k]$ at radius $1-\frac{A_1}n$, and whenever $A_1\ge k+3$ \textup(automatic in the $2$-power line rounds with $a\ge4$, where $A_1=k+a$\textup) and $\binom N{\ell_1}\ge|\B|\,q$, the contrapositive of \cref{thm:B} gives $\eca\bigl(C,\,1-\frac{A_1}n+\frac2n,\,1-\rho-\frac1n\bigr)\ge\frac1{2n}$, as in \cref{prop:slacked}. Finally, the two-tier summary of \cref{rem:twotier} acquires a circle column: the deployed Mersenne-$31$ circle family carries the same kind of near-capacity negative band as the multiplicative rows --- gap $2^{-7}$ at the deployed sizes, $2^{-9}$--$2^{-11}$ universally --- over every challenge field $q<2^{256}$ containing the relevant base field.
\end{remark}

\subsection{The lacunarity obstruction for genus-one (ECFFT) domains}
\label{sec:ecfft}

\begin{remark}[what blocks the elliptic case]\label{rem:ecfft-obstruction}
The proof of \cref{lem:phi-fiber} works because $\varphi$ is a \emph{polynomial}: in the expansion $\prod_{b\in A}(\varphi-b)=\sum_j(-1)^je_j(A)\,\varphi^{\ell-j}$, consecutive terms drop by a full $a=\deg\varphi$ in degree, so the received word occupies the top two degree slots and everything else is a codeword sitting $2a$ degrees down --- gap $\approx2/N$. Consider instead a rational folding map $\psi=f/g$ of degree $a$, with $e:=\deg g<\deg f=a$ and $g$ nonvanishing on $D$, whose fibers inside $D$ are complete and of uniform size $a$. The natural locators are $\Lambda_A:=\prod_{b\in A}(f-b\,g)$, of degree $a\ell$, vanishing on the $\psi$-fibers over $A$; their expansion $\sum_j(-1)^je_j(A)\,f^{\ell-j}g^{\,j}$ has $j$-th term of degree $a\ell-j(a-e)$, so the slack-two window has width $2(a-e)$ in degree and the certified gap is $2(a-e)/n$. Polynomial maps have $e=0$ and recover gap $2/N$. The FFTree domains of ECFFT \cite{BCKL21} fold by $x$-coordinate maps of $2$-isogenies of elliptic curves, which are rational of degree type $(a,e)=(2,1)$; every composition of such maps is again the $x$-map of an isogeny, of type $(a,a-1)$. For these, $a-e=1$ at every scale: the window collapses to width $2$, and the construction certifies only the first staircase step --- a large list at agreement $k+2$, hence an $\eca$ floor at gap $2/n$ --- a nonvacuous but microscopic band, far from the $2^{-9}$--$2^{-11}$ bands above. The circle escapes this because its folding tower becomes polynomial after the $x$-projection: $T_2$ fixes $\infty$ and is totally ramified there. The elliptic $x$-maps also fix $\infty$ but are \emph{not} totally ramified there ($\psi^{-1}(\infty)=\{\infty,x(T)\}$ for a $2$-isogeny with kernel generator $T$), and the locator expansion sees exactly this defect as $e=a-1$. Obtaining a macroscopic universal cap for ECFFT rows therefore requires either a lacunary basis of isogeny locators, a different received-word family, or a matching safe-side theorem showing the collapse is essential. The second exit is realized by the rational graded floor of \cref{prop:graded-rational-floor} \textup(\cref{cor:ecfft-macroscopic}\textup); the residual band form is recorded as Open Problem~\ref{prob:isogeny}.
\end{remark}


\section{Threshold formulation and certificate interface}\label{sec:prize-program}

The Proximity Prize asks for a threshold, not merely for a near-capacity
obstruction.  We therefore use the integer staircase below to state both sides
of the result: lower certificates mark agreements where MCA or list decoding is
unsafe, while upper certificates mark agreements where the corresponding error
is below the target.  The paper does not determine the staircase to one step in
the remaining band, but it reduces the unknown part to a single aperiodic
correlated-agreement problem and gives exact finite certificates for all
printed edges.

\subsection{Established inputs}

For the grand MCA challenge at target error $\varepsilon^*=2^{-128}$, the
unconditional output of \cref{cor:grand} is the following negative certificate:
for every challenge-envelope row satisfying the printed divisor hypothesis,
\[
        \delta\ge 1-\rho-2^{-9}
        \quad(\rho\in\{1/2,1/4,1/8\})
\]
and
\[
        \delta\ge 1-\rho-2^{-10}
        \quad(\rho=1/16)
\]
are impossible MCA radii.  Equivalently, any full determination of
$\delta_C^*(2^{-128})$ for these smooth multiplicative rows must search below
these gaps.  This conclusion uses only the self-contained deep-point conversion,
the quotient-fiber pigeonhole, and support-wise MCA monotonicity; it no longer
rests on importing the Crites--Stewart conversion.

The theorem also distinguishes two different grades of information.  First, the
\emph{threshold grade} is already enough for the prize inequality: the certified
error is $\gg2^{-128}$.  Second, the \emph{failure-profile grade} is still weak:
the present proof gives error of order $1/k$, not error $1-o(1)$.  Closing the
error-one profile at fixed gaps such as $1/64$ is a different, stronger problem.
It would improve our understanding of the failure landscape, but is not needed
to obtain the stated upper bound on $\delta_C^*(2^{-128})$.

The negative certificate also covers a wider row family.  The
same conclusion --- an unsafe near-capacity band at error $\gg2^{-128}$ ---
now holds for mixed-radix multiplicative rows without the integrality
condition $\rho N\in\mathbb Z$ (\cref{cor:mixed-radix}, gaps
$2^{-11}$/$2^{-12}$), for the $x$-coordinate line rounds of circle FRI over
any $p\equiv3\pmod4$, and for the bivariate circle codes over any challenge
field containing $i$ (\cref{cor:circle-grand}, gaps $2^{-9}$--$2^{-11}$),
including the deployed Mersenne-$31$/QM31 family at gap $2^{-7}$
(\cref{cor:circle-deployed}).  These additions use the same three
ingredients; the only new inputs are the exact Chebyshev-fiber geometry of
twin cosets and the torus uniformization, both elementary.

The treatment in \cref{sec:discharge} adds the safe side. \Cref{thm:deep-mca} certifies $\emca(C,\delta)\le(\lfloor\delta n\rfloor+1)/q$ whenever $3\lfloor\delta n\rfloor\le n-k$, for every row at once, and \cref{thm:johnson-list} certifies the interleaved-list challenge up to the Johnson radius uniformly in the arity; \cref{tab:status} records the per-obligation outcome, and \cref{thm:sandwich} assembles the resulting two-sided sandwich.

\subsection{Integer-staircase formulation}

Let
\[
        a(\delta):=\lceil(1-\delta)n\rceil,
        \qquad r(\delta):=n-a(\delta)
\]
be the agreement and integer-radius forms of the same radius.  For a fixed finite-slope line
field $F_{\rm line}$, write $q_{\rm line}=|F_{\rm line}|$ and let
\[
        B_{\rm mca}(a)
\]
be the exact maximum number of MCA-bad line parameters at agreement at least
$a$, for the line family used in the challenge.  Then
\[
        \emca(C,1-a/n)=B_{\rm mca}(a)/q_{\rm line}
\]
under the closed finite-slope convention.  For projective-slope or curve samplers, the same staircase formulation applies after replacing $q_{\rm line}$ and $B_{\rm mca}$ by the corresponding sampler denominator and numerator.  Thus the MCA prize problem is an
integer staircase problem: find the first agreement $a$ for which
\[
        B_{\rm mca}(a)\le 2^{-128}q_{\rm line}.
\]
The present paper supplies lower bounds on $B_{\rm mca}(a)$ near capacity by
routing quotient-fiber list mass through deep points.

\begin{proposition}[one-step threshold certificate]\label{prop:onestep}
Fix a row $C$ and a denominator $q_{\rm line}$.  Suppose that for some agreement
$a_0$ one has a lower certificate
\[
        B_{\rm mca}(a_0)>2^{-128}q_{\rm line}
\]
and an upper theorem
\[
        B_{\rm mca}(a_0+1)\le2^{-128}q_{\rm line}.
\]
Then, under the closed integer-radius convention, radius $1-a_0/n$ is unsafe and
radius $1-(a_0+1)/n$ is safe.  As a real-radius supremum, the threshold is
$1-a_0/n$ and is not attained if the closed ball at that endpoint is counted.
\end{proposition}

\begin{proof}
The function $B_{\rm mca}(a)$ is nonincreasing in $a$, since raising the
required agreement only removes witnesses.  At the closed radius $1-a/n$, the
MCA error is $B_{\rm mca}(a)/q_{\rm line}$.  The two displayed inequalities are
therefore exactly the unsafe/safe transition at consecutive integer agreement
levels.
\end{proof}

\subsection{Conditional MCA one-step criterion}

The following theorem is a deliberately explicit proof blueprint.  Each
hypothesis is a concrete mathematical task, not a black-box ``large field''
assumption.

\begin{theorem}[conditional MCA threshold theorem]\label{thm:conditional-mca}
Fix a smooth multiplicative row $C=\RS[F,D,k]$, a line field $F_{\rm line}$, and
a target $\varepsilon^*=2^{-128}$.  Suppose the following four inputs are proved
for this row, uniformly in the received line:
\begin{enumerate}[label=(\roman*),leftmargin=2.2em]
\item \textbf{Quotient-profile lower and upper ledger.}  For every divisor
$c\mid n$ and every agreement $a=mc+s$ with $0\le s<c$, the quotient-periodic
bad slopes, including remainder supports consisting of $m$ full fibers plus
$s$ residual points, are counted with overlaps across divisors accounted for.
Write the resulting divisor-lattice upper numerator as $B_Q^{\rm all}(a)$;
individual divisor contributions $B_Q(a;c)$ may also be printed, but the upper
ledger must count their union or a certified safe sum rather than only their
maximum.
\item \textbf{Aperiodic residue-line packing.}  After removing tangent and
quotient-periodic locators, the Hankel/residue-line incidence variety has at
most $B_{\rm ap}(a)$ bad line parameters.
\item \textbf{Extension-line accounting.}  If $F_{\rm line}$ is a proper
extension of the field generated by $D$, either all genuinely extension-valued
line witnesses are included in $B_{\rm ap}(a)$, or a separate extension term
$B_{\rm ext}(a)$ is proved.
\item \textbf{No missing line family.}  The line family in the theorem is the
same finite-slope or projective-slope family used by the challenge or protocol
under consideration.
\end{enumerate}
Assume moreover that the resulting upper numerator
\[
        B^+(a):=B_{\rm tan}(a)+B_Q^{\rm all}(a)
                +B_{\rm ap}(a)+B_{\rm ext}(a)
\]
and a corresponding lower numerator $B^-(a)$ satisfy
\[
        B^-(a_0)>\varepsilon^* q_{\rm line},
        \qquad
        B^+(a_0+1)\le\varepsilon^* q_{\rm line}
\]
for some $a_0$.  Then $\delta_C^*(\varepsilon^*)$ is determined at the
one-step accuracy of \cref{prop:onestep}.
\end{theorem}

\begin{proof}
The four hypotheses classify every possible MCA-bad line parameter into one of
four buckets: tangent, quotient-periodic, aperiodic, or genuinely extension-line.
The quotient bucket is counted by the divisor-lattice union ledger $B_Q^{\rm all}$,
so overlaps or simultaneous quotient descriptions cannot invalidate the upper
bound.  The displayed $B^+(a)$ is therefore an upper bound for $B_{\rm mca}(a)$, while
$B^-(a)$ is a certified lower bound.  The conclusion is exactly
\cref{prop:onestep}.
\end{proof}

The present paper supplies part of the lower numerator $B^-(a)$ in the
near-capacity band: quotient-fiber list mass forces MCA mass above the printed
universal caps.  It does not prove the aperiodic upper input (ii), and therefore
does not by itself determine the exact staircase.

\subsection{Conditional interleaved-list one-step criterion}

The grand list-decoding challenge asks for the largest $\delta$ such that an
interleaved list, divided by the field used for the relevant challenge or list
ledger, is below $2^{-128}$.  In the simplest same-field form this is
\[
        |\Lambda(C^{\equiv m},\delta)|\le 2^{-128}|F| .
\]
The quotient-fiber construction in this paper already gives lower floors for
ordinary lists, and ordinary lists embed into interleaved lists by placing the
same received word in one row and zero words in the others.  On the MCA side,
\cref{sec:aperiodic-hankel-certificates} supplies a certified eliminant route
for aperiodic Hankel packing, but it does not prove that the needed eliminants
always exist.  The missing positive list direction is an arbitrary-word
locator-fiber theorem.

\begin{theorem}[conditional interleaved-list threshold theorem]\label{thm:conditional-list}
Fix an interleaving arity $m$.  Suppose that for every received word $U$ and
agreement $a=k+\sigma$ one has a base-code list bound of the form
\[
        |\Lambda(C,1-a/n,U)|
        \le
        L_Q(a,U)+L_{\rm ap}(a),
\]
where $L_Q$ is an explicit quotient-profile list contribution and
$L_{\rm ap}$ is a uniform aperiodic locator-fiber bound.  Suppose also that the
same bound is available at the higher agreement levels that occur in the
codegree recursion, starting with $2a-k$.  Then the interleaved-list threshold is
obtained by applying the codegree recursion row by row and comparing the
resulting numerator to $2^{-128}|F|$.
\end{theorem}

\begin{proof}
For two rows, split the second-row list according to whether its agreement set
has size below or above $2a-k$.  Below $2a-k$, two first-row codewords would
agree on more than $k$ positions and hence are equal; above $2a-k$, use the
base-code list bound at agreement $2a-k$.  Iterating gives the stated $m$-row
recursion.  This is the same codegree mechanism used in the experimental ledger;
the only missing input is the arbitrary-word locator-fiber bound.
\end{proof}

\subsection{Residual ingredients for one-step threshold certificates}

The one-step criteria above consume the following residual ingredients, which we
record with their current status.  The lower-floor half of the periodic-support
ledger is supplied in scanner-ready form by
\cref{lem:heaviest-prefix-locator-floor,thm:quotient-remainder-deep-floor}: every
quotient-plus-remainder agreement $a=mc+s$ has a generated-field prefix-fiber
floor, an image-size fallback, an ambient-box fallback, and a converted CA/MCA
lower numerator.  The support-family upper half is exact at the divisor-lattice
level by \cref{thm:exact-divisor-block-support-ledger}: overlaps between quotient
descriptions at different divisor scales are counted by a common-block
coefficient extraction rather than by a raw safe sum.  The distinct-parameter
branch image is exact inside every declared quotient family by the lcm ledger of
\cref{thm:exact-quotient-image-lcm-ledger,cor:affine-line-lcm-ledger,cor:projective-line-lcm-ledger}; duplicate finite slopes, projective slopes, and
finite curve parameters are coalesced before the numerator is printed.  The
residual-band analysis is split into four certified pieces.  The regular overdetermined nonsingular branch
is canonical by \cref{thm:canonical-affine-rankdrop-ledger,cor:canonical-regular-closed-ball-hankel-packing,thm:canonical-curve-rankdrop-ledger}.  The underdetermined
and rank-drop singular branches are covered by the rank-pivot certificate theorem
\cref{thm:residual-rank-pivot-cover,thm:scanner-checkable-residual-aperiodic-ledger}.  The new compressed-residual layer adds the kernel-compressed eliminant construction of
\cref{thm:kernel-compressed-affine-eliminant,cor:kernel-compressed-projective-curve-eliminants} and the exact arbitrary-residual support-image lcm ledger of
\cref{thm:exact-arbitrary-residual-image-ledger,cor:residual-exact-image-after-paid-removal}.  Thus a residual chart must now be discharged by a small compressed
eliminant, by an exact residual image polynomial, by a Weil-restricted extension
projection bound, or by a named structural obstruction.  Extension-line
accounting is handled on the safe side by
\cref{thm:extension-line-dimension-degree-ledger}: after choosing a $B$-basis of
$F$, each extension chart prints a parameter-projection degree and dimension over
$B$, or an exact zero-dimensional point count.  The curve-sampler ledger is
upgraded for finite-parameter curves and projective lines in parallel: quotient branches have exact gcd/lcm
image ledgers, regular curve rank-drop branches have canonical maximal-minor gcd
ledgers, and residual curve branches have both coefficient-pivot rank charts and
kernel-compressed curve eliminants.  What remains open is no longer a bookkeeping
ambiguity: prove that the compressed residual eliminants are small enough in the
near-capacity rows, or classify any large residual image as quotient,
tangent/common-line, base-field exceptional, genuinely extension-valued, or a new
obstruction.

\begin{enumerate}[label=\textbf{T\arabic*.},leftmargin=2.5em]
\item \textbf{Remainder quotient-profile theorem.}  Generalize the fiber lemma
from full quotient-fiber unions to all agreements $a=mc+s$.  The target is an
exact lower and upper quotient ledger over the divisor lattice of $D$.  The present quotient-ledger package supplies the scanner-ready lower half via heaviest prefix fibers and the
support-union upper half via the exact divisor-block coefficient formula; inside
a declared quotient family, the gcd/lcm image ledger also coalesces duplicate
finite, projective, and curve parameters.  What remains is the structural
exhaustion theorem saying that the declared quotient family captures all
quotient-periodic bad parameters in the intended regime, with the subfield
denominator $|B|$ printed separately from the line denominator $q$.

\item \textbf{Aperiodic Hankel-packing theorem.}  Use the syndrome/Hankel normal
form of \cite{Cho26b} to prove that, after quotient-periodic and tangent
locators are removed, the number of remaining bad line parameters is at most
$B_{\rm ap}(a)$.  The regular overdetermined nonsingular branch is now
canonical by \cref{thm:canonical-affine-rankdrop-ledger,cor:canonical-regular-closed-ball-hankel-packing}: a scanner takes the gcd of all nonzero maximal
Hankel minors and then an lcm over exact agreement levels.  The underdetermined
and rank-drop singular buckets now have the rank-pivot cover of
\cref{thm:residual-rank-pivot-cover}, projective-infinity and curve coefficient
pivots, and the kernel-compressed eliminant ideals of
\cref{thm:kernel-compressed-affine-eliminant,cor:kernel-compressed-projective-curve-eliminants}.  The remaining proof problem is constructive and structural:
produce small eliminants for these compressed residual charts, use the exact
residual image lcm of \cref{thm:exact-arbitrary-residual-image-ledger} when
finite enumeration is feasible, or classify every large residual component as
quotient, tangent/common-line, base-field exceptional, genuinely extension-valued,
or a new obstruction.

\item \textbf{Extension-line lift or counterexample.}  For $B\subset F$ and
$D\subset B$, either prove that genuinely $F$-valued lines are covered by the
aperiodic ledger with a stable numerator, or exhibit explicit extension-valued
counterexamples.  \Cref{cor:Fvalued} shows that deployed failures certified by
this paper must be genuinely extension-valued; \cref{cor:extension-pole-deep-list-floor,cor:extension-pole-quotient-remainder-floor}
provide an explicit lower-side mechanism from any large list by taking the pole
in $F\setminus B$.  On the safe side,
\cref{thm:extension-line-dimension-degree-ledger} gives the certificate format:
Weil-restrict each residual extension chart to $B$, print a parameter-projection
degree/dimension pair or an exact zero-dimensional point count, and divide by the
line field size.  The remaining structural problem is to prove that all genuinely
extension-valued residual charts have small projected dimension/degree or to
exhibit the chart that violates this expectation.

\item \textbf{Arbitrary-word locator-fiber theorem.}  Prove the base-code list
bound needed by \cref{thm:conditional-list}, first for monomial-prefix received
words and quotient-free dimensions, then for arbitrary received words with
quotient cores budgeted.

\item \textbf{Curve and projective-line ledgers.}  Extend the affine-line Hankel
classification to projective slopes and to finite-parameter power curves
$f_0+\gamma f_1+\cdots+\gamma^d f_d$.  The projective and curve ledger package gives exact gcd/lcm ledgers
for projective and finite-curve quotient branches, canonical rank-drop gcds for
regular overdetermined finite-curve buckets, projective-infinity residual charts,
coefficient-pivot curve rank charts, and kernel-compressed projective/curve
eliminant certificates.  The remaining work is to prove smallness or structural
classification for the residual projective and curve charts, including
exceptional vanishing-vector strata and genuinely extension-valued components.

\item \textbf{Certificate scanner.}  Implement a finite staircase scanner that,
for each challenge row, prints tangent, quotient, generated-field entropy,
aperiodic, extension, curve, and interleaved-list numerators at every agreement
$a$.  The desired output is an interval
\[
        a_{\rm unsafe}<a_C^*\le a_{\rm safe}
\]
that shrinks to one step when \cref{thm:conditional-mca} or
\cref{thm:conditional-list} applies.

\item \textbf{Formal finite core.}  Formalize the self-contained deep-point
conversion, the quotient-fiber pigeonhole, support-wise MCA monotonicity, and
the integer staircase lemma.  These are small enough for proof-assistant
checking and directly strengthen a formal certificate package.

\item \textbf{Circle-row completion and genus-one floors.}  Extend the
quotient-support and quotient-image upper ledgers (item~1) and the aperiodic
Hankel classification (item~2) from power locators on multiplicative divisor
lattices to Chebyshev locator towers on $x$-coordinate twin-coset domains,
where the divisor lattice of $D$ is replaced by the lattice of dyadic
Chebyshev scales; and settle the genus-one lower-floor question of
Open Problem~\ref{prob:isogeny}, for which the lacunarity collapse of
\cref{rem:ecfft-obstruction} blocks the present route beyond the first
staircase step.
\end{enumerate}

\subsection{Scope of claims}

We state precisely what this paper establishes and what it does not. Established: a universal unsafe band at capacity for the grand MCA challenge, over smooth multiplicative rows in the stated field-size envelope (\cref{cor:grand,cor:mixed-radix}) and, by the later sections, over Chebyshev, circle, and genus-one rows as well (\cref{cor:circle-grand,cor:ecfft-macroscopic,cor:circle-anyfield}); an unconditional safe region below one third of the distance (\cref{thm:deep-mca}) with a list-side safe region up to the Johnson radius, uniform in the interleaving arity (\cref{thm:johnson-list}); a safe region up to one half of the distance modulo a single classical import (\cref{thm:mca-from-ca,cor:conditional-half}); and the resulting sandwich, which determines the grand threshold of every deployed row within a factor smaller than two modulo that import, and within a factor of three unconditionally (\cref{thm:sandwich,thm:sandwich2}, \cref{tab:scanner2}). Not established, and not claimed: the exact threshold, a positive theorem inside the remaining open band, protocol soundness, or a curve-MCA theorem. What separates the sandwich from a one-step determination is precisely Open Problem~\ref{prob:band}.

\section{Structural ledgers, residual bands, and certificate machinery}\label{sec:discharge}

This section collects the technical infrastructure used by the two-sided threshold statements, organized by mathematical function.  Stabilizer and quotient-descent lemmas isolate the periodic supports and make their numerators field-size independent.  The deep-regime theorem controls MCA directly below one third of the distance and removes the singular Hankel branches there.  The restriction-of-scalars lift shows that extension-valued lines are not a separate class of witnesses.  The Johnson and recursion statements control the list side outside the band.  The curve theorem covers projective and finite-parameter samplers in the deep regime.  Finally, the scanner and finite-certificate grammar package all deployed verdicts as exact integer comparisons.  The only mathematical component not closed by these ledgers is the aperiodic band problem, Open Problem~\ref{prob:band} below.

\subsection{Stabilizer exhaustion and exact quotient descent}\label{sec:discharge-T1}

Throughout this subsection, \emph{multiplicative case} means: $D=\alpha H$ is a multiplicative coset of a cyclic group $H$ of order $n$, acting freely on $D$ by multiplication; for $c\mid n$, $H_c\le H$ is the unique subgroup of order $c$, equal to the full group $\mu_c$ of $c$-th roots of unity in $\B$, and the fibers of $\pi_c:x\mapsto x^c$ on $D$ are exactly the $H_c$-orbits. \emph{Chebyshev case} means: $D=\chi(\mathcal D)$ for a twin coset $\mathcal D=gH\cup g^{-1}H$ of size $2M$ with $\operatorname{ord}(g)\nmid2M$ as in \cref{sec:circle-geometry}; for $c\mid M$, $K_c\le\UU$ is the order-$c$ kernel of the $c$-th power map, and for $S\subseteq D$ we write $\widetilde S:=\chi^{-1}(S)\cap\mathcal D$ for the (automatically inversion-closed) lift. By \cref{lem:cheb-fibers}, the $T_c$-fibers in $D$ are exactly the $\chi$-images of the $\langle K_c,\,u\mapsto u^{-1}\rangle$-orbits on $\mathcal D$. In both cases the \emph{scale lattice} is the divisor lattice of $n$ (resp.\ of $M$).

\begin{definition}[periodicity scale of a support]\label{def:periodicity-scale}
For $S\subseteq D$, the periodicity scale $c(S)$ is $|\{h\in H:\ hS=S\}|$ in the multiplicative case and $|\{\zeta\in K_M:\ \zeta\widetilde S=\widetilde S\}|$ in the Chebyshev case. A support is \emph{aperiodic} if $c(S)=1$.
\end{definition}

\begin{theorem}[stabilizer partition and structural exhaustion]\label{thm:stabilizer-partition}
In either case:
\begin{enumerate}[label=\textup{(\roman*)}]
\item The stabilizer in the definition is a subgroup of a cyclic group, so $c(S)$ divides $n$ \textup(resp.\ $M$\textup), and for every scale $c$,
\[
S\ \text{is a union of complete scale-$c$ fibers}
\quad\Longleftrightarrow\quad
c\mid c(S).
\]
In particular $c(S)$ is the unique maximal such scale, and the supports of each size are partitioned by the value of $c(S)$, with $c(S)=1$ the aperiodic class.
\item Fix a received line and an agreement threshold $a\ge k$, and call a bad parameter \emph{periodically witnessed} if at least one of its witness supports $S$ \textup(in the sense of \cref{def:ca,def:mca}\textup) has $c(S)\ge2$. Then every bad parameter is either periodically witnessed --- and hence counted, with multiplicity one per support, by the family of supports of size $\ge a$ that are unions of complete fibers at some scale $\ge2$ --- or all of its witness supports are aperiodic, in which case it belongs to the aperiodic bucket by definition. The two classes are exhaustive and disjoint.
\end{enumerate}
\end{theorem}

\begin{proof}
(i) The set $\{h\in H:hS=S\}$ is a subgroup of the cyclic group $H$ (resp.\ $\{\zeta\in K_M:\zeta\widetilde S=\widetilde S\}\le K_M$), hence equals $H_{c_0}$ (resp.\ $K_{c_0}$) for a unique divisor $c_0$, which is $c(S)$. In the multiplicative case, the $\pi_c$-fibers in $D$ are the $H_c$-orbits, and a set is a union of complete $H_c$-orbits if and only if it is $H_c$-invariant, i.e.\ $H_c\le H_{c_0}$, i.e.\ $c\mid c_0$. In the Chebyshev case, $S$ is a union of complete $T_c$-fibers if and only if $\widetilde S$ is a union of complete $\langle K_c,\mathrm{inv}\rangle$-orbits; since $\widetilde S$ is inversion-closed automatically, this holds if and only if $\widetilde S$ is $K_c$-invariant, i.e.\ $c\mid c_0$. The partition statement is immediate.

(ii) Exhaustiveness and disjointness are definitional given (i). The counting clause is \cref{lem:one-support-one-line} applied to the support family $\{S:\ |S|\ge a,\ c(S)\ge2\}$: each fixed support pays for at most one finite line parameter, so the number of periodically witnessed bad parameters is at most the size of that family.
\end{proof}

\begin{corollary}[exact size of the periodic support family; $q$-free numerator]\label{cor:periodic-support-count}
Let $P$ be the set of primes dividing $n$ \textup(resp.\ $M$\textup). The number of supports of exact size $A$ with $c(S)\ge2$ is
\[
\Pi(A)\ :=\ \sum_{\emptyset\ne T\subseteq P}(-1)^{|T|+1}\,
\mathbf 1\bigl[\operatorname{lcm}(T)\mid A\bigr]\,
\binom{\,n/\operatorname{lcm}(T)\,}{\,A/\operatorname{lcm}(T)\,},
\]
and consequently, for every received line and every agreement threshold $a\ge k$,
\[
\#\{\text{periodically witnessed bad finite parameters}\}\ \le\ \sum_{A=a}^{n}\Pi(A),
\]
with an extra factor $d$ for degree-$d$ power-curve samplers \textup(\cref{lem:one-support-d-curve}\textup). The numerator is independent of $q$ and of $|\B|$; the subfield order enters only the lower floors, as the periodic-support ledger demands.
\end{corollary}

\begin{proof}
$c(S)\ge2$ if and only if $H_pS=S$ for some $p\in P$, i.e.\ $S$ lies in the union over $p\in P$ of the families of complete-$p$-fiber unions. For $T\subseteq P$, invariance under every $H_p$, $p\in T$, is invariance under the subgroup they generate, which is $H_{\operatorname{lcm}(T)}$; the number of unions of complete $\operatorname{lcm}(T)$-fibers of total size $A$ is $\binom{n/\operatorname{lcm}(T)}{A/\operatorname{lcm}(T)}$ when $\operatorname{lcm}(T)\mid A$ and $0$ otherwise. Inclusion--exclusion gives $\Pi(A)$; the parameter bound is \cref{thm:stabilizer-partition}(ii) \textup(resp.\ \cref{lem:one-support-d-curve}\textup). The Chebyshev case is identical with $K_p\le K_M$.
\end{proof}

This complements \cref{thm:exact-divisor-block-support-ledger}: that formula counts a \emph{declared} remainder-profile family exactly; \cref{cor:periodic-support-count} counts the \emph{intrinsic} periodic family exactly, and \cref{thm:stabilizer-partition} proves that nothing quotient-periodic escapes it. What the periodic-support ledger still lacked was the passage from support-level to image-level counting. The next theorem supplies it exactly on the class of received data for which the quotient family was designed.

\begin{theorem}[exact quotient descent for fiber-constant data]\label{thm:fiber-descent}
Work in the multiplicative or Chebyshev case, write $\varphi_c$ for the scale-$c$ map \textup($X^c$, resp.\ $T_c$\textup) and $Q_c:=\varphi_c(D)$, $N_c:=n/c$. Let $c$ be a scale, let $p\in\F[X]_{<k}$, let $S\subseteq D$ be a union of complete $c$-fibers with $|S|\ge k$, and suppose $p$ agrees on $S$ with a fiber-constant word $W\circ\varphi_c$. Then there is a unique $g\in\F[Y]$ with
\[
p=g\circ\varphi_c,\qquad \deg g<\lceil k/c\rceil,
\]
and $g$ agrees with $W$ on $\varphi_c(S)$, a set of $|S|/c$ points of $Q_c$. Consequently, for a received line $f+zg_0$ with $f$ and $g_0$ both constant on $c$-fibers, and any agreement threshold $a\ge k$: a finite slope $z$ is MCA-bad with some witness support that is a union of complete $c$-fibers if and only if $z$ is MCA-bad for the quotient line $(f_c,g_{0,c})$ on $\RS[\F,Q_c,\lceil k/c\rceil]$ at the corresponding agreement $\ge\lceil a/c\rceil$, with witness $\varphi_c(S)$; and the same holds verbatim for lists: $p\mapsto g$ is a bijection between scale-$c$-witnessed list elements and quotient list elements.
\end{theorem}

\begin{proof}
\emph{Multiplicative case.} For $h\in H_c=\mu_c$ and $x\in S$ we have $hx\in S$ and $\varphi_c(hx)=\varphi_c(x)$, so $p(hx)=W(\varphi_c(hx))=W(\varphi_c(x))=p(x)$. Thus $p(hX)-p(X)$, of degree $<k$, has at least $|S|\ge k$ roots, hence vanishes identically: $p(hX)=p(X)$ for every $h\in\mu_c$. Comparing coefficients, $a_ih^i=a_i$ for all $h\in\mu_c$, so $a_i=0$ whenever $c\nmid i$ (take $h$ a generator of $\mu_c$), i.e.\ $p=g(X^c)$ with $c\deg g=\deg p<k$, so $\deg g<\lceil k/c\rceil$; uniqueness is clear. For $x\in S$, $g(\varphi_c(x))=p(x)=W(\varphi_c(x))$, and $\varphi_c(S)$ has $|S|/c$ elements since $S$ is a union of complete $c$-fibers.

\emph{Chebyshev case.} Work over the compositum $\F':=\F\cdot\F_{p^2}$ and put $P(u):=p(\chi(u))\in\F'[u,u^{-1}]$, a Laurent polynomial with exponents in $[-(k-1),k-1]$. For $\zeta\in K_c=\mu_c$ and $u\in\widetilde S$ we have $\zeta u\in\widetilde S$ and $T_c(\chi(\zeta u))=\chi(\zeta^cu^c)=T_c(\chi(u))$, so $P(\zeta u)=P(u)$ on $\widetilde S$; since $|\widetilde S|=2|S|\ge2k$ exceeds the degree $2(k-1)$ of $u^{k-1}\bigl(P(\zeta u)-P(u)\bigr)$, the identity $P(\zeta u)=P(u)$ holds in $\F'[u,u^{-1}]$. Hence only exponents divisible by $c$ occur, and $P$ is also inversion-symmetric because $\chi$ is, so $P$ lies in the symmetric Laurent ring $\F'[u^c+u^{-c}]$: $P(u)=G(u^c+u^{-c})$ for a unique $G$. Since $u^c+u^{-c}=2\,T_c(\chi(u))$, setting $g(Y):=G(2Y)$ gives $p(\chi(u))=g(T_c(\chi(u)))$ identically, and injectivity of the ring map $r(x)\mapsto r((u+u^{-1})/2)$ yields $p=g\circ T_c$ in $\F'[X]$; as $g$ is the unique solution of an $\F$-linear system with data in $\F$, in fact $g\in\F[Y]$, and $c\deg g=\deg p<k$ as before. Agreement on $\varphi_c(S)=T_c(S)$ follows as above, with $|T_c(S)|=|S|/c$ by \cref{lem:cheb-fibers}.

\emph{Lines and lists.} If $z$ has an MCA witness $S$ that is a union of complete $c$-fibers, the explaining codeword for the point $f+zg_0$ (a fiber-constant word) descends by the above; conversely a quotient explanation composes with $\varphi_c$. For noncontainment, if the pair $(f,g_0)$ were explained on $S$ by $(p_1,p_2)$, both $p_i$ descend (each agrees on $S$ with a fiber-constant word), so the quotient pair would be explained on $\varphi_c(S)$; and conversely quotient explanations compose. Hence badness with a scale-$c$ witness corresponds exactly. For lists, $p\mapsto g$ is well defined and injective (uniqueness of $g$), and $g\mapsto g\circ\varphi_c$ is its inverse from the quotient list, since the full agreement set of $g\circ\varphi_c$ with $W\circ\varphi_c$ is the full $\varphi_c$-preimage of the agreement set of $g$ with $W$.
\end{proof}

\begin{corollary}[structural exhaustion is exact]\label{cor:T1-status}
With the intrinsic \cref{def:periodicity-scale}, the structural exhaustion is a theorem: every bad parameter is quotient-periodic \textup(periodically witnessed, hence counted exactly at the support level by \cref{cor:periodic-support-count}, and exactly at the image level by \cref{thm:fiber-descent} whenever the received line is fiber-constant\textup) or belongs to the aperiodic bucket. The subfield order $|\B|$ appears only in lower floors; all upper numerators above are $q$-free. The entire residual content of the periodic-support ledger is thereby the aperiodic bucket, treated next.
\end{corollary}

\subsection{The deep regime, singular pencils, and the residual band}\label{sec:discharge-T2}

The following theorem is the safe-side counterpart of the paper's unsafe caps. It is self-contained, uses nothing beyond interpolation and double counting, and holds for arbitrary linear codes --- in particular for every multiplicative row, every Chebyshev line-round row, every bivariate circle code over any field \textup(including fields not containing $i$\textup), and every ECFFT row.

\begin{theorem}[deep-regime mutual correlated agreement, unconditional]\label{thm:deep-mca}
Let $C\subseteq\F^n$ be any $\F$-linear code of minimum Hamming weight $w$, let $q=|\F|$, let $\delta\in(0,1)$, and put $r:=\lfloor\delta n\rfloor$. Assume
\[
3r\ \le\ w-1 .
\]
Then for every pair $(f_1,f_2)\in\F^n\times\F^n$, at least one of the following alternatives holds:
\begin{enumerate}[label=\textup{(\alph*)}]
\item \textup{(Tangent pair.)} There are codewords $c_1,c_2\in C$ and a set $T'\subseteq\{1,\dots,n\}$ with $|T'|\le r$ such that $f_i=c_i$ off $T'$ for $i=1,2$; in this case $(f_1,f_2)$ has no CA-bad slope at radii $(\delta,\delta_{\mathrm{int}})$ for any $\delta_{\mathrm{int}}\ge\delta$, and at most $|T'|\le r$ MCA-bad slopes at radius $\delta$, namely a subset of $\{-e_{1,j}/e_{2,j}:\,j\in T',\ e_{2,j}\ne0\}$ where $e_i:=f_i-c_i$.
\item \textup{(Generic pair.)} $\#\{\gamma\in\F:\ \dist(f_1+\gamma f_2,\,C)\le\delta\}\ \le\ r+1$.
\end{enumerate}
Consequently
\[
\eca(C,\delta)\ \le\ \frac{r+1}q,
\qquad
\emca(C,\delta)\ \le\ \frac{r+1}q .
\]
For $C=\RS[\F,D,k]$ on any evaluation set $D$ of size $n$, the hypothesis reads $3\lfloor\delta n\rfloor\le n-k$.
\end{theorem}

\begin{proof}
Write $S_{\rm cl}:=\{\gamma:\dist(f_1+\gamma f_2,C)\le\delta\}$; note $\dist\le\delta$ means disagreement on at most $r$ coordinates, by integrality. If $|S_{\rm cl}|\le1$, alternative (b) holds ($r+1\ge1$). Otherwise fix $\gamma_1\ne\gamma_2$ in $S_{\rm cl}$ with codewords $p_1,p_2$ and error sets $E_1,E_2$ ($|E_i|\le r$), and set
\[
c_2:=\frac{p_1-p_2}{\gamma_1-\gamma_2},
\qquad
c_1:=p_1-\gamma_1c_2,
\qquad
e_i:=f_i-c_i .
\]
Off $T:=E_1\cup E_2$ ($|T|\le2r$) we can solve the two linear equations $f_1+\gamma_if_2=p_i$ coordinatewise, giving $f_2=c_2$ and $f_1=c_1$ there; hence $e_1,e_2$ are supported in $T$. Let $T':=\{j:(e_{1,j},e_{2,j})\ne(0,0)\}\subseteq T$.

\emph{Step 1: every close slope rides the recovered line.} Let $\gamma\in S_{\rm cl}$ with codeword $p$ and error set $E$, $|E|\le r$. Off $T\cup E$ we have $p=f_1+\gamma f_2=c_1+\gamma c_2$, so the codeword $p-(c_1+\gamma c_2)$ has weight at most $|T|+|E|\le3r\le w-1<w$ and is therefore zero. Hence
\[
f_1+\gamma f_2-(c_1+\gamma c_2)\ =\ e_1+\gamma e_2
\quad\text{has weight at most } r
\qquad\text{for every }\gamma\in S_{\rm cl}. \tag{$*$}
\]

\emph{Step 2: dichotomy.} Suppose $|T'|\ge r+1$. For $j\in T'$, the affine function $\gamma\mapsto e_{1,j}+\gamma e_{2,j}$ is not identically zero, so it vanishes for at most one $\gamma$. By $(*)$, each $\gamma\in S_{\rm cl}$ needs at least $|T'|-r\ge1$ vanishing coordinates inside $T'$. Double counting vanishings,
\[
|S_{\rm cl}|\,\bigl(|T'|-r\bigr)\ \le\ |T'|,
\qquad\text{hence}\qquad
|S_{\rm cl}|\ \le\ \frac{|T'|}{|T'|-r}\ \le\ r+1,
\]
which is alternative (b). Otherwise $|T'|\le r$ and $(f_1,f_2)$ is a tangent pair as in (a) with this $T'$.

\emph{Step 3: bad slopes of tangent pairs.} Let $(f_1,f_2)$ be tangent with data $(c_1,c_2,T')$, $|T'|\le r$. First, $\dist_2((f_1,f_2),C^{\equiv2})\le|T'|/n\le r/n\le\delta\le\delta_{\mathrm{int}}$, so no slope is CA-bad at radii $(\delta,\delta_{\mathrm{int}})$. For MCA, note $2r\le w-1$ (from $3r\le w-1$ and $r\ge0$), so two codewords agreeing on at least $n-2r$ coordinates are equal. Let $\gamma$ be MCA-bad with witness $S$, $|S|\ge n-r$. On $S\setminus T'$, which has at least $n-2r$ elements, $f_1+\gamma f_2=c_1+\gamma c_2$; hence the codeword explaining the point on $S$ equals $c_1+\gamma c_2$, and thus $e_1+\gamma e_2$ vanishes on $S\cap T'$. Moreover, any pair explanation $(p'_1,p'_2)$ on $S$ satisfies $p'_i=c_i$ on the $\ge n-2r$ points of $S\setminus T'$, forcing $p'_i=c_i$; so the pair fails to be explained on $S$ exactly when $S\cap T'\ne\emptyset$ (every $j\in T'$ has $(e_{1,j},e_{2,j})\ne0$). Therefore $\gamma$ is MCA-bad if and only if $e_{1,j}+\gamma e_{2,j}=0$ for some $j\in T'$: necessity was just shown; for sufficiency take $S:=(\{1,\dots,n\}\setminus T')\cup\{j\}$, of size $\ge n-r$, on which $c_1+\gamma c_2$ explains the point while the pair is not explained. Each $j\in T'$ contributes at most one slope, giving at most $|T'|\le r$ MCA-bad slopes.

Finally, generic pairs have at most $r+1$ close slopes, hence at most $r+1$ CA- or MCA-bad slopes (an MCA witness of size $\ge n-r$ implies $\dist\le r/n\le\delta$); tangent pairs have $0$ CA-bad and at most $r$ MCA-bad slopes. Taking the maximum over pairs gives the two displayed error bounds. For $\RS[\F,D,k]$, $w=n-k+1$, so $3r\le w-1$ is $3r\le n-k$.
\end{proof}

\begin{remark}[two radii; scope]
The CA statement holds in the two-radius form: for any $\delta_{\mathrm{int}}\ge\delta_{\mathrm{fld}}=\delta$ with $3\lfloor\delta n\rfloor\le w-1$, $\eca(C,\delta,\delta_{\mathrm{int}})\le(r+1)/q$, since the tangent alternative kills the far condition at every $\delta_{\mathrm{int}}\ge\delta$. The constant $\frac13$ can be improved toward $\frac12$ of the distance by the methods of \cite{BCHKS25}; we deliberately keep the three-line-per-step, formalization-ready proof --- it feeds the formal finite core --- and record the import separately. Note also that \cref{thm:deep-mca} applies to the bivariate circle code over $\F_p$ itself \textup(minimum weight $n_c-k_c+1$, by extending scalars to $\F_{p^2}$ and using \cref{lem:circle-rs}\textup), so the deep regime of Open Problem~\ref{prob:isogeny}\textup{(b)} is closed as well.
\end{remark}

We next resolve a named obstruction of the residual-band analysis --- the singular Hankel buckets --- in the regime where \cref{thm:deep-mca} applies. First, the structural normal form of the identically singular branch.

\begin{lemma}[identically singular pencils are locator families]\label{lem:singular-pencil}
In the setting of \cref{sec:aperiodic-hankel-certificates}, fix an exact agreement $A$ with $t_A\ge j_A+1$ and set $M(Z):=H_{t_A,j_A}(u)+ZH_{t_A,j_A}(v)$. The following are equivalent:
\begin{enumerate}[label=\textup{(\roman*)}]
\item every $(j_A+1)\times(j_A+1)$ minor of $M(Z)$ vanishes identically in $Z$;
\item there is a nonzero polynomial vector $\widetilde v(Z)=\sum_{i=0}^{d}Z^iv_i\in\F[Z]^{j_A+1}$ with $M(Z)\widetilde v(Z)\equiv0$, equivalently a chain
\[
H_{t_A,j_A}(u)\,v_0=0,\qquad
H_{t_A,j_A}(u)\,v_{i}=-H_{t_A,j_A}(v)\,v_{i-1}\ (1\le i\le d),\qquad
H_{t_A,j_A}(v)\,v_d=0 .
\]
\end{enumerate}
In that case, writing $\lambda(Z)$ for the last coordinate of $\widetilde v(Z)$, every $Z=\gamma$ with $\lambda(\gamma)\ne0$ admits a monic degree-$j_A$ syndrome annihilator for $u+\gamma v$, i.e.\ every such slope passes the rank test at exact agreement $A$; the at most $\deg\lambda$ remaining slopes pass it after column truncation. The rank test is necessary for badness at agreement $A$ but not sufficient \textup(the annihilator must in addition be $D$-split\textup), so membership in the singular bucket never certifies badness by itself.
\end{lemma}

\begin{proof}
(i) says $\operatorname{rank}_{\F(Z)}M(Z)\le j_A$, which holds if and only if the kernel of $M(Z)$ over the rational function field is nonzero; clearing denominators and dividing by common factors produces a nonzero polynomial kernel vector, and expanding $M(Z)\widetilde v(Z)\equiv0$ in powers of $Z$ gives the chain. The converse direction is immediate. Specializing $Z=\gamma$ with $\lambda(\gamma)\ne0$ and normalizing the last coordinate to $1$ produces the vector $\widetilde\ell$ of the chart atlas with $A(\ell)+\gamma B(\ell)=0$, i.e.\ a monic annihilator; $\lambda\ne0$ has at most $\deg\lambda$ roots. The final sentence restates the convention of \cref{sec:aperiodic-hankel-certificates}.
\end{proof}

\begin{corollary}[the singular buckets are resolved in the deep regime]\label{cor:singular-deep}
Let $C=\RS[\F,D,k]$ and let $A$ be an exact agreement with $3(n-A)\le n-k$. Then for every received line --- in particular for every line in the underdetermined or identically singular buckets of \cref{sec:aperiodic-hankel-certificates} --- the number of finite slopes that are bad at agreement threshold $A$ is at most $n-A+1$ if the pair is not $\frac{n-A}n$-tangent, and at most $n-A$ MCA-bad \textup(zero CA-bad\textup) slopes if it is. In particular, in this regime the singular bucket contributes at most $n-A+1$ to $B_{\rm ap}(A)$ per line, and no line contributes $\Theta(q)$: the ``totally bad line'' scenario that the singular bucket was guarding against does not occur below one third of the distance.
\end{corollary}

\begin{proof}
Immediate from \cref{thm:deep-mca} with $r=n-A$ and $\delta=r/n$: badness at agreement threshold $A$ implies distance at most $r/n$.
\end{proof}

\begin{remark}[calibration: near-capacity unsafety is quotient-borne]\label{rem:quotient-borne}
The unsafe mass produced by this paper does not press on the aperiodic bucket. The fiber witnesses of \cref{lem:fiber,lem:phi-fiber} have agreement supports that are unions of complete fibers, hence periodicity scale $\ge a\ge2$; by \cref{thm:stabilizer-partition} they are charged to the quotient side of the ledger, and by \cref{thm:fiber-descent} even at the image level. So a polynomially bounded aperiodic numerator in the band remains consistent with every theorem in this paper. This motivates stating the entire residual band as a single conjecture.
\end{remark}

\begin{problem}[aperiodic band conjecture --- the residual kernel]\label{prob:band}
Fix a multiplicative or Chebyshev row $C=\RS[\F,D,k]$. Prove or refute: there is a polynomial $P$, independent of $q$, such that for every agreement $a$ in the open band
\[
k+\tfrac{2n}N\ <\ a\ <\ n-\Bigl\lfloor\frac{n-k}3\Bigr\rfloor ,
\]
the number of MCA-bad finite slopes of any received line, all of whose witness supports are aperiodic in the sense of \cref{def:periodicity-scale}, is at most $P(n)$. By \cref{cor:T1-status,thm:weil-lines} \textup(structural exhaustion and the extension-line lift\textup) and \cref{thm:deep-mca} \textup(deep regime discharged\textup), this conjecture is exactly the missing input \textup{(ii)} of \cref{thm:conditional-mca} throughout the band, for every challenge row.
\end{problem}

\subsection{The exact extension-line lift}\label{sec:discharge-T3}

\begin{theorem}[restriction of scalars for lines]\label{thm:weil-lines}
Let $\B\subseteq\F$ with $e:=[\F:\B]$ and $D\subseteq\B$, and fix a $\B$-basis $1=\beta_0,\beta_1,\dots,\beta_{e-1}$ of $\F$. Decompose words as $w=\sum_i\beta_iw_i$ with $w_i\in\B^n$, and let $\mu_{ij}(\gamma)\in\B$ be the multiplication tensor, $\gamma\beta_i=\sum_j\mu_{ij}(\gamma)\beta_j$, which is $\B$-linear and injective in $\gamma$. Then:
\begin{enumerate}[label=\textup{(\roman*)}]
\item A word $c$ lies in $\RS[\F,D,k]$ if and only if every component $c_j$ lies in $\RS[\B,D,k]$, and for all words $w$,
\[
\dist\bigl(w,\RS[\F,D,k]\bigr)\ =\ \dist_e\bigl((w_0,\dots,w_{e-1}),\,\RS[\B,D,k]^{\equiv e}\bigr).
\]
More generally, a set $S\subseteq D$ explains $w$ over $\F$ if and only if $S$ column-explains the component tuple over $\B$.
\item For a pair $(f_1,f_2)$ over $\F$ and $\gamma\in\F$, the components of $f_1+\gamma f_2$ are
\[
(f_1+\gamma f_2)_j\ =\ f_{1,j}+\sum_{i}\mu_{ij}(\gamma)\,f_{2,i},
\]
an $e$-dimensional $\B$-linear \textup(in $\gamma$\textup) family of $e$-interleaved words over $\B$.
\item For every radius \textup(and pair of radii\textup), $\gamma$ is CA-bad \textup(resp.\ MCA-bad\textup) for $(f_1,f_2)$ over $\F$ if and only if the parameter $\gamma$ of the sampler in \textup{(ii)} is bad in the corresponding interleaved sense over $\B$, with pair noncontainment read on the $2e$-tuple $(f_{1,0},\dots,f_{1,e-1},f_{2,0},\dots,f_{2,e-1})$ against $\RS[\B,D,k]^{\equiv2e}$, on the same supports. For $\gamma\in\B$ the sampler restricts to the diagonal $e$-interleaved line $(f_{1,j}+\gamma f_{2,j})_j$.
\end{enumerate}
Consequently the extension term $B_{\rm ext}(a)$ of \cref{thm:conditional-mca} is not an independent quantity: it equals, definitionally, a base-field interleaved-sampler numerator. In particular no ``extension-valued counterexample'' can exist outside the base ledgers: every genuinely $\F$-valued line is charged, exactly, to a base bucket.
\end{theorem}

\begin{proof}
(i) is the componentwise decomposition already used in the proof of \cref{lem:confine} \textup(valid for any $D\subseteq\B$, cf.\ \cref{cor:circle-Fvalued}\textup): the coefficients of a polynomial over $\F$ decompose in the basis, evaluation at $D\subseteq\B$ commutes with the decomposition, and pointwise agreement over $\F$ is simultaneous componentwise agreement, i.e.\ column agreement. (ii) is the definition of $\mu$. (iii) The point condition, the pair condition, and every support-restricted condition in \cref{def:ca,def:mca} are, by (i), equalities of components on sets of coordinates, and these translate verbatim; injectivity of $\gamma\mapsto\mu(\gamma)$ makes the parameterizations match one-to-one. The diagonal claim is $\mu_{ij}(\gamma)=\gamma\delta_{ij}$ for $\gamma\in\B$.
\end{proof}

\begin{corollary}[status of the extension-line lift]\label{cor:T3-status}
The extension-line question is discharged on the ``lift'' side of its dichotomy: extension lines are covered --- exactly, not merely up to a stable numerator --- by base-field ledgers, namely the interleaved-sampler ledgers of the list and curve sides via \cref{thm:weil-lines}. In the deep regime nothing at all remains: \cref{thm:deep-mca} applies over $\F$ directly, uniformly in the field of definition of the pair. On the lower side, \cref{cor:extension-pole-deep-list-floor,cor:extension-pole-quotient-remainder-floor} already produce genuinely extension-valued bad lines from any supplied list, and \cref{lem:confine} \textup(with \cref{cor:Fvalued,cor:circle-Fvalued}\textup) shows the certified deployed failures must be of this kind. What the extension-line lift leaves behind is exactly the band input of the list and curve ledgers, i.e.\ Open Problem~\ref{prob:band}.
\end{corollary}
\subsection{Uniform interleaved-list bounds, exact recursion, and calibration}\label{sec:discharge-T4}

\begin{theorem}[elementary Johnson list bound, uniform in the interleaving arity]\label{thm:johnson-list}
Let $C=\RS[\F,D,k]$ on any evaluation set $D$ of size $n$, let $m\ge1$, and let $a$ be an agreement threshold with $a^2>(k-1)n$. Then for every received $m$-tuple $U$,
\[
\bigl|\Lst\bigl(C^{\equiv m},\,1-\tfrac an,\,U\bigr)\bigr|
\ \le\
\frac{n\,(a-k+1)}{a^2-(k-1)n}\,,
\]
where the list is taken in the column distance $\dist_m$. The same bound holds for any code family in which two distinct codewords column-agree on at most $k-1$ coordinates; in particular it holds for interleaved bivariate circle codes with $k$ replaced by $k_c$.
\end{theorem}

\begin{proof}
Let $\mathbf c^{(1)},\dots,\mathbf c^{(L)}$ be distinct listed codewords with column-agreement sets $S_1,\dots,S_L\subseteq D$, $|S_i|\ge a$. Two distinct interleaved codewords column-agree exactly where all their rows agree; some pair of rows consists of distinct polynomials of degree $<k$, so $|S_i\cap S_j|\le k-1$ for $i\ne j$. Let $\deg(x):=\#\{i:x\in S_i\}$, so $\sum_x\deg(x)\ge La$ and
\[
\sum_{x\in D}\binom{\deg(x)}2=\sum_{i<j}|S_i\cap S_j|\le\binom L2(k-1).
\]
By convexity of $t\mapsto t(t-1)$, $\sum_x\deg(x)(\deg(x)-1)\ge n\cdot\frac{La}n\bigl(\frac{La}n-1\bigr)$ \textup(the right side is nonpositive when $La\le n$, and the average-value bound applies otherwise\textup). Hence $La(La-n)\le L(L-1)(k-1)n$, i.e.\ $a(La-n)\le(L-1)(k-1)n$, i.e.
\[
L\bigl(a^2-(k-1)n\bigr)\ \le\ an-(k-1)n\ =\ n(a-k+1),
\]
which is the claim when $a^2>(k-1)n$.
\end{proof}

\begin{theorem}[exact divisor-lattice list recursion for fiber-constant words]\label{thm:list-recursion}
Multiplicative case; let $c_0\mid n$, let $U=W\circ\pi_{c_0}$ be constant on $c_0$-fibers, and let $a\ge k$. For $p\in\Lst(C,1-a/n,U)$ let $S_p:=\{x\in D:p(x)=U(x)\}$ be the full agreement set and put $c_p:=\gcd\bigl(c(S_p),c_0\bigr)$. For $c\mid c_0$ write $W_c(y):=W\bigl(y^{c_0/c}\bigr)$ on $Q_c$, so $U=W_c\circ\pi_c$. Then, for each $c\mid c_0$, the descent map of \cref{thm:fiber-descent} is a bijection from $\{p\in\Lst(C,1-a/n,U):\ c_p=c\}$ onto
\[
\Bigl\{\,g\in\Lst\bigl(\RS[\F,Q_c,\lceil k/c\rceil],\,1-\lceil a/c\rceil/N_c,\,W_c\bigr)\ :\ \gcd\bigl(c(\bar S_g),\,c_0/c\bigr)=1\,\Bigr\},
\]
where $\bar S_g$ is the full agreement set of $g$ with $W_c$ and its scale is computed in the quotient lattice. In particular
\[
\bigl|\Lst(C,1-\tfrac an,U)\bigr|=\sum_{c\mid c_0}\#\{\cdots\},
\]
and if moreover $a^2>kn$, then every bucket obeys the level-$c$ Johnson bound and
\[
\bigl|\Lst(C,1-\tfrac an,U)\bigr|\ \le\ \sum_{c\mid c_0}\frac{n\,(a-k+2c-1)}{a^2-kn}\ \le\ \tau(c_0)\,\frac{n\,(a-k+2c_0-1)}{a^2-kn}\,,
\]
$\tau$ the number-of-divisors function.
\end{theorem}

\begin{proof}
Fix $p$ with $c_p=c$. Then $c\mid c(S_p)$, so $S_p$ is a union of complete $c$-fibers by \cref{thm:stabilizer-partition}(i), and $U=W_c\circ\pi_c$ is constant on $c$-fibers; \cref{thm:fiber-descent} (with $S=S_p$, $|S_p|\ge a\ge k$) gives $p=g\circ\pi_c$ with $g$ in the displayed quotient list, and $S_p=\pi_c^{-1}(\bar S_g)$ because $p(x)=U(x)\iff g(x^c)=W_c(x^c)$. The power map is equivariant, $\pi_c(hx)=h^c\pi_c(x)$, and $Q_c$ is a coset of $H^{(c)}:=\{h^c:h\in H\}$ with kernel $H_c$; hence $hS_p=S_p\iff h^c\bar S_g=\bar S_g$ and $c(S_p)=c\cdot c(\bar S_g)$. Since $c\mid c_0$, $\gcd(c\cdot c(\bar S_g),c_0)=c\cdot\gcd(c(\bar S_g),c_0/c)$, so $c_p=c$ is exactly the displayed coprimality. Conversely, for $g$ in the displayed set, $p:=g\circ\pi_c$ has full agreement set $\pi_c^{-1}(\bar S_g)$ of size $c|\bar S_g|\ge c\lceil a/c\rceil\ge a$, and the same computation returns $c_p=c$; the two maps are mutually inverse. Summing the buckets gives the exact formula. For the Johnson clause: at level $c$ the code has dimension $\lceil k/c\rceil$ and length $N_c=n/c$, and
\[
\Bigl\lceil\frac ac\Bigr\rceil^2-\Bigl(\Bigl\lceil\frac kc\Bigr\rceil-1\Bigr)\frac nc\ \ge\ \frac{a^2}{c^2}-\frac{kn}{c^2}\ =\ \frac{a^2-kn}{c^2}\ >\ 0,
\]
using $\lceil k/c\rceil-1<k/c$, so \cref{thm:johnson-list} applies at every level; with $\lceil a/c\rceil-\lceil k/c\rceil+1\le(a-k+2c-1)/c$ the level-$c$ bound is at most $n(a-k+2c-1)/(a^2-kn)$.
\end{proof}

\begin{proposition}[calibration: the quotient list term is necessarily word-dependent and exponential]\label{prop:list-calibration}
Let the hypotheses of \cref{lem:phi-fiber}\textup{(ii)} hold and let $U=u_z$ be the pigeonholed received word at agreement $A_2$. Every listed codeword $c_A=-R_A(\varphi)$ lies in the scale-$a$ quotient image $\{g\circ\varphi:\deg g\le\ell_2-2\}$. Hence in any bound of the shape demanded by the interleaved-list ledger,
\[
|\Lst(C^+,1-\tfrac{A_2}n,U)|\ \le\ L_Q(A_2,U)+L_{\rm ap}(A_2),
\]
in which $L_Q$ upper-bounds the quotient-image sublist and $L_{\rm ap}$ the rest, one necessarily has
\[
L_Q(A_2,u_z)\ \ge\ \binom N{\ell_2}\Big/|\B| .
\]
No uniform subexponential $L_Q$ exists in the band; only the aperiodic term $L_{\rm ap}$ can be uniformly polynomial, which is consistent with Open Problem~\ref{prob:band} and with \cref{rem:quotient-borne}.
\end{proposition}

\begin{proof}
Immediate from \cref{lem:phi-fiber}: the $\ge\binom N{\ell_2}/|\B|$ listed codewords are of the displayed composed form, hence are counted by $L_Q$.
\end{proof}

\begin{corollary}[the grand interleaved-list staircase outside the band]\label{cor:list-sandwich}
Fix any row satisfying the hypotheses of \cref{cor:grand}, \cref{cor:mixed-radix}, or \cref{cor:circle-grand}, and any interleaving arity $m\ge1$.
\begin{enumerate}[label=\textup{(\roman*)}]
\item \textup{(Unsafe near capacity.)} For every agreement $a\le k+n/N$ \textup(radius $\ge1-\rho-1/N$\textup),
\[
\bigl|\Lst\bigl(C^{\equiv m},\,1-\tfrac an\bigr)\bigr|\ \ge\ \binom N{\ell_1}\Big/|\B|\ >\ 2^{-128}\,q ,
\]
by \cref{lem:phi-fiber}\textup{(i)} and the entropy counts of the corresponding corollary, together with the embedding $c\mapsto(c,0,\dots,0)$ of base lists into interleaved lists.
\item \textup{(Safe in the Johnson regime.)} For every agreement $a$ with $a^2>(k-1)n$ and $n(a-k+1)/(a^2-(k-1)n)\le2^{-128}q$,
\[
\bigl|\Lst\bigl(C^{\equiv m},\,1-\tfrac an\bigr)\bigr|\ \le\ 2^{-128}\,q ,
\]
by \cref{thm:johnson-list}, uniformly in $m$.
\end{enumerate}
Thus the interleaved-list staircase is determined outside the same Johnson-to-capacity band as the MCA staircase; inside the band, \cref{prop:list-calibration} shows that any resolution of the interleaved-list ledger must carry an exponential, word-dependent quotient term, and \cref{thm:list-recursion} shows that for fiber-constant words this term is exactly a divisor-lattice recursion.
\end{corollary}

\begin{proof}
(i) At agreement $A_1\le k+n/N$ the base-field list floor is $\binom N{\ell_1}/|\B|$ by \cref{lem:phi-fiber}\textup{(i)}; list sizes are nondecreasing as the agreement threshold decreases, and the zero-padding embedding preserves column agreement, so the interleaved list is at least as large. The entropy counts of the cited corollaries give $\binom N{\ell_1}\ge2^{513}\ge|\B|\cdot2^{-128}q\cdot2^{0}$ throughout the challenge envelope, since $|\B|\,q\,2^{-128}\le q^2\,2^{-128}<2^{384}$. (ii) is \cref{thm:johnson-list}.
\end{proof}

\subsection{Curves and projective slopes in the deep regime}\label{sec:discharge-T5}

\begin{theorem}[deep-regime theorem for degree-$d$ power-curve samplers]\label{thm:deep-curve}
Let $C\subseteq\F^n$ be $\F$-linear with minimum weight $w$, let $d\ge1$, let $W_\gamma=f_0+\gamma f_1+\cdots+\gamma^df_d$ be a power-curve sampler, let $\delta\in(0,1)$, $r:=\lfloor\delta n\rfloor$, and assume
\[
(d+2)\,r\ \le\ w-1 .
\]
Then either there are codewords $c_0,\dots,c_d\in C$ and a set $T'$ with $|T'|\le r$ such that $f_i=c_i$ off $T'$ for every $i$ \textup(curve-tangent tuple; at most $d\,|T'|\le dr$ MCA-bad and no CA-bad parameters\textup), or
\[
\#\{\gamma\in\F:\ \dist(W_\gamma,C)\le\delta\}\ \le\ d\,(r+1).
\]
In particular the curve analogues of $\eca$ and $\emca$ \textup(with noncontainment read on the coefficient tuple against $C^{\equiv(d+1)}$, as in \cref{lem:one-support-d-curve}\textup) are at most $d(r+1)/q$ at radius $\delta$. Adjoining the projective parameter at infinity \textup(the leading word $f_d$\textup) adds at most one to every count. For $C=\RS[\F,D,k]$ the hypothesis reads $(d+2)\lfloor\delta n\rfloor\le n-k$.
\end{theorem}

\begin{proof}
Let $S_{\rm cl}$ be the displayed set. If $|S_{\rm cl}|\le d$ the second alternative holds. Otherwise pick distinct $\gamma_0,\dots,\gamma_d\in S_{\rm cl}$ with codewords $p_j$ and error sets $E_j$ ($|E_j|\le r$), and let $c_i:=\sum_j\lambda_{ij}p_j\in C$, where $(\lambda_{ij})$ is the inverse of the Vandermonde matrix $(\gamma_j^i)$; off $T:=\bigcup_jE_j$ ($|T|\le(d+1)r$) the same inversion applied to values gives $f_i=c_i$. Put $e_i:=f_i-c_i$, supported in $T$, and $T':=\{x:(e_{0,x},\dots,e_{d,x})\ne0\}$. For any $\gamma\in S_{\rm cl}$ with codeword $p$ and error set $E$: off $T\cup E$, $p=W_\gamma=\sum_i\gamma^ic_i$, so the codeword $p-\sum_i\gamma^ic_i$ has weight at most $(d+2)r<w$, hence is zero, and
\[
W_\gamma-\sum_i\gamma^ic_i=\sum_i\gamma^ie_i
\quad\text{has weight at most }r .
\]
For $x\in T'$ the nonzero polynomial $\gamma\mapsto\sum_ie_{i,x}\gamma^i$ has at most $d$ roots. If $|T'|\ge r+1$, each $\gamma\in S_{\rm cl}$ needs at least $|T'|-r$ vanishing coordinates in $T'$, so $|S_{\rm cl}|(|T'|-r)\le d|T'|$ and $|S_{\rm cl}|\le d|T'|/(|T'|-r)\le d(r+1)$. Otherwise $|T'|\le r$ and the tuple is curve-tangent. In the tangent case, $2r\le w-1$ makes two codewords agreeing on $\ge n-2r$ coordinates equal; arguing exactly as in Step~3 of \cref{thm:deep-mca} \textup(with $n-(d+2)r\ge n-w+1$ forcing both the point codeword and any tuple explanation to coincide with the $c_i$ off $T'$\textup), a parameter is MCA-bad if and only if $\sum_ie_{i,x}\gamma^i=0$ for some $x\in T'$, giving at most $d|T'|$ parameters, and the coefficient tuple is $\delta$-close to $C^{\equiv(d+1)}$ so no parameter is CA-bad. MCA-bad parameters are close, so the counts cover both notions; the parameter at infinity is a single additional value. For Reed--Solomon, $w=n-k+1$.
\end{proof}

Together with \cref{thm:weil-lines}, \cref{thm:deep-curve} also covers the multi-parameter samplers into which extension lines decompose, since those are handled over $\F$ directly by \cref{thm:deep-mca}; what the curve-sampler ledger still lacks is only the band-regime aperiodic pivot atlas, i.e.\ the curve instance of Open Problem~\ref{prob:band}, together with the exceptional vanishing-vector strata already flagged in \cref{sec:aperiodic-hankel-certificates}.

\subsection{The certificate scanner, its soundness, and deployed output}\label{sec:discharge-T6}

\begin{proposition}[small-field degeneracy]\label{prop:small-field}
Let $C\subsetneq\F^n$ be any linear code and $q=|\F|$. For every $\delta\in(0,1)$, $\emca(C,\delta)\ge1/q$. Consequently $\delta^*_C(\varepsilon^*)=0$ whenever $\varepsilon^*<1/q$; in particular every row over a field of size $q<2^{128}$ --- including the deployed QM31 circle rows, $q=p^4<2^{124}$ --- has a degenerate grand-challenge threshold $\delta^*_C(2^{-128})=0$, and the meaningful staircase question there is posed at a larger target such as $\varepsilon^*=2^{-100}$. At equality $1/q=\varepsilon^*$ the construction below gives only a floor meeting the target exactly, not strict unsafety; all degeneracy statements therefore require $1/q>\varepsilon^*$.
\end{proposition}

\begin{proof}
Pick $f_2\notin C$, any $c_0\in C$ and $\gamma_0\in\F$, and set $f_1:=c_0-\gamma_0f_2$. With $S=D$: the point $f_1+\gamma_0f_2=c_0$ is explained on $S$, while the pair is explained on $S=D$ only if both words lie in $C$, which fails. Hence $\gamma_0$ is MCA-bad at every radius, and $\emca\ge1/q$.
\end{proof}

\begin{definition}[staircase scanner]\label{def:scanner}
The scanner takes a row header $(D\text{-type},n,k,|\B|,q,q_{\rm line}=q)$, a target $\varepsilon^*$, and the applicable divisor data, and for each agreement $a\in\{k+1,\dots,n\}$ evaluates the following verdicts, each tagged by its certifying statement:
\begin{enumerate}[label=\textup{(V\arabic*)},leftmargin=2.6em]
\item \textbf{Cap-unsafe.} If $a\le A_2$ for an instantiated cap \textup(\cref{thm:main,thm:phi-cap} via \cref{cor:grand,cor:mixed-radix,cor:circle-grand,cor:circle-deployed}\textup) and $\frac1{2k}(1-\frac nq)>\varepsilon^*$: mark $a$ MCA-unsafe.
\item \textbf{Floor-unsafe.} If a remainder floor applies at $a$ \textup(\cref{thm:quotient-remainder-deep-floor}\textup) with converted numerator exceeding $\varepsilon^*q$: mark $a$ MCA-unsafe; if a fiber list floor at $a$ exceeds $\varepsilon^*q$ \textup(\cref{lem:fiber,lem:phi-fiber} with the entropy count\textup): mark $a$ list-unsafe.
\item \textbf{Deep-safe.} If $3(n-a)\le n-k$ and $n-a+1\le\varepsilon^*q$: mark $a$ MCA-safe \textup(\cref{thm:deep-mca}\textup); with $(d+2)(n-a)\le n-k$ and $d(n-a+1)\le\varepsilon^*q$: curve-$d$-safe \textup(\cref{thm:deep-curve}\textup).
\item \textbf{Johnson-list-safe.} If $a^2>(k-1)n$ and $n(a-k+1)/(a^2-(k-1)n)\le\varepsilon^*q$: mark $a$ list-safe for every interleaving arity \textup(\cref{thm:johnson-list}\textup).
\end{enumerate}
The output per challenge type is the pair $a_{\rm unsafe}:=\max\{a\ \text{marked unsafe}\}$ and $a_{\rm safe}:=\min\{a\ \text{marked safe}\}$, together with the open interval between them. If $\varepsilon^*<1/q$ the scanner instead reports the degeneracy of \cref{prop:small-field}.
\end{definition}

\begin{theorem}[scanner soundness]\label{thm:scanner-sound}
Every verdict printed by the scanner of \cref{def:scanner} is correct: agreements marked unsafe satisfy $\emca(C,1-a/n)>\varepsilon^*$ \textup(resp.\ list size $>\varepsilon^*q$\textup), agreements marked safe satisfy $\emca(C,1-a/n)\le\varepsilon^*$ \textup(resp.\ list size $\le\varepsilon^*q$, for every $m$\textup), and consequently $a_{\rm unsafe}<a^*_C\le a_{\rm safe}$ in the staircase formulation of \cref{prop:onestep}.
\end{theorem}

\begin{proof}
Each verdict cites a theorem of this paper whose hypotheses are verified numerically by the scanner; monotonicity across agreements is \cref{lem:mca-monotone} for MCA and radius-monotonicity of lists. The bracketing of $a^*_C$ is \cref{prop:onestep}.
\end{proof}

\begin{table}[ht]
\centering
\footnotesize
\setlength{\tabcolsep}{4pt}
\begin{tabular}{@{}lllll@{}}
\toprule
Row / challenge & Radius range $\delta$ & Verdict & Certified numerator & Source \\
\midrule
\multicolumn{5}{@{}l}{\emph{KoalaBear sextic}: $n=2^{21}$, $k=2^{20}$, $q=p^6>2^{185.9}$, $\varepsilon^*=2^{-128}$, $\varepsilon^*q\approx2^{57.9}$}\\
\addlinespace[2pt]
MCA & $[\tfrac12-2^{-7},\,\tfrac12)$ & unsafe & $>2^{-22}\,q\approx2^{164}$ & \cref{cor:deployed} \\
MCA & $\delta\le349525/2^{21}\approx0.16666$ & safe & $\le349526<2^{-167}q$ & \cref{thm:deep-mca} \\
MCA & $(0.16667,\,0.49218)$ & open & --- & Open Problem~\ref{prob:band} \\
List (any $m$) & $\delta\ge\tfrac12-2^{-8}$ & unsafe & $\ge\binom{256}{129}/p\ge2^{216}$ & \cref{lem:fiber}(i) \\
List (any $m$) & $\delta\le0.29289$ \ ($a\ge1482910$) & safe & $\le2^{20}$ & \cref{thm:johnson-list} \\
List (any $m$) & $(0.29290,\,0.49609)$ & open & --- & Open Problem~\ref{prob:band} \\
\midrule
\multicolumn{5}{@{}l}{\emph{M31 circle line round}: $n=2^{21}$, $k=2^{20}$, $q=p^4<2^{124}$; at $2^{-128}$: $\delta^*=0$ (\cref{prop:small-field}); below: $\varepsilon^*=2^{-100}$}\\
\addlinespace[2pt]
MCA & $[\tfrac12-2^{-7},\,\tfrac12)$ & unsafe & $>2^{-22}\,q\approx2^{102}$ & \cref{cor:circle-deployed} \\
MCA & $\delta\le0.16666$ & safe & $\le349526<2^{-105}q$ & \cref{thm:deep-mca} \\
List (any $m$) & $\delta\ge\tfrac12-2^{-8}$ & unsafe & $\ge2^{216}$ & \cref{lem:phi-fiber}(i) \\
List (any $m$) & $\delta\le0.29289$ & safe & $\le2^{20}\le2^{-100}q$ & \cref{thm:johnson-list} \\
\bottomrule
\end{tabular}
\caption{Baseline scanner output (\cref{def:scanner}) for the two deployed rows, before the pincer and the optimized graded-prefix floors; it is retained to display the self-contained deep-safe and Johnson-list-safe certificates, and the final deployed threshold bands are the narrowed bands of \cref{tab:scanner2}. The open MCA band for the circle row at $\varepsilon^*=2^{-100}$ is the same interval $(0.16667,\,0.49218)$ as above.}
\label{tab:scanner}
\end{table}

The table realizes the interval output demanded by the staircase interface: for the KoalaBear MCA staircase, $a_{\rm unsafe}=1{,}064{,}960$ \textup($\delta=\frac12-2^{-7}$\textup) and $a_{\rm safe}=1{,}747{,}627$ \textup($\delta\approx0.16666$\textup); the interval shrinks to one step exactly when Open Problem~\ref{prob:band} is settled, via \cref{thm:conditional-mca} whose other inputs are now discharged \textup(\cref{rem:conditional-status}\textup).

\subsection{The formal finite core}\label{sec:discharge-T7}

\begin{proposition}[finite certificates and decidability]\label{prop:finite-certificate}
Fix a row $(C,q,q_{\rm line})$ and an agreement $a$. The numerator $B_{\rm mca}(a)$ of \cref{sec:prize-program} is a well-defined integer, computable by finite enumeration over pairs, slopes, supports, and codewords; hence the one-step certificate of \cref{prop:onestep} is a finite object. A verifiable prize certificate consists of: the row header; the integer $a_0$; a \emph{lower witness} --- a pair $(f_1,f_2)$, a set of more than $\varepsilon^*q_{\rm line}$ slopes, and for each slope a codeword and a support, checkable by evaluation --- or a theorem tag from this paper with its numerically verified hypotheses; and an \emph{upper witness} --- a theorem tag with numerically verified hypotheses certifying $B_{\rm mca}(a_0+1)\le\varepsilon^*q_{\rm line}$. Verification of a certificate is polynomial in its length.
\end{proposition}

\begin{proof}
All quantifiers in \cref{def:ca,def:mca} range over the finite sets $\F^n$, $\F$, $2^D$, and $C$; the maximum defining $B_{\rm mca}(a)$ is therefore attained and computable. Checking a lower witness is evaluation and counting; checking a theorem tag is checking finitely many integer inequalities.
\end{proof}

\begin{table}[ht]
\centering
\small
\begin{tabular}{@{}ll@{}}
\toprule
Core statement & Proof ingredients \\
\midrule
\cref{thm:A} (deep-point conversion) & polynomial division, pigeonhole over $\F\setminus D$ \\
\cref{lem:fiber}, \cref{lem:phi-fiber} (fiber floors) & elementary symmetric functions, subfield pigeonhole \\
\cref{fact:chain}, \cref{lem:mca-monotone} & set manipulation \\
\cref{prop:onestep} (staircase) & monotone integer staircase \\
\cref{thm:deep-mca} (deep MCA) & linear interpolation, double counting of vanishings \\
\cref{thm:johnson-list} (Johnson lists) & convexity double counting of incidences \\
\cref{thm:stabilizer-partition}, \cref{thm:fiber-descent} & cyclic group actions, root counting \\
\cref{thm:weil-lines} (restriction of scalars) & basis decomposition of $\F/\B$ \\
\cref{def:scanner}, \cref{thm:scanner-sound} & finite case analysis over the statements above \\
\bottomrule
\end{tabular}
\caption{The formal finite core: each statement quantifies over explicitly bounded finite sets and uses only the listed elementary ingredients; the dependency order is top to bottom. Formalizing this table in a proof assistant is the outstanding transcription task.}
\label{tab:core}
\end{table}

The list deliberately includes the new safe-side theorems: a certificate package that contains both an unsafe certificate \textup(\cref{cor:grand} and its relatives\textup) and a machine-checked safe region \textup(\cref{thm:deep-mca,thm:johnson-list}\textup) documents the staircase from both sides in the sense of \cref{prop:onestep}.

\subsection{Tower transfer and the genus-one first step}\label{sec:discharge-T8}

The tower-transfer ledger half is already discharged: \cref{thm:stabilizer-partition,thm:fiber-descent,cor:periodic-support-count} were stated and proved for the Chebyshev case alongside the multiplicative case, with the divisor lattice of $D$ replaced by the dyadic scale lattice of the twin coset, exactly as the tower transfer requires; and Open Problem~\ref{prob:band} is likewise stated for both cases. It remains to treat the genus-one floor. We first extend \cref{def:map-smooth} to rational maps.

\begin{definition}[rationally smooth domain]\label{def:rational-smooth}
Let $\B\subseteq\F$, let $\psi=f/g$ with $f,g\in\B[X]$ coprime, $a:=\deg f>e:=\deg g\ge0$, and let $D\subseteq\B$ with $g$ nonvanishing on $D$. The set $D$ is \emph{$(\psi,a)$-smooth} if every nonempty fiber $\psi^{-1}(w)\cap D$, $w\in\psi(D)$, has exactly $a$ elements; then $N:=|D|/a=|\psi(D)|$. The FFTree domains of ECFFT \cite{BCKL21} are $(\psi,2)$-smooth for the $x$-coordinate map $\psi$ of each $2$-isogeny in the chain, of degree type $(a,e)=(2,1)$.
\end{definition}

\begin{proposition}[rational-map fiber floor; unified $(a,e)$ form]\label{prop:rational-floor}
Let $D\subseteq\B$ be $(\psi,a)$-smooth of size $n$, with $\psi=f/g$ of degree type $(a,e)$, and let $2\le\ell\le N-1$. Put
\[
k'\ :=\ a\ell-2(a-e),
\qquad
u_z\ :=\ \bigl(f(x)^{\ell}+z\,f(x)^{\ell-1}g(x)\bigr)_{x\in D}\in\B^n .
\]
Then there exists $z\in\B$ with
\[
\bigl|\Lst\bigl(\RS[\F,D,k'+1],\,1-\tfrac{a\ell}n,\,u_z\bigr)\bigr|\ \ge\ \binom N{\ell}\Big/|\B| ,
\]
i.e.\ a list at agreement $a\ell=k'+2(a-e)$. For $e=0$ this is \cref{lem:phi-fiber}\textup{(ii)} \textup(window $2a$, gap $\approx2/N$\textup); for $e=a-1$ the window is $2$ and the certified agreement is $k'+2$ --- the first staircase step, in accordance with the collapse of \cref{rem:ecfft-obstruction}.
\end{proposition}

\begin{proof}
For $A\subseteq\psi(D)$ with $|A|=\ell$ put $\Lambda_A:=\prod_{b\in A}(f-bg)$, of degree $a\ell$ with leading coefficient $\mathrm{lc}(f)^\ell\ne0$, vanishing at every $x\in D$ with $\psi(x)\in A$ --- exactly $a\ell$ points by smoothness \textup(and $g(x)\ne0$\textup). Expanding,
\[
\Lambda_A=\sum_{j=0}^{\ell}(-1)^je_j(A)\,f^{\ell-j}g^{\,j},
\]
with $j$-th term of degree $a\ell-j(a-e)$. Set $z_A:=-\eone(A)\in\B$ and
\[
c_A:=\Bigl(-\sum_{j\ge2}(-1)^je_j(A)\,f^{\ell-j}(x)\,g^{\,j}(x)\Bigr)_{x\in D},
\]
a polynomial of degree at most $a\ell-2(a-e)=k'$, so $c_A\in\RS[\B,D,k'+1]$; on the $a\ell$ fiber points, $u_{z_A}=c_A$. For injectivity at fixed $z$: if $c_A=c_{A'}$ as words then as polynomials \textup($\deg\le k'\le a(N-1)<n$\textup), so $\sum_{j\ge2}(-1)^j\bigl(e_j(A)-e_j(A')\bigr)f^{\ell-j}g^{\,j}=0$; dividing by $g^{\ell}$ in $\F(X)$ gives a polynomial identity in $\psi$, and $\psi$ is nonconstant, hence transcendental over $\F$ inside $\F(X)$, so all coefficients vanish: $e_j(A)=e_j(A')$ for $j\ge2$, and $\eone$ agrees via $z$; then $A=A'$ from the root sets of $\prod_{b}(Y-b)$. Pigeonholing $z_A$ over $\B$ finishes as in \cref{lem:phi-fiber}.
\end{proof}

\begin{corollary}[every ECFFT row has an unsafe first staircase step]\label{cor:ecfft-onestep}
Let $D$ be an ECFFT domain that is $(\psi,2)$-smooth \textup(\cref{def:rational-smooth}\textup) over $\B$, let $C=\RS[\F,D,k]$ with $2\mid k$, $\rho=k/n\in[\tfrac1{16},\tfrac12]$, $n<q=|\F|<2^{256}$, and $n\ge2^{12}$. Then, with $\ell=(k+2)/2\le N-1$,
\[
\binom{n/2}{(k+2)/2}\ \ge\ |\B|\Bigl(\frac qk+1\Bigr),
\]
and consequently
\[
\eca\Bigl(C,\,1-\rho-\tfrac2n\Bigr)\ >\ \frac1{2k}\Bigl(1-\frac nq\Bigr),
\qquad
\emca(C,\delta)\ >\ \frac1{2k}\Bigl(1-\frac nq\Bigr)\ \ \text{for every }\delta\in\Bigl[1-\rho-\tfrac2n,\,1-\rho\Bigr),
\]
so $\delta^*_C(\varepsilon^*)\le1-\rho-2/n$ for every $\varepsilon^*\le\frac1{2k}(1-\frac nq)$; and, provided $k<2^{128}$, the list challenge is likewise unsafe at that step, since $\binom{n/2}{(k+2)/2}/|\B|\ge q/k>2^{-128}q$. For odd $k$ the same argument with $\ell=\lfloor k/2\rfloor+1$ gives the step at $1-\rho-1/n$.
\end{corollary}

\begin{proof}
Apply \cref{prop:rational-floor} with $(a,e)=(2,1)$: $k'=2\ell-2=k$, so the list lives in $\RS[\F,D,k+1]=C^+$ at radius $\delta_0:=1-(k+2)/n$, and $\lfloor\delta_0n\rfloor=n-k-2\le n-k-1$ is admissible for \cref{thm:A}. The entropy count: $\ell/(n/2)=(k+2)/n\in(\rho,\rho+2^{-11}]\subseteq[\tfrac1{16},0.502]$, so $\Hb(\ell/(n/2))\ge\Hb(\tfrac1{16})\ge0.3372$ by the endpoint argument of \cref{cor:mixed-radix}, and
\[
\log_2\binom{n/2}{\ell}\ \ge\ \frac n2\,\Hb\Bigl(\frac{\ell}{n/2}\Bigr)-\log_2\Bigl(\frac n2+1\Bigr)\ \ge\ 0.3372\cdot\frac n2-\log_2\Bigl(\frac n2+1\Bigr)\ \ge\ 690-12\ \ge\ 513,
\]
since the middle expression is increasing in $n$ and equals $690.6\ldots-11.0\ldots$ at $n=2^{12}$, while $|\B|(q/k+1)\le q^2+q<2^{513}$. The contradiction with \cref{thm:A} at $\eta=\frac12$, then \cref{fact:chain} and \cref{lem:mca-monotone}, run exactly as in \cref{thm:phi-cap}. The list-challenge clause follows because $\binom{n/2}{\ell}/|\B|\ge q/k+1>2^{-128}q$ whenever $k<2^{128}$, which covers every challenge-envelope row.
\end{proof}
\subsection{The two-sided picture and what remains}\label{sec:discharge-status}

\begin{theorem}[two-sided sandwich for the grand MCA challenge]\label{thm:sandwich}
Let $C=\RS[\F,D,k]$ be a row for which one of the capacity caps of this paper applies at target $\varepsilon^*=2^{-128}$ --- \cref{cor:grand}, \cref{cor:mixed-radix}, \cref{cor:circle-grand} or \cref{cor:circle-deployed}, or \cref{cor:ecfft-onestep} --- with printed gap $G$ \textup(respectively $2^{-9}$/$2^{-10}$, $2^{-11}$/$2^{-12}$, $2^{-9}$--$2^{-11}$ or $2^{-7}$, and $2/n$\textup). Put
\[
r^*\ :=\ \min\Bigl(\Bigl\lfloor\frac{n-k}3\Bigr\rfloor,\ \bigl\lfloor2^{-128}q\bigr\rfloor-1\Bigr).
\]
If $q>2^{128}$ then
\[
\frac{r^*}n\ \le\ \delta^*_C(2^{-128})\ \le\ 1-\rho-G ,
\]
so, whenever $\lfloor2^{-128}q\rfloor-1\ge\lfloor(n-k)/3\rfloor$ --- e.g.\ whenever $q\ge2^{128}n$, as at every deployed multiplicative row --- the grand MCA threshold of such a row is determined up to a multiplicative factor of three, up to the microscopic gap $G$; for $q$ nearer $2^{128}$ the lower edge degrades to $r^*/n$ and no factor-of-three claim is made. If $q<2^{128}$ then $\delta^*_C(2^{-128})=0$ \textup(\cref{prop:small-field}\textup); at the boundary $q=2^{128}$ the pair of \cref{prop:small-field} gives $\emca\ge2^{-128}$ at every radius, which meets the threshold condition with equality and does not force degeneracy.
\end{theorem}

\begin{proof}
Lower bound: $3r^*\le n-k$ and $r^*+1\le\lfloor2^{-128}q\rfloor\le2^{-128}q$, so \cref{thm:deep-mca} gives $\emca(C,r^*/n)\le(r^*+1)/q\le2^{-128}$, and $r^*/n<(n-k)/n=1-\rho$, so $\delta^*\ge r^*/n$ by \cref{def:dstar} (the case $r^*=0$ being trivial). Upper bound: the applicable cap certifies $\emca(C,\delta)>2^{-128}$ for every $\delta\in[1-\rho-G,\,1-\rho)$, so no such $\delta$ belongs to the supremum set, and by \cref{lem:mca-monotone} neither does any larger sub-capacity radius; hence $\delta^*\le1-\rho-G$. The degenerate clause is \cref{prop:small-field} with $\varepsilon^*=2^{-128}<1/q$; the boundary remark is the same construction, whose error floor $1/q$ equals the target exactly at $q=2^{128}$.
\end{proof}

For $\rho=\tfrac12$ and $q\ge2^{128}n$ --- e.g.\ the KoalaBear-sextic row, $q>2^{185}$ --- the sandwich reads $\delta^*_C(2^{-128})\in[0.1666\ldots,\,\tfrac12-G]$: the first two-sided determination of a grand-challenge MCA threshold, exact to within a factor of three.

\begin{table}[ht]
\centering
\small
\begin{tabular}{@{}p{2.7cm}p{7.9cm}l@{}}
\toprule
Component & Outcome & Where \\
\midrule
Periodic-support ledger & Stabilizer exhaustion is exact at the support level; for fiber-constant lines it descends to an exact quotient-image ledger; upper numerators are independent of $q$ and $|\B|$. & \cref{sec:discharge-T1} \\
Deep MCA core & MCA is safe below one third of the distance; singular and rank-drop Hankel branches are removed in that range; the remaining band is Open Problem~\ref{prob:band}, with rank-pivot and compressed-eliminant certificates. & \cref{sec:discharge-T2,thm:deep-mca} \\
Extension-line lift & Extension-valued lines are handled exactly by restriction of scalars; there is no independent extension bucket. & \cref{sec:discharge-T3} \\
Interleaved-list side & Johnson safety is uniform in the interleaving arity; fiber-constant recursion is exact; the quotient term is necessarily exponential in the band. & \cref{sec:discharge-T4} \\
Curve/projective samplers & The deep regime is controlled for degree-$d$ curves and projective parameters; residual curve charts are part of the same aperiodic band problem. & \cref{sec:discharge-T5} \\
Scanner & The finite staircase scanner is specified, proved sound, and run on the deployed rows. & \cref{sec:discharge-T6} \\
Certificate core & The finite certificate format, decidability statement, and dependency chart are supplied; machine formalization is separate implementation work. & \cref{sec:discharge-T7} \\
Tower transfer & The ledger half is transferred to the Chebyshev tower; the genus-one first-step cap is proved, later strengthened to a macroscopic cap by the rational graded floors. & \cref{sec:discharge-T8,sec:answers-isogeny} \\
\bottomrule
\end{tabular}
\caption{Structural and certificate ingredients used in the two-sided threshold statements.}
\label{tab:status}
\end{table}

\begin{remark}[dependency status of the one-step criteria]\label{rem:conditional-status}
The one-step MCA criterion of \cref{thm:conditional-mca} has the following dependency status.  The periodic-support input is supplied by \cref{cor:periodic-support-count} at the support level and by \cref{thm:fiber-descent} at the image level for fiber-constant lines.  The deep-regime input is supplied by \cref{thm:deep-mca,cor:singular-deep}; outside that range it is exactly Open Problem~\ref{prob:band}, represented by the residual rank-pivot and kernel-compressed certificate language of \cref{thm:residual-rank-pivot-cover,thm:kernel-compressed-affine-eliminant,thm:exact-arbitrary-residual-image-ledger}.  The extension-line input is supplied by \cref{thm:weil-lines}.  The sampler input is definitional once the protocol's line sampler is fixed.  For the one-step list criterion of \cref{thm:conditional-list}, the required base-code bound is unconditional throughout the Johnson regime and uniform in $m$ \textup(\cref{thm:johnson-list}, making the codegree recursion unnecessary there\textup), is exactly recursive for fiber-constant words \textup(\cref{thm:list-recursion}\textup), and is provably exponential in the band \textup(\cref{prop:list-calibration}\textup), so any resolution must take that shape.
\end{remark}

After these ingredients are assembled, the remaining mathematical gap has a compact form.  The central residue is Open Problem~\ref{prob:band}, the aperiodic band conjecture, now represented by the finite rank-pivot, compressed-eliminant, and exact-image certificate language of \cref{thm:residual-rank-pivot-cover,thm:kernel-compressed-affine-eliminant,thm:exact-arbitrary-residual-image-ledger}; this is the analytic kernel separating the two-sided interval of \cref{thm:sandwich} from a one-step threshold determination.  The other residues are narrower: the failure-profile upgrade from error $1/k$ to error $1-o(1)$ \textup(Open Problem~\ref{prob:errorone}\textup) and explicit deployed-size certifying pairs \textup(Open Problem~\ref{prob:explicit}\textup); the genus-one band, formerly a separate residue, is settled at macroscopic scale in \cref{sec:answers-isogeny}.  The next section attacks the main band directly, from both sides at once.

\section{Safe-side pincer and scale-optimized unsafe floors}\label{sec:pincer}

The residue identified at the end of \cref{sec:discharge} is a single open band per row. This section attacks the band from both sides and narrows it substantially, with three new results. On the safe side, we prove that \emph{mutual} correlated agreement coincides with plain correlated agreement, up to an explicit tangent term, all the way up to half the minimum distance (\cref{thm:mca-from-ca}); this eliminates the support-mismatch layer of the problem on that range and, combined with the imported unique-decoding proximity gap of Ben-Sasson--Carmon--Ishai--Kopparty--Saraf \cite{BCIKS20}, moves the certified safe frontier of every deployed row from one third to one half of the distance. We also prove a self-contained correlated-agreement bound at half the Johnson radius (\cref{thm:elementary-ca}), which improves the \emph{unconditional} safe frontier for every row of rate below $\tfrac14$. On the unsafe side, we extend the graded locator-prefix pigeonhole of \cref{lem:quotient-remainder-prefix} to every map-smooth scale --- in particular to the Chebyshev scales of the circle tower --- and optimize it over the scale (\cref{prop:graded-prefix-floor}); at deployed sizes this widens the certified unsafe band by a factor of four with a hand-checkable certificate, and by a factor of about sixteen with an exactly verified one, saturating an intrinsic entropy--subfield envelope that we identify in closed form (\cref{rem:entropy-frontier}). \Cref{thm:sandwich2} and \cref{tab:scanner2} assemble the narrowed sandwich.

\subsection{Mutual equals correlated agreement up to half the distance}\label{sec:pincer-half}

The deep-regime theorem (\cref{thm:deep-mca}) controls $\emca$ directly, but only below one third of the distance: its alternatives identify every close point-codeword with a value of a single interpolated line, and that identification costs a factor of three in the radius. The next theorem decouples the two layers of the problem. It shows that below \emph{half} the distance, the support-mismatch phenomenon that separates mutual from plain correlated agreement is carried entirely by the tangent slopes of \cref{thm:deep-mca}\textup{(a)}, of which there are at most $\lfloor\delta n\rfloor$; everything else is already visible to $\eca$.

\begin{theorem}[mutual from correlated agreement, up to half the distance]\label{thm:mca-from-ca}
Let $C\subseteq\F^n$ be any $\F$-linear code of minimum Hamming weight $w$, let $q=|\F|$, let $\delta\in(0,1)$, and put $r:=\lfloor\delta n\rfloor$. Assume
\[
2r\ \le\ w-1 .
\]
Then for every pair $(f_1,f_2)\in\F^n\times\F^n$:
\begin{enumerate}[label=\textup{(\alph*)}]
\item if $\dist_2\bigl((f_1,f_2),C^{\equiv2}\bigr)\le\delta$, the pair has \emph{no} CA-bad slope at radius $\delta$, and at most $r$ MCA-bad slopes at radius $\delta$: writing $(p_1,p_2)$ for a pair of codewords column-agreeing with $(f_1,f_2)$ outside a set $E$ with $|E|\le r$ and $e_i:=f_i-p_i$, the MCA-bad slopes form a subset of $\{-e_{1,j}/e_{2,j}:\ j\in E,\ e_{2,j}\ne0\}$;
\item if $\dist_2\bigl((f_1,f_2),C^{\equiv2}\bigr)>\delta$, every MCA-bad slope at radius $\delta$ is CA-bad at radii $(\delta,\delta)$.
\end{enumerate}
Consequently
\[
\emca(C,\delta)\ \le\ \max\Bigl(\,\eca(C,\delta),\ \frac rq\,\Bigr).
\]
For $C=\RS[\F,D,k]$ on any evaluation set $D$ of size $n$, the hypothesis reads $2\lfloor\delta n\rfloor\le n-k$: mutual and plain correlated agreement coincide, up to the explicit tangent term $r/q$, on the entire unique-decoding range of the interleaved code $C^{\equiv2}$.
\end{theorem}

\begin{proof}
Part \textup{(b)} is immediate from the definitions: an MCA-bad slope $\gamma$ has a witness set $S$ with $|S|\ge\lceil(1-\delta)n\rceil=n-r$ on which $f_1+\gamma f_2$ agrees with a codeword, so $\dist(f_1+\gamma f_2,C)\le r/n\le\delta$; if additionally the pair is $\delta$-far in column distance, $\gamma$ is CA-bad at radii $(\delta,\delta)$ by \cref{def:ca}.

For part \textup{(a)}, fix codewords $(p_1,p_2)$ column-agreeing with $(f_1,f_2)$ outside the column error set $E:=\{j:\ (e_{1,j},e_{2,j})\ne(0,0)\}$, $e_i:=f_i-p_i$, $|E|\le r$. The pair is $\delta$-close, so it has no CA-bad slope. Let $\gamma$ be MCA-bad with witness $(S,c)$: $|S|\ge n-r$, $c\in C$, $c=f_1+\gamma f_2$ on $S$, and no pair of codewords of $C$ explains $(f_1,f_2)$ on $S$.

First, $c=p_1+\gamma p_2$. Indeed, $z:=c-(p_1+\gamma p_2)\in C$ vanishes on $S\setminus E$: there $f_i=p_i$, so $f_1+\gamma f_2=p_1+\gamma p_2$, while $c=f_1+\gamma f_2$ on all of $S$. Since $|S\setminus E|\ge n-2r$, the codeword $z$ has weight at most $2r\le w-1$, hence $z=0$.

Next, $S\not\subseteq D\setminus E$: otherwise $(p_1,p_2)$ would explain the pair on $S$. So there is $j\in S\cap E$. On $S$ we have $0=(f_1+\gamma f_2)-c=(e_1+\gamma e_2)$, so at this $j$,
\[
e_{1,j}+\gamma\,e_{2,j}=0,\qquad (e_{1,j},e_{2,j})\ne(0,0).
\]
If $e_{2,j}=0$ the left side is the nonzero constant $e_{1,j}$, which cannot vanish; hence $e_{2,j}\ne0$ and $\gamma=-e_{1,j}/e_{2,j}$. The MCA-bad slopes of the pair therefore form a subset of the at most $|E|\le r$ listed values. Taking the maximum over pairs gives the displayed bound.
\end{proof}

\begin{remark}[scope and sharpness]\label{rem:half-scope}
Three comments. \textup{(i)} The hypothesis $2r\le w-1$ is exactly the unique-decoding radius of the interleaved code $C^{\equiv2}$ in column weight, and the proof uses it twice, both times through the statement ``a codeword of weight $\le 2r$ vanishes''; beyond half the distance a $\delta$-close pair can admit several column explanations with incomparable agreement sets, and witness codewords need not ride any explanation line, which is precisely where the band begins. \textup{(ii)} The theorem is an unconditional \emph{implication}, with no list bound and no field-size hypothesis; instantiated with the self-contained $\eca$ bound of \cref{thm:deep-mca} it returns the deep theorem (with the marginally better error $\max((r+1)/q,r/q)$ replaced by $(r+1)/q$), and its value is that it accepts \emph{any} stronger correlated-agreement input on the range $(w-1)/3<2r\le w-1$, where no self-contained input is available. \textup{(iii)} The same proof gives the two-radius statement: for $\delta_{\mathrm{int}}\ge\delta$, every MCA-bad slope of a $\delta_{\mathrm{int}}$-far pair is CA-bad at radii $(\delta,\delta_{\mathrm{int}})$, while a $\delta_{\mathrm{int}}$-close pair that is also $\delta$-close contributes at most $r$ slopes; only pairs in the annulus between the two radii require the no-loss form.
\end{remark}

\begin{corollary}[the band conjecture loses its support-mismatch layer below half the distance]\label{cor:band-reduction}
Let $C=\RS[\F,D,k]$ and let $\delta$ satisfy $2\lfloor\delta n\rfloor\le n-k$. Then any correlated-agreement bound $\eca(C,\delta)\le\varepsilon$ implies $\emca(C,\delta)\le\max(\varepsilon,\lfloor\delta n\rfloor/q)$. In particular, on the agreement subrange $a\ge\lceil(n+k)/2\rceil$ of Open Problem~\ref{prob:band}, the aperiodic band conjecture reduces to its correlated-agreement form: column-close pairs contribute at most $\lfloor\delta n\rfloor$ explicit tangent slopes, and what must be bounded polynomially is only, for each column-far pair, the number of slopes $\gamma$ for which $f_1+\gamma f_2$ has an aperiodically witnessed agreement set of size at least $\lceil(1-\delta)n\rceil$.
\end{corollary}

\begin{proof}
Immediate from \cref{thm:mca-from-ca}: case \textup{(a)} pairs contribute at most $r=\lfloor\delta n\rfloor$ slopes, each explicitly a tangent ratio; case \textup{(b)} identifies every remaining MCA-bad slope, together with its witness support, as a CA-bad slope with the same witness.
\end{proof}

\subsection{A self-contained correlated-agreement bound at half the Johnson radius}\label{sec:pincer-ca}

\Cref{thm:mca-from-ca} converts the safe side of the band problem into a demand for correlated-agreement inputs. The strongest known inputs are the proximity-gap theorems of \cite{BCIKS20} (unique-decoding range with error $n/q$; Johnson range with error $\operatorname{poly}(n)/q$), whose proofs run through a Berlekamp--Welch decoder over the rational function field $\F(Z)$ and an induction over curves, and the recent strengthening \cite{BCHKS25}. Those we import. But the certificate interface of \cref{sec:prize-program} also prices \emph{self-contained} inputs, since the formal core (\cref{tab:core}) is built from elementary counting; the following theorem is the elementary member of the family. Its only ingredients are the pair-interleaved Johnson bound of \cref{thm:johnson-list} and a rider double count, and it reaches half the Johnson radius with an explicit list-size error.

\begin{theorem}[elementary correlated agreement at half the Johnson radius]\label{thm:elementary-ca}
Let $C=\RS[\F,D,k]$ on any evaluation set $D$ of size $n$, let $\delta\in(0,1)$, put $r:=\lfloor\delta n\rfloor$, and assume
\[
n-2r\ >\ 0
\qquad\text{and}\qquad
(n-2r)^2\ >\ (k-1)\,n .
\]
Then, with
\[
L_2\ :=\ \frac{n\,\bigl(n-2r-k+1\bigr)}{(n-2r)^2-(k-1)n}
\]
the pair-interleaved Johnson bound of \cref{thm:johnson-list} at agreement $n-2r$,
\[
\eca(C,\delta)\ \le\ \frac{1+(r+1)\,L_2}{q}\,,
\]
and the same bound holds for $\eca(C,\delta,\delta_{\mathrm{int}})$ for every $\delta_{\mathrm{int}}\ge\delta$.
\end{theorem}

\begin{proof}
Fix a pair $(f_1,f_2)$ with $\dist_2((f_1,f_2),C^{\equiv2})>\delta_{\mathrm{int}}\ge\delta$; only such pairs have CA-bad slopes. Let $B\subseteq\F$ be the set of CA-bad slopes; for each $\gamma\in B$ fix a codeword $c_\gamma\in C$ and an agreement set $S_\gamma$ with $|S_\gamma|\ge n-r$ and $c_\gamma=f_1+\gamma f_2$ on $S_\gamma$. If $|B|\le1$ we are done, so fix $\gamma_1\in B$ and consider any $\gamma\in B\setminus\{\gamma_1\}$. Define
\[
P_2:=\frac{c_{\gamma_1}-c_{\gamma}}{\gamma_1-\gamma}\in C,
\qquad
P_1:=c_{\gamma_1}-\gamma_1P_2\in C,
\]
so that, identically as codewords,
\begin{equation}\label{eq:rider-identity}
P_1+\gamma_1P_2=c_{\gamma_1},
\qquad
P_1+\gamma P_2=c_{\gamma}.
\end{equation}
On $S_{\gamma_1}\cap S_{\gamma}$, which has size at least $n-2r$, the two coordinatewise linear equations $f_1+\gamma_1f_2=c_{\gamma_1}$ and $f_1+\gamma f_2=c_{\gamma}$ have the unique solution $f_2=P_2$, $f_1=P_1$; hence the codeword pair $(P_1,P_2)$ column-agrees with $(f_1,f_2)$ on at least $n-2r$ coordinates, i.e.
\[
(P_1,P_2)\ \in\ \mathcal L\ :=\ \Lst\Bigl(C^{\equiv2},\,1-\tfrac{n-2r}{n},\,(f_1,f_2)\Bigr),
\qquad
|\mathcal L|\le L_2
\]
by \cref{thm:johnson-list} with $m=2$ and $a=n-2r$. Say that $\gamma$ \emph{rides} a member $(P_1,P_2)\in\mathcal L$ if $c_\gamma=P_1+\gamma P_2$ as codewords; by \eqref{eq:rider-identity}, every $\gamma\in B\setminus\{\gamma_1\}$ rides some member of $\mathcal L$.

Fix $(P_1,P_2)\in\mathcal L$ and let $E:=\{j:\ (f_{1,j},f_{2,j})\ne(P_{1,j},P_{2,j})\}$ be its column error set. Since the pair is $\delta$-far, $|E|\ge r+1$. If $\gamma$ rides $(P_1,P_2)$, then $(f_1-P_1)+\gamma(f_2-P_2)=(f_1+\gamma f_2)-c_\gamma$ vanishes on $S_\gamma$, hence on at least $|E|-r\ge1$ coordinates of $E$; and at each $j\in E$ the map $\gamma\mapsto(f_{1,j}-P_{1,j})+\gamma(f_{2,j}-P_{2,j})$ is a nonzero affine function, vanishing for at most one slope. Double counting vanishing incidences between riders of $(P_1,P_2)$ and coordinates of $E$,
\[
\#\{\text{riders of }(P_1,P_2)\}\cdot\bigl(|E|-r\bigr)\ \le\ |E|,
\qquad\text{so}\qquad
\#\{\text{riders}\}\ \le\ \frac{|E|}{|E|-r}\ \le\ r+1 .
\]
Summing over $\mathcal L$ and adding back $\gamma_1$: $|B|\le1+(r+1)L_2$. Taking the maximum over pairs completes the proof.
\end{proof}

\begin{remark}[provenance and comparison]\label{rem:ca-provenance}
The argument is elementary and in the folklore spirit ``a list bound implies a proximity gap with a list-size loss''; we do not claim novelty of the mechanism, only of the exact constants in this normalization and of its place in the formal core: its only inputs are \cref{thm:johnson-list} and counting, so it enters the dependency chart of \cref{tab:core} without enlarging the trusted base. The theorems of \cite{BCIKS20,BCHKS25} are strictly stronger where they apply --- they reach the full Johnson radius $1-\sqrt\rho$ with error $\operatorname{poly}(n)/q$ --- at the price of a substantially deeper proof. Note also the built-in factor of two: the recovered pair $(P_1,P_2)$ lives at column radius $2r$, so the list bound is invoked at agreement $n-2r$, which is why the theorem stops at half the Johnson radius rather than at the Johnson radius.
\end{remark}

\begin{corollary}[self-contained safe frontier at half the Johnson radius]\label{cor:self-contained-safe}
Let $C=\RS[\F,D,k]$ with $n\ge k+2$, and let $\delta$, $r$, $L_2$ be as in \cref{thm:elementary-ca}, so $n-2r>0$ and $(n-2r)^2>(k-1)n$. Then
\[
\emca(C,\delta)\ \le\ \frac{1+(r+1)L_2}{q}\,.
\]
For the challenge rows this improves the self-contained safe frontier exactly when $\rho<\tfrac14$: the frontier moves from $(1-\rho)/3$ to $\approx(1-\sqrt\rho\,)/2$, i.e.\ from $0.2917$ to $0.3232$ at $\rho=\tfrac18$ and from $0.3125$ to $0.3750$ at $\rho=\tfrac1{16}$, while at $\rho\in\{\tfrac12,\tfrac14\}$ the deep frontier $(1-\rho)/3$ remains the better self-contained certificate. At the deployed KoalaBear row the half-Johnson certificate reads: $r=307121$, $\delta=r/n\approx0.14645$, $L_2\le1001282$, error at most $(1+307122\cdot1001282)/q<2^{-147}$ --- dominated by \cref{thm:deep-mca} at this rate, but recorded as an independent certificate of the same verdict.
\end{corollary}

\begin{proof}
The hypothesis gives $n-2r>0$ and $(n-2r)^2>(k-1)n$, hence $n-2r>\sqrt{(k-1)n}$; since $n\ge k+2$ we also have $\sqrt{(k-1)n}\ge k$, so $n-2r>k$, i.e.\ $2r\le n-k-1$: the radius lies below half the distance. \Cref{thm:mca-from-ca} therefore applies, and with the input of \cref{thm:elementary-ca},
\[
\emca(C,\delta)\le\max\Bigl(\tfrac{1+(r+1)L_2}{q},\ \tfrac rq\Bigr)=\tfrac{1+(r+1)L_2}{q}.
\]
The frontier comparison is the elementary inequality $(1-\sqrt\rho)/2>(1-\rho)/3\iff\sqrt\rho<\tfrac12$; the KoalaBear numbers are $n=2^{21}$, $k=2^{20}$, $(n-2r)^2-(k-1)n=909700$ at $r=307121$, and $L_2=\lfloor n(n-2r-k+1)/909700\rfloor\le1001282$, with $(1+307122\cdot1001282)<2^{38.2}$ and $q>2^{185.9}$.
\end{proof}

\begin{corollary}[conditional safe frontier at half the distance]\label{cor:conditional-half}
Let $C=\RS[\F,D,k]$, $\rho=k/n$. Import the unique-decoding proximity gap of \cite[Thm.~4.1]{BCIKS20}: in the normalization of \cref{def:ca}, $\eca(C,\delta)\le n/q$ for every $\delta\le(1-\rho)/2$. (Only $\eca(C,\delta)\le\operatorname{poly}(n)/q$ is used below, so the conclusion is robust to the exact constant.) Then for every $\delta$ with $2\lfloor\delta n\rfloor\le n-k$,
\[
\emca(C,\delta)\ \le\ \max\Bigl(\frac nq,\ \frac{\lfloor\delta n\rfloor}q\Bigr)\ =\ \frac nq\,.
\]
Deployed instantiations: for the KoalaBear sextic row, $\emca(C,\delta)\le n/q<2^{-164}<2^{-128}$ for every $\delta\le\tfrac14$ (the endpoint $r=2^{19}$, $2r=n-k$, is included); for the circle line-round row over $\F_{p^4}$, $p=2^{31}-1$, $\emca\le n/q<2^{-102}<2^{-100}$ for every $\delta\le\tfrac14$. In both cases the certified safe frontier moves from one third to one half of the distance.
\end{corollary}

\begin{proof}
Combine the import with \cref{thm:mca-from-ca}; $\lfloor\delta n\rfloor\le n$ makes the maximum $n/q$. For the numbers: $n/q=2^{21}/p^6<2^{-164.9}$ with $p=2^{31}-2^{24}+1$, and $n/q=2^{21}/(2^{31}-1)^4<2^{-102.9}$. The circle case uses the diagonal equivalence of \cref{lem:circle-rs}, which preserves both correlated-agreement errors exactly, so the import applies verbatim to the equivalent Reed--Solomon code on the torus twin coset.
\end{proof}

\subsection{Widening the unsafe band: graded prefix floors on every scale}\label{sec:pincer-floors}

The deployed caps of \cref{cor:deployed,cor:circle-deployed} certify the unsafe band $[\tfrac12-2^{-7},\tfrac12)$ using only the top \emph{two} slots of the locator expansion at the single top scale ($N=256$). The quotient-remainder ledger already prices deeper prefixes at the multiplicative level (\cref{lem:quotient-remainder-prefix}); two observations remain to be made. First, the graded pigeonhole runs verbatim on \emph{every} map-smooth scale --- the locator expansion $\prod_{b\in M}(\varphi-b)=\sum_j(-1)^je_j(M)\varphi^{m-j}$ is graded by $\deg\varphi$ for any polynomial $\varphi$, in particular for the Chebyshev scales of the circle tower --- which extends the tower ledger from the first staircase step to every depth. Second, the floor can be \emph{optimized over the scale}: small scales trade prefix cost $|\B|$ per slot against fiber entropy $\binom{n/c}{m}$, and the optimum widens the certified unsafe band by an order of magnitude at deployed sizes. The route has an intrinsic ceiling, which we identify exactly in \cref{rem:entropy-frontier}.

\begin{proposition}[graded prefix floor for map-smooth scales]\label{prop:graded-prefix-floor}
Let $\B\subseteq\F$ be a subfield, let $D\subseteq\B$ with $|D|=n$, and let $\varphi\in\B[X]$ of degree $c\ge2$ whose fibers within $D$ are complete and of uniform size $c$ (the map-smooth hypothesis of \cref{lem:phi-fiber}); put $N:=n/c$ and $Q:=\varphi(D)\subseteq\B$, so $|Q|=N$. Let $1\le K\le n$ and let $m$ satisfy $\lceil K/c\rceil\le m\le N$; put
\[
w\ :=\ m-\lceil K/c\rceil\ \ge\ 0 .
\]
Then there is a $\B$-valued received word $U:D\to\B$ with
\[
\bigl|\Lst\bigl(\RS[\F,D,K],\,1-\tfrac{mc}{n},\,U\bigr)\bigr|
\ \ge\
\left\lceil\frac{\binom Nm}{|\B|^{\,w}}\right\rceil .
\]
For $\varphi=X^c$ on a multiplicative coset this is \cref{lem:quotient-remainder-prefix} with $s=0$ (note $w=\lfloor(mc-K)/c\rfloor=m-\lceil K/c\rceil$); the content here is that the same graded pigeonhole runs on every polynomial scale, in particular on the Chebyshev scales $T_c$ of the circle tower (\cref{lem:cheb-fibers}), whose locators are not lacunary in the monomial basis but are perfectly graded in powers of $T_c$.
\end{proposition}

\begin{proof}
For $M\subseteq Q$ with $|M|=m$, the locator $\Lambda_M:=\prod_{b\in M}(\varphi-b)\in\B[X]$ vanishes exactly on the union of the $\varphi$-fibers over $M$ inside $D$, a set of size $mc$, and expands as $\Lambda_M=\sum_{j=0}^m(-1)^je_j(M)\,\varphi^{m-j}$ with $e_j(M)\in\B$ the elementary symmetric functions of $M\subseteq\B$. For $z=(z_1,\dots,z_w)\in\B^w$ define the $\B$-valued word
\[
U_z\ :=\ \Bigl(\varphi^m+\textstyle\sum_{j=1}^{w}z_j\,\varphi^{m-j}\Bigr)\Big|_D .
\]
Assign to each $M$ the vector $z(M):=\bigl((-1)^je_j(M)\bigr)_{1\le j\le w}\in\B^w$; by pigeonhole some $z^*$ is assigned by at least $\binom Nm/|\B|^w$ of the $\binom Nm$ sets. For each such $M$ put
\[
c_M\ :=\ U_{z^*}-\Lambda_M\big|_D\ =\ -\sum_{j=w+1}^{m}(-1)^je_j(M)\,\varphi^{m-j}\Big|_D ,
\]
a polynomial of degree at most $c(m-w-1)=c\bigl(\lceil K/c\rceil-1\bigr)\le K-1$, hence a codeword of $\RS[\F,D,K]$; it agrees with $U_{z^*}$ on the $mc$ zeros of $\Lambda_M$ in $D$, so it lies in the displayed list. The map $M\mapsto c_M$ is injective: if $c_M=c_{M'}$ then $\sum_{j>w}(-1)^j\bigl(e_j(M)-e_j(M')\bigr)\varphi^{m-j}=0$ as a polynomial of degree $<n$ vanishing on $D$, hence identically; the powers $\varphi^{m-j}$ have pairwise distinct degrees $c(m-j)$ and are therefore linearly independent over $\F$, so $e_j(M)=e_j(M')$ for $j>w$, while for $j\le w$ equality holds by the choice of $z^*$; thus $\prod_{b\in M}(Y-b)=\prod_{b\in M'}(Y-b)$ and $M=M'$.
\end{proof}

\begin{corollary}[widened unsafe band at deployed sizes]\label{cor:widened-deployed}
\begin{enumerate}[label=\textup{(\roman*)}]
\item \textup{(KoalaBear sextic, MCA.)} With $\B$, $\F$, $D$, $n=2^{21}$, $k=2^{20}$ as in \cref{cor:deployed},
\[
\emca\bigl(\RS[\F_{p^6}\!,D,2^{20}],\,\delta\bigr)\ >\ 2^{-22}
\qquad\text{for every}\qquad
\delta\in\Bigl[\tfrac{15}{32},\,\tfrac12\Bigr),
\]
a certified gap of $2^{-5}$, four times the gap $2^{-7}$ of \cref{cor:deployed}, by a hand-checkable certificate at scale $c=16$.
\item \textup{(KoalaBear sextic, grand list challenge.)} For every $\delta\in[\tfrac{15}{32},1-\rho)$ there is a $\B$-valued word $U$ with $|\Lst(C,\delta,U)|>2^{-128}q$.
\item \textup{(Circle line-round row.)} With $p'=2^{31}-1$, $\F=\F_{p'^4}$, and the twin-coset $x$-domain of size $n=2^{21}$, $k=2^{20}$ as in \cref{cor:circle-deployed}, the same two statements hold verbatim on $[\tfrac{15}{32},\tfrac12)$ resp.\ $[\tfrac{15}{32},1-\rho)$, at the Chebyshev scale $\varphi=T_{16}$, against the target $2^{-100}$.
\end{enumerate}
\end{corollary}

\begin{proof}
All four claims instantiate \cref{prop:graded-prefix-floor} at scale $c=16$, $N=n/c=2^{17}$, $m=69632$, so $mc=k+2^{16}$ and $m/N=\tfrac{17}{32}$, and convert the floor as in \cref{cor:deployed}. The common entropy estimate is
\[
\log_2\binom{N}{m}\ \ge\ N\,\Hb\bigl(\tfrac{17}{32}\bigr)-\log_2(N+1)\ \ge\ 131072\cdot0.99718-17\ >\ 130{,}685 .
\]
\textup{(i)} Take $K=k+1$, so $w=m-\lceil(k+1)/16\rceil=69632-65537=4095$. The threshold for the deep-point conversion (\cref{thm:A} with $\eta=\tfrac12$, admissible since $n-mc\le n-k-1$) is a floor exceeding $q/k+1$:
\[
w\log_2p+\log_2(q/k+1)\ <\ 4095\cdot30.989+166\ <\ 127{,}116\ <\ 130{,}685 ,
\]
with over $3500$ bits of room. So $\eca(C,1-mc/n)>\tfrac1{2k}(1-n/q)>2^{-22}$, and \cref{fact:chain} with \cref{lem:mca-monotone} extend the same lower bound to $\emca(C,\delta)$ for every $\delta\in[1-mc/n,\tfrac12)=[\tfrac{15}{32},\tfrac12)$.
\textup{(ii)} Take $K=k$, so $w=4096$; the list threshold is $\binom Nm/|\B|^w>2^{-128}q$, i.e.\ $w\log_2p+\log_2q-128<4096\cdot30.989+58<127{,}000<130{,}685$; lists only grow with the radius, giving the full range.
\textup{(iii)} $T_{16}$ is a polynomial of degree $16$ over $\B=\F_{p'}$ with complete uniform fibers on the twin-coset $x$-domain at every dyadic scale (\cref{lem:cheb-fibers}), so \cref{prop:graded-prefix-floor} applies with the same $N$, $m$, $w$; the thresholds are weaker ($\log_2(q/k+1)<104$ and $\log_2q-100<24$), and the conversion runs through the diagonal equivalence of \cref{lem:circle-rs} exactly as in \cref{cor:circle-deployed}.
\end{proof}

\begin{remark}[exactly verified frontier]\label{rem:exact-frontier}
The hand-checkable rows above deliberately stop at the round gap $2^{-5}$. Optimizing $(c,m)$ under exact integer verification of $\binom{N}{m}>p^{\,w}(q/k+1)$ resp.\ $\binom Nm\cdot2^{128}>p^{\,w}q$ --- a finite certificate in the sense of \cref{def:scanner,prop:finite-certificate} --- widens the band further, to the following edges (all verified by exact arithmetic; $\Delta:=mc-k$):
\begin{center}\footnotesize
\begin{tabular}{@{}llllll@{}}
\toprule
row / challenge & scale $c$ & $m$ & $w$ & $\Delta$ & certified unsafe range \\
\midrule
KoalaBear, MCA & $16$ & $69748$ & $4211$ & $67392$ & $\delta\in[\,\tfrac{15331}{32768},\tfrac12)=[0.46786\ldots,\tfrac12)$ \\
KoalaBear, list & $32$ & $34874$ & $2106$ & $67392$ & $\delta\in[\,\tfrac{15331}{32768},1-\rho)$ \\
circle $p'{=}2^{31}{-}1$, MCA & $32$ & $34873$ & $2104$ & $67360$ & $\delta\in[\,\tfrac{30663}{65536},\tfrac12)=[0.46787\ldots,\tfrac12)$ \\
circle $p'{=}2^{31}{-}1$, list & $32$ & $34874$ & $2106$ & $67392$ & $\delta\in[\,\tfrac{15331}{32768},1-\rho)$ \\
\bottomrule
\end{tabular}
\end{center}
The certified gap is thus $\approx0.03214\approx2^{-4.96}$ throughout, versus $2^{-7}$ before.
\end{remark}

\begin{remark}[the entropy--subfield frontier]\label{rem:entropy-frontier}
The optimized route has a closed-form ceiling. At scale $c$ the floor beats the conversion threshold iff
\[
\log_2\binom{n/c}{m}\ \gtrsim\ w\log_2|\B|+\log_2(q/k),
\]
and writing $mc=:(\rho+g)n$ --- so $m/N=\rho+g$ and $w=gN+O(1)$ --- and dividing by $N$, this reads
\[
\Hb(\rho+g)\ \ge\ g\,\log_2|\B|+o(1)
\]
for challenge-envelope field sizes ($\log_2q\le256$, so the additive terms are $O(1)$ against $N\to\infty$). Define
\[
g^*(\rho,\beta)\ :=\ \sup\bigl\{\,g\in(0,1-\rho]:\ \Hb(\rho+g)\ge\beta g\,\bigr\},
\qquad
\beta:=\log_2|\B| .
\]
Every gap $g<g^*$ is certified at large enough deployed sizes, and \emph{no} gap beyond $g^*$ is certifiable by graded prefix floors at any scale: the frontier is intrinsic to the route, not to our optimization. For $\beta=31$ (both deployed base primes): $g^*(\tfrac12)\approx0.03216$, $g^*(\tfrac14)\approx0.02748$, $g^*(\tfrac18)\approx0.01920$, $g^*(\tfrac1{16})\approx0.01238$. The deployed KoalaBear row achieves $0.03214$ at $n=2^{21}$: the envelope is saturated to three digits already at deployed sizes. Two structural consequences: the left edge of the band in Open Problem~\ref{prob:band} moves from $k+2n/N$ to $k+g^*n-o(n)$; and the frontier depends on the base field only through $\beta$, so larger base fields \emph{narrow} the certifiable unsafe band --- a quantitative form of the subfield-confinement phenomenon of \cref{sec:subfield}.
\end{remark}

\subsection{The narrowed sandwich}\label{sec:pincer-sandwich}

\begin{theorem}[narrowed two-sided sandwich]\label{thm:sandwich2}
Let $\delta^*_C$ be the challenge threshold of \cref{def:dstar}.
\begin{enumerate}[label=\textup{(\roman*)}]
\item \textup{(KoalaBear sextic MCA row, $\varepsilon^*=2^{-128}$.)} Self-contained:
\[
\delta^*_C(2^{-128})\ \in\ \Bigl[\,\tfrac{349525}{2^{21}},\ \tfrac{15331}{32768}\,\Bigr]\ \approx\ [\,0.16666,\ 0.46787\,] .
\]
With the imported unique-decoding proximity gap \textup{(\cref{cor:conditional-half})}:
\[
\delta^*_C(2^{-128})\ \in\ \Bigl[\,\tfrac14,\ \tfrac{15331}{32768}\,\Bigr],
\]
an open band of width $<0.218$, down from $0.326$ in \cref{thm:sandwich}.
\item \textup{(Circle line-round row, $\varepsilon^*=2^{-100}$.)} Self-contained: $\delta^*_C\in[\tfrac{349525}{2^{21}},\tfrac{30663}{65536}]$; with the import: $\delta^*_C\in[\tfrac14,\tfrac{30663}{65536}]$.
\item \textup{(General envelope rows.)} For any row with $q\ge2^{128}(n+1)$ and a map-smooth scale lattice: the safe edge is $\max\bigl(\lfloor(n-k)/3\rfloor,\,r_{\mathrm{hJ}}\bigr)/n$ self-contained, where $r_{\mathrm{hJ}}$ is the largest radius passing \cref{cor:self-contained-safe}, and $\lfloor(n-k)/2\rfloor/n$ with the import; the unsafe edge is $1-\rho-g$ for every exactly certified $g$, with envelope $g^*(\rho,\log_2|\B|)$ of \cref{rem:entropy-frontier}.
\end{enumerate}
\end{theorem}

\begin{proof}
Lower edges: \cref{thm:deep-mca} at $r=349525=\lfloor(n-k)/3\rfloor$ (errors $349526/q<2^{-167}$ resp.\ $<2^{-105}$, below the respective targets, as in \cref{tab:scanner}); \cref{cor:self-contained-safe} where it exceeds the deep edge; \cref{cor:conditional-half} for the imported edge $\tfrac14$ (errors $2^{-164}$ resp.\ $2^{-102}$, below the respective targets). Upper edges: \cref{cor:widened-deployed} and \cref{rem:exact-frontier} certify $\emca>2^{-22}\ge\varepsilon^*$ (resp.\ $>2^{-100}$ via the circle thresholds) on $[\,\text{edge},\tfrac12)$, so $\delta^*_C$ cannot exceed the printed edge. Part \textup{(iii)} collects the corresponding general statements.
\end{proof}

\begin{table}[t]
\centering\footnotesize
\setlength{\tabcolsep}{4pt}
\begin{tabular}{@{}llllll@{}}
\toprule
row & challenge & unsafe \textup{(widened)} & safe, self-contained & safe, with import & open band \\
\midrule
KoalaBear sextic & MCA, $2^{-128}$ & $[\,0.46787,\tfrac12)$ & $\delta\le0.16666$ & $\delta\le0.25$ & $(0.25,\ 0.46787)$ \\
KoalaBear sextic & list, $2^{-128}$ & $\delta\ge0.46787$ & $\delta\le0.29289$ & --- & $(0.29290,\ 0.46787)$ \\
circle, $p'{=}2^{31}{-}1$ & MCA, $2^{-100}$ & $[\,0.46788,\tfrac12)$ & $\delta\le0.16666$ & $\delta\le0.25$ & $(0.25,\ 0.46788)$ \\
circle, $p'{=}2^{31}{-}1$ & list, $2^{-100}$ & $\delta\ge0.46787$ & $\delta\le0.29289$ & --- & $(0.29290,\ 0.46787)$ \\
\bottomrule
\end{tabular}
\caption{Scanner output after the pincer, superseding \cref{tab:scanner}. ``Unsafe (widened)'' is certified by \cref{cor:widened-deployed} and \cref{rem:exact-frontier} (edges rounded outward to five digits: the exact edges are $\tfrac{15331}{32768}$ and $\tfrac{30663}{65536}$). ``Safe, self-contained'' is \cref{thm:deep-mca} resp.\ \cref{thm:johnson-list} as in \cref{tab:scanner}; at rates below $\tfrac14$, \cref{cor:self-contained-safe} would replace the MCA entries. ``Safe, with import'' is \cref{cor:conditional-half}; the list rows need no import, their safe edge being the self-contained Johnson bound. Open bands are stated for the stronger certificate in each column.}
\label{tab:scanner2}
\end{table}

The reduced kernel, restated once. After the pincer, the open MCA band of each deployed row is $(\tfrac14,\ \text{edge})$ with edge $\approx0.4679$, modulo one import whose own proof is independent of everything here; self-contained, it is $(0.1667,\ \text{edge})$, and at rates below $\tfrac14$ the left endpoint improves to half the Johnson radius. Below half the distance the mutual and plain correlated-agreement questions are now provably the same question up to an explicit tangent term (\cref{thm:mca-from-ca}), so Open Problem~\ref{prob:band} survives only above half the distance, and there the certified unsafe mass presses down to the entropy--subfield envelope $g^*$ (\cref{rem:entropy-frontier}), which the deployed rows already saturate. Closing the band therefore requires one of exactly two kinds of new mathematics: a floor construction whose received words carry more than a graded locator prefix --- anything beating $g^*$ --- or a correlated-agreement theorem crossing half the distance toward capacity, to which \cref{thm:mca-from-ca} would immediately add the mutual layer. This is, for the first time in the program, a two-sided quantitative target with both envelopes made explicit.

\section{Further refinements and finite certificate closure}\label{sec:answers}

This section records three groups of consequences that are not part of the main unsafe/safe sandwich but complete the account.  First, graded rational floors extend the unsafe construction to genus-one scale maps, so the ECFFT obstruction of \cref{rem:ecfft-obstruction} is a one-slot obstruction rather than a macroscopic barrier (\cref{sec:answers-isogeny}), and a stereographic model treats circle codes over challenge fields that do not contain $i$, completing the circle-code transfer (\cref{sec:answers-stereo}).  Second, explicit head words, simple-pole certifying pairs, and the optimized failure profile clarify what is constructive in the failure mechanism and what remains beyond the present method (\cref{sec:answers-explicit,sec:answers-profile}).  Third, \cref{sec:answers-T7} closes the finite certificate grammar over every verdict used in the narrowed sandwich, prints the deployed certificates as integer tuples, bounds verification complexity, and isolates the single imported input to one optional clause.

\subsection{Graded floors on rational scales: the genus-one band is macroscopic}\label{sec:answers-isogeny}

The graded prefix floor of \cref{prop:graded-prefix-floor} pays one subfield factor per locator slot and buys $c$ degrees of agreement per slot on a polynomial scale of degree $c$. On a rational scale of degree type $(a,e)$ the slot spacing is $a-e$ rather than $a$; nothing else changes.

\begin{proposition}[graded prefix floor for rational scales]\label{prop:graded-rational-floor}
Let $\B\subseteq\F$ be a subfield and let $D\subseteq\B$ with $|D|=n$ be $(\psi,a)$-smooth in the sense of \cref{def:rational-smooth}, with $\psi=f/g$ of degree type $(a,e)$, $a>e\ge0$; put $N:=n/a$ and $Q:=\psi(D)\subseteq\B$. Let $1\le K\le n$ and let $m$ satisfy
\[
am\ \ge\ K,\qquad em\ \le\ K,\qquad m\ \le\ N,
\qquad\text{and put}\qquad
w\ :=\ \Bigl\lfloor\frac{am-K}{a-e}\Bigr\rfloor\ \le\ m .
\]
Then there is a $\B$-valued received word $U:D\to\B$ with
\[
\bigl|\Lst\bigl(\RS[\F,D,K],\,1-\tfrac{am}{n},\,U\bigr)\bigr|
\ \ge\
\left\lceil\frac{\binom Nm}{|\B|^{\,w}}\right\rceil .
\]
For $e=0$ this is \cref{prop:graded-prefix-floor} \textup(then $w=m-\lceil K/a\rceil$\textup); for $w\le1$ it recovers the slack-two window of \cref{prop:rational-floor}.
\end{proposition}

\begin{proof}
For $M\subseteq Q$ with $|M|=m$ put $\Lambda_M:=\prod_{b\in M}(f-bg)$. Each factor has degree exactly $a$ with leading coefficient $\mathrm{lc}(f)$, so $\deg\Lambda_M=am$ with leading coefficient $\mathrm{lc}(f)^m\ne0$; and $\Lambda_M\in\B[X]$ vanishes at $x\in D$ if and only if $\psi(x)\in M$, i.e.\ at exactly $am$ points of $D$, by smoothness and $g\ne0$ on $D$. Expanding,
\[
\Lambda_M\ =\ \sum_{j=0}^{m}(-1)^je_j(M)\,f^{\,m-j}g^{\,j},
\]
with $j$-th term of degree $am-j(a-e)$, strictly decreasing in $j$. For $z=(z_1,\dots,z_w)\in\B^w$ define the $\B$-valued word
\[
U_z\ :=\ \Bigl(f(x)^m+\sum_{j=1}^{w}z_j\,f(x)^{\,m-j}g(x)^{\,j}\Bigr)_{x\in D},
\]
which makes sense because $w\le m$: indeed $am-K\le m(a-e)$ is the hypothesis $em\le K$. By pigeonhole over the prefix map $M\mapsto z(M):=\bigl((-1)^je_j(M)\bigr)_{j\le w}\in\B^w$ there is $z^*$ attained by at least $\lceil\binom Nm/|\B|^w\rceil$ sets $M$. For each such $M$ set
\[
c_M\ :=\ U_{z^*}-\Lambda_M|_D\ =\ \Bigl(-\sum_{j>w}(-1)^je_j(M)\,f^{\,m-j}g^{\,j}\Bigr)_{x\in D},
\]
of degree at most $am-(w+1)(a-e)\le K-1$ by the choice of $w$, so $c_M\in\RS[\B,D,K]\subseteq\RS[\F,D,K]$; and $U_{z^*}$ agrees with $c_M$ on the $am$ fiber points of $M$, i.e.\ within relative radius $1-am/n$. Injectivity of $M\mapsto c_M$ on the pigeonhole class: equality of words of degree $\le K-1\le n-1$ forces equality of polynomials, hence $\sum_{j>w}(-1)^j\bigl(e_j(M)-e_j(M')\bigr)f^{\,m-j}g^{\,j}=0$; dividing by $g^m$ in $\F(X)$ gives a polynomial identity in $\psi$, and $\psi$ is nonconstant, hence transcendental over $\F$ inside $\F(X)$, so all coefficients vanish, exactly as in the proof of \cref{prop:rational-floor}. Thus $e_j(M)=e_j(M')$ for $j>w$; for $j\le w$ they agree through $z^*$; hence $\prod_{b\in M}(Y-b)=\prod_{b\in M'}(Y-b)$ and $M=M'$. The $c_M$ are therefore distinct listed codewords for the single word $U_{z^*}$.
\end{proof}

\begin{corollary}[macroscopic universal cap for genus-one rows]\label{cor:ecfft-macroscopic}
Let $D$ be $(\psi,2)$-smooth over $\B$ with $\psi$ of degree type $(2,1)$ \textup(\cref{def:rational-smooth}; in particular every FFTree domain of ECFFT \cite{BCKL21}\textup), let $C=\RS[\F,D,k]$ with $\rho:=k/n\in[\tfrac1{16},\tfrac12]$, $n\ge2^{14}$, $n<q=|\F|<2^{256}$, and $\beta:=\log_2|\B|\le64$. Put $m:=\lceil(k+1+2^{-9}n)/2\rceil$ and $N:=n/2$. Then
\[
\binom Nm\ >\ |\B|^{\,2m-k-1}\Bigl(\frac qk+1\Bigr)
\qquad\text{and}\qquad
\binom Nm\ >\ |\B|^{\,2m-k}\,2^{-128}q ,
\]
and consequently:
\begin{enumerate}[label=\textup{(\roman*)}]
\item \textup(MCA/CA cap\textup) $\eca\bigl(C,\delta\bigr)>\frac1{2k}\bigl(1-\frac nq\bigr)$ for every integer-admissible $\delta\in[1-\frac{2m}n,\,1-\rho-\frac1n]$, and $\emca(C,\delta)$ obeys the same lower bound for every $\delta\in[1-\rho-2^{-9},\,1-\rho)$;
\item \textup(list cap\textup) there is a $\B$-valued word $U$ with $|\Lst(C,1-\frac{2m}n,U)|>2^{-128}q$, so the grand list challenge is unsafe at target $2^{-128}$ on the same band.
\end{enumerate}
The gap $2^{-9}$ coincides with the universal multiplicative gap of \cref{cor:grand}: genus-one rows are no longer behind the polynomial-scale rows.
\end{corollary}

\begin{proof}
Write $w_1:=2m-k-1$ and $w_0:=2m-k$ for the two prefix depths ($K=k+1$ and $K=k$). From $k+1+2^{-9}n\le2m\le k+2^{-9}n+3$: first, $w_0\le2^{-9}n+3=N/256+3$, so $\beta w_0\le N/4+192$; second, $m/N\in[\rho+2^{-9},\,\rho+2^{-9}+2^{-12}]$ using $3/n\le3\cdot2^{-14}<2^{-12}$. On the interval $[\tfrac1{16}+2^{-9},\,\tfrac12+2^{-9}+2^{-12}]$ the concave $\Hb$ is minimized at an endpoint, and $\Hb(\tfrac1{16}+2^{-9})\ge0.344$, $\Hb(\tfrac12+2^{-9}+2^{-12})\ge0.999$, so $\Hb(m/N)\ge0.344$. Hence
\[
\log_2\binom Nm\ \ge\ N\,\Hb(m/N)-\log_2(N+1)\ \ge\ 0.344\,N-\log_2(N+1),
\]
while both right-hand sides are at most $2^{\beta w_0+256}\le2^{N/4+448}$ crudely \textup(using $q/k+1\le q<2^{256}$ and $2^{-128}q<2^{128}$\textup). Since $0.094\,N\ge0.094\cdot2^{13}=770>448+\log_2(N+1)$ for all $N\ge2^{13}$, both displayed inequalities hold. The hypotheses of \cref{prop:graded-rational-floor} hold for both $K\in\{k,k+1\}$: $2m\ge k+1\ge K$; $em=m\le k\le K$ because $2^{-9}n+3\le k$ for $\rho\ge\tfrac1{16}$, $n\ge2^{14}$; and $m\le N$ because $2m\le k+2^{-9}n+3\le n$. Part \textup{(ii)} is the floor at $K=k$. For \textup{(i)}, the floor at $K=k+1$ exceeds $q/k+1\ge\lceil2q\,\eca\rceil$ whenever $\eca(C,\delta)\le\frac12(\frac1k-\frac n{kq})$, contradicting \cref{thm:A} at every integer-admissible $\delta\ge\delta_0:=1-2m/n$ \textup(admissible since $n-2m\le n-k-1$\textup); the assembly through \cref{fact:chain} at $\delta_0$ and \cref{lem:mca-monotone} up to $1-\rho$ is verbatim the proof of \cref{thm:phi-cap}. The band containment is $1-2m/n\le1-\rho-2^{-9}$, i.e.\ $2m\ge k+2^{-9}n$.
\end{proof}

\begin{remark}[the trichotomy of \cref{rem:ecfft-obstruction}, resolved]\label{rem:trichotomy}
\Cref{rem:ecfft-obstruction} offered three exits from the lacunarity collapse: a lacunary basis of isogeny locators, a different received-word family, or a matching safe-side theorem showing the collapse essential. The second exit is the one that exists, and the collapse was not essential: the one-slot analysis priced the window width $2(a-e)$ as the \emph{reach} of the route, when it only sets the \emph{exchange rate} --- $a-e$ degrees of agreement per subfield factor, against an entropy budget of $\Hb$ bits per fiber. Optimizing the number of slots gives the closed-form frontier
\[
\Hb(\rho+g^*)\ =\ g^*\cdot\beta\cdot\frac a{a-e},
\]
so the degree type enters only through the ramification-defect ratio $a/(a-e)$: polynomial scales ($e=0$) have ratio $1$ and recover \cref{rem:entropy-frontier}; genus-one $x$-maps have $e=a-1$, ratio $a$, optimized at the first fold $a=2$, i.e.\ an effective subfield cost of $2\beta$. Numerically $g^*(\tfrac12,62)\approx0.0161\approx2^{-5.96}$ at $\beta=31$ and $g^*(\tfrac12,128)\approx0.0078\approx2^{-7}$ at $\beta=64$: macroscopic in every case, against the $2/n$ of the one-slot route. The safe side needs no new work: \cref{thm:deep-mca,thm:mca-from-ca,thm:elementary-ca,cor:conditional-half} are domain-agnostic, so every $(\psi,2)$-smooth row carries the full narrowed sandwich of \cref{thm:sandwich2}'s shape --- self-contained safe frontier at $(1-\rho)/3$, safe at $(1-\rho)/2$ modulo the single import, unsafe from gap $2^{-9}$ \textup(and, row by row, from the optimized frontier $g^*(\rho,2\beta)$\textup). This answers Open Problem~\ref{prob:isogeny}\textup{(a)} as posed.
\end{remark}

\subsection{The stereographic model: circle codes over every challenge field}\label{sec:answers-stereo}

The torus uniformization of \cref{lem:circle-rs} needs $i\in\F$. The circle is a conic with the rational point $(1,0)$, so it also has a uniformization defined over $\F_p$ itself --- the stereographic projection --- and this one transports the code theory to \emph{every} challenge field. The price is that the resulting Reed--Solomon domain is no longer a multiplicative coset; but it is rationally smooth of the same $(2,1)$ type as the genus-one rows, so \cref{sec:answers-isogeny} covers it.

\begin{lemma}[stereographic uniformization of circle codes; no $i$ required]\label{lem:stereographic}
Let $p\equiv3\pmod4$, let $\mathcal D=gH\cup g^{-1}H\subseteq\UU$ be a twin coset of size $2M$ with $\operatorname{ord}(g)\nmid2M$ and $4\mid M$, let $E:=\iota^{-1}(\mathcal D)\subseteq\mathcal C(\F_p)$, and let $\F\supseteq\F_p$ be \emph{any} finite field. Define $s(x,y):=y/(1+x)$ and $T:=s(E)$.
\begin{enumerate}[label=\textup{(\roman*)}]
\item The four torsion points $(\pm1,0),(0,\pm1)$ lie in $\iota^{-1}(H)$ and hence not in $E$; $s$ is injective on $E$ with inverse $s\mapsto\bigl(\frac{1-s^2}{1+s^2},\frac{2s}{1+s^2}\bigr)$; and $T\subseteq\F_p\setminus\{0,\pm1\}$ with $|T|=2M$. \textup(Over $\F_p$ one has $1+s^2\ne0$ throughout, since $p\equiv3\pmod4$.\textup)
\item With the diagonal vector $d:=\bigl((1+s(P)^2)^{-w}\bigr)_{P\in E}$ and the coordinate bijection $s|_E:E\to T$,
\[
\mathcal C_w(\F,E)\ =\ d\circ\RS[\F,T,2w+1].
\]
Consequently, by \cref{lem:diag-invariance}, the bivariate circle code on $E$ and $\RS[\F,T,2w+1]$ have identical list sizes at every radius and identical $\eca$ and $\emca$ at every \textup(pair of\textup) radii --- with no hypothesis on $i\in\F$.
\item $T$ is $(\psi,2)$-smooth over $\B:=\F_p$ in the sense of \cref{def:rational-smooth} for the tangent-halving map $\psi(s):=\frac{s^2-1}{2s}$, of degree type $(2,1)$: the $\psi$-fibers within $T$ are the pairs $\{s(P),s(P\tau)\}=\{s,-1/s\}$, $\tau:=(-1,0)$, complete because $\tau\in\iota^{-1}(H)$ and of uniform size $2$ because $s^2\ne-1$ in $\F_p$; and $N=M$, $Q=\psi(T)\subseteq\F_p$.
\end{enumerate}
\end{lemma}

\begin{proof}
(i) $H$ is the unique subgroup of order $M$ of the cyclic group $\UU$ of order $p+1$; since $4\mid M$ it contains all elements of order dividing $4$, whose $\iota^{-1}$-images are exactly the four torsion points. From $\operatorname{ord}(g)\nmid2M$ we get $g,g^{-1}\notin H$, so $gH\cap H=g^{-1}H\cap H=\emptyset$ and $\mathcal D\cap H=\emptyset$; in particular $(-1,0)\notin E$, so $s$ is defined on $E$, and $(1,0),(0,\pm1)\notin E$, so $0,\pm1\notin T$. For the inverse formula: from $x^2+y^2=1$, $1+s^2=\frac{(1+x)^2+y^2}{(1+x)^2}=\frac2{1+x}$, whence $\frac{1-s^2}{1+s^2}=(1+x)-1=x$ and $\frac{2s}{1+s^2}=s(1+x)=y$; and $1+s^2\ne0$ for $s\in\F_p$ because $-1$ is a nonsquare when $p\equiv3\pmod4$.

(ii) Every class of degree $\le w$ has the canonical form $f_0(x)+y\,f_1(x)$ with $\deg f_0\le w$, $\deg f_1\le w-1$, a space of dimension $2w+1$ \textup(proof of \cref{lem:circle-rs}\textup). Substituting the inverse of (i) and clearing denominators,
\[
(1+s^2)^w\bigl(f_0(x)+y\,f_1(x)\bigr)\Big|_{x=\frac{1-s^2}{1+s^2},\,y=\frac{2s}{1+s^2}}
\]
is a polynomial in $s$ of degree at most $2w$: the monomial $x^j$ contributes $(1-s^2)^j(1+s^2)^{w-j}$ and $y\,x^j$ contributes $2s\,(1-s^2)^j(1+s^2)^{w-1-j}$. These $2w+1$ polynomials form a basis of $\F[s]_{\le2w}$: the first family is even, the second odd, so it suffices to check each family separately; and in the variable $v:=s^2$ the first family is $\{(1-v)^j(1+v)^{w-j}\}_{j\le w}$, the image of the monomial basis $\{Y^jZ^{w-j}\}$ under the invertible substitution $(Y,Z)=(1-v,1+v)$ \textup($p$ odd\textup), and similarly for the odd family. Hence $F\mapsto(1+s^2)^wF$ is a linear isomorphism from the circle space onto $\F[s]_{\le2w}$, and evaluating at $P\in E$ gives $F(P)=(1+s(P)^2)^{-w}\,f(s(P))$ with $f\in\F[s]_{\le2w}$ ranging over all of $\RS[\F,T,2w+1]$: the displayed code identity. The vector $d$ is $\F_p$-valued and everywhere nonzero, so \cref{lem:diag-invariance} applies.

(iii) The group law is $(x_1,y_1)\cdot(x_2,y_2)=(x_1x_2-y_1y_2,\,x_1y_2+x_2y_1)$, so $P\tau=(-x,-y)$ and
\[
s(P)\,s(P\tau)\ =\ \frac y{1+x}\cdot\frac{-y}{1-x}\ =\ \frac{-y^2}{1-x^2}\ =\ -1,
\]
using $x^2+y^2=1$ and $y\ne0$ on $E$. Since $\tau\in\iota^{-1}(H)$ \textup($2\mid M$\textup) and $\mathcal D H=\mathcal D$ coset-wise, $E\tau=E$, so the involution $s\mapsto-1/s$ preserves $T$. The fibers of $\psi$: $\psi(s)=b$ if and only if $s^2-2bs-1=0$, whose two roots multiply to $-1$; hence every fiber is $\{s,-1/s\}$, of size exactly $2$ since $s^2=-1$ is impossible in $\F_p$, and complete within $T$ by the closure just shown. Finally $f=X^2-1$ and $g=2X$ are coprime of degrees $2>1$, and $g$ is nonvanishing on $T$ because $0\notin T$.
\end{proof}

\begin{corollary}[circle codes over every challenge field; Open Problem~\ref{prob:isogeny}\textup{(b)}]\label{cor:circle-anyfield}
Let $E$ be a standard-position circle domain as in \cref{lem:stereographic} with $M=2^{21}$ \textup(the deployed shape: $n_c=2M=2^{22}$, $k_c=2^{21}+1$\textup), and let $\F\supseteq\F_p$, $p=2^{31}-1$, be any challenge field, $q=|\F|$.
\begin{enumerate}[label=\textup{(\roman*)}]
\item If $q<2^{100}$ --- in particular $\F=\F_p$ and $\F=\F_{p^3}$, hence every field without $i$ among the extensions $\F_{p^r}$, $r\le4$, that deployed verifiers sample from, since $i\in\F_{p^r}$ if and only if $2\mid r$ --- the row is degenerately unsafe at every deployed target by \cref{prop:small-field}, and Open Problem~\ref{prob:isogeny}\textup{(b)} is out of scope there as posed.
\item If $i\in\F$, \cref{cor:circle-grand,cor:circle-deployed} apply verbatim through \cref{lem:circle-rs}.
\item If $i\notin\F$ and $q\ge2^{100}$ \textup(the first such field is $\F_{p^5}$\textup), the stereographic model of \cref{lem:stereographic} transports the row to $\RS[\F,T,k_c]$ on the $(\psi,2)$-smooth tangent domain $T$, and \cref{prop:graded-rational-floor} with
\[
m=2^{20}+2^{15}+1,\qquad w=2m-(k_c+1)=2^{16},\qquad
\binom{2^{21}}{m}\ >\ p^{\,2^{16}}\cdot2^{256}
\]
\textup(verified exactly in integer arithmetic; the margin exceeds $2^{63000}$\textup) gives, for every $q\in(n_c,2^{256})$ at once, an MCA and list cap on
\[
\delta\ \in\ \bigl[\,1-\rho_c-2^{-6},\ 1-\rho_c\bigr),
\qquad
\text{error}\ >\ \tfrac1{2k_c}\bigl(1-\tfrac{n_c}q\bigr)\ >\ 2^{-23},
\]
by the assembly of \cref{cor:ecfft-macroscopic} \textup(here $\beta=\log_2p$, effective cost $2\beta$, asymptotic frontier $g^*\approx2^{-5.96}$\textup). The safe side is generic: \cref{thm:deep-mca} applies to the circle code directly \textup(as already noted after its proof\textup), and \cref{thm:mca-from-ca,thm:elementary-ca,cor:conditional-half} apply to the diagonal model. Every bivariate circle row over every challenge field therefore carries a two-sided sandwich, and the deep-point conversion \cref{thm:A} does extend to the bivariate setting --- through the stereographic model rather than a new proof.
\end{enumerate}
\end{corollary}

\begin{proof}
(i) is \cref{prop:small-field} with $q=p<2^{31}$ resp.\ $q=p^3<2^{93}$: $\emca\ge1/q>2^{-100}$ at every radius, so $\delta^*_C(\varepsilon^*)=0$ for every deployed target $\varepsilon^*\le2^{-100}$; the parity criterion is $\operatorname{ord}(i)=4\mid p^r-1$ if and only if $r$ is even, as $p\equiv3\pmod4$. (ii) is quoted. (iii): the transport is \cref{lem:stereographic}\textup{(ii)}, the smoothness is \textup{(iii)}, and the hypotheses of \cref{prop:graded-rational-floor} hold with $K\in\{k_c,k_c+1\}$: $2m=k_c+1+2^{16}\ge K$, $em=m\le k_c\le K$, $m\le N=2^{21}$. The single displayed integer inequality dominates both thresholds for every $q<2^{256}$, since $q/k_c+1\le q<2^{256}$ and $2^{-100}q<2^{156}$; the edge satisfies $1-2m/n_c=1-\rho_c-2^{-6}-2^{-21}+2^{-22}\le1-\rho_c-2^{-6}$, and $\lfloor\delta_0 n_c\rfloor=n_c-2m\le n_c-k_c-1$, so \cref{thm:A}, \cref{fact:chain}, and \cref{lem:mca-monotone} assemble exactly as in \cref{cor:ecfft-macroscopic}; the error numerics use $q\ge p^5$, so $\frac1{2k_c}\bigl(1-\frac{n_c}q\bigr)>2^{-22}(1-2^{-21})(1-2^{-133})>2^{-23}$.
\end{proof}

\subsection{Explicit witness words and explicit certifying pairs}\label{sec:answers-explicit}

Every floor so far selects its received word by pigeonhole over a prefix. At the first staircase step the pigeonhole can be removed entirely: if the support is assembled from \emph{antipodal pairs}, its first elementary symmetric function vanishes identically, and the received word collapses to a pure power.

\begin{theorem}[explicit head-and-pairs floor]\label{thm:explicit-head-floor}
Let $D$ be $(\varphi,c)$-smooth over $\B$ \textup(\cref{def:map-smooth}\textup) with $N:=n/c$ even, and suppose $Q:=\varphi(D)$ satisfies $-Q=Q$ and $0\notin Q$, so that $Q$ is partitioned into $N/2$ antipodal classes $\{y,-y\}$. Let $K\le n$ and let $m$ satisfy $2\le m\le N-2$ and $cm\le K-1+2c$.
\begin{enumerate}[label=\textup{(\roman*)}]
\item If $m$ is even, the \emph{pure power word} $u:=\varphi^m|_D$ satisfies
\[
\bigl|\Lst\bigl(\RS[\F,D,K],\,1-\tfrac{cm}n,\,u\bigr)\bigr|\ \ge\ \binom{N/2}{m/2},
\]
the listed codewords being $c_M:=u-\Lambda_M|_D$ for $M$ a union of $m/2$ antipodal classes.
\item If $m$ is odd, then for every $t\in Q$ the \emph{one-head word} $u:=\bigl(\varphi^m-t\,\varphi^{m-1}\bigr)|_D$ satisfies
\[
\bigl|\Lst\bigl(\RS[\F,D,K],\,1-\tfrac{cm}n,\,u\bigr)\bigr|\ \ge\ \binom{N/2-1}{(m-1)/2},
\]
with $M=\{t\}\cup(\text{$(m-1)/2$ antipodal classes distinct from $\{t,-t\}$})$.
\end{enumerate}
In both cases the word, the family, and every listed codeword are explicit closed forms, with no pigeonhole and no subfield division. Both deployed towers qualify: on a multiplicative coset, $Q$ is a coset of the even-order subgroup $\mu_N$, so $-Q=Q$ and $0\notin Q$; on a Chebyshev domain $D=\chi(\mathcal D)$ with $2c\mid M$, the antipodal classes of $Q=T_c(D)$ are exactly the $T_2$-fibers inside $Q$, complete of size $2$ by \cref{lem:cheb-fibers} at scales $c$ and $2c$ \textup(two elements with equal $T_2$-value satisfy $y'=-y$, and $0\notin Q$ since a $T_2$-fiber $\{0\}$ would have size $1$\textup).
\end{theorem}

\begin{proof}
In case \textup{(i)}, $M$ is a union of antipodal pairs, so $e_1(M)=0$ and $\Lambda_M=\varphi^m+\sum_{j\ge2}(-1)^je_j(M)\varphi^{m-j}$, whence $c_M:=u-\Lambda_M|_D=-\sum_{j\ge2}(-1)^je_j(M)\varphi^{m-j}|_D$ has degree at most $c(m-2)\le K-1$ and agrees with $u$ on the $cm$ fiber points of $M$. In case \textup{(ii)}, $e_1(M)=t+\sum(\text{pair sums})=t$, so the slot-$1$ coefficient of $\Lambda_M$ is $-t\varphi^{m-1}$ for every $M$ in the family, and $c_M:=u-\Lambda_M|_D$ again consists of the slots $j\ge2$ only, of degree at most $c(m-2)\le K-1$; note that $e_2(M)$ and beyond are permitted to vary, since they land in the codeword rather than in the received word. Injectivity: $c_M=c_{M'}$ as words forces $\Lambda_M=\Lambda_{M'}$ on all of $D$, and both are polynomials of degree $cm<cN=n$, hence equal as polynomials, with equal root multisets, hence $M=M'$ through the fibers. The family counts are the number of choices of antipodal classes.
\end{proof}

\begin{corollary}[explicit witnesses at the deployed rows]\label{cor:explicit-deployed}
At the KoalaBear row of \cref{cor:deployed} \textup($n=2^{21}$, $k=2^{20}$, $\B=\F_p$, $q=p^6$\textup) and at the Chebyshev line-round row of \cref{cor:circle-deployed}\textup{(a)} \textup($n=2^{21}$, $k=2^{20}$, $q=p'^4$, $p'=2^{31}-1$, target $2^{-100}$\textup):
\begin{enumerate}[label=\textup{(\roman*)}]
\item \textup(MCA, gap $2^{-8}$\textup) With $c=2^{12}$, $m=k/c+2=258$, $K=k+1$: the explicit word $u=X^{k+2^{13}}|_D$ \textup(resp.\ $u=T_c^{258}|_D$\textup) has list at agreement $k+2^{13}$ at least $\binom{256}{129}\ge2^{251}$, exceeding $q/k+1<2^{166}$ \textup(resp.\ $<2^{104}$\textup); by \cref{thm:A} with $\eta=\frac12$, exactly as in \cref{thm:main}, $\eca(C,\delta)>\frac1{2k}(1-\frac nq)$ on the integer-admissible part of $[\frac12-2^{-8},\frac12)$ and $\emca$ obeys the same bound throughout it.
\item \textup(lists, gap $2^{-7}$\textup) With $c=2^{14}$, $m=k/c+1=65$, $K=k$: for any $t\in Q$ the explicit word $u=\bigl(X^{k+2^{14}}-t\,X^{k}\bigr)|_D$ \textup(resp.\ its $T_c$ analogue\textup) has list at agreement $k+2^{14}$ at least $\binom{63}{32}=916312070471295267>2^{59}$, exceeding $2^{-128}q<2^{58}$ \textup(resp.\ $2^{-100}q<2^{24}$\textup): the first staircase step of the grand list challenge, at its original printed edge $\frac12-2^{-7}$, is certified by a fully explicit word.
\end{enumerate}
All four inequalities are verified exactly in integer arithmetic.
\end{corollary}

\begin{proof}
The hypotheses of \cref{thm:explicit-head-floor} hold with $N=512$ resp.\ $128$: parity as displayed, $cm\le K-1+2c$ \textup(with equality in \textup{(i)}\textup), and the Chebyshev scale conditions $2c\mid M=2^{21}$. The threshold comparisons are the displayed integer inequalities; the MCA conversion and monotone assembly are as quoted.
\end{proof}

The words above certify lists explicitly, and \cref{thm:A} converts large lists into CA failure --- but by contradiction, without exhibiting the failing pair. Open Problem~\ref{prob:explicit}, however, asks for the \emph{pair}. The proof of \cref{thm:A} is more generous than its statement: it converts through the completely explicit simple-pole lines $\bigl(U/(x-\alpha),\,-1/(x-\alpha)\bigr)$ --- precisely the residue-line normal form with extension denominator that \cref{cor:Fvalued} predicts --- and its collision count can be read as a lower bound holding for \emph{most} poles.

\begin{theorem}[explicit certifying pairs, up to the choice of a pole]\label{thm:explicit-pairs}
At the KoalaBear row of \cref{cor:deployed}, let $u:=X^{k+2^{13}}|_D$ be the pure power word of \cref{cor:explicit-deployed}\textup{(i)}, let $L_0:=\lceil(q-n)/k\rceil<2^{166}$, and fix any explicit subfamily of $L_0$ of its listed codewords $P_M\in\F[X]_{<k+1}$ \textup(e.g.\ the lexicographically least $L_0$ antipodal-class sets; $L_0\le\binom{256}{129}$\textup). For $\alpha\in\Omega:=\F\setminus D$ define the explicit pair
\[
f_\alpha(x):=\frac{u(x)}{x-\alpha},
\qquad
g_\alpha(x):=-\frac1{x-\alpha}
\qquad(x\in D).
\]
Then for at least half of all $\alpha\in\Omega$, the pair $(f_\alpha,g_\alpha)$ has at least
\[
\Bigl\lceil\frac{q-n}{3k}\Bigr\rceil
\]
distinct CA-bad slopes at radius $\delta=\frac12-2^{-8}$, namely among the explicit values $\{P_M(\alpha)\}$; its CA-bad-slope density exceeds $2^{-21.6}>2^{-22}$ there and at every larger integer-admissible radius, and its MCA-bad-slope density exceeds the same bound throughout $[\frac12-2^{-8},\frac12)$. All but at most $p$ of these poles have $\alpha\notin\B$, and for those the pair is genuinely $\F$-valued, in accordance with \cref{cor:Fvalued}. The same construction on the circle line-round row gives density $>2^{-21.6}>2^{-23}$ for at least half of all poles.
\end{theorem}

\begin{proof}
Each $P_M$ agrees with $u$ on at least $a=k+2^{13}$ points, so, as in the proof of \cref{thm:A}, for every $\alpha\in\Omega$ and every $M$ the slope $P_M(\alpha)$ carries the point $f_\alpha+P_M(\alpha)g_\alpha$ to within radius $1-a/n=\frac12-2^{-8}$ of $C$, while the pair is farther than every $\delta<1-k/n=\frac12$ from $C^{\equiv2}$ by the $(X-\alpha)G+1$ argument there; so every distinct value $P_M(\alpha)$ is a CA-bad slope at every integer-admissible $\delta\in[\frac12-2^{-8},\frac12)$, hence \textup(as failing CA is failing MCA, \cref{fact:chain}\textup) an MCA-bad slope throughout that interval. It remains to count distinct values for many $\alpha$ at once. Distinct sets $M$ give distinct polynomials $P_M$ \textup(\cref{thm:explicit-head-floor}\textup), and $P_M-P_{M'}$ is nonzero of degree at most $k$, so the total number of colliding pairs over all $\alpha\in\Omega$ is at most $k\binom{L_0}2$. Let $x^*:=\lceil2k\binom{L_0}2/(q-n)\rceil$; by Markov, fewer than $(q-n)/2$ poles have collision count at least $x^*$, so at least half have collision count at most $x^*-1$. For such $\alpha$, Cauchy--Schwarz as in \cref{thm:A} gives
\[
M(\alpha)\ \ge\ \frac{L_0^2}{L_0+2(x^*-1)}\ \ge\ \frac{L_0}{1+2k(L_0-1)/(q-n)}\ >\ \Bigl\lceil\frac{q-n}{3k}\Bigr\rceil-1,
\]
so, since $M(\alpha)$ is an integer,
$M(\alpha)\ge\lceil(q-n)/(3k)\rceil$. The last comparison and
$\lceil(q-n)/3k\rceil/q>2^{-21.6}$ are verified exactly in integer arithmetic
at both rows. The subfield count is $|\Omega\cap\B|\le|\B|=p$; for
$\alpha\notin\B$ the word $g_\alpha$ is not $\B$-valued \textup(its value at any
$x\in D\subseteq\B$ lies in $\B$ only if $x-\alpha\in\B$\textup), and no common
scalar repairs this while keeping $f_\alpha$ a $\B$-multiple, so the pair falls
outside the inert class of \cref{lem:confine}, as \cref{cor:Fvalued} requires
of any pair with this density.
\end{proof}

\begin{remark}[what remains of Open Problem~\ref{prob:explicit}]\label{rem:explicit-scope}
\Cref{thm:explicit-pairs} answers Open Problem~\ref{prob:explicit} at $\delta=\frac12-2^{-8}$, on the upper half of the stated band, in the predicted normal form, and for both densities the problem asks about; the non-constructiveness has collapsed from a received word --- an object of $n\log_2q$ bits with no known certificate --- to the choice of one field element in an explicit majority subset of $\F\setminus D$. Two residues remain. First, a certified \emph{individual} pair: derandomizing the pole. Second, the left half of the band: at gap exactly $2^{-7}$ the explicit family at scale $2^{13}$ has only $\binom{128}{65}<2^{125}$ members, short of the $\approx2^{166}$ the counting needs, so explicitness there --- and throughout the widened band of \cref{cor:widened-deployed}, where the floors are pigeonhole-selected deep prefixes --- would require explicit families with prescribed deep prefixes, i.e.\ explicit subsets of $\mu_N$ with many prescribed elementary symmetric functions. We record this as the surviving form of Open Problem~\ref{prob:explicit}.
\end{remark}

\subsection{The failure profile: the deep-point ceiling}\label{sec:answers-profile}

\Cref{thm:A} was used throughout with $\eta=\frac12$, giving error floors of $\frac1{2k}(1-\frac nq)$. The dichotomy improves when the list is far above threshold, and the improvement is exactly quantifiable.

\begin{proposition}[optimized failure profile]\label{prop:profile}
Let $C=\RS[\F,D,k]$, let $\delta$ be integer-admissible \textup($\lfloor\delta n\rfloor\le n-k-1$, $q>n$\textup), and put $T:=\frac1k\bigl(1-\frac nq\bigr)$. If some received word has $\bigl|\Lst(C^+,\delta,u)\bigr|\ge1+\kappa qT$ for a real $\kappa>0$, then
\[
\eca(C,\delta)\ \ge\ \frac{\kappa}{\kappa+1}\,T .
\]
\end{proposition}

\begin{proof}
Write $x:=\eca(C,\delta)$; if $x\ge T$ there is nothing to prove, so assume $x<T$ and set $\eta:=x/T\in[0,1)$. Then $x=\eta T$, so \cref{thm:A} applies and gives $L\le\lceil qx/(1-\eta)\rceil\le qx/(1-\eta)+1$, i.e.\ $\kappa qT\le L-1\le qx/(1-x/T)$. Rearranging, $\kappa(T-x)\le x$, i.e.\ $x\ge\kappa T/(\kappa+1)$.
\end{proof}

\begin{corollary}[deployed profile at the first staircase step]\label{cor:profile-deployed}
At the KoalaBear row, the fiber floor of \cref{lem:phi-fiber} at scale $2^{13}$ gives $L\ge\lceil\binom{256}{130}/p\rceil>2^{220}$ at agreement $k+2^{14}$, hence $\kappa\ge2^{54}$ and
\[
\eca\bigl(C,\delta\bigr)\ \ge\ \bigl(1-2^{-54}\bigr)\,\frac1k\Bigl(1-\frac nq\Bigr)
\]
for every integer-admissible $\delta\in[\frac12-2^{-7},\frac12-2^{-21}]$, with the same lower bound for $\emca(C,\delta)$ throughout $[\frac12-2^{-7},\frac12)$: twice the printed constant of \cref{cor:deployed}, saturating the dichotomy's ceiling to relative error $2^{-54}$. On the circle line-round row the same floor gives $\kappa\ge2^{116}$ and the factor $1-2^{-116}$.
\end{corollary}

\begin{proof}
$\kappa=\lfloor(L-1)k/(q-n)\rfloor\ge2^{54}$ \textup(resp.\ $2^{116}$\textup) is verified exactly; $\kappa/(\kappa+1)\ge1-2^{-54}$; the list persists at every larger radius, and \cref{fact:chain} with \cref{lem:mca-monotone} extend as usual.
\end{proof}

\begin{remark}[the ceiling, and what keeps Open Problem~\ref{prob:errorone} open]\label{rem:profile-ceiling}
Two ceilings frame Open Problem~\ref{prob:errorone}. First, the route ceiling: as $\kappa\to\infty$ the profile of \cref{prop:profile} tends to $T=\frac1k(1-\frac nq)$ and no list, however large, certifies more through \cref{thm:A} --- error beyond $\approx\frac1k$ requires strengthening the dichotomy itself, which is exactly the unnamed problem following Open Problem~\ref{prob:errorone}. Second, at the specific radius $1-\rho-\frac1{64}$ of Open Problem~\ref{prob:errorone} over \emph{prime} fields $2^{150}\le p\le2^{256}$, the present floors certify nothing at all: the only subfield is $\B=\F_p$, so every prefix frontier obeys $g^*\le1/\beta=1/\log_2p\le\frac1{150}<\frac1{64}$, and the radius lies strictly outside the entropy--subfield envelope. \textup(There is also no confinement barrier there: \cref{lem:confine} caps $\B$-valued pairs at $|\B|/q$, which is vacuous when $\B=\F$.\textup) So Open Problem~\ref{prob:errorone} is genuinely beyond both mechanisms of this paper: it needs either non-prefix floors at macroscopic gaps over prime fields, or conversions past $\frac1k$. The certified state after this section: error at least $(1-2^{-54})\frac1k(1-\frac nq)$ at gap $2^{-7}$, error at least $\frac1{2k}(1-\frac nq)$ at the widened gaps, and error one still open everywhere.
\end{remark}

\subsection{Certificate closure of the two-sided sandwich}\label{sec:answers-T7}

\Cref{prop:finite-certificate} supplied the initial certificate format and \cref{tab:core} the formal core.  The pincer and the refinements above add new verdict types.  This subsection closes the certificate grammar over the statements used by the two-sided sandwich, prints the deployed certificates, and isolates the single imported proximity-gap input.

\begin{definition}[certificate grammar, second edition]\label{def:certificate-v2}
Extend \cref{def:scanner} by the following verdict types, each a tuple of integers together with a theorem tag:
\begin{enumerate}[label=\textup{(V\arabic*)},leftmargin=2.6em,start=5]
\item \textbf{Prefix-floor \textup(PF\textup) / rational-floor \textup(RF\textup).} A scale datum \textup(polynomial or Chebyshev $c$, or rational type $(a,e)$\textup), integers $(m,w,K)$, and the single inequality $\binom Nm>|\B|^w\cdot\Theta$, $\Theta\in\{q/k+1,\ \varepsilon^*q\}$: mark the band $[1-am/n,1-\rho)$ MCA-unsafe \textup(via \cref{prop:graded-prefix-floor,prop:graded-rational-floor}, \cref{thm:A}\textup) resp.\ $1-am/n$ list-unsafe.
\item \textbf{Explicit-head \textup(XH\textup).} Integers $(c,m,K)$ with $cm\le K-1+2c$ and the inequality $\binom{N/2-\epsilon}{\lfloor m/2\rfloor}>\Theta$, $\epsilon\in\{0,1\}$ by parity: as \textup{(V5)}, with the received word printed in closed form \textup(\cref{thm:explicit-head-floor}\textup).
\item \textbf{Mutual-from-correlated \textup(MC\textup).} An integer $r$ with $2r\le w_{\min}-1$ and a CA-verdict handle: transfer a CA-safe verdict at radius $r/n$ to MCA-safe with error $\max\{\varepsilon_{\rm CA},\,r/q\}$ \textup(\cref{thm:mca-from-ca}\textup). The CA handle is either \textbf{HJ} --- integers $(r,L_2)$ with $n-2r>0$, $(n-2r)^2>(k-1)n$, $(L_2+1)\bigl((n-2r)^2-(k-1)n\bigr)>n(n-2r-k+1)$, and $(1+(r+1)L_2)/q\le\varepsilon^*$, certifying the self-contained half-Johnson CA bound of \cref{thm:elementary-ca} --- or \textbf{CH} --- the inequality $2r\le n-k$ with error $n/q\le\varepsilon^*$, tagged \textup{\textsc{import}}:~\cite[Thm.~4.1]{BCIKS20} \textup(\cref{cor:conditional-half}\textup).
\item \textbf{Profile \textup(PR\textup).} A floor handle and an integer $\kappa$ with $L\ge1+\kappa q T$: certify error $\ge\frac{\kappa}{\kappa+1}T$ \textup(\cref{prop:profile}\textup).
\end{enumerate}
\end{definition}

\begin{theorem}[certificate closure of the two-sided sandwich]\label{thm:T7-closure}
With the grammar of \cref{prop:finite-certificate} extended by \cref{def:certificate-v2}:
\begin{enumerate}[label=\textup{(\roman*)}]
\item \textup(Soundness\textup) Every verdict of \cref{def:certificate-v2} is implied by its cited theorem once its integer inequalities hold; \cref{thm:scanner-sound} extends verbatim.
\item \textup(Completeness for the printed results\textup) Every clause of \cref{thm:sandwich2}, of \cref{tab:scanner2}, of \cref{cor:ecfft-macroscopic}, of \cref{cor:circle-anyfield}, and of \cref{cor:explicit-deployed,cor:profile-deployed} is certified by a certificate of \cref{prop:finite-certificate} or \cref{def:certificate-v2}; the deployed certificates are printed in \cref{tab:certs}.
\item \textup(Complexity\textup) A \textup{(V5)}/\textup{(V6)} certificate is checked with $O(N)$ multiplications of integers of $O(N+\log q)$ bits \textup(the binomial by its product formula, the power $|\B|^w$ by repeated squaring\textup); every other certificate with $O(1)$ arithmetic operations on integers of $O(\log q+\log n)$ bits. A full row is checked in $O(n)$ big-integer multiplications.
\item \textup(Import isolation\textup) Exactly one imported statement occurs anywhere in the grammar: the \textup{\textsc{import}} tag of the CH handle. Deleting it degrades only the safe edge, from $(1-\rho)/2$ to $\max\{(1-\rho)/3,\ \text{half-Johnson}\}$, and leaves every other verdict intact. The remaining statements --- \cref{tab:core} extended by the rows of \cref{tab:core2} --- quantify over explicitly bounded finite sets and use only the listed elementary ingredients; their dependency closure is import-free.
\end{enumerate}
\end{theorem}

\begin{proof}
(i) is quotation: each verdict's hypotheses are the hypotheses of its theorem, specialized to integers. (ii) is inspection of the quoted proofs against \cref{tab:certs}: each proof consists of its printed integer inequalities plus the cited assembly steps, all of which are verdicts or core statements. (iii): $\binom Nm=\prod_{j=1}^m\frac{N-m+j}j$ uses $2m\le2N$ multiplications and divisions of integers of at most $N+O(\log N)$ bits; thresholds and powers are $O(\log w)$ squarings of $O(w\log|\B|+\log q)$-bit integers, and $w\log|\B|=O(N)$ in every certificate above; the remaining verdicts compare $O(1)$ products. (iv): the only \textup{\textsc{import}} tag is CH; \cref{thm:sandwich2} lists its conditional clause separately, and the self-contained clauses cite only \cref{thm:deep-mca,thm:elementary-ca,thm:mca-from-ca} on the safe side. The new core rows are in \cref{tab:core2}; each proof above is elementary in the stated ingredients, as the reader may verify line by line.
\end{proof}

\begin{table}[ht]
\centering
\small
\begin{tabular}{@{}ll@{}}
\toprule
Added core statement & Proof ingredients \\
\midrule
\cref{thm:mca-from-ca} \textup(mutual from correlated\textup) & minimum weight, two-point interpolation \\
\cref{thm:elementary-ca} \textup(half-Johnson CA\textup) & pair-list double counting, convexity \\
\cref{prop:graded-prefix-floor}, \cref{prop:graded-rational-floor} & symmetric functions, graded pigeonhole \\
\cref{thm:explicit-head-floor} \textup(explicit heads\textup) & antipodal pairing, no pigeonhole \\
\cref{thm:explicit-pairs} \textup(explicit pairs\textup) & \cref{thm:A}'s collision count, Markov \\
\cref{prop:profile} \textup(optimized profile\textup) & one-line rearrangement of \cref{thm:A} \\
\cref{lem:stereographic} \textup(stereographic model\textup) & conic parametrization, basis change \\
\cref{def:certificate-v2}, \cref{thm:T7-closure} & finite case analysis over the above \\
\bottomrule
\end{tabular}
\caption{The formal finite core, continued: statements added by the pincer and the refinements, extending \cref{tab:core}. Each row quantifies over explicitly bounded finite sets; the dependency order is top to bottom, below the rows of \cref{tab:core}. The one import (\textsc{CH}) is not a core row.}
\label{tab:core2}
\end{table}

\begin{table}[ht]
\centering
\small
\begin{tabular}{@{}llll@{}}
\toprule
Row / verdict & Certificate & Verified inequality & \\
\midrule
KB MCA, PF & $(c,m,w)=(16,\,69748,\,4211)$ & $\binom{2^{17}}{69748}>p^{4211}(q/k+1)$ & exact \\
KB list, PF & $(32,\,34874,\,2106)$ & $\binom{2^{16}}{34874}>p^{2106}\,2^{-128}q$ & exact \\
Circle MCA, PF & $(32,\,34873,\,2104)$ & $\binom{2^{16}}{34873}>p'^{2104}(q/k+1)$ & exact \\
Circle list, PF & $(32,\,34874,\,2106)$ & $\binom{2^{16}}{34874}>p'^{2106}\,2^{-100}q$ & exact \\
KB MCA, XH & $(2^{12},\,258)$, $u=X^{k+2^{13}}|_D$ & $\binom{256}{129}>q/k+1$: $2^{251.6}$ vs $2^{165.94}$ & exact \\
KB list, XH & $(2^{14},\,65)$, one head & $\binom{63}{32}>2^{-128}q$: $2^{59.66}$ vs $2^{57.94}$ & exact \\
Circle MCA, XH & $(2^{12},\,258)$, $u=T_c^{258}|_D$ & $\binom{256}{129}>q/k+1$: $2^{251.6}$ vs $2^{104}$ & exact \\
Circle list, XH & $(2^{14},\,65)$, one head & $\binom{63}{32}>2^{-100}q$: $2^{59.66}$ vs $2^{24}$ & exact \\
Both rows, MC & $r=\lfloor(n-k)/2\rfloor=2^{19}$ & $2r\le n-k=2^{20}$ & --- \\
KB safe, HJ & $(r,L_2)=(307121,\,1001282)$ & $(1+(r+1)L_2)/q<2^{-147}$ & exact \\
KB safe, CH \textsc{import} & $r=2^{19}$ & $n/q<2^{-164}$ & --- \\
Circle safe, CH \textsc{import} & $r=2^{19}$ & $n/q<2^{-102}$ & --- \\
KB profile, PR & $\kappa=2^{54}$ & $\lceil\binom{256}{130}/p\rceil>1+2^{54}q T$ & exact \\
$i$-free circle, RF & $(m,w)=(2^{20}{+}2^{15}{+}1,\,2^{16})$ & $\binom{2^{21}}{m}>p'^{\,2^{16}}2^{256}$ & exact \\
\bottomrule
\end{tabular}
\caption{Printed certificates for the deployed rows (\cref{def:certificate-v2}; $T=\frac1k(1-\frac nq)$). Each line is one integer comparison; ``exact'' means it was verified in exact integer arithmetic. The two \textup{\textsc{import}} lines are optional imported half-distance clauses --- the only non-self-contained entries.}
\label{tab:certs}
\end{table}

\medskip
\noindent\textbf{Remaining residues.}
The refinements above reduce the remaining landscape to a small number of explicit residues.  The genus-one and circle-code questions are reduced to the same aperiodic-band shape after the rational-floor and stereographic transports of \cref{cor:ecfft-macroscopic,cor:circle-anyfield}.  The explicit-witness question is resolved at the level of closed-form words and residue-line pairs in \cref{cor:explicit-deployed,thm:explicit-pairs}, with the remaining scope stated in \cref{rem:explicit-scope}.  The error-one profile remains open: \cref{cor:profile-deployed} raises the certified profile to $(1-2^{-54})\frac1k(1-\frac nq)$, while \cref{rem:profile-ceiling} explains why reaching $1-o(1)$ requires new input.  The central open question remains Open Problem~\ref{prob:band}, in the narrowed CA-form shape of \cref{cor:band-reduction} and \cref{thm:sandwich2}.  The certificate material above is finite and checkable; a proof-assistant transcription would be implementation work rather than a new mathematical ingredient.


\section{Discussion}\label{sec:discussion}

\paragraph{Values versus fibers.}
Above the per-rate value-set reach of \cite{Cho26b} (roughly $2^{150}$--$2^{162}$), the value route --- showing that the slope multiset $\eone(\ell^\wedge Q)$ covers a $2^{-128}$ fraction of $\F$ at a fixed large prime --- runs into the per-fiber collision problem (\cite[Rem.\ ``The remaining band'']{Cho26b}): characteristic-zero locator values can collide arbitrarily badly under reduction at a worst-case prime. \cite{Cho26b} therefore closed the band at the $\delta^*$ grade only conditionally and with weaker constants --- error $2^{-42}$ at gap $2^{-11}$, assuming $|\F|\ge2n$, via quotient-core lists and the same conversion \cref{thm:A} --- leaving the \emph{error-one} grade open. The present route also never counts values, and improves each parameter of that cap: pigeonhole guarantees one heavy fiber regardless of how the values $\eone(A)$ distribute, \cref{thm:A} converts heavy fibers into correlated-agreement error directly, and the subfield pigeonhole carries the construction to extension fields with no condition relating $n$ and $|\F|$. The cost relative to the error-one results of \cite{Cho26a,Cho26b} below $2^{150}$ is quantitative: error $\approx\frac1{2k}$ instead of $1-o(1)$, and gap $2^{-9}$ (uniformly $2^{-10}$ across all four rates) instead of $\frac1{64}$. For the prize question --- is $\emca\le2^{-128}$ achievable above the stated gaps? --- this cost is immaterial.

\paragraph{Limits of the method.}
Hypothesis \eqref{eq:hyp} needs $N\,\Hb(\rho)\gtrsim2\log_2q-\log_2k$, so the smallest certifiable gap is
\[
\frac2N\;\approx\;\frac{2\,\Hb(\rho)}{2\log_2q-\log_2k}\;\approx\;\frac{\Hb(\rho)}{\log_2 q},
\]
and the certified error is capped at $\approx\frac1k$ by the shape of \cref{thm:A} (any $\eta<1$). Since $q$ enters only as the size of the field from which the line challenge $\gamma$ is drawn, this is also a protocol-facing floor: enlarging the challenge field cannot push a deep-point-conversion failure below gap $\approx\Hb(\rho)/\log_2 q_{\mathrm{line}}$, which quantitatively bounds the ``larger challenge field'' fallback in certificate frameworks. The construction also needs $N\le n$, i.e.\ $\log_2q\lesssim\Hb(\rho)\,n/2$; this covers the entire cryptographic regime $\mathrm{poly}(n)\le q\le2^{o(n)}$, but for $q\ge2^{\Hb(\rho)n/2}$ the method is vacuous and the value-route caps of \cite{Cho26a,Cho26b} are the only ones known. Within the band, the \emph{error-one} question at gap $\frac1{64}$ (Open Problem~\ref{prob:errorone}) remains exactly as open as before: the present result bounds $\delta^*$, not the failure profile near capacity.

\paragraph{Comparison with the cyclotomic sieve.}
The sieve of \cite{Cho26a} produces, for infinitely many primes, failures of error $(\log p)^{-6}$ at gap $\Theta(1/\log p)$, but its prime set is given by Siegel--Walfisz and is ineffective. \Cref{cor:grand} trades error strength for universality: every field below $2^{256}$, fully effective, and elementary given \cref{thm:A}. The two results triangulate the failure landscape from different directions.

\paragraph{Isogeny-smooth reach and the lacunarity dichotomy.}
The counting side of the method extends beyond multiplicative smoothness. The fiber lemma needs only a polynomial folding map with complete uniform fibers inside the domain (\cref{def:map-smooth,lem:phi-fiber}), and the divisibility hypothesis $a\mid k$ is removed by rounding, so the universal cap now covers mixed-radix multiplicative rows without integral $\rho N$ (\cref{cor:mixed-radix}), the $x$-coordinate line rounds of circle FRI over $p\equiv3\pmod4$, and --- through the torus uniformization, over any challenge field containing $i$ --- the bivariate circle codes themselves, at the deployed Mersenne-$31$/QM31 parameters with error $>2^{-23}$ at gap $2^{-7}$ (\cref{cor:circle-grand,cor:circle-deployed}). The boundary of the method is now a clean dichotomy in the degree type $(a,e)$ of the folding map: polynomial towers ($e=0$: powers, Chebyshev) give gap $\approx2/N$, while elliptic isogeny towers ($e=a-1$: ECFFT) collapse the slack-two window to gap $2/n$ (\cref{rem:ecfft-obstruction}). Subfield confinement and the extension-pole mechanism also carry over to circle rows (\cref{cor:circle-Fvalued}), so the deployed circle failures are again fiber-borne, genuinely extension-valued, and constructively reachable from any supplied list via extension poles.

\paragraph{Pressure on the packing conjecture.}
By the exact normal form of \cite{Cho26b}, $\emca(C,\delta)=\max_t\Lambda^{\mathrm{NC}}_t(\delta)/q$, where $\Lambda^{\mathrm{NC}}$ is a noncontained residue-line packing number. \Cref{cor:deployed} therefore forces $\max_t\Lambda^{\mathrm{NC}}\ge(q-n)/2k\approx2^{164}$ at gap $2^{-7}$ on the deployed sextic --- superpolynomially many noncontained residue lines at $n=2^{21}$. This puts no direct pressure on the final MCA conjecture of \cite[Conj.\ ``Final floor- and quotient-corrected MCA asymptotic'']{Cho26b}: at gap $2^{-7}$ both of its corrected-reserve hypotheses fail ($\sigma=2^{14}$ lies below $Cn/\log_2n\approx2^{16.6}$, and $\sigma\log_2 q_{\mathrm{gen}}=2^{14}\cdot31\approx2^{19}$ is far below $\log_2\binom n{k+\sigma}\approx2^{21}$), so the conjectured ceiling of the form $\bigl(n^{1+o(1)}+2^{(\beta/\Hb)\,Q_{\Hb}(\eta)}\bigr)/q$ is simply not asserted there --- the packing mass extracted through \cref{thm:A} lives in the region the framework already marks as below-floor. The genuine structural question this paper raises for the framework of \cite{Cho26b} is whether conversion-extracted fiber mass persists \emph{above} the corrected reserve: whether, at reserves where the conjecture does apply, the fiber route still produces packings beyond the $n^{1+o(1)}$ aperiodic term plus the quotient-periodic term, which would force an additional fiber-borne term into the conjectured ceiling. This is open.

\paragraph{Verification caveat.}
The main cap no longer depends on importing the Crites--Stewart theorem: \cref{thm:A} proves the needed deep-point list-to-CA conversion directly, including the integer-radius condition $f\le n-k-1$.  \Cref{thm:B} remains a separate slacked fallback imported through the BCHKS/ABF route, and readers using that fallback should consult the underlying source directly.  The main composition is deliberately shallow: it uses the self-contained conversion, the locator-fiber pigeonhole, and support-wise MCA monotonicity.

\paragraph{What remains after the pincer and certificate closure.}
The structural ledgers, extension-line accounting, scanner, and certificate grammar are closed at the level needed for the printed sandwich.  The deep regime and the Johnson list regime are theorems; the near-capacity region is unsafe by the quotient and graded floors; and the genus-one and circle field-of-definition extensions are covered by \cref{cor:ecfft-macroscopic,cor:circle-anyfield}.  What remains is the aperiodic band conjecture of Open Problem~\ref{prob:band}, the error-one profile question of Open Problem~\ref{prob:errorone}, and the explicit deployed-pair strengthening of Open Problem~\ref{prob:explicit}.

\paragraph{Formerly open questions resolved above.}
The genus-one macroscopic-cap and circle field-of-definition questions are resolved by \cref{cor:ecfft-macroscopic,cor:circle-anyfield}.  The explicit-witness question is resolved at the first staircase step and for simple-pole certifying pairs at gap $2^{-8}$ \textup(\cref{cor:explicit-deployed,thm:explicit-pairs}\textup); its full-gap strengthening is restated below.

\paragraph{Remaining open problems.}

\begin{problem}[explicit deployed witnesses at the full deployed gap]\label{prob:explicit}
For the KoalaBear-sextic deployed row, exhibit explicit pairs $(f,g)$ with MCA-bad-slope density $>2^{-22}$ at the gap $2^{-7}$ certified by \cref{cor:deployed}; on the integer-admissible subinterval, exhibit the corresponding CA-bad-slope density as well.  By \cref{cor:Fvalued} such pairs exist and are not $\B$-rational, and the natural candidates are residue-line normal forms with denominator $E\in\F[X]$ not defined over $\B$.  The present paper gives explicit received words and simple-pole certifying pairs of exactly this normal form, $E=X-\alpha$, at the first staircase step and at gap $2^{-8}$ \textup(\cref{cor:explicit-deployed,thm:explicit-pairs}\textup); the remaining task is to push the explicit pair construction to the full deployed gap, derandomize the choice of the pole $\alpha$, and print the deployed-size bucket tables.
\end{problem}

\begin{problem}\label{prob:errorone}
Error one in the band: does $\emca(C,1-\rho-\frac1{64})=1-o(1)$ hold for all primes $2^{150}\le p\le2^{256}$? This is the per-fiber collision problem proper; \cref{cor:grand} caps $\delta^*$ there but says nothing above error $\approx\frac1{2k}$.  \Cref{cor:profile-deployed} raises the certified profile to $(1-2^{-54})\frac1k$ at gap $2^{-7}$, and \cref{rem:profile-ceiling} shows the question lies strictly beyond both mechanisms of this paper; it remains open.
\end{problem}

\begin{problem}
Improve the threshold of \cref{thm:A}: any strengthening of the hypothesis range $\eca\le\eta/k$ toward $\eca\le\eta\cdot\mathrm{polylog}(q)/k$, or of the list conclusion, immediately tightens the certified error of \cref{thm:main} toward $1/\mathrm{poly}\log$; the gap, by contrast, is governed by the binomial in \eqref{eq:hyp} and would need a denser fiber construction to shrink.  Cf.\ \cref{prop:profile}: optimizing the dichotomy already attains $(1-o(1))/k$, so passing $1/k$ is exactly what a polylog strengthening would buy.
\end{problem}

\begin{problem}[aperiodic band on genus-one and circle tangent domains]\label{prob:isogeny}
Extend the aperiodic band conjecture, in the narrowed CA-form of \cref{cor:band-reduction}, to the tangent-domain Reed--Solomon models arising from genus-one isogeny towers and to the stereographic circle models over challenge fields not containing $i$.  The macroscopic unsafe floors on these rows are supplied by \cref{cor:ecfft-macroscopic,cor:circle-anyfield}, and the generic safe side applies verbatim; what remains is the same safe-side band problem as in Open Problem~\ref{prob:band}, transported through the rational scale map.
\end{problem}

\paragraph{Acknowledgements.} This paper builds on the Crites--Stewart conversion perspective \cite{CS25}, the BCHKS theorem \cite{BCHKS25}, and the locator framework of \cite{Cho26a,Cho26b}, following the synthesis of Arnon--Boneh--Fenzi \cite{ABF26}; it proves the main conversion needed here directly, sharpens the conditional universal-cap theorem, and refines the subfield-confinement theorem of \cite{Cho26b}.


\begin{thebibliography}{BCHKS25}

\bibitem[PP26]{ProximityPrize}
The Proximity Prize.
\newblock The Proximity Prize: \$1M Reed--Solomon Challenge.
\newblock Preliminary release, 2026.
\newblock \url{https://proximityprize.org/}.


\bibitem[ABF26]{ABF26}
G.~Arnon, D.~Boneh, and G.~Fenzi.
\newblock Open problems in list decoding and correlated agreement.
\newblock IACR Cryptology ePrint Archive, Report 2026/680, 2026.
\newblock \url{https://eprint.iacr.org/2026/680}.

\bibitem[BCHKS25]{BCHKS25}
E.~Ben-Sasson, D.~Carmon, U.~Hab{\"o}ck, S.~Kopparty, and S.~Saraf.
\newblock On proximity gaps for Reed--Solomon codes.
\newblock IACR Cryptology ePrint Archive, Report 2025/2055, 2025.
\newblock \url{https://eprint.iacr.org/2025/2055}.


\bibitem[BCIKS20]{BCIKS20}
E.~Ben-Sasson, D.~Carmon, Y.~Ishai, S.~Kopparty, and S.~Saraf.
\newblock Proximity gaps for Reed--Solomon codes.
\newblock In \emph{Proc.\ 61st IEEE Symposium on Foundations of Computer Science (FOCS)}, 2020.
\newblock Extended manuscript: IACR Cryptology ePrint Archive, Report 2020/654.

\bibitem[BCKL21]{BCKL21}
E.~Ben-Sasson, D.~Carmon, S.~Kopparty, and D.~Levit.
\newblock Elliptic curve fast Fourier transform (ECFFT) part~I: Fast polynomial algorithms over all finite fields.
\newblock arXiv:2107.08473, 2021.

\bibitem[Cho26a]{Cho26a}
P.~Chojecki.
\newblock Capacity-edge obstructions to Reed--Solomon mutual correlated agreement over smooth multiplicative domains.
\newblock Manuscript, 2026.

\bibitem[Cho26b]{Cho26b}
P.~Chojecki.
\newblock Slack, quotient cores, and the entropy gap for smooth-domain Reed--Solomon codes: list fibers, cyclotomic rigidity, Galois-amplified collisions, and the corrected slack mutual-correlated-agreement theory.
\newblock Manuscript, merged corrected edition, 2026.

\bibitem[CS25]{CS25}
E.~Crites and A.~Stewart.
\newblock On Reed--Solomon proximity gaps conjectures.
\newblock IACR Cryptology ePrint Archive, Report 2025/2046, 2025.
\newblock \url{https://eprint.iacr.org/2025/2046}.

\bibitem[HLP24]{HLP24}
U.~Hab{\"o}ck, D.~Levit, and S.~Papini.
\newblock Circle STARKs.
\newblock IACR Cryptology ePrint Archive, Report 2024/278, 2024.
\newblock \url{https://eprint.iacr.org/2024/278}.

\end{thebibliography}

\end{document}
